% Options for packages loaded elsewhere
\PassOptionsToPackage{unicode}{hyperref}
\PassOptionsToPackage{hyphens}{url}
\documentclass[
  ukrainian,
]{article}
\usepackage{xcolor}
\usepackage[margin=2.5cm]{geometry}
\usepackage{amsmath,amssymb}
\setcounter{secnumdepth}{-\maxdimen} % remove section numbering
\usepackage{iftex}
\ifPDFTeX
  \usepackage[T1]{fontenc}
  \usepackage[utf8]{inputenc}
  \usepackage{textcomp} % provide euro and other symbols
\else % if luatex or xetex
  \usepackage{unicode-math} % this also loads fontspec
  \defaultfontfeatures{Scale=MatchLowercase}
  \defaultfontfeatures[\rmfamily]{Ligatures=TeX,Scale=1}
\fi
\usepackage{lmodern}
\ifPDFTeX\else
\setmainfont{DejaVu Serif}
\setsansfont{DejaVu Sans}
\setmonofont{DejaVu Sans Mono}
\setmathfont{STIX Two Math}
\fi
\ifPDFTeX\else
  % xetex/luatex font selection
\fi
% Use upquote if available, for straight quotes in verbatim environments
\IfFileExists{upquote.sty}{\usepackage{upquote}}{}
\IfFileExists{microtype.sty}{% use microtype if available
  \usepackage[]{microtype}
  \UseMicrotypeSet[protrusion]{basicmath} % disable protrusion for tt fonts
}{}
\makeatletter
\@ifundefined{KOMAClassName}{% if non-KOMA class
  \IfFileExists{parskip.sty}{%
    \usepackage{parskip}
  }{% else
    \setlength{\parindent}{0pt}
    \setlength{\parskip}{6pt plus 2pt minus 1pt}}
}{% if KOMA class
  \KOMAoptions{parskip=half}}
\makeatother
\usepackage{longtable,booktabs,array}
\newcounter{none} % for unnumbered tables
\usepackage{calc} % for calculating minipage widths
% Correct order of tables after \paragraph or \subparagraph
\usepackage{etoolbox}
\makeatletter
\patchcmd\longtable{\par}{\if@noskipsec\mbox{}\fi\par}{}{}
\makeatother
% Allow footnotes in longtable head/foot
\IfFileExists{footnotehyper.sty}{\usepackage{footnotehyper}}{\usepackage{footnote}}
\makesavenoteenv{longtable}
\usepackage{graphicx}
\graphicspath{{/Users/kalinicenkovitalijmikolajovic/.codex/worktrees/5c15/Script_for_translate_book/Output_Final/}}
\setkeys{Gin}{width=0.82\linewidth,height=0.82\textheight,keepaspectratio}
\makeatletter
\newsavebox\pandoc@box
\newcommand*\pandocbounded[1]{% scales image to fit in text height/width
  \sbox\pandoc@box{#1}%
  \Gscale@div\@tempa{\textheight}{\dimexpr\ht\pandoc@box+\dp\pandoc@box\relax}%
  \Gscale@div\@tempb{\linewidth}{\wd\pandoc@box}%
  \ifdim\@tempb\p@<\@tempa\p@\let\@tempa\@tempb\fi% select the smaller of both
  \ifdim\@tempa\p@<\p@\scalebox{\@tempa}{\usebox\pandoc@box}%
  \else\usebox{\pandoc@box}%
  \fi%
}
% Set default figure placement to htbp
\def\fps@figure{htbp}
\makeatother
\ifLuaTeX
\usepackage[bidi=basic,shorthands=off,provide=*]{babel}
\else
\usepackage[bidi=default,shorthands=off,provide=*]{babel}
\fi
\ifLuaTeX
  \usepackage{selnolig} % disable illegal ligatures
\fi
\setlength{\emergencystretch}{3em} % prevent overfull lines
\providecommand{\tightlist}{%
  \setlength{\itemsep}{0pt}\setlength{\parskip}{0pt}}
\usepackage{bookmark}
\IfFileExists{xurl.sty}{\usepackage{xurl}}{} % add URL line breaks if available
\urlstyle{same}
\hypersetup{
  pdflang={uk},
  hidelinks,
  pdfcreator={LaTeX via pandoc}}

\author{}
\date{}

\begin{document}

\section{РОЗВ'ЯЗАННЯ ДИФЕРЕНЦІАЛЬНИХ РІВНЯНЬ У ПОХІДНИХ ПОХІДНИХ ЗА
ДОПОМОГОЮ МЕТОДУ КОНФОРМНИХ
ВІДОБРАЖЕНЬ}\label{ux440ux43eux437ux432ux44fux437ux430ux43dux43dux44f-ux434ux438ux444ux435ux440ux435ux43dux446ux456ux430ux43bux44cux43dux438ux445-ux440ux456ux432ux43dux44fux43dux44c-ux443-ux43fux43eux445ux456ux434ux43dux438ux445-ux43fux43eux445ux456ux434ux43dux438ux445-ux437ux430-ux434ux43eux43fux43eux43cux43eux433ux43eux44e-ux43cux435ux442ux43eux434ux443-ux43aux43eux43dux444ux43eux440ux43cux43dux438ux445-ux432ux456ux434ux43eux431ux440ux430ux436ux435ux43dux44c}

МЕТА ЛЕКЦІЇ: Показати, як деякі двовимірні граничні задачі можна
перетворити на інші, простіші задачі за допомогою конформного
відображення. Припустимо, наприклад, що ми хочемо розв'язати рівняння
Лапласа

\begin{verbatim}
\begin{cases} u(0, t) = 0, \\ u_x(1, t) = -u(1, t), \\ u(x, 0) = x, \\ u_t(x, 0) = 0. \end{cases}
\end{verbatim}

з певною граничною умовою в області комплексної форми (площина змінних x
і y). Цю задачу можна перетворити на нову, в якій потрібно знайти
розв'язок рівняння Лапласа

у простішій області (площина змінних і і і).

\pandocbounded{\includegraphics[keepaspectratio]{images/_page_353_Figure_9.jpeg}}

\begin{verbatim}
У конформному відображенні рівняння Лапласа \varphi_{xx}+\varphi_{yy}=0 у координатній площині (x,y) повертається у рівняння Лапласа \varphi_{nn}\dashv -\varphi_{vo}=0 у координатній площині (u,v). Інші рівняння, такі як \varphi_{xx}+2\varphi_{xy}+\varphi_{yy}=0 , не мають такої властивості. Після розв'язання рівняння Лапласа \varphi(u,v) у простій області (
\end{verbatim}

\begin{verbatim}
наприклад, у колі, напівплощині, квадраті) знайдено, достатньо підставити вирази U=U(X, Y), V=V(X, Y) у цей розв'язок, і ми отримаємо розв'язок нашої задачі \phi[u(x, y), v(x, y)] , вираженого через початкові змінні.
\end{verbatim}

У цій лекції ми розглянемо:

\begin{enumerate}
\def\labelenumi{\arabic{enumi})}
\item
  про властивості конформних відображень,
\item
  як їх будувати (для деяких районів),
\item
  як розв'язувати задачі методом конформних відображень.
\end{enumerate}

Однією з головних труднощів, що виникають при розв'язанні задач меж, є
складна форма кордонів. Навіть межі відносно простої форми часто роблять
задачу дуже складною для розв'язання. Звісно, можна спробувати
розв'язати задачу за допомогою методів теорії збурень, але теорія
збурень працює добре лише у випадках, коли межа близька до простої.

Однак існує спосіб розв'язати двовимірне рівняння Лапласа в областях з
комплексною межею. Цей метод базується на конформних відображеннях.

Перш ніж почати розв'язувати задачі цим методом, давайте трохи
ознайомимося з поняттям конформного відображення та комплексних змінних
загалом.

\section{Комплексні функції та конформні
відображення}\label{ux43aux43eux43cux43fux43bux435ux43aux441ux43dux456-ux444ux443ux43dux43aux446ux456ux457-ux442ux430-ux43aux43eux43dux444ux43eux440ux43cux43dux456-ux432ux456ux434ux43eux431ux440ux430ux436ux435ux43dux43dux44f}

У цій лекції ми ознайомимося лише з тими поняттями з теорії функцій
комплексної змінної, які нам знадобляться пізніше. Представимо
комплексне число z = x + iy як точку на площині комплексного xy (рис.
47.1).

Щоб ввести поняття конформного відображення, спершу ознайомимося з
поняттям функції комплексної змінної

.

Тут z --- це комплексна змінна (з деякої області площини z), а w ---
нова комплексна змінна, яка отримає значення відповідно до формули w =
\texttt{f(z)} у новій комплексній площині w.

Наприклад, комплексна функція, визначена в \emph{першому квадранті}
комплексної z-площини, \texttt{w=z\^{}2} відображає цю область у всю
півплощину v\textgreater0 у комплексній площині w.

Таблиця 47 2

{\def\LTcaptype{none} % do not increment counter
\begin{longtable}[]{@{}
  >{\raggedright\arraybackslash}p{(\linewidth - 4\tabcolsep) * \real{0.2342}}
  >{\raggedright\arraybackslash}p{(\linewidth - 4\tabcolsep) * \real{0.2252}}
  >{\raggedright\arraybackslash}p{(\linewidth - 4\tabcolsep) * \real{0.5405}}@{}}
\toprule\noalign{}
\begin{minipage}[b]{\linewidth}\raggedright
Площина d
\end{minipage} & \begin{minipage}[b]{\linewidth}\raggedright
\texttt{z\ \textbackslash{}longrightarrow\ z\^{}2}
\end{minipage} & \begin{minipage}[b]{\linewidth}\raggedright
Площина S
\end{minipage} \\
\midrule\noalign{}
\endhead
\bottomrule\noalign{}
\endlastfoot
0 & & 0 \\
t1 → t & & -1(2) \\
Вісь x Перший квадрант & & Позитивна дійсна піввісь верхньої
півплощини \\
\end{longtable}
}

Таблиця 47.2 показує відповідність, яка встановлюється шляхом
відображення \texttt{w=z\^{}2} між деякими точками на площинах z і w.

\pandocbounded{\includegraphics[keepaspectratio]{images/_page_355_Figure_5.jpeg}}

Рис. 47.1. Площина комплексної змінної
\texttt{\textbackslash{}boldsymbol\{z\}} та некольорових формул бігунів
Р И С У Н О К

\pandocbounded{\includegraphics[keepaspectratio]{images/_page_355_Figure_7.jpeg}}

47.2. Відображення першого квадранта площини
\texttt{\{\textbackslash{}bf\ z\}} у верхню півплощину
\texttt{\{\textbackslash{}bf\ w\}} .

Інакше кажучи, якщо ми говоримо про комплексісні функції змінної w =
\texttt{f(z)}, то можна розглянути питання, як криві площини z
перетворюються на криві в площині

Це питання не ставиться в теорії функцій дійсної змінної, з якою читач
уже добре знайомий.

Перш ніж рухатися далі, давайте розглянемо, як \texttt{w=z\^{}2} складне
відображення типів у еквівалентній дійсній формі. Дійсна форма показує,
як початкові координати (x,y) у площині z перетворюються на нові
координати (u,v) у площині w \texttt{w=z\^{}2} .

\begin{verbatim}
\begin{cases} u(x, 0) = x, \\ u_t(x, 0) = 0. \end{cases}
\end{verbatim}

І давайте окремо прирівняємо реальні та уявні частини в цьому
співвідношенні. У результаті ми отримуємо

\begin{verbatim}
\begin{array}{ll} u=x^2-y^2,\\ v=2xy, \end{array}
\end{verbatim}

(действительная форма отображения \texttt{w=z\^{}2} ).

Ця форма нам знадобиться пізніше. У цьому прикладі це цікаво, оскільки
показує, як гіперболи в площині z переходять до координатних прямих u =
const і v = const у площині w.

Тепер перейдемо до вивчення одного конкретного типу функцій комплексної
змінної, які виконують конформні відображення.

\section{Определение конформного
отображения}\label{ux43eux43fux440ux435ux434ux435ux43bux435ux43dux438ux435-ux43aux43eux43dux444ux43eux440ux43cux43dux43eux433ux43e-ux43eux442ux43eux431ux440ux430ux436ux435ux43dux438ux44f}

Відображення w=\texttt{f(z)} комплексної площини z у комплексну площину
w називається конформним у точці \texttt{z\_\{\textbackslash{}bullet\}}
площини z, якщо похідна
\texttt{f\textquotesingle{}(z\_\{\textbackslash{}bullet\})\ \textbackslash{}neq\ 0}
. Відображення \texttt{f(z)} називається конформним у регіоні D, якщо
\texttt{f\textquotesingle{}(z)\ \textbackslash{}neq\ 0} в кожній точці
регіону D.

Наприклад, відображення \texttt{f(z)\ =\ z\^{}2} конформно всюди, крім z
= 0, оскільки
\texttt{z\textquotesingle{}(z)\ =\ 2z\ \textbackslash{}neq\ 0} для всіх
\texttt{z\ \textbackslash{}neq\ 0} . З іншого боку, відображення
\texttt{e\^{}z} конформно по всій площині z, оскільки скрізь
\texttt{f\textquotesingle{}(z)\ =\ e\^{}z\ \textbackslash{}neq\ 0} . Яке
призначення конформних відображень? Відповідь полягає в тому, що при
розв'язанні рівняння Лапласа
\texttt{\textbackslash{}phi\_\{xx\}\ +\ \textbackslash{}phi\_\{yy\}\ =\ 0}
в межах деякої області площини змінних x і y, ми можемо розглядати цю
площину як площину комплексної змінної z. Розглянемо комплексне
відображення w = \texttt{f(z)} площини z у площину w, у якій введені
координати (u, v). Коли w = \texttt{f(z}\phi\emph{\{xy\} + \phi}\{uy\} =
0\texttt{)} відображається, рівняння Лапласа
\texttt{\textbackslash{}varphi\_\{xx\}\ +\ \textbackslash{}varphi\_\{yy\}\ =\ 0}
перетворюється у нове диференціальне рівняння в похідних похідних, яке
залежить від нових координат і v. To --- рівняння, яке
\texttt{\textbackslash{}phi\_\{xx\}\ +\ \textbackslash{}phi\_\{uu\}\ =\ 0}
повернеться у рівняння
\texttt{\textbackslash{}varphi\_\{nn\}\ +\ \textbackslash{}varphi\_\{nn\}\ =\ 0}
. Таким чином, ми придумаємо ідею знайти конформне відображення, яке
переносить область з комплексною межею у область з простою межею
(нагадаємо, що відображення \texttt{w=z\^{}*} переносить межу першого
квадранта площини z у фактичну вісь площини w).

Ось кілька прикладів, які покажуть, як можна розв'язати диференціальні
рівняння у похідних за допомогою конформних відображень.

Уравнение Лапласа в верхней полуплоскости.

Припустимо, що нам потрібно розв'язати наступну задачу Діріхле у верхній
половині видимості (рис. 47.3):

(47.1)

\begin{verbatim}
можно найти стандартным методом разделения переменных. При этом оказывается, что решение   u\left(x,\,t\right)  , представленное в виде ряда
\end{verbatim}

Один із способів розв'язання цієї проблеми --- застосувати перетворення
Фур'є до змінної x, але набагато краще конформно відобразити верхню
половину площини у площину (u, v). Читачеві не буде складно побачити, що
конформне відображення

является суперпозицией простейших колебаний. Пространственная часть
\texttt{X\_n}(x) каждого элементарного колебания является решением
(собственной функцией) некоторой краевой задачи Штурма---Лиувилля.

\begin{verbatim}
Итак, мы научились решать волиовое уравнение   u_{II}=c^2u_{xx}   на всей прямой   -\infty < x < \infty  . Формула Даламбера представляет решение в виде суммы двух бегущих волн, распространяющихся в противоположные стороны. Если же рассматривать это же уравнение на ограниченном промежутке 0 < x < L, то бегущих волн уже не станет, так как они будут взаимодействовать с границами. Вместо них возникают другие волны, которые называют стоячими волнами. Посмотрим, например, что случится, если гитарную струну, закрепленную в точках x=0, x=L, привести в движение. Для ответа на этот вопрос нам необходимо решить следующую смешанную задачу гиперболического типа:
\end{verbatim}

Вона має такі властивості (див. RNS. 47.4).

\begin{enumerate}
\def\labelenumi{\arabic{enumi}.}
\tightlist
\item
  Вона відображає верхню півплощину площини z у смуги
  \texttt{-\textbackslash{}infty\ \textless{}\ u\ \textless{}\ \textbackslash{}infty}
  , \texttt{0\ \textless{}\ v\ \textless{}\ \textbackslash{}pi} , у
  площині w.
\end{enumerate}

\pandocbounded{\includegraphics[keepaspectratio]{images/_page_357_Figure_11.jpeg}}

Фіг. 47.3 Проблема граничного значення у верхній половині площі.

\begin{itemize}
\tightlist
\item
  \begin{enumerate}
  \def\labelenumi{\arabic{enumi}.}
  \setcounter{enumi}{1}
  \tightlist
  \item
    Ребро прямої y=0, -1 \textless{} x \textless{} 1, у площині z (де
    потенціал дорівнює 1), переходить у пряму
    \texttt{v=\textbackslash{}pi,} -
    \texttt{\textbackslash{}infty\ \textless{}\ u\ \textless{}\ \textbackslash{}infty}
    , у площині w. - 3. Два промені x\textgreater1, y=0, і x\textless-1,
    y=0, у площині z проходять відповідно у позитивний і негативний
  \end{enumerate}
\end{itemize}

Читач повинен бути впевнений, що може підтвердити ці твердження.

\pandocbounded{\includegraphics[keepaspectratio]{images/_page_358_Figure_2.jpeg}}

Рис. 47.4. Перетворення складної задачі на легку за допомогою
конформного відображення.

Важливість цих властивостей полягає в тому, що вони перетворюють
початкову вадаху (47.1) на дуже просту задачу

(47.2)

\begin{verbatim}
Розв'язок очевидний: \varphi(u, v) = \frac{1}{\pi}v .
\end{verbatim}

\begin{verbatim}
Тепер, щоб знайти розв'язок початкової задачі, потрібно пояснити v через x і y і підставити цей вираз на q(u,v)=v'\pi . Після цього ми отримуємо
\end{verbatim}

Виявляється, що в цьому випадку хвилі, що рухаються, відбиваються від
меж так, що отримані коливання стають непоширюючими, а зберігають форму
в одному положенні, тобто Ф і г

\pandocbounded{\includegraphics[keepaspectratio]{images/_page_147_Picture_5.jpeg}}

. 20.1. Три типові стоячі хвилі: \texttt{X(x)} \texttt{T(t)}.

Декілька прикладів стоячих хвиль показано на рис. 20.1. Якщо профілі цих
стоячих хвиль відомі \texttt{X\_n(x)} і природа коливань кожної хвилі
\texttt{T\_n(t)} , то розв'язок проблеми коливань гітарних струн є
суперпозицією найпростіших коливань
\texttt{X\_n(x)\textbackslash{},T\_n(t)} і

сывается в виде

\begin{verbatim}
= \ln\left[\frac{z-1}{z+1}\right] + i \arg\left[\frac{z-1}{z+1}\right]
\end{verbatim}

\texttt{Залишилося\ лише\ обрати\ шанси\ \textquotesingle{}c\_n\textquotesingle{}\ так,\ щоб\ при\ t\ =\ 0\ початкові\ умови\ \textquotesingle{}u\textbackslash{}left(x,\textbackslash{}\ 0\textbackslash{}right)=f\textbackslash{}left(x\textbackslash{}right)\textquotesingle{}\ і\ \textquotesingle{}u\_t\textbackslash{}left(x,\textbackslash{}\ 0\textbackslash{}right)=g\textbackslash{}left(x\textbackslash{}right)\textquotesingle{}\ були\ виконані.}

Задачу о колебаниях гитарной струны мы решим методом разделения
переменных.

\section{Розв'язання проблеми коливань обмеженої струни методом
розділення
змінних}\label{ux440ux43eux437ux432ux44fux437ux430ux43dux43dux44f-ux43fux440ux43eux431ux43bux435ux43cux438-ux43aux43eux43bux438ux432ux430ux43dux44c-ux43eux431ux43cux435ux436ux435ux43dux43eux457-ux441ux442ux440ux443ux43dux438-ux43cux435ux442ux43eux434ux43eux43c-ux440ux43eux437ux434ux456ux43bux435ux43dux43dux44f-ux437ux43cux456ux43dux43dux438ux445}

Чтобы решить смешанную задачу

(УЧП)

\begin{verbatim}
v = \arg\left[\frac{z-1}{z+1}\right] = \arg\left[\frac{x+iy-1}{x+1+iy}\right] =
= \arg\left[\frac{2y}{x^2+y^2-1}\right]
\varphi(x, y) = \frac{1}{\pi} \arctan\left[\frac{2y}{x^2 + y^2 - 1}\right].
\begin{cases} u(0, t) = 0, \\ u(L, t) = 0, \end{cases}
x^{2} + y^{2} = 1,
\end{verbatim}

якщо потенціал на внутрішньому колі дорівнює одинці, а на зовнішньому
колі дорівнює двом (рис. 47.5). Іншими словами, нам потрібно розв'язати
проблему

(47.3)

Проблема полягає в тому, щоб знайти конформне відображення, яке
перетворює цю область на щось простіше, де нашу проблему легко
розв'язати. У цьому випадку очевидно, що ця проста ділянка має бути
кільцем. Відповідна експозиція Р И С У Н О К

\pandocbounded{\includegraphics[keepaspectratio]{images/_page_359_Figure_10.jpeg}}

47.5. Конформне відображення на кільце.

Це вже не так очевидно, але існують довідники, де можна знайти дисплей
із вказаними властивостями. Одним із найзагальних є «Довідник конформних
відображень»

G. Kober 1), яка містить сотні конформних відображень між різними
регіонами. Якщо читач витратить трохи часу і краще ознайомиться з
конформними таблицями відображення, він знайде відповідне конформне
відображення для задачі (47.3). У нашому конкретному випадку це буде
образ виду

(47.4)

де s=-0,146, а t=-6,85. У своїй еквівалентній дійсній формі це буде

(47.5)

\begin{verbatim}
\begin{cases} u(x, 0) = f(x), \\ u_t(x, 0) = g(x), \end{cases}
v = \gamma [y(x-t)],
\gamma = 2t/[(x-t)^2 + y^2].
\end{verbatim}

найдем сначала стоячне волны, т. е. решения вида

\begin{verbatim}
\begin{aligned} \varphi_{uu} + \varphi_{vv} &= 0, & \text{в кольче,} \\ (\Gamma Y) & & & & & & & \\ (\Gamma Y) & & & & & & \\ (\Gamma Y) & & & & & \\ & & & & & \\ (\Gamma Y) & & & & & \\ & & & & \\ & & & & \\ & & & & \\ & & & & \\ & & & \\ & & & \\ & & & \\ & & & \\ & & & \\ & & & \\ & & \\ & & \\ & & \\ & & \\ & & \\ & & \\ & & \\ & & \\ & & \\ & & \\ & & \\ & & \\ & & \\ & & \\ & & \\ & & \\ & & \\ & & \\ & & \\ & & \\ & & \\ & & \\ & & \\ & & \\ & & \\ & & \\ & & \\ & & \\ & & \\ & & \\ & & \\ & & \\ & & \\ & & \\ & & \\ & & \\ & & \\ & & \\ & & \\ & & \\ & & \\ & & \\ & & \\ & & \\ & & \\ & & \\ & & \\ & & \\ & & \\ & & \\ & & \\ & & \\ & & \\ & & \\ & & \\ & & \\ & & \\ & & \\ & & \\ & & \\ & & \\ & & \\ & & \\ & & \\ & & \\ & & \\ & & \\ & & \\ & & \\ & & \\ & & \\ & & \\ & & \\ & & \\ & & \\ & & \\ & & \\ & & \\ & & \\ & & \\ & & \\ & & \\ & & \\ & & \\ & & \\ & & \\ & & \\ & & \\ & & \\ & & \\ & & \\ & & \\ & & \\ & & \\ & & \\ & & \\ & & \\ & & \\ & & \\ & & \\ & & \\ & & \\ & & \\ & & \\ & & \\ & & \\ & & \\ & & \\ & & \\ & & \\ & & \\ & & \\ & & \\ & & \\ & & \\ & & \\ & & \\ & & \\ & & \\ & & \\ & & \\ & & \\ & & \\ & & \\ & & \\ & & \\ & & \\ & & \\ & & \\ & & \\ & & \\ & & \\ & & \\ & & \\ & & \\ & & \\ & & \\ & & \\ & & \\ & & \\ & & \\ & & \\ & & \\ & & \\ & & \\ & & \\ & & \\ & & \\ & & \\ & & \\ & & \\ & & \\ & & \\ & & \\ & & \\ & & \\ & & \\ & & \\ & & \\ & & \\ & & \\ & & \\ & & \\ & & \\ & & \\ & & \\ & & \\ & & \\ & & \\ & & \\ & & \\ & & \\ & & \\ & & \\ & & \\ & & \\ & & \\ & & \\ & & \\ & & \\ & & \\ & & \\ & & \\ & & \\ & & \\ & & \\ & & \\ & & \\ & & \\ & & \\ & & \\ & & \\ & & \\ & & \\ & & \\ & & \\ & & \\ & & \\ & & \\ & & \\ & & \\ & & \\ & & \\ & & \\ & & \\ & & \\ & & \\ & & \\ & & \\ & & \\ & & \\ & & \\ & & \\ & & \\ & & \\ & & \\ & & \\ & & \\ & & \\ & & \\ & & \\ & & \\ & & \\ & & \\ & & \\ & & \\ & & \\ & & \\ & & \\ & & \\ & & \\ & & \\ & & \\ & & \\ & & \\ & & \\ & & \\ & & \\ & & \\ & & \\ & & \\ & & \\ & & \\ & & \\ & & \\ & & \\ & & \\ & & \\ & & \\ & & \\ & & \\ & & \\ & & \\ & & \\ & & \\ & & \\ & & \\ & & \\ & & \\ & & \\ & & \\ & & \\ & & \\ & & \\ & & \\ & & \\ & & \\ & & \\ & & \\ & & \\ & & \\ & & \\ & & \\ & & \\ & & \\ & & \\ & & \\ & & \\ & & \\ & & \\ & & \\ & & \\ & & \\ & & \\ & & \\ & & \\ & & \\ & & \\ & & \\ & & \\ & & \\ & & \\ & & \\ & & \\ & & \\ & & \\ & & \\ & & \\ & & \\ & & \\ & & \\ & & \\ & & \end{aligned}
\end{verbatim}

Підставляючи цей вираз у рівняння (20.1) і ділячи змінні, отримуємо дві
звичайні диференціальні рівняння

\begin{verbatim}
\varphi(r) = a \ln r + b, \quad r = u^2 + v^2.
\varphi(u, v) = 0.57 \ln(u^2 + v^2) + 1
\end{verbatim}

\begin{verbatim}
где константа   \lambda   — пока что любое действительное число —   \infty < \lambda < \infty   .
\end{verbatim}

\begin{verbatim}
Результаты исследования решений этих двух уравнений при различных значениях   \lambda   приведены на рис 20.2.
\end{verbatim}

\begin{verbatim}
Теперь мы должны *опбросить* те решения, которые либо неограничены при   t \to \infty  , либо становятся тождественно равными нулю, если потребовать выполнения граничных условий u(0, t) = u(L, t) = 0.
\end{verbatim}

\begin{verbatim}
В качестве упражнения читателю предлагается проверить, что нетривиальные и ограниченные решения можно получить только при отрицательных значениях параметра   \lambda  . Таким образом, наша задача— найти такие значения констант A, B, C, D и отрица-
\end{verbatim}

\begin{verbatim}
тельной константы   \lambda = -\beta^2  , чтобы выражение
\end{verbatim}

(20.2)

\begin{verbatim}
\varphi(x, y) = 0.57 \ln(u^2 + v^2) + 1,
\end{verbatim}

удовлетворяло граничным условиям. Если это удастся, то мы будем
располагать множеством простейших колебаний струны, каждое из которых
удовлетворяет уравнению с частными произ- Р И С

\pandocbounded{\includegraphics[keepaspectratio]{images/_page_149_Figure_5.jpeg}}

20.2. Построение стоячих воли при раздичных значениях параметра λ.

водными и граничным условиям. Затем останется только найти такую
линейную комбинацию этих колебаний, чтобы сумма удовлетворяла начальным
условиям при t=0.

Подстановка (20.2) в граничные условия \texttt{u(0,\ t)} =
\texttt{u(L,\ t)} = 0 дает

\begin{verbatim}
u(0, t) = X(0)T(t) = D\bigl[A\sin(\alpha\beta t) + B\cos(\alpha\beta t)\bigr] = 0 \Rightarrow D = 0,
\end{verbatim}

и

\begin{verbatim}
u(L, t) = X(L)T(t) = C\sin(\beta L)\bigl[A\sin(\alpha\beta t) + B\cos(\alpha\beta t)\bigr] = 0 \Rightarrow \sin(\beta L) = 0.
\end{verbatim}

\begin{verbatim}
Другими словами, константа разделения \beta должна удовлетворять уравнению \sin(\beta L) = 0, откуда
\end{verbatim}

\begin{verbatim}
\beta_n = \frac{n\pi}{L}, \qquad n = 0, 1, 2, \ldots
\end{verbatim}

\begin{verbatim}
Заметим, что если во втором уравнении (20.3) положить C=0, то получится тривиальное решение X(x)T(t) \equiv 0. Следовательно, мы нашли последовательность элементарных колебаний струны:
\end{verbatim}

\begin{verbatim}
u_n(x, t) = X_n(x)T_n(t) = \sin\!\left(\frac{n\pi x}{L}\right)\left[a_n \sin\!\left(\frac{n\pi \alpha t}{L}\right) + b_n \cos\!\left(\frac{n\pi \alpha t}{L}\right)\right].
\end{verbatim}

(20.4) або

\begin{verbatim}
u_n(x, t) = R_n \sin\!\left(\frac{n\pi x}{L}\right)\cos\!\left(\frac{n\pi \alpha (t - \delta_n)}{L}\right).
\end{verbatim}

\begin{verbatim}
Тут a_n, b_n, R_n і \delta_n — довільні сталі. Кожне таке елементарне коливання є стоячою хвилею.
\end{verbatim}

\pandocbounded{\includegraphics[keepaspectratio]{images/_page_150_Picture_5.jpeg}}

Рис. 20.3. Стоячі хвилі \texttt{u\_n}(x, t) =
\texttt{X\_n}(x)\texttt{T\_n}(t).

В нашей задаче уравнение и граничные условия линейны и однородны,
поэтому произвольная сумма элементарных стоячих волн также является
решением волнового уравнения и удовлетворяет граничным условиям. Мы
должны теперь найти такую суперпозицию стоячих волн, чтобы она
удовлетворяла также и начальным условиям, тогда получим решение нашей
задачи. Подстановка суммы

Зверніть увагу на зображення.

\pandocbounded{\includegraphics[keepaspectratio]{images/_page_362_Figure_1.jpeg}}

\pandocbounded{\includegraphics[keepaspectratio]{images/_page_362_Figure_2.jpeg}}

\begin{enumerate}
\def\labelenumi{\arabic{enumi}.}
\setcounter{enumi}{3}
\tightlist
\item
  Розв'язати змішану задачу Діріхле-Неймана з кутовим порошком за
  допомогою розчину під кутом 45°
\end{enumerate}

\pandocbounded{\includegraphics[keepaspectratio]{images/_page_362_Figure_5.jpeg}}

ПРИМІТКА. Комплексна функція
\texttt{w\ =\ \textbackslash{}ln\ z\ =\ \textbackslash{}ln\ \textbar{}z\textbar{}\ +\ i\ \textbackslash{}arg(z)}
відображає промінь \texttt{\textbackslash{}theta\ =\ c\_1} у площині z у
пряму \texttt{v\ =\ c\_1} , а коло \texttt{r\ =\ c\_2} у пряму
\texttt{u\ =\ \textbackslash{}ln\ c\_2} у площині w (див. рисунки).

\pandocbounded{\includegraphics[keepaspectratio]{images/_page_362_Figure_7.jpeg}}

\section{ТАБЛИЦІ ІНТЕГРАЛЬНИХ
ПЕРЕТВОРЕНЬ}\label{ux442ux430ux431ux43bux438ux446ux456-ux456ux43dux442ux435ux433ux440ux430ux43bux44cux43dux438ux445-ux43fux435ux440ux435ux442ux432ux43eux440ux435ux43dux44c}

Таблиця A: Експоненціальне перетворення Фур'є,

Таблиця B: Перетворення синуса Фур'є.

Таблиця St Cosine Фур'є-перетворення.

Таблиця Dr скінченне перетворення Фур'є, таблиця та скінченне косинусне
перетворення Фур'є.

Таблиця D: Трансформація Лапласа.

\section{\texorpdfstring{Прийняті позначення
\texttt{\textbackslash{}delta\textbackslash{}left(x\textbackslash{}right)}
---
дельта-функція,}{Прийняті позначення \textbackslash delta\textbackslash left(x\textbackslash right) --- дельта-функція,}}\label{ux43fux440ux438ux439ux43dux44fux442ux456-ux43fux43eux437ux43dux430ux447ux435ux43dux43dux44f-deltaleftxright-ux434ux435ux43bux44cux442ux430-ux444ux443ux43dux43aux446ux456ux44f}

в начальные условия

\begin{verbatim}
\begin{cases} 0, x < a, \\ 1, x \geqslant a, \end{cases}
\end{verbatim}

\begin{verbatim}
u_t(x, \theta) = g(x)  приводит к двум уравнениям
\end{verbatim}

\begin{verbatim}
H(x-a) =
\end{verbatim}

Функція Хевісайду,

\begin{verbatim}
H(a-x) = \begin{cases} 1, & x \leqslant a, \\ 0, & x > a, \end{cases}
\end{verbatim}

дзеркальною функцією Хевісайду.

ТАБЛИЦЯ A. Експоненціальне перетворення Фур'є

Воспользовавшись соотношениями ортогональности

1. f'(x) \textbar{} iωF (ω) \textbar{}

\textbar------------------------------\textbar-------------------------------------------------\textbar{}

\begin{verbatim}
| 2. f''(x) | -\omega^2 F(\omega) |
\end{verbatim}

\begin{verbatim}
| 3. fn (x), (n-та похідна) | (i\omega)^n F(\omega) |
\end{verbatim}

\begin{verbatim}
| 4. f(ax), a > 0 | \frac{1}{a} F\left(\frac{\omega}{a}\right) |
\end{verbatim}

\begin{verbatim}
| 5. f(x-a) | e^{-ia\omega}F(\omega) |
\end{verbatim}

\begin{verbatim}
| 6. e^{-a^2x^2} | \frac{1}{a\sqrt{2}}e^{-\omega^2/4a^2} |
\end{verbatim}

\begin{verbatim}
| 7. e^{-a_1 x} | \sqrt{\frac{2}{\pi}} \frac{a}{a^2 + \omega^2} |
\end{verbatim}

\begin{enumerate}
\def\labelenumi{\arabic{enumi}.}
\setcounter{enumi}{19}
\item
  гріх (ax)
\item
\item
\end{enumerate}

\begin{verbatim}
\begin{cases} 0, & m \neq n, \\ L/2, & m = n, \end{cases}
\end{verbatim}

\begin{enumerate}
\def\labelenumi{\arabic{enumi}.}
\setcounter{enumi}{9}
\tightlist
\item
\item
\end{enumerate}

находим коэффициенты \texttt{a\_n} и \texttt{b\_n} (20.5)

\begin{enumerate}
\def\labelenumi{\arabic{enumi}.}
\setcounter{enumi}{11}
\tightlist
\item
\end{enumerate}

Итак, задача решена. Ее решение записывается в виде

(20.6)

\begin{enumerate}
\def\labelenumi{\arabic{enumi}.}
\setcounter{enumi}{12}
\tightlist
\item
\end{enumerate}

а шанси \texttt{a\_n} і \texttt{b\_n} визначаються за формулами (20,5).
Перед завершенням лекції дозвольте зробити кілька зауважень. \# ПОВАГА
1. Якщо початкова \emph{швидкість} струни дорівнює нулю, тоді розв'язок
(20.6) має вигляд

\begin{enumerate}
\def\labelenumi{\arabic{enumi}.}
\setcounter{enumi}{13}
\tightlist
\item
\end{enumerate}

І це можна інтерпретувати так. Припустимо, ми розклали початкове
зміщення струни

\begin{enumerate}
\def\labelenumi{\arabic{enumi}.}
\setcounter{enumi}{14}
\tightlist
\item
\end{enumerate}

в ряд по синусам

\begin{enumerate}
\def\labelenumi{\arabic{enumi}.}
\setcounter{enumi}{15}
\tightlist
\item
\end{enumerate}

Тогда каждый член такого разложения должен совершать колебания

\begin{enumerate}
\def\labelenumi{\arabic{enumi}.}
\setcounter{enumi}{16}
\tightlist
\item
\end{enumerate}

Якщо ми підсумуємо всі ці коливання, отримаємо розв'язок нашої проблеми.
Наприклад, припустимо, що

\texttt{Початкове\ переміщення\ струни\ задається\ функцією\ \textquotesingle{}f(x)\ =\ \textbackslash{}sin(\textbackslash{}pi\ x/L)\ +\ 0.5\ \textbackslash{}sin(3\textbackslash{}pi\ x/\ L)\ +\ 0,25\ \textbackslash{}sin(5\textbackslash{}pi\ x/L)\textquotesingle{}\ .}

Повна відповідь системи на початкову умову такого типу дорівнює сумі
відповідей на кожен терм, тобто

\begin{enumerate}
\def\labelenumi{\arabic{enumi}.}
\setcounter{enumi}{17}
\item
\item
  Член із числом n у розкладі (20.6) називається ** модою осциляції**
  або ** p-тою гармонікою**. Використовуючи тригонометричні формули, ці
  гармоніки можна записати як
\item
\end{enumerate}

де \texttt{R\_n} і \texttt{\textbackslash{}delta\_n} --- нові довільні
константи (амплітуда та фазовий кут). Ця нова форма представлення n-го
моду виявляється більш корисною для аналізу коливань. Частота коливань
od (рад/с) та моду визначаються відповідно до формули

\begin{enumerate}
\def\labelenumi{\arabic{enumi}.}
\setcounter{enumi}{9}
\tightlist
\item
\end{enumerate}

де T --- натяг струни, p --- лінійна густина струни. Зверніть увагу, що
частота n-ї гармоніки у n разів більша за основну частоту (n=1). Ця
властивість частоти не характерна для всіх типів коливань. Приємний звук
струни скриму або гітари зумовлений множинністю обертонів основної
частоти. У барабані частоти гармонік високого порядку не є кратними
основної частоти коливань. \# ЗАВДАННЯ 1. Знайти розв'язок задачі
(20.1), якщо початкові умови задаються формулами

\begin{enumerate}
\def\labelenumi{\arabic{enumi}.}
\setcounter{enumi}{9}
\tightlist
\item
\end{enumerate}

\texttt{u\_t(x,\ 0)\ =\ 0.} Побудувати графіки розв'язків для різних
моментів часу. Чи є це рішення періодичним у часі? Який у нього період?

\begin{enumerate}
\def\labelenumi{\arabic{enumi}.}
\setcounter{enumi}{1}
\item
  Знайти розв'язок проблеми коливань струн (20.1) за початкових умов
  вигляду
\item
\end{enumerate}

\texttt{\textquotesingle{}u\_t(x,\ 0)\ =\ \textbackslash{}sin(3\textbackslash{}pi\ x/L).\textquotesingle{}\ Как\ выглядит\ график\ решения\ в\ различные\ моменты\ времени?}

\begin{itemize}
\tightlist
\item
  \begin{enumerate}
  \def\labelenumi{\arabic{enumi}.}
  \setcounter{enumi}{2}
  \tightlist
  \item
    Показати це, якщо на рис. 20.2 значення
    \texttt{\textbackslash{}lambda\ \textbackslash{}geqslant\ 0} , тоді
    розв'язки \texttt{X(x)} \texttt{T(t)} будуть або необмеженими, або
    ідентичними нулю.
  \end{enumerate}
\item
  Знайти розв'язок проблеми коливань струн за початкових умов вигляду
\end{itemize}

\begin{enumerate}
\def\labelenumi{\arabic{enumi}.}
\setcounter{enumi}{11}
\tightlist
\item
\end{enumerate}

\texttt{\textquotesingle{}u\_1(x,\ 0)\ =\ (3\textbackslash{}pi\ \textbackslash{}alpha/L)\textbackslash{}sin(3\textbackslash{}pi\ x/L).\textquotesingle{}\ 5.\ Гітарна\ струна\ довжиною\ L\ =\ 1\ тягнеться\ за\ середню\ точку\ значенням\ h\ (див.\ рис.)\ Початкове\ положення\ струни\ можна\ встановити\ як}

\begin{enumerate}
\def\labelenumi{\arabic{enumi}.}
\setcounter{enumi}{12}
\tightlist
\item
\item
\item
\item
\end{enumerate}

\begin{verbatim}
\begin{array}{llll} (\mathrm{Y} \Pi \Pi) & u_{tt} = \alpha^2 u_{xx} - \beta u_t, & 0 < x < 1, & 0 < t < \infty, \\ (\Gamma \mathrm{Y}) & \begin{cases} u \ (0, \ t) = 0, \\ u \ (1, \ t) = 0, \end{cases} & 0 < t < \infty, \\ (\mathrm{H} \mathrm{Y}) & \begin{cases} u \ (x, \ 0) = f \ (x), \\ u_1 \ (x, \ 0) = 0. \end{cases} & 0 \leqslant x \leqslant 1, \end{array}
\end{verbatim}

\begin{enumerate}
\def\labelenumi{\arabic{enumi}.}
\setcounter{enumi}{16}
\tightlist
\item
\item
\item
\end{enumerate}

\begin{verbatim}
\begin{array}{ll} (\mathrm{Y} \mathrm{H} \Pi) & u_{tt} = \alpha^2 u_{xx} + Kx, \quad 0 < x < 1, \ 0 < t < \infty, \\ (\mathrm{\Gamma} \mathrm{Y}) & \begin{cases} u \ (0, \ t) = 0, \\ u \ (1, \ t) = 0, \end{cases} & 0 < t < \infty, \\ (\mathrm{H} \mathrm{Y}) & \begin{cases} u \ (x, \ 0) = f \ (x), \\ u_t \ (x, \ 0) = 0. \end{cases} & 0 \leqslant x \leqslant 1, \end{array}
\end{verbatim}

\begin{enumerate}
\def\labelenumi{\arabic{enumi}.}
\setcounter{enumi}{18}
\tightlist
\item
\item
\end{enumerate}

\section{ОСЦИЛЯЦІЇ ПУЧКА (РІВНЯННЯ ЧЕТВЕРТОГО ПОРЯДКУ В ЧАСТИННИХ
ПОХІДНИХ)}\label{ux43eux441ux446ux438ux43bux44fux446ux456ux457-ux43fux443ux447ux43aux430-ux440ux456ux432ux43dux44fux43dux43dux44f-ux447ux435ux442ux432ux435ux440ux442ux43eux433ux43e-ux43fux43eux440ux44fux434ux43aux443-ux432-ux447ux430ux441ux442ux438ux43dux43dux438ux445-ux43fux43eux445ux456ux434ux43dux438ux445}

МЕТА-ЛЕКЦІЯ: Показати, як проблема осциляції пучка веде до рівняння
четвертого порядку і як ця задача розв'язується у випадку вільно
спокійних кінців. Проводиться порівняння коливань променя з вібраціями
струни скрипки.

\pandocbounded{\includegraphics[keepaspectratio]{images/_page_154_Figure_5.jpeg}}

Головна різниця між поперечними коливаннями тонкої балки та поперечними
коливаннями струни полягає в тому, що промінь чинить опір згину. Не
заглиблюючись у механіку тонких балок, припустимо, що, враховуючи опір
згинання (замість рівняння хвилі) веде до рівняння четвертого порядку

(21.1)

\begin{verbatim}
# ТАБЛИЦЯ B. Перетворення Спнуса-Фур'є i \sqrt{\frac{\pi}{2}} [\delta(\omega+a) - \delta(\omega-a)]
\end{verbatim}

\texttt{rge\ \textquotesingle{}\textbackslash{}alpha\^{}2\ =\ K/\textbackslash{}rho\textquotesingle{}\ ,}

K---модуль сдвига (чем больше K, тем жестче балка, тем выше частота
колебаний),

р --- линейная плотность балки (масса/ед. длины).

Оскільки це перший випадок, коли читач стикається з рівнянням, вищим за
другий порядок у цій книзі, буде корисно розв'язати типову задачу
коливань пучка. Пізніше ми ознайомимося з іншими проблемами теорії
пучків.

\section{Свободно опирающаяся
балка}\label{ux441ux432ux43eux431ux43eux434ux43dux43e-ux43eux43fux438ux440ux430ux44eux449ux430ux44fux441ux44f-ux431ux430ux43bux43aux430}

Розглянемо невеликі коливання тонкої балки, кінці якої вільно лежать на
двох опорах. Коли ми кажемо «вільно відпочивати», ми маємо на увазі, що
кінці балки не рухаються, але

Клони пучка на кінцевих точках можна змінювати (кінці пучка фіксуються
за допомогою штифтових пристроїв, рис. 21.1).

Совершенно ясно, что на концах балки должны выполняться граничные
условия

\begin{verbatim}
F(\omega) = \frac{?}{\pi} \int_{0}^{\infty} f(x) \sin(\omega x) dx
\end{verbatim}

\texttt{u(1,\ t)\ =\ 0,} але зовсім не очевидно, що на кінцях балки
мають бути виконані ще дві граничні умови:

\begin{verbatim}
0 < \omega < x
\end{verbatim}

\texttt{u\_\{xx\}(1,\ t)\ =\ 0.} Використовуючи теорію тонких балок,
можна показати, що момент згинання в балці пропорційний значенню
\texttt{u\_\{xx\}} , а момент згинання на вільно лежачому кінці має
дорівнювати P I C

\pandocbounded{\includegraphics[keepaspectratio]{images/_page_155_Figure_8.jpeg}}

.21.1. Свободно оппрающаяся балка.

Нуль вен. Відповідно, зображення, зображене на рис. 21.1 Осцилюючий
пучок описується такою змішаною задачею (для простоти значення o
вважається рівним одиниці):

(21.2)

f'\,'(0)

\begin{enumerate}
\def\labelenumi{\arabic{enumi}.}
\setcounter{enumi}{1}
\tightlist
\item
  \texttt{f(\textbackslash{}alpha\ x)}
\end{enumerate}

\begin{verbatim}
\begin{aligned} & (\text{Y} \cdot \Pi) \quad u_{tt} = -u_{xxxx}, \quad 0 < x < 1, \quad 0 < t < \infty, \\ & \begin{cases} u \cdot (0, \ t) = 0, \\ u_{xx} \cdot (0, \ t) = 0, \\ u \cdot (1, \ t) = 0, \\ u_{xx} \cdot (1, \ t) = 0, \end{cases} \quad 0 < t < \infty, \\ & \begin{cases} u \cdot (x, \ 0) = f(x), \\ u_{1}(x, \ 0) = g(x). \end{cases} \quad 0 \leqslant x \leqslant 1, \end{aligned}
\end{verbatim}

\begin{enumerate}
\def\labelenumi{\arabic{enumi}.}
\setcounter{enumi}{2}
\item
  \texttt{e\^{}\{-\textbackslash{}alpha\ x\}}
\item
\end{enumerate}

Щоб розв'язати цю задачу, ми використаємо метод розділення змінних. Ми
будемо шукати лише періодичні розв'язки, тобто флуктуації

(21.3)

\begin{verbatim}
[2/\pi\omega]^{1/2}
\end{verbatim}

Зверніть увагу, що вибір розв'язку вигляду (21.3) фактично означає, що
ми взяли константу розділення у методі розділення змінних як від'ємну.

\texttt{Подставим\ теперь\ (21.3)\ в\ уравнение\ колебаний\ балки\ и\ получим\ уравнение\ для\ \textquotesingle{}X\textbackslash{}left(x\textbackslash{}right)\textquotesingle{}}

\begin{verbatim}
\frac{2}{\pi\omega} [1-\cos(\omega a)]
\end{verbatim}

общее решение которого записывается в виде

1

Щоб визначити коефіцієнти C, D, E та F, підставляємо загальний розв'язок
у граничні умови

\begin{verbatim}
e^{-1}a + \omega
\frac{1}{2}e^{-\alpha}\sin\omega
\begin{split} u\left(0,\ t\right) &= 0 \Rightarrow X\left(0\right)T\left(t\right) = 0 \Rightarrow X\left(0\right) = 0 \Rightarrow C + E = 0 \\ u_{xx}\left(0,\ t\right) &= 0 \Rightarrow X''\left(0\right)T\left(t\right) = 0 \Rightarrow X''\left(0\right) = 0 \Rightarrow -C + E = 0 \\ u\left(1,\ t\right) &= 0 \Rightarrow D\sin\sqrt{\omega} + F\sin\sqrt{\omega} = 0, \\ u_{xx}\left(1,\ t\right) &= 0 \Rightarrow -D\sin\sqrt{\omega} + F\sin\sqrt{\omega} = 0. \end{split}
\frac{1-e^{-a\omega}}{\omega}
\frac{2}{\pi}F''(\omega)
F \sinh \sqrt{\omega} = 0,
f(x) = \int_{0}^{\infty} F(\omega) \cos(\omega x) d\omega \qquad \qquad F(x) = \frac{2}{\pi} \int_{0}^{\infty} f(x) \cos(\omega x) dx
D \sin \sqrt{\omega} = 0.
\end{verbatim}

\begin{enumerate}
\def\labelenumi{\arabic{enumi}.}
\tightlist
\item
\end{enumerate}

\begin{verbatim}
F = 0
\end{verbatim}

\section{ТАБЛИЦЯ D. Скінченне перетворення синуса
Фур'є}\label{ux442ux430ux431ux43bux438ux446ux44f-d.-ux441ux43aux456ux43dux447ux435ux43dux43dux435-ux43fux435ux440ux435ux442ux432ux43eux440ux435ux43dux43dux44f-ux441ux438ux43dux443ux441ux430-ux444ux443ux440ux454}

\begin{verbatim}
У скінченному перетворенні синуса Фур'є функція f(x), визначена в інтервалі 0 \le x \le \pi , перетворюється у послідовність S_n , n = 1, 2 ... За формулою
\end{verbatim}

\begin{verbatim}
\omega_n = (n\pi)^2,
\end{verbatim}

\begin{verbatim}
Коли йдеться про таблиці, зручно розглядати змінну x як належну до інтервалу [0, l]. Будь-який інший чисельний інтервал [a, b] можна перетворити на [0, \pi] за допомогою формули
\end{verbatim}

\begin{verbatim}
u_n(x, t) = X_n(x) T_n(t) = [a_n \sin(n\pi)^2 t + b_n \cos(n\pi)^2 t] \sin(n\pi x).
\end{verbatim}

\begin{verbatim}
У скінченному синусоїдальному перетворенні немає нічого загадкового: елементи послідовності \mathcal{S}_n збігаються з коефіцієнтами в \sin(nx) у розкладі функції f(x) у ряд Фур'є, тобто
\end{verbatim}

\begin{verbatim}
u(x, t) = \sum_{n=1}^{\infty} \left[ a_n \sin(n\pi)^2 t + b_n \cos(n\pi)^2 t \right] \sin(n\pi x)
f(x) = \sum_{n=1}^{\infty} S_n \sin(nx)
u(x, 0) = f(x) = \sum_{n=1}^{\infty} b_n \sin(n\pi x),
0 \le x \le \pi
u_t(x, 0) = g(x) = \sum_{n=1}^{\infty} (n\pi)^2 a_n \sin(n\pi x).
\end{verbatim}

\begin{enumerate}
\def\labelenumi{\arabic{enumi}.}
\tightlist
\item
\end{enumerate}

\begin{verbatim}
a_n = \frac{2}{(n\pi)^2} \int_0^1 g(x) \sin(n\pi x) dx,
\end{verbatim}

\begin{enumerate}
\def\labelenumi{\arabic{enumi}.}
\setcounter{enumi}{8}
\tightlist
\item
\end{enumerate}

\begin{verbatim}
b_n = 2 \int_0^1 f(x) \sin(n\pi x) dx.
\end{verbatim}

\begin{verbatim}
\frac{n}{n^2 + a^2} [1 - (-1)^n e^{ax}] 10. \frac{\sin a (\pi - x)}{\sin a\pi} \frac{2n}{\pi (n^2 + a^2)} ТАБЛИЦЯ E. Скінченне перетворення Косніуса Фур'є
\end{verbatim}

\begin{verbatim}
У скінченному косинусному перетворенні Фур'є функція f(x), визначена на інтервалі 0 \le x \le \pi , перетворюється у послідовність C_n , n=0,\ 1,\ 2,\ \ldots , за формулою
\end{verbatim}

\begin{verbatim}
u(x, 0) = \sin(\pi x) + 0.5 \sin(3\pi x),
\end{verbatim}

\begin{verbatim}
Коли йдеться про таблиці, зручно розглядати змінну x як таку, що належить до інтервалу [0,\pi] . Будь-який інший числовий інтервал [a,b] можна перетворити на [0,\pi] за допомогою формули
\end{verbatim}

\begin{verbatim}
a_n=0
\end{verbatim}

\begin{verbatim}
У скінченному косинусному перетворенні немає нічого загадкового: елементи послідовності C_n збігаються з коефіцієнтами, коли \cos{(nx)} у косинусному розкладі функції f(x) у ряд Фур'є, тобто
\end{verbatim}

\begin{verbatim}
u(x, t) = \cos(\pi^2 t) \sin(\pi x) + 0.5 \cos(9\pi^2 t) \sin(3\pi x)
f(x) = \frac{C_0}{2} + \sum_{n=1}^{\infty} C_n \cos(nx) \qquad C_n = \frac{2}{\pi} \int_{0}^{\pi} f(x) \cos nx \, dx
u(x, t) = \cos(\pi t) \sin(\pi x) + 0.5 \cos(3\pi t) \sin(3\pi x)
\end{verbatim}

\begin{enumerate}
\def\labelenumi{\arabic{enumi}.}
\tightlist
\item
\end{enumerate}

\begin{verbatim}
u(0, t) = 0,
-n^{2}C_{n} - \frac{2}{\pi} \{f'(0) - (-1)^{n} f'(n)\}
u(x, t) = \sum_{n=1}^{\infty} X_n(x) \left[ a_n \sin(\omega_n t) + b_n \cos(\omega_n t) \right],
\begin{cases} 2, &n=0 \\ 0, &n=1, 2, \dots \end{cases}
(\text{УЧП}) \quad u_{tt} + u_{xxxx} = 0, \quad 0 < x < 1, \quad 0 < t < \infty,
\begin{cases} 1, & n=m \\ 0, & n\neq m \end{cases}
(\text{ГУ}) \quad
\end{verbatim}

x

\begin{verbatim}
\begin{cases} u(0, t) &= 0, \\ u_x(0, t) &= 0, \\ u_{xx}(1, t) &= 0, \\ u_{xxx}(1, t) &= 0, \end{cases}
(\text{HV}) \quad
\begin{cases} u(x, 0) &= f(x), \\ u_t(x, 0) &= g(x). \end{cases}
\quad 0 \leqslant x \leqslant 1,
\end{verbatim}

6 M 601

ПРИМІТКА. Хоча власні функції, \texttt{X\_n(x)} в цій задачі, вже не є
звичайними синусами, згідно з загальною теорією Штурма-Лівілля, вони
будуть ортогональними на сегменті {[}0, 1{]}.

\begin{enumerate}
\def\labelenumi{\arabic{enumi}.}
\setcounter{enumi}{1}
\tightlist
\item
  Знайти розв'язок для вільно лежачого балки на обох кінцях за
  початкових умов
\end{enumerate}

\begin{verbatim}
u(x, 0) = \sin(\pi x),
\end{verbatim}

\texttt{u\_t(x,\ 0)\ =\ \textbackslash{}sin(\textbackslash{}pi\ x).}
\texttt{0\ \textbackslash{}le\ x\ \textbackslash{}le\ 1,} 3. Решите
задачу 2 с начальными условиями

\begin{verbatim}
\begin{array}{l} u(x, 0) = 1 - x^2, \\ u_t(x, 0) = 0. \end{array}
\quad 0 \le x \le 1,
\end{verbatim}

\begin{enumerate}
\def\labelenumi{\arabic{enumi}.}
\setcounter{enumi}{3}
\tightlist
\item
  Пусть левый конец (x=0) балки жестко закреплен в стену, а правый конец
  (x=1) свободно опирается, как показано на рис. 21.3. Решите задачу для
  балки с такими граничными условиями и объясните, как найти собственные
  частоты колебаний в этом случае. Знать собственные частоты балок очень
  важно, чтобы избежать резонанса с внешними воздействиями.
\end{enumerate}

Лекция 22

\section{ПЕРЕХОД К БЕЗРАЗМЕРНЫМ ПЕРЕМЕННЫМ МЕТА ЛЕКЦІЇ: Показать, каким
образом краевые, начальные и другие виды физических задач можно записать
в безразмерном виде. В конкретных областях знания (физике, химии,
биологии или экономике) уравнения записываются
по-разному.}\label{ux43fux435ux440ux435ux445ux43eux434-ux43a-ux431ux435ux437ux440ux430ux437ux43cux435ux440ux43dux44bux43c-ux43fux435ux440ux435ux43cux435ux43dux43dux44bux43c-ux43cux435ux442ux430-ux43bux435ux43aux446ux456ux457-ux43fux43eux43aux430ux437ux430ux442ux44c-ux43aux430ux43aux438ux43c-ux43eux431ux440ux430ux437ux43eux43c-ux43aux440ux430ux435ux432ux44bux435-ux43dux430ux447ux430ux43bux44cux43dux44bux435-ux438-ux434ux440ux443ux433ux438ux435-ux432ux438ux434ux44b-ux444ux438ux437ux438ux447ux435ux441ux43aux438ux445-ux437ux430ux434ux430ux447-ux43cux43eux436ux43dux43e-ux437ux430ux43fux438ux441ux430ux442ux44c-ux432-ux431ux435ux437ux440ux430ux437ux43cux435ux440ux43dux43eux43c-ux432ux438ux434ux435.-ux432-ux43aux43eux43dux43aux440ux435ux442ux43dux44bux445-ux43eux431ux43bux430ux441ux442ux44fux445-ux437ux43dux430ux43dux438ux44f-ux444ux438ux437ux438ux43aux435-ux445ux438ux43cux438ux438-ux431ux438ux43eux43bux43eux433ux438ux438-ux438ux43bux438-ux44dux43aux43eux43dux43eux43cux438ux43aux435-ux443ux440ux430ux432ux43dux435ux43dux438ux44f-ux437ux430ux43fux438ux441ux44bux432ux430ux44eux442ux441ux44f-ux43fux43e-ux440ux430ux437ux43dux43eux43cux443.}

Переход к безразмерным переменным позволяет привести уравнения к одному
и тому же виду. Именно по этой причине при изучении уравнений с частными
производными в математике стараются отвлекаться от физического смысла
входящих в них параметров. Вот почему химик, физик или биолог должен
преобразовать свое уравнение к той форме, которая принята в нашей книге.

Основная идея введения новых (безразмерных) переменных состоит в том,
что после перехода к безразмерным переменным вадача становится чисто
математической и уже не содержит

карактерных физических констант. Именно таким способом многие различные
уравнения физики, биологии и химии, содержащие всякие нюансы, связанные
с физическими параметрами, приводятся к одной и той же простой форме
(см. рис. 22.1).

Познакомимся с этим преобразованием на простом примере.

\pandocbounded{\includegraphics[keepaspectratio]{images/_page_161_Figure_4.jpeg}}

Рис. 22.1. Некоторые задачи, приводящиеся к одной и той же безразмерной
форме.

\section{Приведение диффузионной задачи к безразмерному
виду}\label{ux43fux440ux438ux432ux435ux434ux435ux43dux438ux435-ux434ux438ux444ux444ux443ux437ux438ux43eux43dux43dux43eux439-ux437ux430ux434ux430ux447ux438-ux43a-ux431ux435ux437ux440ux430ux437ux43cux435ux440ux43dux43eux43cux443-ux432ux438ux434ux443}

\begin{verbatim}
Начнем со смешанной задачи для отрезка с начальным условием на температуру  u\left(x,\,0\right)=\sin\left(\pi x/L\right)  и граничными температурами на концах  T_1  и  T_2 . Другими словами, мы решаем задачу
\end{verbatim}

(22.1)

\begin{verbatim}
\begin{array}{ccccccccccccccccccccccccccccccccccc
\end{verbatim}

(см. рис. 22.2).

\pandocbounded{\includegraphics[keepaspectratio]{images/_page_161_Figure_10.jpeg}}

Рис. 22.2. Задача теплопроводности в плоскости переменных (х, т).

Наша цель --- дать новую эквивалентную (22.1) постановку задачи таким
образом, чтобы

в новой задаче не было никаких физических параметров (вроде α);

\begin{enumerate}
\def\labelenumi{\arabic{enumi})}
\setcounter{enumi}{1}
\tightlist
\item
  начальные и граничные условия стали проще.
\end{enumerate}

\begin{verbatim}
Для осуществления этой цели, мы введем три новых безразмерных переменных U,   \xi   и   \tau  , чтобы заменить прежние u, x и t по схеме
\end{verbatim}

\begin{verbatim}
u \to U
\end{verbatim}

\begin{verbatim}
(безразмерная температура),   x \to \xi   (безразмерная длина),   t \to \tau   (безразмерное время).
\end{verbatim}

Мы проведем все три преобразования одновременно.

\begin{verbatim}
# Преобразование зависимой переменной  u \rightarrow U
\end{verbatim}

Определим функцию \texttt{U(x,\ t)} по формуле

\begin{verbatim}
U(x, t) = \frac{u(x, t) - T_1}{T_2 - T_1}
\end{verbatim}

.

\begin{verbatim}
Ясно, что новая температура   U\left(x,\,t\right)   безразмерна, поскольку мы делим градусы Цельсия на градусы Цельсия. Новые граничные условия для   U\left(x,\,t\right)   при x=0 и x=L примут вид   U\left(0,\,t\right)=0   и   U\left(L,\,t\right)=1  . Давайте поставим задачу для новой функции   U\left(x,\,t\right)  . После несложных преобразований задача (22.1) переходит в новую задачу
\end{verbatim}

задачу

\begin{verbatim}
\begin{array}{ll} \text{(УЧП)} & U_t = \alpha^2 U_{xx}, \quad 0 < x < L, \quad 0 < t < \infty, \\ U_t = 0, \quad 0 < t < \infty, \\ U_t = 0, \quad 0 < t < \infty, \\ U_t = 0, \quad 0 < t < \infty, \\ U_t = 0, \quad 0 < t < \infty, \\ U_t = 0, \quad 0 < t < \infty, \\ U_t = 0, \quad 0 < t < \infty, \\ U_t = 0, \quad 0 < t < \infty, \\ U_t = 0, \quad 0 < t < \infty, \\ U_t = 0, \quad 0 < t < \infty, \\ U_t = 0, \quad 0 < t < \infty, \\ U_t = 0, \quad 0 < t < \infty, \\ U_t = 0, \quad 0 < t < \infty, \\ U_t = 0, \quad 0 < t < \infty, \\ U_t = 0, \quad 0 < t < \infty, \\ U_t = 0, \quad 0 < t < \infty, \\ U_t = 0, \quad 0 < t < \infty, \\ U_t = 0, \quad 0 < t < \infty, \\ U_t = 0, \quad 0 < t < \infty, \\ U_t = 0, \quad 0 < t < \infty, \\ U_t = 0, \quad 0 < t < \infty, \\ U_t = 0, \quad 0 < t < \infty, \\ U_t = 0, \quad 0 < t < \infty, \\ U_t = 0, \quad 0 < t < \infty, \\ U_t = 0, \quad 0 < t < \infty, \\ U_t = 0, \quad 0 < t < \infty, \\ U_t = 0, \quad 0 < t < \infty, \\ U_t = 0, \quad 0 < t < \infty, \\ U_t = 0, \quad 0 < t < \infty, \\ U_t = 0, \quad 0 < t < \infty, \\ U_t = 0, \quad 0 < t < \infty, \\ U_t = 0, \quad 0 < t < \infty, \\ U_t = 0, \quad 0 < t < \infty, \\ U_t = 0, \quad 0 < t < \infty, \\ U_t = 0, \quad 0 < t < \infty, \\ U_t = 0, \quad 0 < t < \infty, \\ U_t = 0, \quad 0 < t < \infty, \\ U_t = 0, \quad 0 < t < \infty, \\ U_t = 0, \quad 0 < t < \infty, \\ U_t = 0, \quad 0 < t < \infty, \\ U_t = 0, \quad 0 < t < \infty, \\ U_t = 0, \quad 0 < t < \infty, \\ U_t = 0, \quad 0 < t < \infty, \\ U_t = 0, \quad 0 < t < \infty, \\ U_t = 0, \quad 0 < t < \infty, \\ U_t = 0, \quad 0 < t < \infty, \\ U_t = 0, \quad 0 < t < \infty, \\ U_t = 0, \quad 0 < t < \infty, \\ U_t = 0, \quad 0 < t < \infty, \\ U_t = 0, \quad 0 < t < \infty, \\ U_t = 0, \quad 0 < t < \infty, \\ U_t = 0, \quad 0 < t < \infty, \\ U_t = 0, \quad 0 < t < \infty, \\ U_t = 0, \quad 0 < t < \infty, \\ U_t = 0, \quad 0 < t < \infty, \\ U_t = 0, \quad 0 < t < \infty, \\ U_t = 0, \quad 0 < t < \infty, \\ U_t = 0, \quad 0 < t < \infty, \\ U_t = 0, \quad 0 < t < \infty, \\ U_t = 0, \quad 0 < t < \infty, \\ U_t = 0, \quad 0 < t < \infty, \\ U_t = 0, \quad 0 < t < \infty, \\ U_t = 0, \quad 0 < t < \infty, \\ U_t = 0, \quad 0 < t < \infty, \\ U_t = 0, \quad 0 < t < \infty, \\ U_t = 0, \quad 0 < t < \infty, \\ U_t = 0, \quad 0 < t < \infty, \\ U_t = 0, \quad 0 < t < \infty, \\ U_t = 0, \quad 0 < t < \infty, \\ U_t = 0, \quad 0 < t < \infty, \\ U_t = 0, \quad 0 < t < \infty, \\ U_t = 0, \quad 0 < t < \infty, \\ U_t = 0, \quad 0 < t < \infty, \\ U_t = 0, \quad 0 < t < \infty, \\ U_t = 0, \quad 0 < t < \infty, \\ U_t = 0, \quad 0 < t < \infty, \\ U_t = 0, \quad 0 < t < \infty, \\ U_t = 0, \quad 0 < t < \infty, \\ U_t = 0, \quad 0 < t < \infty, \\ U_t = 0, \quad 0 < t < \infty, \\ U_t = 0, \quad 0 < t < \infty, \\ U_t = 0, \quad 0 < t < \infty, \\ U_t = 0, \quad 0 < t < \infty, \\ U_t
\end{verbatim}

\begin{verbatim}
Если угодно, то здесь можно остановиться, решить задачу для   U\left(x,\, t\right)   и затем найти   u\left(x,\,t\right)   по формуле
\end{verbatim}

\begin{verbatim}
u(x, t) = T_1 + (T_2 - T_1) U(x, t).
\end{verbatim}

Однако мы пойдем дальше и займемся преобразованием независимых
переменных x и t.

\begin{verbatim}
# Преобразование пространственной координаты  x \rightarrow \xi
\end{verbatim}

\begin{verbatim}
Представляется совершенно очевидным, как следует выбрать безразмерную переменную   \xi  . Поскольку   0 \leqslant x \leqslant L  , полагаем
\end{verbatim}

\begin{verbatim}
\xi = x/L
\end{verbatim}

.

После вычисления производных

\begin{verbatim}
U_x = U_{\xi} \xi_x = \frac{1}{L} U_{\xi},
U_{xx} = \frac{1}{L^2} U_{\xi\xi}
\end{verbatim}

наступна задача стає очевидною (у змінних U, \texttt{\textbackslash{}xi}
, t):

\begin{verbatim}
\begin{array}{lll} (\text{YM}\Pi) & U_t = (\alpha/L)^2 \, U_{\xi\xi}, \quad 0 < \xi < 1, \quad 0 < t < \infty, \\ (22.3) & (\Gamma \text{Y}) & \begin{cases} U(0,\ t) = 0, \\ U(1,\ t) = 1, \\ U(\xi,\ 0) = \frac{\sin\left(\pi\xi\right) - T_1}{T_2 - T_1}, \quad 0 \leqslant \xi \leqslant 1. \end{cases} \end{array}
\end{verbatim}

Ми вже пройшли дві третини шляху. Останній крок пов'язаний із введенням
безрозмірного часу \texttt{\textbackslash{}tau} для виключення
коефіцієнта \texttt{{[}\textbackslash{}alpha/L{]}\^{}2} з
диференціального рівняння.

\section{\texorpdfstring{\texttt{t\ \textbackslash{}rightarrow\ \textbackslash{}tau}
Перетворення
часу}{t \textbackslash rightarrow \textbackslash tau Перетворення часу}}\label{t-rightarrow-tau-ux43fux435ux440ux435ux442ux432ux43eux440ux435ux43dux43dux44f-ux447ux430ux441ux443}

Як обрати час без розмірів? Це не так очевидно, як у випадку перших двох
змінних. Однак, оскільки наша мета --- виключити сталу
\texttt{{[}\textbackslash{}alpha/L{]}\^{}2} з рівняння (22.3), ми
діятимемо наступним чином.

\begin{itemize}
\item
  \begin{enumerate}
  \def\labelenumi{\arabic{enumi}.}
  \tightlist
  \item
    Спробуємо замінити вигляд \texttt{\textbackslash{}tau\ =\ ct} , де c
    --- невідома константа.
  \end{enumerate}

  \begin{itemize}
  \tightlist
  \item
    \begin{enumerate}
    \def\labelenumi{\arabic{enumi}.}
    \setcounter{enumi}{1}
    \tightlist
    \item
      Давайте розрахуємо
      \texttt{u\_t\ =\ u\_\{\textbackslash{}tau\}\ \textbackslash{}tau\_t\ =\ c\ u\_\{\textbackslash{}tau\}}
      .
    \end{enumerate}
  \item
    \begin{enumerate}
    \def\labelenumi{\arabic{enumi}.}
    \setcounter{enumi}{2}
    \tightlist
    \item
      Підставимо цю похідну у рівняння (22.3) і отримаємо
    \end{enumerate}
  \end{itemize}
\end{itemize}

\begin{verbatim}
cu_{\mathbf{t}} = [\alpha/L]^2 u_{\xi\xi}.
\end{verbatim}

Тому потрібно обрати \texttt{c\ =\ {[}\textbackslash{}alpha/L{]}\^{}2} ,
тобто новий час слід вводити за формулою

\begin{verbatim}
\tau = [\alpha/L]^2 t.
\end{verbatim}

Застосовуючи це перетворення до попередньої задачі (22.3), ми можемо
повністю перейти до безрозмірних змінних U, \texttt{\textbackslash{}xi}
та \texttt{\textbackslash{}tau\_i}

\begin{verbatim}
\begin{array}{cccccccccccccccccccccccccccccccccccc
\end{verbatim}

Новая безразмерная система обладает следующими свойствами:

\begin{enumerate}
\def\labelenumi{\arabic{enumi}.}
\item
  Уравнение не содержит параметров.
\item
  Граничные условия имеют простой вид.
\item
  Початкова умова суттєво не змінилася (тобто містить лише відомі
  функції).
\item
  Завдання стало простішим і компактнішим порівняно з оригінальним.
  Розв'язок цієї проблеми можна побудувати раз і назавжди, тож якщо
  читач перетворює свою початкову задачу (22.1) у безрозмірну форму
  (22.4) і знаходить її розв'язок
  \texttt{U(\textbackslash{}xi,\ \textbackslash{}tau)} у підручнику чи
  науковому журналі, то він може знайти розв'язок початкової задачі
  (22.1) простим перерахуванням
\end{enumerate}

\begin{verbatim}
u(x, t) = T_1 + (T_2 - T_1) U(x/L, \alpha^2 t/L^2).
\end{verbatim}

На цьому завершується наша дискусія проблеми перетворення задач у
безрозмірний вигляд. Не існує універсальних правил, за якими слід
вводити нові змінні, тут слід покладатися на фізичну інтуїцію і
пробувати різні варіанти.

Завершимо цю лекцію простим прикладом перетворення задачі на
безрозмірний вигляд, розв'язання нової задачі і повернення до початкових
змінних.

\section{Приклад перетворення гіперболічної задачі на безрозмірний
вигляд}\label{ux43fux440ux438ux43aux43bux430ux434-ux43fux435ux440ux435ux442ux432ux43eux440ux435ux43dux43dux44f-ux433ux456ux43fux435ux440ux431ux43eux43bux456ux447ux43dux43eux457-ux437ux430ux434ux430ux447ux456-ux43dux430-ux431ux435ux437ux440ux43eux437ux43cux456ux440ux43dux438ux439-ux432ux438ux433ux43bux44fux434}

Рассмотрим колеблющуюся струну

(22.5)

\begin{verbatim}
\begin{array}{llll} (\text{УЧП}) & u_{tt} = \alpha^2 u_{xx}, & 0 < x < L, & 0 < t < \infty, \\ u(0, t) = 0, & 0 < t < \infty, \\ u(L, t) = 0, & 0 < t < \infty, \\ u(X, 0) = \sin(\pi x/L) + (0.5)\sin(3\pi x/L), \\ u_t(x, 0) = 0. & \end{array}
\end{verbatim}

Перетворюючи пояснювальні змінні у безрозмірний вигляд (немає потреби
трансформувати вас)

\begin{verbatim}
\xi = x/L
\end{verbatim}

, \texttt{\textbackslash{}tau\ =\ {[}\textbackslash{}alpha/L{]}\ t}
получаем новую задачу

\begin{verbatim}
(\text{YYII}) \qquad u_{\tau\tau} = u_{\xi\xi}, \quad 0 < \xi < 1, \quad 0 < \tau < \infty, 1
(\text{\GammaY}) \quad
\begin{cases} u(0, \tau) = 0, \\ u(1, \tau) = 0, \\ u(\xi, 0) = \sin(\pi\xi) + 0.5\sin(3\pi\xi), \\ u_{\xi}(\xi, 0) = 0, \end{cases}
\quad 0 < \xi \leqslant 1,
(\text{HY}) \quad
\begin{cases} u(\xi, 0) = \sin(\pi\xi) + 0.5\sin(3\pi\xi), \\ u_{\xi}(\xi, 0) = 0, \end{cases}
\quad 0 \leqslant \xi \leqslant 1,
\end{verbatim}

решением которой является функция

\begin{verbatim}
u(\xi, \tau) = \cos(\pi \tau) \sin(\pi \xi) + 0.5 \cos(3\pi \tau) \sin(3\pi \xi).
\end{verbatim}

Якщо повернутися до координат x і t, розв'язок задачі (22.5) набуває
вигляду

\begin{verbatim}
u(x, t) = \cos(\pi \alpha t/L)\sin(\pi x/L) + 0.5\cos(3\pi \alpha t/L)\sin(3\pi x/L).
\end{verbatim}

\section{ПОВАГА - 1. Аналіз розмірності особливо важливий у чисельних
методах, оскільки більшість програм написані для розв'язання загальних
математичних задач і не передбачають використання великої кількості
фізичних параметрів. Ті, хто хоче користуватися цими програмами, повинні
перетворити свою задачу у відповідну форму, розв'язати її на комп'ютері,
а потім довести числові результати до потрібного
розміру.}\label{ux43fux43eux432ux430ux433ux430---1.-ux430ux43dux430ux43bux456ux437-ux440ux43eux437ux43cux456ux440ux43dux43eux441ux442ux456-ux43eux441ux43eux431ux43bux438ux432ux43e-ux432ux430ux436ux43bux438ux432ux438ux439-ux443-ux447ux438ux441ux435ux43bux44cux43dux438ux445-ux43cux435ux442ux43eux434ux430ux445-ux43eux441ux43aux456ux43bux44cux43aux438-ux431ux456ux43bux44cux448ux456ux441ux442ux44c-ux43fux440ux43eux433ux440ux430ux43c-ux43dux430ux43fux438ux441ux430ux43dux456-ux434ux43bux44f-ux440ux43eux437ux432ux44fux437ux430ux43dux43dux44f-ux437ux430ux433ux430ux43bux44cux43dux438ux445-ux43cux430ux442ux435ux43cux430ux442ux438ux447ux43dux438ux445-ux437ux430ux434ux430ux447-ux456-ux43dux435-ux43fux435ux440ux435ux434ux431ux430ux447ux430ux44eux442ux44c-ux432ux438ux43aux43eux440ux438ux441ux442ux430ux43dux43dux44f-ux432ux435ux43bux438ux43aux43eux457-ux43aux456ux43bux44cux43aux43eux441ux442ux456-ux444ux456ux437ux438ux447ux43dux438ux445-ux43fux430ux440ux430ux43cux435ux442ux440ux456ux432.-ux442ux456-ux445ux442ux43e-ux445ux43eux447ux435-ux43aux43eux440ux438ux441ux442ux443ux432ux430ux442ux438ux441ux44f-ux446ux438ux43cux438-ux43fux440ux43eux433ux440ux430ux43cux430ux43cux438-ux43fux43eux432ux438ux43dux43dux456-ux43fux435ux440ux435ux442ux432ux43eux440ux438ux442ux438-ux441ux432ux43eux44e-ux437ux430ux434ux430ux447ux443-ux443-ux432ux456ux434ux43fux43eux432ux456ux434ux43dux443-ux444ux43eux440ux43cux443-ux440ux43eux437ux432ux44fux437ux430ux442ux438-ux457ux457-ux43dux430-ux43aux43eux43cux43fux44eux442ux435ux440ux456-ux430-ux43fux43eux442ux456ux43c-ux434ux43eux432ux435ux441ux442ux438-ux447ux438ux441ux43bux43eux432ux456-ux440ux435ux437ux443ux43bux44cux442ux430ux442ux438-ux434ux43e-ux43fux43eux442ux440ux456ux431ux43dux43eux433ux43e-ux440ux43eux437ux43cux456ux440ux443.}

\begin{itemize}
\tightlist
\item
  \begin{enumerate}
  \def\labelenumi{\arabic{enumi}.}
  \setcounter{enumi}{1}
  \tightlist
  \item
    Розмірний аналіз дозволяє математикам працювати з диференціальним
    рівнянням у похідних похідних, не виключаючи багатьох параметрів і
    констант, які не впливають на математичну суть задачі.
  \end{enumerate}
\item
  \begin{enumerate}
  \def\labelenumi{\arabic{enumi}.}
  \setcounter{enumi}{2}
  \tightlist
  \item
    Не завжди потрібно конвертувати всі змінні у безрозмірний вигляд,
    іноді достатньо перетворити одну або дві змінні. \# ЗАВДАННЯ
  \end{enumerate}
\end{itemize}

\begin{enumerate}
\def\labelenumi{\arabic{enumi}.}
\tightlist
\item
  Шляхом заміни змінних
\end{enumerate}

\begin{verbatim}
\xi = x/L
\end{verbatim}

, \texttt{\textbackslash{}tau\ =\ {[}\textbackslash{}alpha/L{]}\ t}
преобразуйте задачу (22.5) к безразмерному виду (22.6).

\begin{enumerate}
\def\labelenumi{\arabic{enumi}.}
\setcounter{enumi}{1}
\tightlist
\item
  Найдите безразмерную форму для задачи
\end{enumerate}

\begin{verbatim}
\begin{array}{lll} (\text{Y} \, \Pi) & u_t = \alpha^2 u_{xx}, & 0 < x < L, \\ (\Gamma \, \text{Y}) & \left\{ \begin{array}{ll} u \, (0, \, t) = T_1, \\ u \, (L, \, t) = 0, \\ u \, (x, \, 0) = T_2, & 0 \leqslant x \leqslant L. \end{array} \end{array}
\end{verbatim}

\begin{verbatim}
\right. \\ \end{array}
\end{verbatim}

\begin{enumerate}
\def\labelenumi{\arabic{enumi}.}
\setcounter{enumi}{2}
\tightlist
\item
  С помощью замены
\end{enumerate}

\begin{verbatim}
U(x, t) = \frac{u(x, t) - T_1}{T_2 - T_1}
\end{verbatim}

преобразуйте задачу (22.1) в задачу (22.2).

\begin{verbatim}
4. Почему замена переменной   \tau = \alpha t   приводит к исключению параметра   \alpha   из волнового уравнения   u_{tt} = \alpha^2 u_{xx}  ? В чем физический смысл этой замены? Напоминаем вам, что   \alpha  —скорость
\end{verbatim}

распространения волн. При поиске новых координат в большинстве случаев
главную роль играет интуиция.

Как выбрать новую пространственную координату ξ, чтобы исключить и из
уравнения

\begin{verbatim}
u_t + vu_x = 0?
\end{verbatim}

Лекция 23

\section{КЛАССИФИКАЦИЯ УРАВНЕНИЙ С ЧАСТНЫМИ ПРОИЗВОДНЫМИ (КАНОНИЧЕСКАЯ
ФОРМА ГИПЕРБОЛИЧЕСКОГО
УРАВНЕНИЯ)}\label{ux43aux43bux430ux441ux441ux438ux444ux438ux43aux430ux446ux438ux44f-ux443ux440ux430ux432ux43dux435ux43dux438ux439-ux441-ux447ux430ux441ux442ux43dux44bux43cux438-ux43fux440ux43eux438ux437ux432ux43eux434ux43dux44bux43cux438-ux43aux430ux43dux43eux43dux438ux447ux435ux441ux43aux430ux44f-ux444ux43eux440ux43cux430-ux433ux438ux43fux435ux440ux431ux43eux43bux438ux447ux435ux441ux43aux43eux433ux43e-ux443ux440ux430ux432ux43dux435ux43dux438ux44f}

МЕТА ЛЕКЦІЇ: Показать, что линейное уравнение второго порядка с двумя
независимыми переменными

\begin{verbatim}
Au_{xx} + Bu_{xy} + Cu_{yy} + Du_x + Eu_y + Fu = G
\end{verbatim}

(A, B, C, D, E, F и G являются функциями x и y или могут быть
константами) относится к одному из следующих типов:

\begin{itemize}
\tightlist
\item
  \begin{enumerate}
  \def\labelenumi{\arabic{enumi})}
  \tightlist
  \item
    гиперболическому (если B\^{}2 4AC \textgreater{} 0 ),
  \end{enumerate}
\item
  \begin{enumerate}
  \def\labelenumi{\arabic{enumi})}
  \setcounter{enumi}{1}
  \tightlist
  \item
    параболическому (если B\^{}2 4AC = ),
  \end{enumerate}
\item
  \begin{enumerate}
  \def\labelenumi{\arabic{enumi})}
  \setcounter{enumi}{2}
  \tightlist
  \item
    эллиптическому (если 0 B\^{}2 4AC \textless{} 0 ).
  \end{enumerate}
\end{itemize}

\begin{verbatim}
Показать, какие новые переменные   \xi = \xi(x, y)   и   \eta = \eta(x, y)   можно ввести вместо x и y, чтобы упростить исходное уравнение. В новых переменных   \xi   и   \eta   уравнение с частными производными приводится к одному из следующих видов (в вависимости от того, будет ли величина   B^2 - 4AC   положительна, равна нулю или отрицательна):
\end{verbatim}

\begin{itemize}
\tightlist
\item
  \begin{enumerate}
  \def\labelenumi{\arabic{enumi}.}
  \tightlist
  \item
  \end{enumerate}
\end{itemize}

\begin{verbatim}
\begin{cases} u_{\xi\xi} u_{\eta\eta} = \Psi(\xi,  \eta, u, u_{\xi}, u_{\eta}) і \text{(дві канонічні форми гіперболічного рівняння),} \end{cases}
\end{verbatim}

\begin{verbatim}
- 2.   u_{\eta\eta} = \Phi(\xi, \eta, u, u_{\xi}, u_{\eta})   (каноническая форма параболического уравнения),
\end{verbatim}

\begin{verbatim}
- 3.   u_{\xi\xi} + u_{\eta\eta} = \Phi(\xi, \eta, u, u_{\xi}, u_{\eta})   (каноническая форма эллиптического уравнения),
\end{verbatim}

\begin{verbatim}
где   \Phi   и   \Psi  —функции от частных производных первого порядка, зависимой переменной u и новых независимых переменных   \xi   и   \eta  . Конкретный вид функций   \Phi   и   \Psi   зависит от исходного уравнения.
\end{verbatim}

У читателя может возникнуть мысль, что глава, посвященная классификации
уравнений, должна размещаться в самом начале книги. Вероятно, это
правильно, и многие книги действительно начинаются с этого материала.
Однако правда и то, что многие студенты не могут с энтузиазмом изучать
нечто такое, о чем у них нет никакого представления. Именно по этой
причине мы отложили изучение темы, посвященной классификации уравнений с
частными производными, до настоящего момента.

Наша задача - классифицировать все уравнения вида

(23 1)

\begin{verbatim}
Au_{xx} + Bu_{xy} + Cu_{yy} + Du_x + Eu_y + Fu = G,
\end{verbatim}

где A, B, C, D, E, F и G---вообще говоря, функции {[}x и y. Согласно
предлагаемой схеме классификации уравнение называется

\begin{enumerate}
\def\labelenumi{\arabic{enumi})}
\item
  гиперболическим в точке (\texttt{x\_0}, \texttt{y\_0}) , если
  B\^{}2(\texttt{x\_0}, \texttt{y\_0}) - 4A(\texttt{x\_0},
  \texttt{y\_0}) \помножити на \помножити на C(\texttt{x\_0},
  \texttt{y\_0}) \textgreater{} 0 2) параболнческим в точке
  (\texttt{x\_0}, \texttt{y\_0}) , если B\^{}2(\texttt{x\_0},
  \texttt{y\_0}) - 4A(\texttt{x\_0}, \texttt{y\_0}) \помножити на
  C(\texttt{x\_0}, \texttt{y\_0}) = 0 ,
\item
  эллиптическим в точке (\texttt{x\_0}, \texttt{y\_0}) , если
  B\^{}2(\texttt{x\_0}, \texttt{y\_0}) - 4A(\texttt{x\_0},
  \texttt{y\_0}) \помножити на \помножити на C(\texttt{x\_0},
  \texttt{y\_0}) \textless{} 0 и в зависимости от типа его можно
  привести к одной из канонических форм. Для того чтобы было легче
  разобраться в классификационной схеме, мы приведем четыре примера
  гиперболических, параболических и эллиптических уравнений.
\end{enumerate}

\section{Примеры гиперболических, параболических и эллиптических
уравнений
10}\label{ux43fux440ux438ux43cux435ux440ux44b-ux433ux438ux43fux435ux440ux431ux43eux43bux438ux447ux435ux441ux43aux438ux445-ux43fux430ux440ux430ux431ux43eux43bux438ux447ux435ux441ux43aux438ux445-ux438-ux44dux43bux43bux438ux43fux442ux438ux447ux435ux441ux43aux438ux445-ux443ux440ux430ux432ux43dux435ux43dux438ux439-10}

\begin{enumerate}
\def\labelenumi{\arabic{enumi}.}
\tightlist
\item
  У равнение теплопроводности \texttt{u\_t\ =\ u\_\{xx\}}--- линейное
  уравнение второго порядка вида (23.1) с коэффициентами
\end{enumerate}

\begin{verbatim}
A=1
,   B=0  ,   C=0  ,   D=0  ,   E=-1  ,   F=0  ,   G=0  ,
\end{verbatim}

так что B\^{}2-4AC=0 при всех x и t. Следовательно, это уравнение будет
параболическим при всех x и t. В общем уравнении (23 1) переменная y
играет роль времени t. Результат не изменится, если в общем уравнении
поменять местами x и y.

\begin{enumerate}
\def\labelenumi{\arabic{enumi}.}
\setcounter{enumi}{1}
\tightlist
\item
  Волновое уравнение \texttt{u\_\{1t\}\ =\ u\_\{xx\}}также принадлежит
  классу уравнений вида (23.1) с коэффициентами
\end{enumerate}

\begin{verbatim}
A=1
,   B=0  ,   C=-1  ,   D=E=F=G=0  .
\end{verbatim}

Так как B\^{}2 - 4AC = 4 для всех x и t, то это уравнение является
еиперболическим при всех x и t.

1) Следует указать, что приведенные в примерах уравнения являются
частными случаями более общих параболических, гиперболических в
эллиптических уравнений, а основные их свойства характерны и для общего
случая.

\begin{enumerate}
\def\labelenumi{\arabic{enumi}.}
\setcounter{enumi}{2}
\tightlist
\item
  У равнение Лапласа \texttt{u\_\{xx\}} +
  \texttt{u\_\{yy\}\ =\ 0}является эллиптическим при всех x и y, так как
  B\^{}2 - 4AC = - 4 \textless{} 0 .
\end{enumerate}

\begin{verbatim}
4. Линейное уравнение   xu_{xx} + u_{yy} = \sin x   с переменными коэффициентами также принадлежит классу уравнений вида (23.1), но для него   B^2-4AC=-4x   и, следовательно, уравнение будет эллиптическим при x > 0,
\end{verbatim}

лараболическим при x=0, гиперболическим при x \textless{} 0.

Этот пример иллюстрирует тот факт, что тип уравнения о переменными
коэффициентами может изменяться от точки к точке.

Читатель должен обратить внимание на следующее: тип уравнения (23.1)
определяется только коэффициентами при вторых првизводных и никак не
зависит ни от коэффициентов при первых производных и самой функции, ни
от свободного члена.

\begin{verbatim}
Обратимся теперь к главной теме нашей лекции — приведению гиперболических уравнений к каноническому виду. Известно, что если в данной области пространства уравнение является гиперболическим, то вместо x и y можно ввести координаты   \xi   и   \eta  (характеристические переменные) таким образом, что уравнение примет простейшую форму
\end{verbatim}

(23.2)

\begin{verbatim}
u_{\xi\eta} = \Phi(\xi, \eta, u, u_{\xi}, u_{\eta}).
\end{verbatim}

\begin{verbatim}
Это уравнение содержит только одну производную второго порядка   u_{\xi\eta}  , а функция   \Phi(\xi,) \eta,u,u_{\xi},u_{\eta})   зависит только от новых независимых переменных § и η, зависимой переменной и и первых производных из и ип. Конкретный вид функции Ф зависит, конечно, от исходного уравнения и формул перехода к новым координатам § и п. Научимся отыскивать эти функции.
\end{verbatim}

\section{Каноническая форма гиперболического
уравнения}\label{ux43aux430ux43dux43eux43dux438ux447ux435ux441ux43aux430ux44f-ux444ux43eux440ux43cux430-ux433ux438ux43fux435ux440ux431ux43eux43bux438ux447ux435ux441ux43aux43eux433ux43e-ux443ux440ux430ux432ux43dux435ux43dux438ux44f}

Рассмотрим уравнение общего вида

(23.3)

\begin{verbatim}
Au_{xx} + Bu_{xy} + Cu_{yy} + Du_x + Eu_y + Fu = G
\end{verbatim}

и пусть в интересующей нас области B\^{}2 - 4AC \textgreater{} 0 . Нам
нужно ввести новые координаты 1)

\begin{verbatim}
\xi = \xi(x, y),
\eta = \eta(x, y)
\end{verbatim}

\texttt{тому\ в\ рівнянні\ (23.3)\ залишається\ лише\ одна\ часткова\ похідна\ другого\ порядку\ \textquotesingle{}u\_\{\textbackslash{}xi\textbackslash{}eta\}\textquotesingle{}\ (не\ слід\ шукати\ заміну\ змінним,\ за\ допомогою\ якої\ гіперболічне\ вирівнювання}

1) Необхідно, щоб це перетворення координат було локально оборотним і
подвійно диференційованим. --- Ед. Ед.

У цьому випадку це неможливо).

Вычислим сначала частные производные

\begin{verbatim}
\begin{array}{c} u_x = u_\xi \xi_x + u_\eta \eta_x, \\ u_y = u_\xi \xi_y + u_\eta \eta_y, \\ (23.4) \ u_{xx} = u_{\xi\xi} \xi_x^2 + 2u_{\xi\eta} \xi_x \eta_x + u_{\eta\eta} \eta_x^2 + u_{\xi} \xi_{xx} + u_{\eta} \eta_{xx}, \\ u_{xy} = u_{\xi\xi} \xi_x^2 + u_{\xi\eta} (\xi_x \eta_y + \xi_y \eta_x) + u_{\eta\eta} \eta_x \eta_y + u_{\xi} \xi_{xy} + u_{\eta} \eta_{xy}, \\ u_{yy} = u_{\xi\xi} \xi_y^2 + 2u_{\xi\eta} \xi_y \eta_y + u_{\eta\eta} \eta_y^2 + u_{\xi} \xi_{yy} + u_{\eta} \eta_{yy}. \end{array}
\end{verbatim}

Підставляючи ці відношення до початкового рівняння (23.3), після
простих, але громіздких обчислень отримуємо

(23.5)

\begin{verbatim}
\overline{A}u_{\xi\xi} + \overline{B}u_{\xi\eta} + \overline{C}u_{\eta\eta} + \overline{D}u_{\xi} + \overline{E}u_{\eta} + \overline{F}u = \overline{G},
\overline{A} = A\xi_{x}^{2} + B\xi_{x}\xi_{y} + C\xi_{y}^{2},
\overline{B} = 2A\xi_{x}\eta_{y} + B(\xi_{x}\eta_{y} + \xi_{y}\eta_{x}) + 2C\xi_{y}\eta_{y},
\overline{C} = A\eta_{x}^{2} + B\eta_{x}\eta_{y} + C\eta_{y}^{2},
\overline{D} = A\xi_{xx} + B\xi_{xy} + C\xi_{yy} + D\xi_{x} + E\xi_{y},
\overline{E} = A\eta_{xx} + B\eta_{xy} + C\eta_{yy} + D\eta_{x} + E\eta_{y},
\overline{F} = F,
\overline{G} = G.
\end{verbatim}

Читателю предоставится возможность проделать эти выкладки при решении
задач.

\begin{verbatim}
Наш следующий шаг—выбрать функции   \xi = \xi(x, y)   и   \eta = \eta(x, y)   таким образом, чтобы коэффициенты   \overline{A}   и   \overline{C}   обратились в нуль. Это позволит нам привести исходное уравиение к канонической форме. Итак, необходимо потребовать, чтобы
\end{verbatim}

\begin{verbatim}
\bar{A} = A\xi_x^2 + B\xi_x\xi_y + C\xi_y^2 = 0,
\end{verbatim}

\begin{verbatim}
\bar{C} = A\eta_x^2 + B\eta_x\eta_y + C\eta_y^2 = 0.  Эти уравнения можно представить в виде
\end{verbatim}

\begin{verbatim}
\begin{split} &A\left[\xi_{x}/\xi_{y}\right]^{2}+B\left[\xi_{x}/\xi_{y}\right]+C=0,\\ &A\left[\eta_{x}/\eta_{y}\right]^{2}+B\left[\eta_{x}/\eta_{y}\right]^{2}+C=0. \end{split}
\end{verbatim}

\begin{verbatim}
Разрешив их относительно  [\xi_x/\xi_y]  и  [\eta_x/\eta_y] , получаем
\end{verbatim}

(23.7)

\begin{verbatim}
\begin{bmatrix} \xi_x/\xi_y \end{bmatrix}
= \frac{-B + \sqrt{B^2 - 4AC}}{2A}
\end{verbatim}

(характеристические

\begin{verbatim}
[\eta_x/\eta_y] = \frac{-B - \sqrt{B^2 - 4AC}}{2A}
\end{verbatim}

уравнения).

Кожне квадратичне рівняння для
\texttt{{[}\textbackslash{}xi\_x/\textbackslash{}xi\_y{]}} і
\texttt{{[}\textbackslash{}eta\_x/\textbackslash{}eta\_y{]}} має два
ядра, але залишаємо лише одне, щоб вони відрізнялися.

\pandocbounded{\includegraphics[keepaspectratio]{images/_page_170_Figure_3.jpeg}}

Фіг. 23.1. Характеристики \texttt{\textbackslash{}xi(x,\ y)\ =\ c} та
\texttt{\textbackslash{}eta(x,\ y)\ =\ c} .

Задача була зведена до знаходження двох функцій
\texttt{\textbackslash{}xi(x,\ y)} і \texttt{\textbackslash{}eta(x,\ y)}
, так що відношення
\texttt{{[}\textbackslash{}xi\_x/\textbackslash{}xi\_y{]}} і
\texttt{{[}\textbackslash{}eta\_x/\textbackslash{}eta\_y{]}}
задовольняють рівняння (23.7). Дуже легко знайти такі функції 1) якщо
звернути увагу на рис. 23.1. Щоб зрозуміти, як це можна зробити,
розглянемо просте рівняння

\begin{verbatim}
u_{xx} - 4u_{yy} + u_x = 0.
\end{verbatim}

Его характеристики определяются из уравнений

\begin{verbatim}
\begin{aligned}
\frac{dy}{dx} &= -\frac{\xi_x}{\xi_y} = \frac{B - \sqrt{B^2 - 4AC}}{2A} = -2, \\
\frac{dy}{dx} &= -\frac{\eta_x}{\eta_y} = \frac{B + \sqrt{B^2 - 4AC}}{2A} = 2.
\end{aligned}
\end{verbatim}

Решая эти уравнения относительно у, получаем

\begin{verbatim}
y = -2x + c_1,
y = 2x + c_2.
\end{verbatim}

\begin{verbatim}
Чтобы найти   \xi   и   \eta  , разрешим полученные уравнения относительно констант   c_1   и   c_2  , оставляя их в правой части и перенося остальные слагаемые в левую.
\end{verbatim}

Выражения, стоящие в левых частях, примем за искомые функции § и пр

\begin{verbatim}
\xi = y + 2x = c_1,
\eta = y - 2x = c_2.
\end{verbatim}

1) См., например, Тихонов А. Н., Самарский А. А. Уравнения
математической физики. --- М.: Наука, 1972 и др. издания. --- Прим. ред.

Очевидно, що функції \texttt{\textbackslash{}xi} і
\texttt{\textbackslash{}eta} введені таким чином, задовольняють рівняння
характеристик. Ці нові координати показані на рис. 23.2. На цьому наше
обговорення того, як знайти нові координати, завершується. Останнє, що
залишилося --- визначити новий тип початкового рівняння.

\pandocbounded{\includegraphics[keepaspectratio]{images/_page_171_Figure_3.jpeg}}

Фіг. 23.2. Характеристичні координати рівняння
\texttt{u\_\{xx\}\ -\ 4u\_\{uu\}\ +\ u\_x\ =\ 0} .

Це досить просто: знайти канонічну форму диференціального рівняння з
частковими похідними, підставити нові координати
\texttt{\textbackslash{}xi(x,\ y)} і
\texttt{\textbackslash{}eta(x,)\ y)} у рівняння

\begin{verbatim}
\overline{A}u_{\xi\xi} + \overline{B}u_{\xi\eta} + \overline{C}u_{\eta\eta} + \overline{D}u_{\xi} + \overline{E}u_{\eta} + \overline{F}u = \overline{G}
\end{verbatim}

,

де \texttt{\textbackslash{}overline\{A\}} ,
\texttt{\textbackslash{}overline\{B\}} ,
\texttt{\textbackslash{}overline\{C\}} ,
\texttt{\textbackslash{}overline\{D\}} ,
\texttt{\textbackslash{}overline\{E\}} ,
\texttt{\textbackslash{}overline\{F\}} та
\texttt{\textbackslash{}overline\{G\}} визначаються формулами (23.6).

Прежде чем закончить лекцию, давайте на конкретном примере посмотрим,
как «работает» общий метод приведения к каноническому виду.

Приведение гиперболического уравнения
\texttt{y\^{}2u\_\{xx\}-x\^{}2u\_\{yy\}=0} к канонической форме

Рассмотрим уравнение

\begin{verbatim}
y^2u_{xx}-x^2u_{yy}=0
\end{verbatim}

, \texttt{x\textgreater{}} , \texttt{0\ y\textgreater{}0} ,

що є гіперболічним у першому квадранті. Знайдемо нові координати так,
щоб початкове рівняння набуло канонічної форми.

ШАГ 1. Решаем уравнения характеристик

\begin{verbatim}
\frac{dy}{dx} = \frac{B + \sqrt{B^2 - 4AC}}{2A} = \frac{x}{y}
\end{verbatim}

(напомним, что этот шаг

\begin{verbatim}
\frac{dy}{dx} = \frac{B - \sqrt{B^2 - 4AC}}{2A} = -\frac{x}{y}.
\end{verbatim}

обеспечивает
\texttt{\textbackslash{}overline\{A\}\ =\ \textbackslash{}overline\{C\}\ =\ 0}
),

\pandocbounded{\includegraphics[keepaspectratio]{images/_page_172_Figure_2.jpeg}}

Фіг. 23.3. Нові координати карактора.

Інтегруючи ці звичайні диференціальні рівняння методом розділення
змінних, ми отримуємо два неявних співвідношення (за бажанням можна
виразити y через x)

\begin{verbatim}
y^2 - x^2 = \text{const},
y^2 + x^2 = \text{const}.
\end{verbatim}

Следовательно, новые координаты следует вводить по формулам

\begin{verbatim}
\xi = y^2 - x^2
\end{verbatim}

,

\begin{verbatim}
\eta = y^2 + x^2  .
\end{verbatim}

Эти новые координаты изображены на рис. 23.3. Теперь следует получить
новое уравнение. Вычислим все коэффициенты

\begin{verbatim}
\widetilde{A} = 0   (так и должно быть; мы специально выбирали   \xi \xi   и   \eta   так, чтобы этот коэффициент обращался в нуль).
\end{verbatim}

\begin{verbatim}
\vec{B} = 2A\xi_x\eta_x + B(\xi_x\eta_y + \xi_y\eta_x) + 2C\xi_y\eta_y = -16x^2y^2
\end{verbatim}

,

\begin{verbatim}
\overline{C} = 0   (по той же причине, что и   \overline{A} = 0  ),
\end{verbatim}

\begin{verbatim}
\bar{D} = A\xi_{xx} + B\xi_{xy} + C\xi_{yy} + D\xi_x + E\xi_y = -2(x^2 + y^2)
\widetilde{E} = A\eta_{xx} + B\eta_{xy} + C\eta_{yy} + D\eta_x + E\eta_y = 2(y^2 - x^2),
\tilde{F} = F = 0
\end{verbatim}

,

\begin{verbatim}
\overline{G} = G = 0
\end{verbatim}

,

подставим их в уравнение

\begin{verbatim}
\overline{A}u_{\xi\xi} + \overline{B}u_{\xi\eta} + \overline{C}u_{\eta\eta} + \overline{D}u + \overline{E}u_{\eta} + \overline{F}u = \overline{G}
\end{verbatim}

и получим

\begin{verbatim}
u_{\xi\eta} = \frac{-(x^2 + y^2) u_{\xi} + (y^2 - x^2) u_{\eta}}{8x^2y^3}.
\end{verbatim}

\begin{verbatim}
ШАГ 2. Выражая, наконец, x и y через   \xi   и   \eta  , окончательно получаем
\end{verbatim}

(23.8)

\begin{verbatim}
u_{\xi\eta} = \frac{\eta u_{\xi} - \xi u_{\eta}}{2(\xi^2 - \eta^2)}.
\end{verbatim}

Задача решена. Теперь мы знаем, что нужно сделать, чтобы найти новые
координаты и привести уравнение к канонической форме.

\section{ЗАУВАЖЕННЯ}\label{ux437ux430ux443ux432ux430ux436ux435ux43dux43dux44f}

\begin{enumerate}
\def\labelenumi{\arabic{enumi}.}
\tightlist
\item
  На самом деле для гиперболических уравнений существует две
  канонические формы. Вторую можно получить из первов заменой переменных
  вида
\end{enumerate}

\begin{verbatim}
\alpha = \alpha (\xi, \eta) = \xi + \eta,
\end{verbatim}

\begin{verbatim}
\beta = \beta (\xi, \eta) = \xi - \eta.  Воспользуемся этой заменой для преобразования уравнения (23.8). Так как
\end{verbatim}

\begin{verbatim}
\begin{array}{l} u_{\xi} = u_{\alpha}\alpha_{\xi} + u_{\beta}\beta_{\xi} = u_{\alpha} + u_{\beta}. \\ u_{\eta} = u_{\alpha}\alpha_{\eta} + u_{\beta}\beta_{\eta} = u_{\alpha} - u_{\beta}, \\ u_{\xi\eta} = u_{\alpha\alpha}\alpha_{\eta} + u_{\alpha\beta}\beta_{\eta} + u_{\beta\alpha}\alpha_{\eta} + u_{\beta\beta}\beta_{\eta} = u_{\alpha\alpha} - u_{\beta\beta}, \end{array}
\end{verbatim}

TO

(23 9)

\begin{verbatim}
u_{\alpha\alpha} - u_{\beta\beta} = \frac{-\beta u_{\alpha} - \alpha u_{\beta}}{2\alpha\beta}.
\end{verbatim}

\begin{verbatim}
При желании можно выразить   \alpha   и   \beta   через исходные координаты x и y. В результате получаем
\end{verbatim}

\begin{verbatim}
\alpha = \xi + \eta = (y^2 - x^2) + (y^2 + x^2) = 2y^3,
\beta = \xi - \eta = (y^2 - x^2) - (y^2 + x^2) = -2x^2.
\end{verbatim}

\begin{itemize}
\tightlist
\item
  Читач може запитати, чому необхідно класифікувати та зводити їх до
  канонічної форми диференціальних рівнянь у похідних похідних.

  \begin{itemize}
  \tightlist
  \item
    \begin{enumerate}
    \def\labelenumi{\alph{enumi})}
    \tightlist
    \item
      Поділ рівнянь на класи гіперболічних, параболічних та еліптичних
      відповідає поділу фізичних процесів на три основні класи:
      хвильовий, дифузійний і стаціонарний. Математичні особливості
      розв'язків цих трьох типів рівнянь є абсолютно різними.
    \end{enumerate}
  \item
    \begin{enumerate}
    \def\labelenumi{\alph{enumi})}
    \setcounter{enumi}{1}
    \tightlist
    \item
      У більшості праць, присвячених розв'язанню гіперболічних задач,
      припускається, що початкове рівняння записується у канонічній
      формі, тобто у вигляді
    \end{enumerate}
  \end{itemize}
\end{itemize}

\begin{verbatim}
u_{\xi\xi} - u_{\eta\eta} = \Psi(\xi, \, \eta, \, u_{\xi}, \, u, u_{\eta}).
\end{verbatim}

Якщо у нас є рівняння і ми хочемо отримати його розв'язок, ми повинні
перетворити його у канонічну форму і використати відомі результати.

\begin{enumerate}
\def\labelenumi{\alph{enumi})}
\setcounter{enumi}{2}
\tightlist
\item
  Для чисельного розв'язання канонічного гіперболічного рівняння було
  створено багато різних комп'ютерних програм.
\end{enumerate}

Функція \texttt{\textbackslash{}Phi} ( \texttt{\textbackslash{}xi} ,
\texttt{\textbackslash{}xi\ \textbackslash{}eta} , u,
\texttt{u\_\{\textbackslash{}xi\}} , \texttt{u\_\{\textbackslash{}eta\}}
) вводиться у VM \texttt{\textbackslash{}Theta} як підпрограма, тому
рівняння спочатку має бути канонічним. Як тільки розв'язок знайдено в
нових координатах, його завжди можна перетворити назад на початкові
координати. \# ЗАВДАННЯ - 1. Дізнайтеся, чи є наступні рівняння
гіперболічними, параболічними чи еліптичними: - : \textbf{a)}
\texttt{u\_\{xx\}\ u\_\{xy\}\ =\ 0} , - 6)
\texttt{u\_\{tt\}\ =\ u\_\{xx\}\ +\ u\_x\ +\ hu} ,
\texttt{-\ **B)**\ \textquotesingle{}u\_\{xx\}\ +\ 3u\_\{yy\}\ =\ \textbackslash{}sin\ x\textquotesingle{}\ ,}
- r) \texttt{u\_\{xx\}\ +\ u\_\{yy\}\ =\ f(x,\ y)} , -
\texttt{\textbackslash{}mathbf\{D\}\ )\ \textbackslash{}\ u\_\{rr\}\ +\ \textbackslash{}frac\{1\}\{r\}\ u\_r\ +\ \textbackslash{}frac\{1\}\{r\^{}2\}\ u\_\{\textbackslash{}theta\ \textbackslash{}theta\}\ =\ f(r,\ \textbackslash{}\ \textbackslash{}theta).}
- 2. Отримайте співвідношення (23,4), (23,5) та (23,6). - 3.
Переконайтеся, що рівняння

\begin{verbatim}
3u_{xx} + 7u_{xy} + 2u_{yy} = 0
\end{verbatim}

гиперболического типа при всех x и y и найдите
\texttt{x\ a\ p\ a\ \textbackslash{}kappa\ m\ e\ p\ u\ c} - m u + e c
\kappa u + e c \kappa u + e c \kappa u + e c \kappa u + e c \kappa u + e
c \kappa u + e c \kappa u + e c \kappa u + e c \kappa u + e c \kappa u +
e c \kappa u + e c \kappa u + e c \kappa u + e c \kappa u + e c \kappa u
+ e c \kappa u + e c \kappa u + e c \kappa u + e c \kappa u + e e c
\kappa u + e c \kappa u + e c \kappa u + e c \kappa u + e c \kappa u + e
c \kappa u + e c \kappa u + e c \kappa u + e c \kappa u + e c \kappa u +
e c \kappa u + e c \kappa u + e c \kappa u + e c \kappa u + e c \kappa u
+ e c \kappa u + e c \kappa u + e c \kappa u + e c \kappa u + e c \kap u
+ e c \kap

pa u + e c \kappa u + e c \kappa u + e c \kappa u + e c \kappa u + e c
\kappa u + e c \kappa u + e c \kappa u + e c \kappa u + e c \kappa u + e
c \kappa u + e c \kappa u + e c \kappa u + e c \kappa u + e c \kappa u +
e c \kappa u + e c \kappa u + e c \kappa u + e c \kappa u + e c \kappa u
+ e c \kappa u + e c \kappa u + e c \kappa u + e e c \kappa u + e c
\kappa u + e c \kappa u + e c \kappa u + e c \kappa u + e c \kappa u + e
c \kappa u + e c \kappa u + e c \kappa u + e c \kappa u + e c \kappa u +
e c \kappa u + e c \kappa u + e c \kappa u + e c \kappa u + e c \kappa u
+ e c \kappa u + e c \kappa u + e c \kappa u + e c \kappa u + e c
\kappa u + e c \kappa u + e c \kappa u + e c \kappa u + e c \kappa u + e
c \kappa u + e c \kappa u + e c \kappa u + e c \kappa u + e c \kappa u +
e c \kappa u + e c \kappa u + e c \kappa u + e c \kappa u + e c \kappa u
+ e c \kappa u + e c \kappa u + e c \kappa u + e c \kappa u + e c
\kappa u + e u + e c \kappa u + e c \kappa u + e c \kappa u + e c
\kappa u + e c \kappa u + e c \kappa u + e c \kappa u + e c \kappa u + e
c \kappa u + e c \kappa u + e c \kappa u + e c \kappa u + e c \kappa u +
e c \kappa u + e c \kappa u + e c \kappa u + e c \kappa u + e c \kappa u
+ e c \kappa u + e c ~kappa u + e c \kappa u + e c \kappa u + e c
\kappa u + e c \kappa u + e c \kappa u + e c \kappa u + e c \kappa u + e
c \kappa u + e c \kappa u + e c \kappa u + e c \kappa u + e c \kappa u +
e c \kappa u + e c \kappa u + e c \kappa u + e c \kappa u + e c \kappa u
+ e c \kappa u + e e c \kappa u + e c \kappa u + e c \kappa u + e c
\kappa u + e c \kappa u + e c \kappa u + e c \kappa u + e c \kappa u + e
c \kappa u + e c \kappa u + e c \kappa u + e c \kappa u + e c \kappa u +
e c \kappa u + e c \kappa u + e c \kappa u + e c \kappa u + e c \kappa u
+ e c \kappa u + e c \kappa u + e c \kappa u + e c \kappa u + e c
\kappa u + e c \kappa u + e c \kappa u + e c \kappa u + e c \kappa u + e
c \kappa u + e c \kappa u + e c \kappa u + e c \kappa u + e c \kappa u +
e c \kappa u + e c \kappa u + e c \kappa u + e c \kappa u + e c \kappa u
+ e c \kappa u + e c \kappa u + e c \kappa u + e u + e c \kappa u + e c
\kappa u + e c \kappa u + e c \kappa u + e c \kappa u + e c \kappa u + e
c \kappa u + e c \kappa u + e c \kappa u + e c \kappa u + e c \kappa u +
e c \kappa u + e c \kappa u + e c \kappa u + e c \kappa u + e c \kappa u
+ e c \kappa u + e c \kappa u + e c \kappa u + e c ~kappa u + e c
\kappa u + e c \kappa u + e c \kappa u + e c \kappa u + e c \kappa u + e
c \kappa u + e c \kappa u + e c \kappa u + e c \kappa u + e c \kappa u +
e c \kappa u + e c \kappa u + e c \kappa u + e c 4. Приведите уравнение
из предыдущей задачи к каноническому виду

\begin{verbatim}
u_{\xi\eta} = \Phi(\xi, \eta, u, u_1, u_2).
\end{verbatim}

\begin{enumerate}
\def\labelenumi{\arabic{enumi}.}
\setcounter{enumi}{4}
\tightlist
\item
  Продолжите вадачу 4 и приведите уравнение ко второй канонической форме
\end{enumerate}

\begin{verbatim}
u_{\alpha\alpha} - u_{\beta\beta} = \Psi(\alpha, \beta, u, u_{\alpha}, u_{\beta}).
\end{verbatim}

\begin{enumerate}
\def\labelenumi{\arabic{enumi}.}
\setcounter{enumi}{5}
\tightlist
\item
  Найдите характеристики уравнения
\end{enumerate}

\begin{verbatim}
u_{xx} + 4u_{xy} = 0.
\end{verbatim}

Перетворіть рівняння на нові координати, розв'яжіть його і поверніться
до координат x і y, щоб отримати розв'язок початкової здогадки.

\section{ХВИЛЬОВЕ РІВНЯННЯ У ВІЛЬНОМУ ПРОСТОРІ (ДВОВИМІРНІ ТА ТРИВИМІРНІ
ЗАДАЧІ)}\label{ux445ux432ux438ux43bux44cux43eux432ux435-ux440ux456ux432ux43dux44fux43dux43dux44f-ux443-ux432ux456ux43bux44cux43dux43eux43cux443-ux43fux440ux43eux441ux442ux43eux440ux456-ux434ux432ux43eux432ux438ux43cux456ux440ux43dux456-ux442ux430-ux442ux440ux438ux432ux438ux43cux456ux440ux43dux456-ux437ux430ux434ux430ux447ux456}

МЕТАЛЕКЦІЇ: Розв'яжіть задачу з початковими умовами

(УЧП)

\begin{verbatim}
u_{tt} = e^{2} (u_{xx} + u_{yy} + u_{zz}),
\begin{cases} -\infty < x < \infty, \\ -\infty < y < \infty, \\ -\infty < z < \infty, \end{cases}
\end{verbatim}

(НУ)

\begin{verbatim}
\begin{cases} u(x, y, z, 0) = \varphi(x, y, z), \\ u_{t}(x, y, z, 0) = \psi(x, y, z) \end{cases}
\end{verbatim}

в трехмерном пространстве и показать, что полученное решение
удовлетворяет \emph{принципу Гюйгенса}. Для решения соответствующей
двумерной задачи

(УЧП)

\begin{verbatim}
u_{tt} = c^2 (u_{xx} + u_{yy}), \quad
\begin{cases} -\infty < x < \infty, \\ -\infty < y < \infty, \end{cases}
\end{verbatim}

(НУ)

\begin{verbatim}
Mr_1 + K = 2643
\end{verbatim}

використовується метод спуску, і доведено, що двовимірні розв'язки не
задовольняють принцип Гюйгенса. І нарешті, використовуючи метод спуску,
показано, що розв'язок одномірного рівняння виражається формулою
д'Аламбера, з якою ми познайомилися раніше.

У попередніх лекціях ми розглядали задачу з початковими умовами для
нескінченної вібруючої струни і показали, що її розв'язок задається
формулою д'Аламбера. Читачеві має бути зрозуміло, що в тривимірному
просторі одномірне хвильове рівняння описує плоскі хвилі. Наприклад,
звукові або електромагнітні хвилі на достатньо великих відстанях від
джерел можна вважати плоскими хвилями і, отже, описуватися одномірним
рівнянням. Існує така термінологія:

\begin{itemize}
\tightlist
\item
  \begin{enumerate}
  \def\labelenumi{\arabic{enumi})}
  \tightlist
  \item
    одномірні хвилі називаються плоскими хвилями,
  \end{enumerate}
\item
  \begin{enumerate}
  \def\labelenumi{\arabic{enumi})}
  \setcounter{enumi}{1}
  \tightlist
  \item
    Двовимірні хвилі називають циліндричними хвилями.
  \end{enumerate}
\end{itemize}

\begin{enumerate}
\def\labelenumi{\arabic{enumi})}
\setcounter{enumi}{2}
\tightlist
\item
  тривимірні хвилі називаються сферичними хвилями.
\end{enumerate}

Іншими словами, одномірне хвильове рівняння може описувати або плоскі
хвилі в просторах більшої кількості розмірів, або коливання струн.
Завдання лекції --- узагальнити формулу д'Аламбера для випадку двох і
трьох вимірів.

\section{Хвилі у тривимірному
просторі}\label{ux445ux432ux438ux43bux456-ux443-ux442ux440ux438ux432ux438ux43cux456ux440ux43dux43eux43cux443-ux43fux440ux43eux441ux442ux43eux440ux456}

Розглянемо сферичні хвилі у тривимірному просторі за заданих початкових
умов, тобто розв'яжемо задачу з початковими умовами

(24.1)

\begin{verbatim}
\begin{cases} -\infty < x < \infty, \\ -\infty < y < \infty, \\ -\infty < z < \infty, \end{cases}
\end{verbatim}

(24.1)

\begin{verbatim}
\begin{cases} u(x, y, z, 0) = \varphi(x, y, z), \\ u_{t}(x, y, z, 0) = \psi(x, y, z). \end{cases}
P_{2}
\begin{aligned} u_{tt} &= c^2 \Delta u, \\ (HY) & \begin{cases} u(x, y, z, 0) = 0, \\ u_t(x, y, z, 0) = \psi(x, y, z), \end{cases} \end{aligned}
\end{verbatim}

\texttt{где\ экаком\ \textquotesingle{}\textbackslash{}Delta\textquotesingle{}\ обозначен\ дифференциальный\ оператор}

\begin{verbatim}
Цю задачу можна розв'язати методом  \Phi  Юр'єра  ^{1}  ), і виявляється, що розв'язок можна представити у вигляді
\end{verbatim}

\begin{verbatim}
P_{1}
\end{verbatim}

\begin{verbatim}
где  \overline{\psi}  — *среднее* значение начального распределения  \psi  *по сфере* радиуса ct с центром в точке (x, y, z), т. е.
\end{verbatim}

\texttt{\textquotesingle{}+\ ct\ \textbackslash{}cos\ \textbackslash{}varphi\textquotesingle{}\ )\ \textquotesingle{}(ct)\^{}2\ \textbackslash{}sin\ \textbackslash{}varphi\ d\textbackslash{}varphi\ d\textbackslash{}theta\textquotesingle{}\ .}

\texttt{Аргументи\ функції\ \textquotesingle{}\textbackslash{}psi\textquotesingle{}\ проходити\ по\ поверхні\ сфери,\ якщо\ значення\ \textquotesingle{}\textbackslash{}theta\textquotesingle{}\ і\ \textquotesingle{}\textbackslash{}phi\textquotesingle{}\ змінюються\ відповідно\ в\ межах\ \textquotesingle{}{[}0,\ 2\textbackslash{}pi{]}\textquotesingle{}\ і\ \textquotesingle{}{[}0,\ \textbackslash{}pi{]}\textquotesingle{}\ (рисунок\ 24.1).\ Цю\ формулу\ можна\ інтерпретувати\ так:}

См Владимиров В. С. Уравнения математической физики. --- М.: Наука, 1971
и др. издания. --- Прим ред.

дующим образом: каждая точка пространства излучает расходящуюся (со
скоростью c) сферическую волну, через t секунд точка с координатами (x,
y, z) окажется под воздействием начального Р И С

\pandocbounded{\includegraphics[keepaspectratio]{images/_page_177_Picture_3.jpeg}}

. 24 1. Решение является средним значением начатьного распределения на
сфере.

возмущения, сосредоточенного на сфере раднуса ct с центром в заданной
точке (см. рис. 24.2).

Формула (24.3) дозволяє обчислити розв'язок задачі для більшості
початкових умов на комп'ютері. Можливо, читачу буде цікаво знайти
розв'язок для деяких простих функцій f.

\pandocbounded{\includegraphics[keepaspectratio]{images/_page_177_Picture_7.jpeg}}

Фіг. 24.2. Початкове збурення \texttt{\textbackslash{}psi} поширюється з
кожної точки у всіх напрямках.

Щоб отримати повне рішення, розглянемо другу половину проблеми

(24.4)

\begin{verbatim}
\begin{aligned} (\text{YAII}) & u_{tt} = c^2 \Delta u, \\ (\text{HY}) & \begin{cases} u = 0, \\ u_t = \varphi, \end{cases} \end{aligned}
u_h(r, \theta) = \sum_{n=0}^{\infty} r^n [a_n \cos(n\theta) + b_n \sin(n\theta)].
\Delta u = -\frac{r^2}{2} - \frac{r^3}{2} \cos{(2\theta)}
u = \frac{\partial}{\partial t} [t\overline{\varphi}].
u_p(r, \theta) = Ar^4 + Br^4 \cos(2\theta).
u(x, t) = \frac{1}{2c} \int_{x-ct}^{x+ct} \varphi(s) ds.
u_p(r, \theta) = -\frac{r^4}{32} - \frac{r^4}{24} \cos(2\theta).
u_t(x, t) = \frac{1}{2} [\varphi(x + ct) + \varphi(x - ct)].
\sum_{n=0}^{\infty} [a_n \cos(n\theta) + b_n \sin(n\theta)] - \frac{1}{32} - \frac{1}{24} \cos(2\theta) = 0.
(\text{УЧП}) \quad u_{tt} = c^2 \Delta u, \quad (x, y, z) \in \mathbb{R}^3,
u_1(r, \theta) = \frac{1}{32} + \frac{1}{24}r^2\cos(2\theta) - \frac{r^4}{24}\cos(2\theta) - \frac{r^4}{32} = \frac{(r^4 - 1)}{32} - \frac{(r^4 - r^2)}{24}\cos(2\theta).
(\text{HV}) \quad
u = u_0 + u_1 = r \cos \theta - \frac{(r^4 - 1)}{32} - \frac{(r^4 - r^2)}{24} \cos (2\theta).
\begin{cases} u = \varphi, \\ u_t = \psi \end{cases}
P_{2}
\Delta u = 0, \quad 0 < r < 1 + \frac{1}{4} \sin \theta,
\end{verbatim}

записывается в виде

\begin{verbatim}
u \left( 1 + \frac{1}{4} \sin \theta, \ \theta \right) = \cos \theta, \quad 0 \le \theta \le 2\pi.
\end{verbatim}

\begin{verbatim}
где  \overline{\phi}  и  \overline{\psi}  — средние значения функций  \phi  и  \psi  на сфере раднусом ct с центром в точке (x, y, z).
\end{verbatim}

\begin{verbatim}
Цей розв'язок зазвичай називають формулою Пуассона для хвильового рівняння у тривимірному просторі. Це природне узагальнення формули д'Аламбера на тривимірний випадок. Найважливішим у формулі Пуассона є те, що обидва інтеграли в  \bar{\phi}  і  \bar{\psi}  взяті з поверхні сфери. Цей факт дозволяє нам надати таку важливу інтерпретацію рішення.
\end{verbatim}

\begin{verbatim}
У момент  t=t_1  розв'язок u у (x,y,z) залежить лише від початкових розподілів сфери з радіусом  ct_1  центрованим на (x,y,z) (рисунок 24.3). Тепер припустимо, що початкові розподіли  \phi  і  \psi  дорівнюють нулю всюди, крім малої сфери (див. рисунок 24.3). З часом радіус P И С
\end{verbatim}

\pandocbounded{\includegraphics[keepaspectratio]{images/_page_179_Picture_3.jpeg}}

. 24.3. Схема поясияет, как начальное возмущение влияет на точку (к, ч,
г).

\begin{verbatim}
Сфера, що оточує точку (x, y, z), зростає зі швидкістю c, так що через  t_{\hat{z}}  секунди ця сфера починає перетинати область початкового збурення, і тому розв'язок u(x, y, z, t) стає ненульовим. У  t_{\hat{z}} < t < t_{\hat{a}}  розв'язок U у (x, y, z) залишається ненульовим, оскільки сфера продовжує перетинати область початкового збурення, але на  t=t_{\hat{a}}  розв'язок знову стає нульовим. Іншими словами, розв'язок, що поширюється з початкової області збурення, має чітко визначену задню кромку. Цей загальний принцип відомий як принцип Гюйгенса для тривимірного простору. Виявляється, що передній край хвилі завжди чітко окреслений, але задній край чітко окреслений лише в просторах розмірності 3, 5, 7.... Ми вже знаємо з формули д'Аламбера, що початкові умови
\end{verbatim}

\pandocbounded{\includegraphics[keepaspectratio]{images/_page_349_Figure_5.jpeg}}

\begin{verbatim}
u\left(1+\frac{1}{4}\sin\theta,\ \theta\right)=\cos\theta
\end{verbatim}

\begin{verbatim}
приводят, вообще говоря, к нерезкому заднему фронту (поскольку формула Даламбера подразумевает интегрирование функции  \psi(x)  в пределах от x-ct до x+ct).
\end{verbatim}

Покажем теперь, что принцип Гюйгенса неприменим к цилиндрическим волнам.
Такого рода решения описывают волны на воде от точечного источника,
когда задний фронт волны не резкий, а постепенно затухает до нуля.

\section{Двумерное волновое
уравнение}\label{ux434ux432ux443ux43cux435ux440ux43dux43eux435-ux432ux43eux43bux43dux43eux432ux43eux435-ux443ux440ux430ux432ux43dux435ux43dux438ux435}

Чтобы решить двумерную волновую задачу

(УЧП)

\begin{verbatim}
u(1+\epsilon\sin\theta, \theta) = \cos\theta, \quad 0 \le \epsilon \le 1/4.
f(x+h) = f(x) + f'(x)h + f''(x)\frac{h^2}{2!} + \dots,
\begin{cases} -\infty < x < \infty, \\ -\infty < y < \infty, \end{cases}
\end{verbatim}

\pandocbounded{\includegraphics[keepaspectratio]{images/_page_350_Figure_2.jpeg}}

\begin{verbatim}
u(1+\varepsilon\sin\theta,\ \theta)=u(1,\ \theta)+u_r(1,\ \theta)(\varepsilon\sin\theta)+u_{rr}(1,\ \theta)\frac{(\varepsilon\sin\theta)^2}{21}+\ldots
\end{verbatim}

(НУ)

\begin{verbatim}
\begin{array}{l} u_{xx} + u_{yy} = 0, \quad 0 < x < 1, \quad 0 < y < 1, \\ u(x, 1) = 0, \quad 0 < x < 1, \quad 0 < y < 1, \\ u(x, 0) = 0, \quad 0 < x < 1, \quad 0 < y < 1, \\ u(0, y) = 1, \quad \frac{\partial u}{\partial x}(1, y) = 0 \end{array}
\end{verbatim}

Розглянемо розв'язок тривимірної задачі за початкових умов, які залежать
лише від двох змінних X і Y.

Сделав это, получим, что трехмерная формула

\begin{verbatim}
\Delta u = 0, \quad 0 < r < 1 + \varepsilon \sin \theta,
\end{verbatim}

для розв'язку та опише циліндричні хвилі і, отже, дасть розв'язок
двовимірної задачі. Цей метод називається методом спуску. Після всіх
розрахунків ми отримуємо 1) (зовсім не тривіальне)

\begin{verbatim}
u(1, \theta) + u_r(1, \theta) (\varepsilon \sin \theta) + u_{rr}(1, \theta) \frac{(\varepsilon \sin \theta)^2}{2!} + \dots = \cos \theta.
\end{verbatim}

Перед нами стоїть розв'язок двовимірного хвильового рівняння, і хоча,
найімовірніше, його практичне застосування передбачає використання
комп'ютера, цікаво аналізувати його з точки зору принципу Гюйгенса. У
цьому розв'язці обидва інтеграли початкових збурень розглядаються вздовж
внутрішньої частини (тут слово «внутрішнє» відіграє вирішальну роль)
кола з радіусом cte e центр у точці (x, y). Іншими словами, якщо ми
проаналізуємо це рішення так само, як у тривимірному випадку, то
побачимо, що початкове збурення поширюється з чітко визначеним переднім
краєм і розмитим заднім краєм. Отже, принцип Гюйгенса не виконується у
просторі двох вимірів.

Нарешті, припустимо, що початкові збурення \texttt{\textbackslash{}phi}
і \texttt{\textbackslash{}phi} залежать лише від однієї змінної. У цьому
випадку виникають плоскі хвилі, і після застосування методу спуску
отримуємо

\&lt;sup\textgreater2) См., например, Тихонов А. Н., Самарский А. А.
Уравнения математической физики. --- М.: Наука, 1972 и др. издания. ---
Прим. ред.

хорошо известную формулу Даламбера

\begin{verbatim}
(46.8) u = u_0 + \varepsilon u_1 + \varepsilon^2 u_2 + \dots.
\end{verbatim}

Слід зазначити, що в цьому випадку використання методу спуску вимагає
нетривіальних розрахунків. За формулою д'Аламбера початкове зміщення f
призводить до утворення гострого заднього краю, а початкова швидкість f
його розмиває. Іншими словами, одномірний випадок є незвичайним тим, що
принцип Гюйгенса виконується для початкових зсув і не застосовується для
початкових швидкостей. Можна стверджувати, що загалом принцип Гюйгенса
не виконується в одномірному випадку. \# ПОВАГА Ми не описували метод
спуску детально, оскільки неможливо зупинитися на питанні, як зменшити
розмірність інтегралів, включених у формули. Усі ці розрахунки можна
знайти в навчальній літературі, перелік якої наведено наприкінці книги.
Загальна ідея полягає в тому, що при розв'язанні задач у низьковимірному
просторі можна використати відоме розв'язання подібної задачі у просторі
вищої вимірності, а потім спростити його, припускаючи, що початкові та
граничні умови не залежать від частини змінних. Читачеві має бути
зрозуміло, що метод спуску можна застосувати не лише до розглянутої
проблеми. \# ЗАВДАННЯ 1. Показати, що в одномірному випадку розв'язок
задачі

\begin{verbatim}
P_{0}
\begin{cases} \Delta u_{0} = 0, & 0 < r < 1 \text{ (внутри круга),} \\ u_{0}(1, \theta) = \cos \theta, & u_{0}(r, \theta) = r \cos \theta, \end{cases}
\begin{cases} u(x, \ 0) = \varphi(x), \\ u_t(x, \ 0) = 0 \end{cases}
P_{1}
\begin{cases} \Delta u_{1} = 0, & 0 < r < 1 \text{ (внутри круга),} \\ u_{1}(1, \theta) = -\sin \theta \frac{\partial u_{0}(1, \theta)}{\partial r} = -\sin \theta \cos \theta \end{cases}
\begin{aligned} & (\forall \mathsf{Y}\Pi) \quad u_{tt} = c^2 u_{xx}, \\ & (\mathsf{H} \mathsf{Y}) \quad \begin{cases} u\left(x, \ 0\right) = 0, \\ u_t\left(x, \ 0\right) = \phi\left(x\right). \end{cases} \end{aligned}
u = u_0 + \frac{1}{4}u_1 + \left(\frac{1}{4}\right)^2 u_1 + \cdots
u_u + \frac{1}{4} u_t
\end{verbatim}

\begin{enumerate}
\def\labelenumi{\arabic{enumi}.}
\setcounter{enumi}{1}
\tightlist
\item
  Результаты задачи 1 используйте для решения следующей задачи:
\end{enumerate}

(УЧП)

\begin{verbatim}
u_t = (1+x) u_{xx}, \\ u(x, 0) = q(x), \quad -\infty < x < \infty,
\end{verbatim}

(НУ)

\begin{verbatim}
u_1 = (1 + \varepsilon_1) u_{xx}
u = u_1 + \varepsilon u_1 + \varepsilon^2 u_2 + \cdots
\begin{cases} u(x, 0) = x, \\ u_t(x, 0) = 0. \end{cases}
P_{0}
\end{verbatim}

\begin{enumerate}
\def\labelenumi{\arabic{enumi}.}
\setcounter{enumi}{2}
\tightlist
\item
  Проілюструвати та прокоментувати розв'язок тривимірної проблеми
\end{enumerate}

(УЧП)

\begin{verbatim}
,  (x,  y, z) \in R^3  ,  u = 0   u_t =
\end{verbatim}

\begin{verbatim}
g(x) = e^{-x^2}.
\end{verbatim}

\begin{enumerate}
\def\labelenumi{\arabic{enumi}.}
\setcounter{enumi}{3}
\tightlist
\item
  Решите аналогичную (см. задачу 3) задачу для двумерного уравнения
\end{enumerate}

(УЧП)

(НУ)

Зверніть увагу, що в перших двох задачах коефіцієнти є сталими. Читач
повинен розуміти, що параметр b має бути достатньо малим, а необережний
ряд може відрізнятися

\begin{verbatim}
# ЗАДАЧІ - 1. Подставьте разложение (46.8) в задачу (46.7) и получите последовательность задач P_0,\ P_1,\ P_2,\ \dots . 2. Покажите, что нелинейную задачу
\end{verbatim}

\begin{enumerate}
\def\labelenumi{\arabic{enumi}.}
\setcounter{enumi}{4}
\tightlist
\item
  Решите аналогичную (см. задачу 3) задачу для одномерных плоских волн
\end{enumerate}

(УЧП)

(НУ)

\begin{itemize}
\tightlist
\item
  \begin{enumerate}
  \def\labelenumi{\arabic{enumi}.}
  \setcounter{enumi}{5}
  \tightlist
  \item
    Какова физическая природа того факта, что принцип Гюйгенса не
    справедлив в двумерном случае?
  \end{enumerate}
\item
  \begin{enumerate}
  \def\labelenumi{\arabic{enumi}.}
  \setcounter{enumi}{6}
  \tightlist
  \item
    Воспользуйтесь формулой Лейбница
  \end{enumerate}
\end{itemize}

\begin{verbatim}
u(r, \theta) = u_0(r, \theta) + \frac{1}{4}u_1(r, \theta)
\end{verbatim}

и продифференцируйте интеграл

\begin{verbatim}
\Delta u = 0,
\end{verbatim}

по времени t.

\section{СКІНЧЕННІ ПЕРЕТВОРЕННЯ ФУР'Є (ПЕРЕТВОРЕННЯ СИНУСА ТА
КОСИНУСА)}\label{ux441ux43aux456ux43dux447ux435ux43dux43dux456-ux43fux435ux440ux435ux442ux432ux43eux440ux435ux43dux43dux44f-ux444ux443ux440ux454-ux43fux435ux440ux435ux442ux432ux43eux440ux435ux43dux43dux44f-ux441ux438ux43dux443ux441ux430-ux442ux430-ux43aux43eux441ux438ux43dux443ux441ux430}

МЕТАЛЕКЦІЇ: Введіть два нових інтегральних перетворення (скінченні
спіусні та косинусні перетворення):

\begin{verbatim}
q_{xx} + \varphi_{uu} = 0
\end{verbatim}

(конечное синус-преобразование),

\begin{verbatim}
q_{uu} + q_{vv} = 0
\end{verbatim}

(конечное косинус-преобразование),

\pandocbounded{\includegraphics[keepaspectratio]{images/_page_353_Figure_9.jpeg}}

(обратное синус-преобразование),

\begin{verbatim}
w = f(z)
\end{verbatim}

(обратное косинус-преобразование),

та показати, як розв'язувати граничні проблеми (зокрема гетерогенні) за
допомогою цих перетворень. Раніше ми ознайомилися з регулярними
перетвореннями Фур'є та Лапласа і побачили, як застосування цих
перетворень до рівняння з частковими похідними зводить задачу до
розв'язання звичайного диференціального рівняння. Звичайне перетворення
Фур'є означає, що змінна, яку потрібно перетворити, змінюється від
\texttt{-\textbackslash{}infty} до \texttt{+\textbackslash{}infty} ,
тому воно використовується для розв'язання задач у вільному просторі
(без меж). У цій лекції ми покажемо, як розв'язувати задачі граничних
значень (з межами) шляхом перетворення обмежених змінних (це перший раз,
коли ми стикаємося з такою трансформацією).

Давайте на деякий час забудемо про наші завдання і розберемося з самими
трансформаціями. Дамо визначення перетворень, формули циркуляції та
основні властивості. Для чого потрібні трансформації і як їх
використовувати, ми дізнаємося трохи пізніше. Однак, коротко кажучи,
суть інтегральних перетворень полягає у розкладі функцій, включених у
задачу як

частотам, решении семейства задач для каждой частоты и после-

дующем суммировании всех полученных результатов.

Почнемо з функції \texttt{f(x)}, визначеної на інтервалі {[}0, L{]}.
Скінченні синусоїдальні та косинусоїдальні перетворення цієї функції
визначаються формулами

\texttt{Конечное\ синус-преобразование\ \textquotesingle{}S{[}f{]}\ =\ S\_n\ =\ \textbackslash{}frac\{2\}\{L\}\ \textbackslash{}int\_\{\textbackslash{}sigma\}\^{}\{L\}\ f(x)\ \textbackslash{}sin(n\textbackslash{}pi\ x/L)\ dx\textquotesingle{}\ \textquotesingle{}n\ =\ 1,\ 2,\ \textbackslash{}dots\textquotesingle{}\ \textquotesingle{}C{[}f{]}\ =\ C\_n\ =\ \textbackslash{}frac\{2\}\{L\}\ \textbackslash{}int\_\{\textbackslash{}sigma\}\^{}\{L\}\ f(x)\ \textbackslash{}cos(n\textbackslash{}pi\ x/L)\ dx\textquotesingle{}\ \textquotesingle{}n\ =\ 0,\ 1,\ \textbackslash{}dots\textquotesingle{}\ Конечное\ косинус-преобразование}

Читач, ймовірно, помітив, що ці перетворення не відрізняються від формул
для визначення коефіцієнтів Фур'є, коли їх розкладають на ряд синусами
або косинусами. Інверсійні формули для цих перетворень є типовими рядами
Фур'є на синусах і косинусах.

Обратное синус-преобразование

\pandocbounded{\includegraphics[keepaspectratio]{images/_page_355_Figure_5.jpeg}}

\pandocbounded{\includegraphics[keepaspectratio]{images/_page_355_Figure_7.jpeg}}

Обратное косинус-преобразование

Отметим, что в обратном косинус-преобразовании суммирование начинается с
n=0, а в синус-преобразовании --- с n=1.

\section{Приклади синусоїдального
перетворення}\label{ux43fux440ux438ux43aux43bux430ux434ux438-ux441ux438ux43dux443ux441ux43eux457ux434ux430ux43bux44cux43dux43eux433ux43e-ux43fux435ux440ux435ux442ux432ux43eux440ux435ux43dux43dux44f}

Пусть

\begin{verbatim}
u + iv = (x + iy)^2 = x^2 - y^2 + 2ixy
\end{verbatim}

тогда

\begin{verbatim}
\left\{
\end{verbatim}

Графік функції \texttt{f(x)} та її перетворення показані на рис. 25.1.
Обернене перетворення записується як

Чи знаєте ви, як виглядає графік цієї функції з сегментом {[}0, 1{]}?
Подумай про це. Зверніть увагу, що синусоїдальне перетворення функції
\texttt{f(x)} визначається лише для додатних цілих чисел n.~Іншими
словами, скінченний синус і косинус перетворюють функції у чисельні
послідовності.

\pandocbounded{\includegraphics[keepaspectratio]{images/_page_185_Figure_3.jpeg}}

Рис. 25.1. Графики функции \texttt{f(x)} = 1 и ее преобразования.

\section{Свойства
преобразований}\label{ux441ux432ux43eux439ux441ux442ux432ux430-ux43fux440ux435ux43eux431ux440ux430ux437ux43eux432ux430ux43dux438ux439}

Прежде чем приступить к решению задач, мы должны получить некоторые
свойства этих преобразований.

\begin{verbatim}
Якщо you(x, t) — функція  \partial \theta yx  змінних, і ми виконуємо перетворення змінної x, то
\end{verbatim}

(действительная форма отображения \texttt{w=z\^{}2} ).

Ця форма нам знадобиться пізніше. У цьому прикладі це цікаво, оскільки
показує, як гіперболи в площині z переходять до координатних прямих u =
const і v = const у площині w.

Тепер перейдемо до вивчення одного конкретного типу функцій комплексної
змінної, які виконують конформні відображення.

\section{Определение конформного
отображения}\label{ux43eux43fux440ux435ux434ux435ux43bux435ux43dux438ux435-ux43aux43eux43dux444ux43eux440ux43cux43dux43eux433ux43e-ux43eux442ux43eux431ux440ux430ux436ux435ux43dux438ux44f-1}

Відображення w=\texttt{f(z)} комплексної площини z у комплексну площину
w називається конформним у точці \texttt{z\_\{\textbackslash{}bullet\}}
площини z, якщо похідна
\texttt{f\textquotesingle{}(z\_\{\textbackslash{}bullet\})\ \textbackslash{}neq\ 0}
. Відображення \texttt{f(z)} називається конформним у регіоні D, якщо
\texttt{f\textquotesingle{}(z)\ \textbackslash{}neq\ 0} в кожній точці
регіону D.

Наприклад, відображення \texttt{f(z)\ =\ z\^{}2} конформно всюди, крім
точки z = 0, оскільки
\texttt{z\textquotesingle{}(z)\ =\ 2z\ \textbackslash{}neq\ 0} для всіх
\texttt{z\ \textbackslash{}neq\ 0} . З іншого боку, відображення
\texttt{e\^{}z} конформно у всій z-площині, оскільки скрізь
\texttt{f\textquotesingle{}(z)\ =\ e\^{}z\ \textbackslash{}neq\ 0} . Яке
призначення конформних відображень? Відповідь полягає в тому, що,
розв'язуючи рівняння Лапласа
\texttt{\textbackslash{}phi\_\{xx\}\ +\ \textbackslash{}phi\_\{yy\}\ =\ 0}
в певній області площини змінних x і y, ми можемо розглядати цю площину
як площину комплексної змінної z. Розглянемо комплексне відображення w =
\texttt{f(z)} площини z у площину w, у якій введені координати (u, v).
Відображуючи w = \texttt{f(z}\phi\emph{\{xy\} + \phi}\{uy\} =
0\texttt{)}, рівняння Лапласа
\texttt{\textbackslash{}varphi\_\{xx\}\ +\ \textbackslash{}varphi\_\{yy\}\ =\ 0}
перетворюється на нове диференціальне рівняння в частинних похідних, яке
залежить від нових координат і v. До --- рівняння, яке
\texttt{\textbackslash{}phi\_\{xx\}\ +\ \textbackslash{}phi\_\{uu\}\ =\ 0}
повернеться у рівняння
\texttt{\textbackslash{}varphi\_\{nn\}\ +\ \textbackslash{}varphi\_\{nn\}\ =\ 0}
. Отже, ми придумаємо ідею знайти конформне відображення, яке переносить
область з комплексною межею у область з простою межею (нагадаємо, що
відображення \texttt{w=z\^{}*} переносить межу першого квадранта площини
z у дійсну вісь площини w).

Ось кілька прикладів, які покажуть, як можна розв'язати диференціальні
рівняння у похідних за допомогою конформних відображень.

Уравнение Лапласа в верхней полуплоскости.

Припустимо, що нам потрібно розв'язати наступну задачу Діріхле у верхній
половині видимості (рис. 47.3):

(47.1)

(Зверніть увагу, що перетворення було виконано на змінній x, і отримана
послідовність залежить лише від часу t.)

А как быть с производными? Ниже приводятся несколько полезных формул

\begin{verbatim}
\begin{aligned} (Y & \Pi) & \varphi_{xx} + \varphi_{yy} = 0, & -\infty < x < \infty, & 0 < y < \infty, \\ (47.1) & \varphi(x, 0) = \begin{cases} 0, & |x| > 1, \\ 1, & |x| \le 1. \end{cases}
\end{verbatim}

\begin{verbatim}
w = \ln\left\{\frac{z-1}{z+1}\right\}
\begin{split} S\left[u_{t}\right] &= \frac{dS\left[u\right]}{dt}, \quad S\left[u_{tt}\right] = \frac{d^{2}S\left[u\right]}{dt^{2}}, \\ S\left[u_{xx}\right] &= -\left[n\pi/L\right]^{2}S\left[u\right] + \frac{2n\pi}{L^{2}}\left[u\left(0,\ t\right) + (-1)^{n+2}u\left(L,\ t\right)\right], \\ C\left[u_{xx}\right] &= -\left[n\pi/L\right]^{2}C\left[u\right] - \frac{2}{L}\left[u_{x}\left(0,\ t\right) + (-1)^{n+1}u_{x}\left(L,\ t\right)\right]. \end{split}
\end{verbatim}

\pandocbounded{\includegraphics[keepaspectratio]{images/_page_357_Figure_11.jpeg}}

\pandocbounded{\includegraphics[keepaspectratio]{images/_page_358_Figure_2.jpeg}}

\begin{verbatim}
(\text{УЧП}) \quad u_{tt} = u_{xx} + \sin(\pi x), \quad 0 < x < 1, \quad 0 < t < \infty,
\begin{cases} (\text{УЧП}) & \varphi_{uu} + \varphi_{vv} = 0, \quad -\infty < u < \infty, \quad 0 < v < \pi, \\ (\text{47.2}) & \begin{cases} \varphi(u, 0) = 0, \\ \varphi(u, \pi) = 1. \end{cases} \end{cases}
(\text{ГУ}) \quad
w = u + iv = \ln\left[\frac{z-1}{z+1}\right] =
\begin{cases} u(0, t) = 0, \\ u(1, t) = 0, \end{cases}
= \ln\left[\frac{z-1}{z+1}\right] + i \arg\left[\frac{z-1}{z+1}\right]
v = \arg\left[\frac{z-1}{z+1}\right] = \arg\left[\frac{x+iy-1}{x+1+iy}\right] =
= \arg\left[\frac{x^2+y^2-1+i2y}{(x+1)^2+y^2}\right] =
= \arg\left[\frac{2y}{x^2+y^2-1}\right]
\begin{cases} u(x, 0) = 1, \\ u_t(x, 0) = 0, \end{cases}
\varphi(x, y) = \frac{1}{\pi} \arctan\left[\frac{2y}{x^2 + y^2 - 1}\right].
x^{2} + y^{2} = 1,
\end{verbatim}

Ми можемо розв'язати задачу наступними трьома кроками:

ШАГ І (Определение подходящего преобразования.)

Оскільки змінна x змінюється від 0 до 1, ми використаємо фінальне
перетворення, а саме синусоїдальне перетворення; Пізніше стане
зрозуміло, чому. Ми могли б розв'язати цю задачу, використовуючи
перетворення Лапласа на змінній t (цей метод не відрізняється від
скінченного синусоїдального перетворення за інтенсивністю праці).

КРОК 2 (Виконати конвертацію).

Для удобства будем использовать обозначение
\texttt{S\_n(t)\ =\ S{[}u{]}} . Применим синус-преобразование к
исходному уравнению

\begin{verbatim}
(x-1)^{2} + y^{2} = 9,
\end{verbatim}

и воспользуемся тождеством для синус-преобразования, получим

\begin{verbatim}
\begin{aligned} & (\text{УЧП}) & \phi_{xx} + \phi_{yy} = 0, & \text{внутри } D, \\ & (\text{ГУ}) & \begin{cases} \varphi(x, y) = 1 & \text{на } x^2 + y^2 = 1, \\ \varphi(x, y) = 2 & \text{на } (x - 1)^2 + y^2 = 9. \end{cases} \end{aligned}
\begin{split} \frac{d^2S_n\left(t\right)}{dt^2} &= -\left(\pi n\right)^2S_n\left(t\right) + 2n\pi\left[u\left(0,\ t\right) + \left(-1\right)^{n+1}u\left(1,\ t\right)\right] + D_n\left(t\right) = \\ &= -\left(n\pi\right)^2S_n\left(t\right) + D_n\left(t\right), \end{split}
\end{verbatim}

\pandocbounded{\includegraphics[keepaspectratio]{images/_page_359_Figure_10.jpeg}}

\begin{verbatim}
w = 2t \left[ \frac{z-s}{z-t} \right],
\end{verbatim}

где

\begin{verbatim}
u = \gamma [(x-s)(x-t)-\gamma^2],
v = \gamma [y(x-t)],
\gamma = 2t/[(x-t)^2 + y^2].
\end{verbatim}

Якщо тепер перетворити початкові умови задачі на граничне значення,
отримаємо початкові умови для звичайного диференціального рівняння

\begin{verbatim}
\begin{aligned} \varphi_{uu} + \varphi_{vv} &= 0, & \text{в кольче,} \\ (\Gamma Y) & & & & & & & \\ (\Gamma Y) & & & & & & \\ (\Gamma Y) & & & & & \\ & & & & & \\ (\Gamma Y) & & & & & \\ & & & & \\ & & & & \\ & & & & \\ & & & & \\ & & & \\ & & & \\ & & & \\ & & & \\ & & & \\ & & & \\ & & \\ & & \\ & & \\ & & \\ & & \\ & & \\ & & \\ & & \\ & & \\ & & \\ & & \\ & & \\ & & \\ & & \\ & & \\ & & \\ & & \\ & & \\ & & \\ & & \\ & & \\ & & \\ & & \\ & & \\ & & \\ & & \\ & & \\ & & \\ & & \\ & & \\ & & \\ & & \\ & & \\ & & \\ & & \\ & & \\ & & \\ & & \\ & & \\ & & \\ & & \\ & & \\ & & \\ & & \\ & & \\ & & \\ & & \\ & & \\ & & \\ & & \\ & & \\ & & \\ & & \\ & & \\ & & \\ & & \\ & & \\ & & \\ & & \\ & & \\ & & \\ & & \\ & & \\ & & \\ & & \\ & & \\ & & \\ & & \\ & & \\ & & \\ & & \\ & & \\ & & \\ & & \\ & & \\ & & \\ & & \\ & & \\ & & \\ & & \\ & & \\ & & \\ & & \\ & & \\ & & \\ & & \\ & & \\ & & \\ & & \\ & & \\ & & \\ & & \\ & & \\ & & \\ & & \\ & & \\ & & \\ & & \\ & & \\ & & \\ & & \\ & & \\ & & \\ & & \\ & & \\ & & \\ & & \\ & & \\ & & \\ & & \\ & & \\ & & \\ & & \\ & & \\ & & \\ & & \\ & & \\ & & \\ & & \\ & & \\ & & \\ & & \\ & & \\ & & \\ & & \\ & & \\ & & \\ & & \\ & & \\ & & \\ & & \\ & & \\ & & \\ & & \\ & & \\ & & \\ & & \\ & & \\ & & \\ & & \\ & & \\ & & \\ & & \\ & & \\ & & \\ & & \\ & & \\ & & \\ & & \\ & & \\ & & \\ & & \\ & & \\ & & \\ & & \\ & & \\ & & \\ & & \\ & & \\ & & \\ & & \\ & & \\ & & \\ & & \\ & & \\ & & \\ & & \\ & & \\ & & \\ & & \\ & & \\ & & \\ & & \\ & & \\ & & \\ & & \\ & & \\ & & \\ & & \\ & & \\ & & \\ & & \\ & & \\ & & \\ & & \\ & & \\ & & \\ & & \\ & & \\ & & \\ & & \\ & & \\ & & \\ & & \\ & & \\ & & \\ & & \\ & & \\ & & \\ & & \\ & & \\ & & \\ & & \\ & & \\ & & \\ & & \\ & & \\ & & \\ & & \\ & & \\ & & \\ & & \\ & & \\ & & \\ & & \\ & & \\ & & \\ & & \\ & & \\ & & \\ & & \\ & & \\ & & \\ & & \\ & & \\ & & \\ & & \\ & & \\ & & \\ & & \\ & & \\ & & \\ & & \\ & & \\ & & \\ & & \\ & & \\ & & \\ & & \\ & & \\ & & \\ & & \\ & & \\ & & \\ & & \\ & & \\ & & \\ & & \\ & & \\ & & \\ & & \\ & & \\ & & \\ & & \\ & & \\ & & \\ & & \\ & & \\ & & \\ & & \\ & & \\ & & \\ & & \\ & & \\ & & \\ & & \\ & & \\ & & \\ & & \\ & & \\ & & \\ & & \\ & & \\ & & \\ & & \\ & & \\ & & \\ & & \\ & & \\ & & \\ & & \\ & & \\ & & \\ & & \\ & & \\ & & \\ & & \\ & & \\ & & \\ & & \\ & & \\ & & \\ & & \\ & & \\ & & \\ & & \\ & & \\ & & \\ & & \\ & & \\ & & \\ & & \\ & & \end{aligned}
\begin{cases} 4/n\pi, & n = 1, 3, \dots, \\ 0, & n = 2, 4, \dots, \end{cases}
\varphi(r) = a \ln r + b, \quad r = u^2 + v^2.
\varphi(u, v) = 0.57 \ln(u^2 + v^2) + 1
\varphi(x, y) = 0.57 \ln(u^2 + v^2) + 1,
\end{verbatim}

Итак, решим теперь новое семейство задач Кощи:

(ОДУ)

\begin{verbatim}
w = \ln \left[ \frac{z-1}{z+1} \right].
z = re^{i\theta}
\begin{cases} 1, & n = 1, \\ 0, & n = 2, 3, \dots, \end{cases}
\end{verbatim}

\begin{verbatim}
\begin{aligned} (\text{Y} \forall \Pi) & \qquad \phi_{\text{vx}} + \phi_{yy} = 0, \quad 0 < x < \infty, \quad 0 < y < \infty, \\ \left\{ \begin{array}{l} \phi(x, \ 0) = \begin{cases} 1, & 0 < x < 1, \\ 0, & 1 \leqslant x < \infty, \end{cases} \\ \phi(0, \ y) = \begin{cases} 1, & 0 < y < 1, \\ 0, & 1 \leqslant y < \infty. \end{cases} \end{aligned}
\end{verbatim}

\begin{verbatim}
\right.
\end{verbatim}

\pandocbounded{\includegraphics[keepaspectratio]{images/_page_362_Figure_1.jpeg}}

\begin{verbatim}
\begin{cases} S_n(0) = \begin{cases} 4/n\pi, & n = 1, 3, \dots, \\ 0, & n = 2, 4, \dots, \end{cases}
\end{verbatim}

\pandocbounded{\includegraphics[keepaspectratio]{images/_page_362_Figure_2.jpeg}}

\pandocbounded{\includegraphics[keepaspectratio]{images/_page_362_Figure_5.jpeg}}

\pandocbounded{\includegraphics[keepaspectratio]{images/_page_362_Figure_7.jpeg}}

\begin{verbatim}
\begin{cases} \frac{dS_n(0)}{dt} = 0. \end{cases}
H(x-a) =
\begin{cases} 0, x < a, \\ 1, x \geqslant a, \end{cases}
\end{verbatim}

В результате получаем

,

где

\begin{verbatim}
H(x-a) =
\begin{cases} 0, & x < a, \\ 1, & x \geqslant a, \end{cases}
\begin{cases} 0, & n = 2, 4, \dots, \\ (4/n\pi)\cos(n\pi t), & n = 3, 5, 7, \dots \end{cases}
f(x) = \mathcal{F}^{-1}[F] = \frac{1}{\sqrt{2\pi}} \int_{-\infty}^{+\infty} F(\omega) e^{-i\omega x} d\omega \qquad F(\omega) = \mathcal{F}[f] = \frac{1}{\sqrt{2\pi}} \int_{-\infty}^{+\infty} f(x) e^{-i\omega x} dx
\end{verbatim}

Следовательно, решение нашей задачи можно представить в виде

\begin{verbatim}
\begin{cases} 1, |x| < a \\ 0, |x| > a \end{cases}
\end{verbatim}

\section{ЗАУВАЖЕННЯ}\label{ux437ux430ux443ux432ux430ux436ux435ux43dux43dux44f-1}

\begin{enumerate}
\def\labelenumi{\arabic{enumi}.}
\item
  Для решения задачи методом конечного синус- или косинуспреобразования
  краевые условия при x=0 и x=L должны иметь вид
\item
\end{enumerate}

\begin{verbatim}
\begin{cases} 1, |x| < 1 \\ 0, |x| > 1 \end{cases}
\begin{cases} u(0, t) = f(t) \\ u(L, t) = g(t) \end{cases}
\delta(x-a)
\end{verbatim}

\texttt{f(x)} * \texttt{g(x)}

\begin{verbatim}
\begin{array}{ll} (\mathrm{Y} \mathrm{U} \Pi) & u_t = u_{xx}, \quad 0 < x < 1, \quad 0 < t < \infty, \\ (\mathrm{\Gamma} \mathrm{Y} & \left\{ \begin{array}{ll} u_x \left( 0, \ t \right) = 0 \\ u_x \left( 1, \ t \right) = 0, \end{array} \end{array}
\end{verbatim}

\begin{verbatim}
(1+x^2)^{-1}
xe^{-a+x^4}, a > 0
\end{verbatim}

\begin{enumerate}
\def\labelenumi{\arabic{enumi}.}
\setcounter{enumi}{1}
\tightlist
\item
  Решите в общем виде задачу
\end{enumerate}

\texttt{H(x+a)} - \texttt{H(x-a)}

\begin{verbatim}
\frac{a}{x^2+a^2}
\end{verbatim}

\begin{verbatim}
\right. & 0 < t < \infty, \\ (\text{НУ}) & u(x, 0) = 0, \quad 0 \leqslant x \leqslant 1. \end{array}
\end{verbatim}

\begin{verbatim}
\frac{2ax}{(x^2+a^2)^2}
\begin{cases} \cos(ax), |x| > \pi/2a \\ 0, |x| > \pi/2a \end{cases}
\end{verbatim}

\section{ПРИНЦИП СУПЕРПОЗИЦІЇ Є ОСНОВОЮ ТЕОРІЇ ЛІНІЙНИХ СИСТЕМ
МЕТАЛЕКЦІЇ: Ознайомити читача з принципом суперпозиції та показати, як
його можна використати для розкладання початкової задачі на підзадачі,
потім розв'язання всіх підзадач одразу і, підсумуючи результати,
знаходження розв'язку початкової проблеми. Ми також хочемо показати, що
два основні методи розв'язання лінійних рівнянь --- метод розділення
змінних і метод інтегрального перетворення --- використовують принцип
суперпозиції.}\label{ux43fux440ux438ux43dux446ux438ux43f-ux441ux443ux43fux435ux440ux43fux43eux437ux438ux446ux456ux457-ux454-ux43eux441ux43dux43eux432ux43eux44e-ux442ux435ux43eux440ux456ux457-ux43bux456ux43dux456ux439ux43dux438ux445-ux441ux438ux441ux442ux435ux43c-ux43cux435ux442ux430ux43bux435ux43aux446ux456ux457-ux43eux437ux43dux430ux439ux43eux43cux438ux442ux438-ux447ux438ux442ux430ux447ux430-ux437-ux43fux440ux438ux43dux446ux438ux43fux43eux43c-ux441ux443ux43fux435ux440ux43fux43eux437ux438ux446ux456ux457-ux442ux430-ux43fux43eux43aux430ux437ux430ux442ux438-ux44fux43a-ux439ux43eux433ux43e-ux43cux43eux436ux43dux430-ux432ux438ux43aux43eux440ux438ux441ux442ux430ux442ux438-ux434ux43bux44f-ux440ux43eux437ux43aux43bux430ux434ux430ux43dux43dux44f-ux43fux43eux447ux430ux442ux43aux43eux432ux43eux457-ux437ux430ux434ux430ux447ux456-ux43dux430-ux43fux456ux434ux437ux430ux434ux430ux447ux456-ux43fux43eux442ux456ux43c-ux440ux43eux437ux432ux44fux437ux430ux43dux43dux44f-ux432ux441ux456ux445-ux43fux456ux434ux437ux430ux434ux430ux447-ux43eux434ux440ux430ux437ux443-ux456-ux43fux456ux434ux441ux443ux43cux443ux44eux447ux438-ux440ux435ux437ux443ux43bux44cux442ux430ux442ux438-ux437ux43dux430ux445ux43eux434ux436ux435ux43dux43dux44f-ux440ux43eux437ux432ux44fux437ux43aux443-ux43fux43eux447ux430ux442ux43aux43eux432ux43eux457-ux43fux440ux43eux431ux43bux435ux43cux438.-ux43cux438-ux442ux430ux43aux43eux436-ux445ux43eux447ux435ux43cux43e-ux43fux43eux43aux430ux437ux430ux442ux438-ux449ux43e-ux434ux432ux430-ux43eux441ux43dux43eux432ux43dux456-ux43cux435ux442ux43eux434ux438-ux440ux43eux437ux432ux44fux437ux430ux43dux43dux44f-ux43bux456ux43dux456ux439ux43dux438ux445-ux440ux456ux432ux43dux44fux43dux44c-ux43cux435ux442ux43eux434-ux440ux43eux437ux434ux456ux43bux435ux43dux43dux44f-ux437ux43cux456ux43dux43dux438ux445-ux456-ux43cux435ux442ux43eux434-ux456ux43dux442ux435ux433ux440ux430ux43bux44cux43dux43eux433ux43e-ux43fux435ux440ux435ux442ux432ux43eux440ux435ux43dux43dux44f-ux432ux438ux43aux43eux440ux438ux441ux442ux43eux432ux443ux44eux442ux44c-ux43fux440ux438ux43dux446ux438ux43f-ux441ux443ux43fux435ux440ux43fux43eux437ux438ux446ux456ux457.}

Припустимо, що інженеру потрібно знайти відгук лінійної системи на
вхідний ефект f.~Загальний підхід до розв'язання задачі такий:

\begin{itemize}
\tightlist
\item
  \begin{enumerate}
  \def\labelenumi{\arabic{enumi})}
  \tightlist
  \item
    розкласти f на елементарні ефекти
    \texttt{f\ =\ \textbackslash{}sum\ f\_k} ;
  \end{enumerate}
\item
  \begin{enumerate}
  \def\labelenumi{\arabic{enumi})}
  \setcounter{enumi}{1}
  \tightlist
  \item
    знайти реакцію системи \texttt{u\_k} на вплив \texttt{f\_k} ;
  \end{enumerate}
\end{itemize}

\begin{enumerate}
\def\labelenumi{\arabic{enumi})}
\setcounter{enumi}{2}
\tightlist
\item
  сложить отклики
  \texttt{u\ =\ \textbackslash{}sum\_\{i=1\}\^{}\{n\}\ u\_\{i\}} .
\end{enumerate}

Якщо система лінійна, то отримана сума є реакцією системи на вхідний
ефект f.~Це принцип суперпоезії (див. стор. 26.1).

\pandocbounded{\includegraphics[keepaspectratio]{images/_page_189_Picture_10.jpeg}}

Фіг. 26.1. Основна ідея принципу суперпоемації; L --- диференціальний
оператор; f --- вхід; У---виходь.

Ця базова ідея може використовуватися для розв'язання складних змішаних
задач, розбиваючи їх на прості підзавдання, які розв'язуються за
допомогою

кожне з них, а потім підсумовування отриманих результатів (звісно,
диференціальне рівняння і граничні умови мають бути лінійними).

\section{Розкладання змішаної задачі на дві
простіші}\label{ux440ux43eux437ux43aux43bux430ux434ux430ux43dux43dux44f-ux437ux43cux456ux448ux430ux43dux43eux457-ux437ux430ux434ux430ux447ux456-ux43dux430-ux434ux432ux456-ux43fux440ux43eux441ux442ux456ux448ux456}

Предположим, нам необходимо решить линейную задачу (обозначим ее P)

(УЧП)

\begin{enumerate}
\def\labelenumi{\arabic{enumi}.}
\setcounter{enumi}{17}
\tightlist
\item
\end{enumerate}

(Р) (ГУ)

\begin{verbatim}
\begin{cases} 1-|x|, |x| < 1 \\ 0, |x| > 1 \end{cases}
\begin{cases} u(0, t) = 0, \\ u(1, t) = 0, \\ u(x, 0) = \sin(2\pi x), \quad 0 \le x \le 1. \end{cases}
\cos(ax)
\frac{2}{\pi} \frac{\sin a\omega}{\omega}
\end{verbatim}

Перед нами неоднородное уравнение теплопроводности, так что метод
разделения переменных неприменим. Можпо, конечно, воспользоваться
конечным синус-преобразованием по переменной x или преобразованием
Лапласа по переменной t, но вместо этого мы рассмотрим две подзадачи

\begin{verbatim}
\frac{2}{\pi} \frac{\sin a\omega}{(\omega^2+a^2)^2}
\frac{\pi}{2} e^{-a+\omega}
\begin{cases} u(0, t) = 0, \\ u(1, t) = 0, \\ u(1, t) = 0, \end{cases}
\frac{\pi}{2} e^{-a+\omega}
\frac{\pi}{2} e^{-a+\omega}
\frac{\pi}{2} e^{-a+\omega}
\frac{\pi}{2} e^{-a+\omega}
\end{verbatim}

\begin{verbatim}
Каждая из этих задач легко решается и в то же время ясно, что  \mathit{suma}  решений задач  P_1  и  P_2  является решением исходной вадачи, т. е.
\end{verbatim}

\begin{verbatim}
\frac{2ax}{(x^2+a^2)^2}
\end{verbatim}

Решение задачи \texttt{P\_\{i\}} Розв'язання задачі \texttt{P\_i}
Розв'язання задачі \texttt{P\_i} Можна показати, що в загальному випадку
розв'язок задачі

\begin{verbatim}
\frac{\cos(ax), |x| > \pi/2a}{0, |x| > \pi/2a}
\frac{1-|x|, |x| < 1}{0, |x| > 1}
\begin{cases} u_t = u_{xx} + f(x, t), \\ u(0, t) = 0, \\ u(1, t) = 0, \\ u(x, 0) = \varphi(x) \end{cases}
\cos(ax)
\cos(ax)
\frac{\pi}{2} [\delta(\omega+a) + \delta(\omega-a)]
\begin{cases} u_t = u_{xx} + f(x, t), \\ u(0, t) = 0, \\ u(1, t) = 0, \\ u(x, 0) = 0 \end{cases}
f(x) = \int_{0}^{\infty} F(\omega) \sin(\omega x) d\omega
F(\omega) = \frac{?}{\pi} \int_{0}^{\infty} f(x) \sin(\omega x) dx
0 < x < \infty
\begin{cases} u_t = u_{xx}, \\ u(0, t) = 0, \\ u(1, t) = 0, \\ u(x, 0) = \varphi(x). \end{cases}
0 < \omega < x
\end{verbatim}

f'\,'(0)

\begin{verbatim}
-\omega^2 F(\omega) + \frac{2}{\pi} \omega f(0)
\frac{1}{\alpha} F\left(\frac{\omega}{\alpha}\right)
\end{verbatim}

проследим, как оно решается методом конечного синус-преобразования. Вот
что мы должны сделать: разложить еходное воздействие \texttt{f(x,\ t)}
на компоненты, найти отклик \texttt{U\_n} на каждую компоненту, а затем
сложить эти отклики. Математически это может быть и не совсем очевидно,
но давайте внимательно проследим за ходом решения. Разложим входящие в
уравнение с частными производными члены в ряды Фурье по синусам:

\begin{verbatim}
\frac{2\omega}{\pi (\alpha^2 + \omega^2)}
x^{-1/2}
[2/\pi\omega]^{1/2}
\end{verbatim}

where

\begin{verbatim}
\frac{2}{\pi\omega} [1-\cos(\omega a)]
\end{verbatim}

1

\begin{verbatim}
e^{-1}a + \omega
\end{verbatim}

Зверніть увагу, по-перше, що коефіцієнти \texttt{A\_n} , \texttt{B\_n}
\texttt{F\_n} дійсно є функціями часу t, оскільки ми застосували
синусоїдальне перетворення до функції двох змінних x, t, u, а по-друге,
ми розклали вхідний ефект \texttt{f(x,t)} на прості компоненти
\texttt{F\_n}(t) . Тепер знайдіть відповіді \texttt{U\_n}(t) для кожного
компонента \texttt{F\_n}(t) , потім додайте всі відповіді
\texttt{U\_n}(t) і отримайте розв'язок \texttt{u(x,t)}.

Щоб знайти відповіді \texttt{U\_n}(t) , необхідно робити невеликі
розрахунки з коефіцієнтами \texttt{A\_n}(t) і \texttt{B\_n}(t) , щоб у
вирази інтегрування була включена лише функція u, а не її похідні
\texttt{u\_i} і \texttt{u\_\{xx\}} . Інтегруючи на частини, ми отримуємо

\begin{verbatim}
\frac{1}{2}e^{-\alpha}\sin\omega
\frac{1-e^{-a\omega}}{\omega}
\begin{split} A_n(t) &= 2\int_0^1 u_t \sin(n\pi x) \, dx = \frac{d}{dt} \bigg[ 2\int_0^1 u(x, t) \sin(n\pi x) \, dx \bigg] = \frac{dU_n(t)}{dt} \,, \\ B_n(t) &= 2\int_0^1 u_{xx} \sin(n\pi x) \, dx = -(n\pi)^2 \, U_n(t) \,+ \\ &\quad + 2n\pi \, \big[ u(0, t) + (-1)^{n+1} \, u(1, t) \big], \end{split}
\frac{2}{\pi}F''(\omega)
f(x) = \int_{0}^{\infty} F(\omega) \cos(\omega x) d\omega \qquad \qquad F(x) = \frac{2}{\pi} \int_{0}^{\infty} f(x) \cos(\omega x) dx
0 < x < \infty \qquad \qquad 0 < \omega < \infty
\begin{cases} u(0, t) = 0, \\ u(1, t) = 0 \end{cases}
\end{verbatim}

j'\,'(x)

\begin{verbatim}
S_n = \frac{2}{\pi} \int_0^{\pi} f(x) \sin(nx) dx.
\end{verbatim}

в последнее соотношение, получаем

\begin{verbatim}
y = \frac{\pi}{b - a} (x - a).
\end{verbatim}

и, следовательно, (26.1) принимает вид

\begin{verbatim}
f(x) = \sum_{n=1}^{\infty} S_n \sin(nx).
\end{verbatim}

\begin{verbatim}
Поскольку это разложение является тождеством по x, его коэффициенты должны быть равны нулю,  \tau  . е.
\end{verbatim}

\begin{verbatim}
f(x) = \sum_{n=1}^{\infty} S_n \sin(nx)
\end{verbatim}

.

Таким образом, мы получили связь между входом и выходом системы, между
коэффициентами \texttt{F\_n} и \texttt{U\_n} . Прежде чем решать
уравнения для \texttt{U\_n}(t) , обратимся к начальному условию

\begin{verbatim}
S_n = \frac{2}{\pi} \int_{0}^{\pi} f(x) \sin(nx) dx
\end{verbatim}

Если разложить \texttt{u(x,\ 0)} в ряд Фурье по синусам и приравнять его
нулю:

\begin{verbatim}
0 \le x \le \pi
\end{verbatim}

\begin{verbatim}
то определяются начальные условия для функции  U_n\left(t\right)  , т. е.
\end{verbatim}

\begin{verbatim}
n = 1, 2, ...
\end{verbatim}

\begin{verbatim}
,  n = 1,  2, \dots  Мы получили разложение нашей смешанной задачи на ряд простых задач об отклике системы на элементарные входные воздействия:
\end{verbatim}

(OДУ)

f'\,'(x)

(HУ) \texttt{U\_n}(0) = 0, n = 1, 2, \ldots{} Каждое из этих уравнений
легко решается методом интегрирующего множителя (или с помощью
преобразования Лапласа). В любом случае получаем

\begin{verbatim}
\frac{\pi}{2} e^{ax}
\end{verbatim}

Таким образом, мы получили отклики \texttt{U\_n}(t) на простейшие
входные воздействия \texttt{F\_n}(t) . Для получения решения исходной
задачи осталось совершить последний шаг---сложить все отклики

\begin{verbatim}
C_n = \frac{2}{\pi} \int_0^{\pi} f(x) \cos(nx) dx.
\end{verbatim}

\begin{verbatim}
Читатель, должно быть, заметил, что каждый отклик  U_n(t)  умножается на весовой множитель  \sin(n\pi x)  . Конечно, в разложении функции f(x, t) на компоненты  F_n(t)  присутствуют те же весовые множители.
\end{verbatim}

\section{ЗАМЕЧАНИЕ}\label{ux437ux430ux43cux435ux447ux430ux43dux438ux435}

Використання скінченного синусоїдального перетворення дозволяє записати
розв'язок у вигляді ряду, але для більшості інших інтегральних
перетворень розв'язки представлені як інтеграли (неперервний розклад).
Читача можна попросити дати інтерпретацію проблеми

\begin{verbatim}
y = \frac{\pi}{b-a}(x-a).
f(x) = \frac{C_0}{2} + \sum_{n=1}^{\infty} C_n \cos(nx).
f(x) = \frac{C_0}{2} + \sum_{n=1}^{\infty} C_n \cos(nx) \qquad C_n = \frac{2}{\pi} \int_{0}^{\pi} f(x) \cos nx \, dx
0 \le x \le \pi \qquad n = 0, 1, 2, \dots
\end{verbatim}

\begin{verbatim}
\begin{array}{ll} (\mathrm{Y} + \Pi) & u_{tt} = u_{xx} + \sin{(3\pi x)}, & 0 < x < 1, & 0 < t < \infty, \\ (\mathrm{\Gamma} \mathrm{Y}) & \left\{ \begin{array}{ll} u\left(0, \ t\right) = 0, \\ u\left(1, \ t\right) = 0, \\ \end{array}
\end{verbatim}

f'\,'(x)

\begin{verbatim}
-n^{2}C_{n} - \frac{2}{\pi} \{f'(0) - (-1)^{n} f'(n)\}
\begin{array}{ll} u\left(x, \ 0\right) = \sin{(\pi x)}, \\ u_{t}\left(x, \ 0\right) = 0, \end{array}
\end{verbatim}

Для каждой из подзадач выберите метод решения по своему BKYCY,

\begin{enumerate}
\def\labelenumi{\arabic{enumi}.}
\setcounter{enumi}{2}
\tightlist
\item
  Нехай \texttt{u\_1} і \texttt{u\_2} є розв'язками наступних рівнянь:
\end{enumerate}

\begin{enumerate}
\def\labelenumi{\alph{enumi})}
\tightlist
\item
\end{enumerate}

;\\
b) \texttt{u\_t\ =\ u\_\{xx\}\ +\ e\^{}t} ;

\begin{enumerate}
\def\labelenumi{\Alph{enumi})}
\setcounter{enumi}{1}
\tightlist
\item
\end{enumerate}

\begin{verbatim}
5. \cos(mx), m=1, 2, \dots
\end{verbatim}

;\\
F) \texttt{u\_t\ =\ u\_\{xx\}\ +\ u\^{}2} .

\begin{verbatim}
\mathbf{\Pi}  для яких рівнянь функція  u_1 + u_2  також буде розв'язком? Які висновки ви можете зробити зі своїх відповідей?
\end{verbatim}

\begin{enumerate}
\def\labelenumi{\arabic{enumi}.}
\setcounter{enumi}{3}
\tightlist
\item
  Найдите четыре смещанные задачи, сумма решений которых дает решение
  следующей задачи:
\end{enumerate}

\begin{verbatim}
\begin{aligned} u_t &= u_{xx} + f(x, t), & 0 < x < 1, & 0 < t < \infty, \\ & \begin{cases} u(0, t) = g_1(t), \\ u(1, t) = g_2(t), \end{cases} & 0 < t < \infty, \\ & u(x, 0) = \varphi(x), & 0 \le x \le 1. \end{aligned}
\end{verbatim}

\begin{enumerate}
\def\labelenumi{\arabic{enumi}.}
\setcounter{enumi}{4}
\tightlist
\item
  Решите задачу Коши
\end{enumerate}

(OAY)

x

(HY) \texttt{U\_n}(0) = 0. Можете ли вы проверить полученное решение?

\begin{enumerate}
\def\labelenumi{\arabic{enumi}.}
\setcounter{enumi}{5}
\tightlist
\item
  Предположим, что функции \texttt{u\_1} и \texttt{u\_2} удовлетворяют
  линейным однородным граничным условиям
\end{enumerate}

\begin{verbatim}
\begin{cases} u_x(1, t) + h_2 u(1, t) = 0.  Чи відповідає функція u_i + u_i цим умовам? , .
\end{verbatim}

Лекция 27

\section{РІВНЯННЯ ПЕРШОГО ПОРЯДКУ (ХАРАКТЕРИСТИЧНИЙ
МЕТОД)}\label{ux440ux456ux432ux43dux44fux43dux43dux44f-ux43fux435ux440ux448ux43eux433ux43e-ux43fux43eux440ux44fux434ux43aux443-ux445ux430ux440ux430ux43aux442ux435ux440ux438ux441ux442ux438ux447ux43dux438ux439-ux43cux435ux442ux43eux434}

МЕТАЛЕКЦІЇ: Введіть поняття рівняння з похідними в похідних першого
порядку (поки що ми розглядали рівняння другого порядку) і ознайомтеся з
важливим методом розв'язання задач із початковими умовами --- методом
характеристик. Проблема, яку ми зараз розглянемо, виглядає так:

(УЧП)

\_ \textbar, -\infty \textless{} x \textless{} \infty  , 0 \textless{} t
\textless{} \infty  ,

\begin{verbatim}
(HV)   u(x, 0) = \varphi(x), \quad -\infty < x < \infty.
\end{verbatim}

Слід одразу зазначити, що вперше ми розв'яжемо задачу з змінними
коефіцієнтами. Виявляється, що якщо ми рухаємося з координат (x, t) на
нові, характеристичні координати (s, \tau), то наша UFP перетвориться на
звичайне диференціальне рівняння. Тоді можна розв'язати звичайне
диференціальне рівняння, тобто знайти u(s, \tau) , і на останньому кроці
виразити s і t за x і t і отримати \texttt{u(x,\ t)}.

Напомним читателю, что, когда мы решали уравнение дифф

\begin{verbatim}
\Delta u = u_{xx} + u_{yy},   -\infty < x < \infty  ,   0 < t < \infty
\end{verbatim}

Константа \alpha\^{}2 виконувала роль коефіцієнта дифузії, а v ---
швидкість середовища. Очевидно, якщо \alpha=0 (тобто дифузії немає), то
розв'язок буде зсуватися як ціле число вздовж осі x зі швидкістю v
(залишається лише конвекція). Іншими словами, якщо початкова умова ---
\texttt{u(x,0)}=\varphi(x) , то відповідне розв'язання рівняння

\begin{verbatim}
\Delta u = u_{rr} + \frac{1}{r}u_r + \frac{1}{r^2}u_{\theta\theta}
\end{verbatim}

будет иметь вид \texttt{u(x,\ t)} = \varphi(x-vt) .

Це дає нам підстави вважати, що розв'язок рівняння першого порядку

\begin{verbatim}
\Delta u = u_{xx} + u_{yy} + u_{zz}
\end{verbatim}

переносится вдоль потока со скоростью

\begin{verbatim}
\Delta u = u_{rr} + \frac{1}{r}u_r + \frac{1}{r^2}u_{\theta\theta} + u_{zz}.
\end{verbatim}

Звісно, якщо a і b --- константи, то розв'язок --- це хвиля, що
рухається з постійною швидкістю. Якщо \texttt{a(x,\ t)} та
\texttt{b(x,\ t)}

\pandocbounded{\includegraphics[keepaspectratio]{images/_page_196_Figure_10.jpeg}}

Фіг. 27.1. Початкове значення в точці x впливає на розв'язок лише вздовж
характеристики у площині змінних (x,t). Параметр s змінюється вздовж
характеристик від нуля до нескінченності, починаючи з початкової лінії
умови; \tau --- сталий за кожною характеристикою; \tau змінюється вздовж
початкової лінії стану.

зависят от x и t, то скорость потока изменяется как вдоль потока, так и
во времени (читатель сможет увидеть, что начальная кривая сильно
искажается). Здесь так много аналогий с конвекцией!

Вернемся к нашей основной задаче:

(УЧП)

\begin{verbatim}
\Delta u = u_{rr} + \frac{2}{r}u_r + \frac{1}{r^2}u_{\theta\theta} + \frac{\cot\theta}{r^2}
\end{verbatim}

(НУ) \texttt{u(x,\ 0)} = \varphi(x), -\infty \textless{} x \textless{}
\infty.

Розв'язок цього рівняння першого порядку базується на наступному
фізичному факті: початкова умова в певній точці x переноситься на tx
площину вздовж прямої, що

называется характеристи-кой (см. рис. 27.1)

Таким чином, наше рівняння відрізняється від інших рівнянь (наприклад,
рівняння теплопровідності \texttt{u\_t\ =\ u\_\{xx\})},для яких
початкове значення в точці x впливає на розв'язок у всіх точках простору
і в будь-який час. Якщо читач пам'ятає, то початковий P I S

\pandocbounded{\includegraphics[keepaspectratio]{images/_page_197_Figure_5.jpeg}}

\texttt{.\ 27.2.\ Початкове\ збурення\ \textquotesingle{}u\textbackslash{}left(x,0\textbackslash{}right)\textquotesingle{}\ у\ точці\ породжує\ дві\ хвилі.\ Початкове\ збурення\ струни,\ згідно\ з\ хвильовим\ рівнянням,\ поширюється\ за\ двома\ характеристиками.}

Верхня струна в точці X впливає на розв'язок уздовж \часткових ліній syx
у площині TX (що відповідає двом хвилям, що рухаються) (див. рисунок
27.2).

Природно, ідея пропонує ввести дві нові координати s s і \tau (замість x
і t), щоб

\begin{itemize}
\tightlist
\item
  s змінювалися вздовж характеристичних кривих,
\item
  t змінюється вздовж початкової кривої (бажано вздовж лінії t=0).
\end{itemize}

Спочатку розглянемо нову координату s. Якщо s вибрано так, щоб
задовольнити наведену умову, то рівняння

\begin{verbatim}
\Delta u = u_{rr} + \frac{2}{r}u_r + \frac{1}{r^2}u_{\theta\theta} + \frac{\cot\theta}{r^2}u_0 + \frac{1}{r^2\sin\theta}u_{\phi\phi}
\end{verbatim}

превратится в обыкновенное дифференциальное уравнение К о н е О д н

\pandocbounded{\includegraphics[keepaspectratio]{images/_page_370_Figure_1.jpeg}}

ак залишається питання, як знайти ці характеристики. Відповідь проста:
ми оберемо характеристики \{{[}\texttt{x(s)}, \texttt{t(s)}{]}: 0
\textless{} s \textless{} \infty\} так, щоб вони задовольняли систему

\pandocbounded{\includegraphics[keepaspectratio]{images/_page_370_Picture_2.jpeg}}

\pandocbounded{\includegraphics[keepaspectratio]{images/_page_370_Figure_3.jpeg}}

\begin{verbatim}
\begin{cases} \frac{dx}{ds} = a(x, t), \\ \frac{dt}{ds} = b(x, t). \end{cases}
-\infty < x < +\infty
-\infty < y < +\infty
\end{verbatim}

Ясно, что в этом случае

\begin{verbatim}
-\infty < z < +\infty
\end{verbatim}

\begin{verbatim}
Другими словами, на характеристиках  \{[x(s), t(s)]: 0 < s < \infty\}  уравнение с частными производными становится обыкновенным дифференциальным уравнением. Поясним эти идеи на конкретном примере.
\end{verbatim}

\section{ПРИМЕР}\label{ux43fux440ux438ux43cux435ux440}

Припустимо, нам потрібно розв'язати наступну задачу Коші з сталими
коефіцієнтами:

\begin{verbatim}
\begin{array}{lll}
\end{verbatim}

\begin{verbatim}
\begin{array}{llll} (\forall \Pi \Pi) & u_x + u_t + 2u = 0, & -\infty < x < \infty, & 0 < t < \infty \\ (\forall \Psi) & u(x, 0) = \sin x, & -\infty < x < \infty. \end{array}
-\infty & < x < +\infty & r \geqslant 0, \ 0 \leqslant \theta \leqslant 2\pi, \\
\frac{dx}{ds} = 1, \quad \frac{dt}{ds} = 1, \quad 0 < s < \infty.
-\infty & < y < +\infty & -\alpha \leqslant z \leqslant \alpha, \\
x(s) = s + c_1, \quad t(s) = s + c_2.
x(0) = \tau,
-\infty & < z < +\infty & x = r \cos \theta, \\
x = s + \tau
x-t=\tau
y = r \sin \theta, \\
\frac{du}{ds} + 2u = 0, \quad 0 < s < \infty,
u(0) = \sin \tau.
z = z.
u\left(s,\,\tau\right)=\sin\tau e^{-2s}
\end{verbatim}

\begin{verbatim}
\end{array}
\end{verbatim}

\begin{verbatim}
x = s + \tau
s=t
r \geqslant 0
u(x, t) = \sin(x-t)e^{-xt}
\Delta u = 0
(\forall \Pi) \quad a(x, t) u_x + b(x, t) u_t + c(x, t) u = 0, \quad -\infty < x < \infty, \\ 0 < t < \infty,
\Delta u + \lambda^2 u = 0
u(x, 0) = f(x), \quad -\infty < x < \infty.
\Delta u = k
\frac{dx}{ds} = a(x, t), \qquad \frac{dt}{ds} = b(x, t).
\Delta u + k(E - V)u = 0
x = x (s, \tau),
t = t (s, \tau).
\frac{du}{ds} + c[x(s, \tau), t(s, \tau)]u = 0, 0 < s < \infty,
(\text{УЧП})
\frac{dx}{ds} = x
\frac{dt}{ds} = 1
x = \tau e^s,
u_{tt} = c^2 \Delta u
t = s.
u_{tt} = c^2 \Delta u - hu_t
\frac{du}{ds} + su = 0, \quad 0 < s < \infty,
u_t = \alpha^2 u_{xx}
u(0) = F(\tau).
u_t = \alpha^2 u_{xx} - hu_x
u(s, \tau) = F(\tau) e^{-s} \ell^2
u_t = \alpha^2 u_{xx} - ku
s=t
u_t = \alpha^2 u_{xx} + f(x, t)
u(x, t) = F(xe^{-t}) e^{-t^2/2}
\end{verbatim}

Уравнение одномерной диффузии Уравнение конвективной диффузии

Уравнение теплопроводности с поглощением Уравнение теплопроводности с
источником

\section{НЕЛИНЕЙНЫЕ УРАВНЕНИЯ ПЕРВОГО ПОРЯДКА (ЗАКОНЫ
СОХРАНЕНИЯ)}\label{ux43dux435ux43bux438ux43dux435ux439ux43dux44bux435-ux443ux440ux430ux432ux43dux435ux43dux438ux44f-ux43fux435ux440ux432ux43eux433ux43e-ux43fux43eux440ux44fux434ux43aux430-ux437ux430ux43aux43eux43dux44b-ux441ux43eux445ux440ux430ux43dux435ux43dux438ux44f}

МЕТА ЛЕКЦІЇ: Ввести понятие нелинейных уравнений первого порядка и
показать, как некоторые из них (так называемые законы сохранения) можно
использовать для описания физических явлений. Например, закон сохранения

\begin{verbatim}
u_i + g(u)u_x = 0
\end{verbatim}

с начальным условием

\begin{verbatim}
u(x, 0) = \varphi(x)
\end{verbatim}

используется для описания потока автомобилей на автостраде. Этот пример
показывает, что уравнения с частными производными применимы не только в
физике, биологии или технике. Более того, это нелинейное уравнение имеет
разрывные решения, которые описывают ударные волны, распространяющиеся в
потоке автомобилей вдоль дороги.

Одно из наиболее важных уравнений в частных производных --- это закон
сохранения

\begin{verbatim}
u_t + f_x = 0.
\end{verbatim}

Этот закон утверждает, что увеличение некоторой физической величины
\texttt{u\_t} равно изменению потока--- \texttt{f\_x} этой величины
через поперечное сечение (поток, направленный слева направо, считается
положительным). В динамике жидкости величина \texttt{u(x,t)} может
обозначать плотность жидкости в точке x, а величина
\texttt{f(x,t)}---поток жидкости (т. е. количество жидкости, протекающей
через поперечное сечение в точке x в момент времени t). Мы не станем в
этой лекции рассматривать динамику невязкой несжимаемой жидкости. Вместо
этого мы покажем, как закон сохранения можно использовать для прогноза
динамики дорожного движения (нас будет интересовать поток автомобилей, а
не поток молекул воды).

Сначала займемся выводом закона сохранения.

\section{Вывод закона
сохранения}\label{ux432ux44bux432ux43eux434-ux437ux430ux43aux43eux43dux430-ux441ux43eux445ux440ux430ux43dux435ux43dux438ux44f}

Предположим, что автомобили движутся слева направо по скоростной
автостраде, у которой нет боковых съездов и въездов. Обозначим

\texttt{u(x,\ t)} --- плотность автомобилей в точке x (число автомобилей
на единицу длины),

\texttt{f(x,\ t)} --- ноток автомобилей в точке x (число автомобилей,
проходящих через точку x в одну минуту).

Совершенно очевидно, что на участке дороги {[}a, b{]} изменение числа
автомобилей (за единицу времени) дается следующими двумя выражениями:

Изменение числа автомобилей на
\texttt{{[}a,\ b{]}\ =\ \textbackslash{}frac\{d\}\{dt\}\ \textbackslash{}int\_a\^{}b\ u(x,\ t)\ dx}

И

Изменение числа автомобилей на
\texttt{{[}a,b{]}\ =\ f(a,t)\ -\ f(b,t)\ =\ -\textbackslash{}int\_\{0\}\^{}\{b\}\ \textbackslash{}frac\{\textbackslash{}partial\ f\}\{\textbackslash{}partial\ x\}(x,t)dx}
.

Последнее соотношение --- следствие основной теоремы интегрального
исчисления. Приравнивая эти два интеграла, получаем

\begin{verbatim}
\int_{a}^{b} \frac{\partial u}{\partial t}(x, t) dx = -\int_{a}^{b} \frac{\partial f}{\partial x}(x, t) dx.
\end{verbatim}

Поскольку промежуток {[}a,b{]} произволен, делаем вывод, что
подынтегральные функции равны, т е. выполняется закон сохранения

\begin{verbatim}
u_t + f_x = 0.
\end{verbatim}

\section{Применение закона сохранения к задаче о дорожном
движении}\label{ux43fux440ux438ux43cux435ux43dux435ux43dux438ux435-ux437ux430ux43aux43eux43dux430-ux441ux43eux445ux440ux430ux43dux435ux43dux438ux44f-ux43a-ux437ux430ux434ux430ux447ux435-ux43e-ux434ux43eux440ux43eux436ux43dux43eux43c-ux434ux432ux438ux436ux435ux43dux438ux438}

Уравнение \texttt{u\_t\ +\ f\_x\ =\ 0} по форме очень простое. Правда, в
нем две неизвестные функции! Вопрос в том, как воспользоваться этим
уравнением и что это нам даст5

В задачах дорожного движения пользуются экспериментально найденной
зависимостью потока автомобилей \texttt{f(u)} от плотности автомобилей в
данной точке u. Представляется очевидным, что если плотность u растет,
то растет и поток f (по крайней мере в точке). Типичная модель дорожного
движения задается формулой (см. рис. 28.1)

\begin{verbatim}
f(u) = Au (1-u).
\end{verbatim}

Используются и другие модели:

\begin{itemize}
\tightlist
\item
  \begin{enumerate}
  \def\labelenumi{\arabic{enumi})}
  \tightlist
  \item
    \texttt{f(u)} = ku (линейный рост потока),
  \end{enumerate}
\item
  \begin{enumerate}
  \def\labelenumi{\arabic{enumi})}
  \setcounter{enumi}{1}
  \tightlist
  \item
    \texttt{f(u)\ =\ u\^{}2} (квадратичный рост потока).
  \end{enumerate}
\end{itemize}

\pandocbounded{\includegraphics[keepaspectratio]{images/_page_205_Figure_5.jpeg}}

Рис. 28 1. Типичная зависимость потока от плотности.

Теперь мы подготовлены к тому, чтобы применить уравнение
\texttt{u\_i+f\_x=0} к задаче дорожного движения. По правилу
дифференцирования сложной функции
\texttt{f\_x=\textbackslash{}frac\{df\}\{du\}\textbackslash{},u\_x} ;
значит, закон сохранения можно переписать в виде

\begin{verbatim}
u_t + \frac{df}{du} u_x = u_t + g(u) u_x = 0.
\end{verbatim}

Пусть, например, зависимость потока от плотности квадратичная, т. е.
\texttt{f(u)\ =\ u\^{}\{2\}} , тогда закон сохранения принимает
следующую форму:

\begin{verbatim}
u_t + 2uu_x = 0.
\end{verbatim}

Если начальная плотность распределения автомобилей была
\texttt{u(x,\ 0)\ =\ \textbackslash{}varphi(x)} , то, чтобы найти
плотность распределения в произвольный момент времени, необходимо решить
задачу с начальным условием (задачу Коши):

\begin{verbatim}
\begin{array}{llll} (\text{УЧП}) & u_t + 2uu_x = 0, & -\infty < x < \infty, & 0 < t < \infty, \\ (\text{HV}) & u\left(x, \ 0\right) = \varphi\left(x\right), & -\infty < x < \infty. \end{array}
\end{verbatim}

Займемся теперь решением нелинейной задачи Коши.

\section{Нелинейная задача с начальным условием (задача
Коши)}\label{ux43dux435ux43bux438ux43dux435ux439ux43dux430ux44f-ux437ux430ux434ux430ux447ux430-ux441-ux43dux430ux447ux430ux43bux44cux43dux44bux43c-ux443ux441ux43bux43eux432ux438ux435ux43c-ux437ux430ux434ux430ux447ux430-ux43aux43eux448ux438}

Рассмотрим задачу

\begin{verbatim}
\begin{array}{llll} (\mathsf{Y}\mathsf{P}\Pi) & u_t + g\left(u\right)u_x = 0, & -\infty < x < \infty, & 0 < t < \infty, \\ (\mathsf{H}\mathsf{Y}) & u\left(x, \ 0\right) = \varphi\left(x\right), & -\infty < x < \infty. \end{array}
\end{verbatim}

Вспомним предыдущую лекцию, где мы изучали уравнение конвективного
переноса:

\begin{verbatim}
u_t + vu_x = 0
\end{verbatim}

.

У цьому рівнянні \texttt{u(x,\ t)} --- це концентрація матерії в потоку,
що рухається зі швидкістю v

\begin{verbatim}
u_1 + g(u)u_x = 0
\end{verbatim}

и уравнением переноса. Будем считать, что частина воды в гочке
\texttt{x\_0} движется со скоростью \texttt{g(u)} (вдоль по течению или
против него). Тогда через t секунд координата частины будет

\begin{verbatim}
x = x_0 + g(u) t
\end{verbatim}

(уравнение характеристик).

Пам'ятайте, що концентрація \texttt{u\ (x,\ t)} не змінюється вздовж
характеристики. Якщо ми знаємо початкову концентрацію
\texttt{u(x\_a,\ 0)} ,

\pandocbounded{\includegraphics[keepaspectratio]{images/_page_206_Figure_8.jpeg}}

Рис. 28.2. Характеристики уравнения \texttt{u\_t+g} (u) \texttt{u\_x=0}
; a---общий вид характеристики, выходящей из т. ( \texttt{x\_0} , 0),
вдоль характеристики решение остается постоянным и равным величине
\texttt{u\ (} x\_0 \texttt{,\ 0)}; b---пример характеристики,
начинающейся из т. (2, 0), уравнения \texttt{u\_t+3uu\_x=0} с
\texttt{u\ (2,\ 0)}=1.

то уравнение характеристик принимает вид
\texttt{x\ =\ x\_0\ +\ g{[}u(x\_0,\ 0){]}\ t} (уравнение характеристики,
начинающейся в точке \texttt{(x\_0,\ 0)} ).

Пусть, например, мы решаем задачу

(УЧП)

\begin{verbatim}
u_t + 3uu_x = 0
\end{verbatim}

,\\
(НУ) u(x, 0) =

\begin{verbatim}
u_t = Du_{xx} - Vu_x.
\end{verbatim}

Чтобы найти характеристику, выходящую из точки, например (2, 0), мы
должны написать

\begin{verbatim}
x = 2 + g[u(2, 0)]t =
\end{verbatim}

= \texttt{2\ +\ g{[}1{]}t\ =} = \texttt{2\ +\ 3t} .

Можна стверджувати, що розв'язок нашої задачі \texttt{u(x,\ t)} дорівнює
одному на лінін x=2+3t, тобто \texttt{u(2+3t,\ t)}=1 для всіх
t\textgreater0. Фіг. 28.2 показує характеристики загального випадку та
цього конкретного завдання.

\texttt{Тепер\ зрозуміло,\ що,\ знаючи\ характеристики\ в\ кожній\ точці\ і\ знаючи,\ що\ розв\textquotesingle{}язок\ за\ характеристиками\ не\ змінюється,\ можна\ отримати\ розв\textquotesingle{}язок\ \textquotesingle{}u\textbackslash{}left(x,\textbackslash{}\ t\textbackslash{}right)\textquotesingle{}\ у\ будь-який\ момент\ часу\ t.\ Ми\ не\ отримали\ явного\ виразу\ для\ \textquotesingle{}u\textbackslash{}left(x,\textbackslash{}\ t\textbackslash{}right)\textquotesingle{}\ через\ змінні\ x\ і\ t,\ але\ знання\ характеристик\ рівняння\ дозволяє\ розв\textquotesingle{}язати\ деякі\ цікаві\ задачі.}

\section{Задача дорожного
движения}\label{ux437ux430ux434ux430ux447ux430-ux434ux43eux440ux43eux436ux43dux43eux433ux43e-ux434ux432ux438ux436ux435ux43dux438ux44f}

Займемся моделью дорожного движения, в которой поток вычисляется по
формуле \texttt{f(u)=u\^{}2} , а начальная плотность изо- Р И С

\pandocbounded{\includegraphics[keepaspectratio]{images/_page_207_Figure_8.jpeg}}

. 28.3. Начальная плотность автомобилей движущихся слева направо.

бражена на рис. 28.3. Другими словами, нам нужно решить задачу

\begin{verbatim}
(\text{YUII}) \quad u_t + 2uu_x = 0, \quad -\infty < x < \infty, \quad 0 < t < \infty,
(\text{HV}) \quad u(x, \ 0) =
\begin{cases} 1, & x \leq 0, \\ 1 - x, & 0 < x < 1, \\ 0, & 1 \leq x. \end{cases}
\end{verbatim}

Начнем с определения характеристик, выходящих из точек вида
\texttt{(x\_0,\ 0)} . Для \texttt{x\_0\ \textless{}\ 0} получаем

\begin{verbatim}
x = x_0 + g [u(x_0, 0)] t =
\end{verbatim}

= \texttt{x\_0\ +\ g\ {[}1{]}\ t\ =} = \texttt{x\_0\ +\ 2t} .

Разрешая относительно t, находим уравнение характеристик

\begin{verbatim}
t = \frac{1}{2} (x - x_0).
\end{verbatim}

Эти прямые линии изображены на рис. 28.4. Теперь рассмотрим точки
\texttt{0\ \textless{}\ x\_0\ \textless{}\ 1} . Здесь характеристики
определяются следующими соотношениями:

\begin{verbatim}
x = x_0 + g[u(x_0, 0)] t =
= x_0 + g[(1 - x_0)] t =
= x_0 + 2(1 - x_0) t.
\end{verbatim}

Разремая последнее соотношение относительно t, получаем

\begin{verbatim}
t = \frac{x - x_0}{2(1 - x_0)}.
\end{verbatim}

При \texttt{1\ ≤\ x\_0\ \textless{}\ ∞} характеристики

\begin{verbatim}
x = x_0 + g[u(x_0, 0)]t =
\end{verbatim}

= \texttt{x\_0\ +\ g{[}0{]}t\ =} = \texttt{x\_0} приводят к вертикальным
линиям, выходящим из точек хо. Все характеристики задачи изображены на
рис. 28.4.

\pandocbounded{\includegraphics[keepaspectratio]{images/_page_208_Figure_10.jpeg}}

Рис. 28.4. Характеристики уравнения \texttt{u\_t\ +\ 2uu\_x\ =\ 0} .

Ситуація на дорозі при 0 \textless{} t \textless{} 1/2 досить зрозуміла.
Рис. 28.5 показує розв'язок нашої задачі в різні моменти часу. Зверніть
увагу, що характеристики сходяться в одній точці t=1/2. Отже, щоб
розв'язати задачу на t\textgreater1/2, потрібно використовувати інший
метод. Коли характеристики сходяться в одній точці, говоримо про ударну
хвилю (розрив розв'язку). Тепер потрібно відповісти на питання: на якій
швидкості є передній край

ударной волны будет перемещаться вдоль автострады? Хотя это и не
очевидно, но оказывается, что скорость распространения разрыва находится
по формуле

\begin{verbatim}
S = \frac{f(u_R) - f(u_L)}{u_R - u_L},
\end{verbatim}

где \texttt{u\_R} и \texttt{u\_L} ---значения решения справа и слева от
волнового фронта, а \texttt{f(u\_R)} и \texttt{f(u\_L)} ---значения
потока при этих плотностях.

\pandocbounded{\includegraphics[keepaspectratio]{images/_page_209_Figure_5.jpeg}}

Рис. 28.5. Плотность дорожного движения в различные моменты времени.

У нашому прикладі швидкість поширення щілини виявляється рівною одиниці.
Насправді, якщо згадати цей \texttt{f(u)\ =\ u\^{}2} , то одразу
отримуємо

\begin{verbatim}
S = \frac{0-1}{0-1} = 1
\end{verbatim}

.

Это означает, что при t\textgreater1/2 фронт волны будет перемещаться
слева направо со скоростью единица.

Полное решение задачи

\begin{verbatim}
(\text{Y} \, \Pi) \quad u_t + 2uu_x = 0, \quad -\infty < x < \infty,
(\text{HY}) \quad u(x, \ 0) =
\begin{cases} 1, & x \leq 0, \\ 1 - x, & 0 < x < 1, \\ 0, & 1 < x, \end{cases}
\end{verbatim}

изображено на рис. 28.5. \# ЗАУВАЖЕННЯ 1. Ударна хвиля в нашому прикладі
виникає через те, що потік швидко зростає зі збільшенням густини u. Якщо
потік заданий формулою \texttt{f(u)} = u, то рівняння було б у вигляді
\texttt{u\_1+u\_x=0} , і розв'язок цього рівняння представляє хвилю, що
рухається вправо без спотворень. Якщо на мить замислитися над значенням
формули \texttt{f(u)} = u, стає очевидно, що розв'язок рухатиметься
таким чином.

\begin{enumerate}
\def\labelenumi{\arabic{enumi}.}
\setcounter{enumi}{1}
\tightlist
\item
  Прямой подстановкой решения в виде
\end{enumerate}

\begin{verbatim}
u = \varphi(x - g(u) t)
\end{verbatim}

можно убедиться в том, что эта формула неявно определяет решение задачи

\begin{verbatim}
\begin{array}{lll} u_t + g\left(u\right) u_x = 0, & -\infty < x < \infty, & 0 < t < \infty, \\ u\left(x, \ 0\right) = \varphi\left(x\right), & -\infty < x < \infty. \end{array}
\end{verbatim}

Например, решение задачи Коши

\begin{verbatim}
(Y\Pi\Pi) \quad u_t + uu_x = 0,
(HY) \quad u(x, 0) = x
\end{verbatim}

неявно вадается соотношением

\begin{verbatim}
u = \varphi [x - g(u) t] =
= x - g(u) t = x - ut.
\end{verbatim}

У нашому конкретному випадку можливо отримати рішення у явній формі (у
багатьох інших випадках це неможливо)

\begin{verbatim}
u(x, t) = \frac{x}{1+t}.
\end{verbatim}

Проверьте, что эта функция действительно является решением нашей задачи.

\section{\texorpdfstring{\textbf{ЗАДАЧИ}}{ЗАДАЧИ}}\label{ux437ux430ux434ux430ux447ux438}

\begin{enumerate}
\def\labelenumi{\arabic{enumi}.}
\tightlist
\item
  Решите задачу Коши
\end{enumerate}

\begin{verbatim}
\begin{array}{llll} (\text{YUII}) & u_t + uu_x = 0, & -\infty < x < \infty, & 0 < t < \infty, \\ (\text{HV}) & u(x, \ 0) = \begin{cases} 0, & x < 0, \\ x, & 0 \leqslant x. \end{cases} \end{array}
\end{verbatim}

Постройте график решения для различных моментов времени. Как вы
интерпретируете это решение? Как выглядит соотношение между потоком и
плотностью в этой задаче?

\begin{enumerate}
\def\labelenumi{\arabic{enumi}.}
\setcounter{enumi}{1}
\tightlist
\item
  Решите задачу
\end{enumerate}

\begin{verbatim}
\begin{split} & (\text{УЧП}) \quad u_t + u^2 u_x = 0, \quad -\infty < x < \infty, \quad 0 < t < \infty, \\ & (\text{НУ}) \quad u_t(x, 0) = \begin{cases} 0, & x < 0, \\ x, & 0 \leqslant x. \end{cases} \end{split}
\end{verbatim}

Какова здесь связь между потоком и плотностью? Совнадает ли поведение
полученного решения с тем, что вы ожидали? Сравните решения задачи 1 и
задачи 2.

\begin{enumerate}
\def\labelenumi{\arabic{enumi}.}
\setcounter{enumi}{2}
\tightlist
\item
  Предположим, что невязкая жидкость течет по трубе и при этом
  просачивается через стенки по закону \texttt{F(u)} = ku (г/см·с).
  Закон сохранения (скорее закон изменения) в этом случае записывается в
  виде
\end{enumerate}

\begin{verbatim}
u_t + f_x = -F(u)
\end{verbatim}

.

Знайдіть розв'язок цього рівняння, якщо \texttt{f(u)} = u і
\texttt{\textbackslash{}varphi(x)\ =\ u} (x, 0) (початковий розподіл
задається загальним чином). Дайте тлумачення отриманого рішення.

\begin{enumerate}
\def\labelenumi{\arabic{enumi}.}
\setcounter{enumi}{3}
\item
  Знайти розв'язок попереднього рівняння, якщо втрати визначаються
  функцією \texttt{F(x,\ t)} = 1/x. Перевірте розв'язок. Ви розумієте
  його фізичне значення?
\item
  Проверьте, что функция и, неявно заданная соотношением
\end{enumerate}

\begin{verbatim}
u = \varphi[x - g(u) t],
\end{verbatim}

является решением нелинейной задачи

\begin{verbatim}
(\text{УЧП}) \quad \begin{array}{ll} u_1 + g(u) u_x = 0, \\ (\text{НУ}) & u(x, 0) = \varphi(x). \end{array}
\end{verbatim}

\# СИСТЕМЫ УРАВНЕНИЙ С ЧАСТНЫМИ ПРОИЗВОДНЫМИ МЕТА ЛЕКЦІЇ: Познакомить
читателя с тем фактом, что многие физические системы (особенно в
динамике жидкости) нельзя описать одним уравнением. Для описания таких
систем требуются системы связанных уравнений. В эти системы входят
некоторые неизвестные функции, такие, как давление
\texttt{p(x,\ y,\ z,\ t)}, плотность
\texttt{\textbackslash{}rho(x,\ y,\ z,\ t)} , температура
\texttt{T(x,\ y,\ z,\ t)} и их частные производные. Связи между этими
функциями определяются физическими законами и наша задача --- найти все
эти функции одновременно.

В лекции будет показано, как можно решить линейную систему уравнений

\begin{verbatim}
u_t + Au_x = 0
\end{verbatim}

если преобразовать ее в новую систему независимых уравнений

\begin{verbatim}
v_1 + \Lambda v_x = 0,
\end{verbatim}

а затем независимо решить каждое уравнение этой системы.

У багатьох галузях науки існують системи величин, які взаємопов'язані не
одним, а кількома відношеннями одночасно. Наприклад, у гідродинамікі це
чотири рівняння

\begin{verbatim}
u_t + uu_x + vu_y + \frac{1}{\rho} p_x = 0
\end{verbatim}

\texttt{(сохранение\ импульса\ вдоль\ оси\ \textquotesingle{}x\textquotesingle{}\ ),\ (29.1)\ \textquotesingle{}v\_t\ +\ uv\_x\ +\ vv\_y\ +\ \textbackslash{}frac\{1\}\{\textbackslash{}rho\}\ p\_y\ =\ 0\textquotesingle{}\ (сохранение\ импульса\ вдоль\ оси\ \textquotesingle{}y\textquotesingle{}\ ),\ (сохранение\ массы),\ (сохранение\ массы),\ (сохранение\ массы),\ \textquotesingle{}\textbackslash{}left(\textbackslash{}frac\{p\}\{\textbackslash{}rho\^{}\{\textbackslash{}gamma\}\}\textbackslash{}right)\_t\ +\ u\textbackslash{}left(\textbackslash{}frac\{p\}\{\textbackslash{}rho\^{}\{\textbackslash{}gamma\}\}\textbackslash{}right)\_x\ +\ v\textbackslash{}left(\textbackslash{}frac\{p\}\{\textbackslash{}rho\^{}\{\textbackslash{}gamma\}\}\textbackslash{}right)\_y\ =\ 0\textquotesingle{}\ (сохранение\ энергии),\ \textquotesingle{}p(x,\ y,\ z,\ t)\textquotesingle{}\ —\ давление\ в\ жидкости,\ \textquotesingle{}u(x,\ y,\ z,\ t)\textquotesingle{}\ —\ \textquotesingle{}x\textquotesingle{}\ -компонента\ скорости,\ \textquotesingle{}v(x,\ y,\ z,\ t)\textquotesingle{}\ —\ \textquotesingle{}y\textquotesingle{}\ -компонента\ скорости,\ \textquotesingle{}\textbackslash{}rho(x,\ y,\ z,\ t)\textquotesingle{}\ —\ плотность\ жидкости.}

Приведенная нелинейная система уравнений называется уравнениями Эйлера
движения жидкости. Задача состоит в одновременном отыскании всех
неизвестных функций р, и, v и р, которые удовлетворяли бы всем четырем
уравнениям (вместе с начальными и граничными условиями).

Існують й інші причини вивчати системи рівнянь. У теорії звичайних
диференціальних рівнянь (якщо читач пам'ятає) показано, що одне рівняння
другого порядку можна представити як систему з двох рівнянь першого
порядку. Хоча теорія диференціальних рівнянь у похідних не така проста,
дуже часто одне рівняння з похідними у похідних порядків зводиться до
системи диференціальних рівнянь першого порядку. Наприклад, телеграфне
рівняння

\begin{verbatim}
u_{tt} = c^2 u_{xx} + a u_t + b u
\end{verbatim}

путем введения трех новых переменных

\begin{verbatim}
u_1 = u_1
u_2 = u_2
u_3 = u_4
\end{verbatim}

можно свести к системе уравнений первого порядка

(29.2)

\begin{verbatim}
\begin{cases} \frac{\partial u_1}{\partial x} = u_2, \\ \frac{\partial a_1}{\partial t} = u_3, \\ \frac{\partial u_3}{\partial t} = c^2 \frac{\partial u_2}{\partial x} + au_3 + bu_4. \end{cases}
\end{verbatim}

Для того чтобы решать системы типа (29.2), необходимо хорошо знать такие
понятия, как собственные значения и собственные векторы матрицы. По эфой
причине мы предлагаем вниманию читателя краткий обзор линейной алгебры.

\section{Обзор линейной
алгебры}\label{ux43eux431ux437ux43eux440-ux43bux438ux43dux435ux439ux43dux43eux439-ux430ux43bux433ux435ux431ux440ux44b}

Общий случай

Частный случай

Определение. Если А---квадратная матрица порядка n, то собственными
значениями матрицы называются n корней характе-

Пусть

\begin{verbatim}
A =
\begin{bmatrix} 1 & 1 \\ 4 & 1 \end{bmatrix}
\end{verbatim}

.

ристического уравнения

\begin{verbatim}
\det\left(A-\lambda I\right)=0,
\end{verbatim}

\begin{verbatim}
где  \det(A-\lambda I)  — определитель матрицы  A-\lambda I  .
\end{verbatim}

Деякі власні значення можуть збігатися.

Тогда собственные значения являются корнями уравнения

\begin{verbatim}
\det
\begin{bmatrix} 1-\lambda & 1\\ 4 & 1-\lambda \end{bmatrix}
= (1-\lambda)^2 - 4 =
= \lambda^2 - 2\lambda - 3 = 0.
\end{verbatim}

Откуда

\begin{verbatim}
\lambda_1 = -1
\end{verbatim}

\texttt{,\ \textquotesingle{}\textbackslash{}lambda\_2\ =\ 3\textquotesingle{}\ .}

Визначення. Якщо \texttt{\textbackslash{}lambda} --- власне значення
матриці A, то відповідний власний вектор є таким ненульовим вектором, що
задовольняє відношення

\begin{verbatim}
Ax = \lambda x
\end{verbatim}

.

Пример. Вектор X =

\begin{verbatim}
u_t = -Vu_x
\end{verbatim}

\texttt{является\ собственным\ вектором,\ соответствующим\ собственному\ значению}
\lambda\_2 = 3 , так как

\begin{verbatim}
\begin{bmatrix} 1 & 1 \\ 4 & 1 \end{bmatrix}
\begin{bmatrix} 1 \\ 2 \end{bmatrix}
= 3
\begin{bmatrix} 1 \\ 2 \end{bmatrix}
\end{verbatim}

.

Определение. Матрица \texttt{A\^{}\{-1\}} называется обратной к матрице
A, если

\begin{verbatim}
AA^{-1} = A^{-1}A = I_1
\end{verbatim}

Пример, Обратной к

\begin{verbatim}
A =
\begin{bmatrix} 2 & -2 \\ 1 & 1 \end{bmatrix}
\end{verbatim}

является матрица

\begin{verbatim}
A^{-1} =
\begin{bmatrix} 0.25 & 0.5 \\ -0.25 & 0.5 \end{bmatrix}
\end{verbatim}

,

де I - одинична матриця.

поскольку

\begin{verbatim}
AA^{-1} = A^{-1}A =
\begin{bmatrix} 1 & 0 \\ 0 & 1 \end{bmatrix}
\end{verbatim}

.

\section{Діагональна
матриця}\label{ux434ux456ux430ux433ux43eux43dux430ux43bux44cux43dux430-ux43cux430ux442ux440ux438ux446ux44f}

\begin{verbatim}
Нехай A — квадратна матриця порядку n і нехай всі її власні значення  \lambda_1  ,  \lambda_2  ...  \lambda_n  різні. Побудуємо матрицю P так, щоб у неї були *стовпці* координати власних векторів (k-й стовпець містить власний вектор  X_k  , що відповідає власному значенню  \lambda_k  ), тобто
\end{verbatim}

\begin{verbatim}
P =
\begin{bmatrix} X_1 & X_{\tilde{2}} & \dots & X_n \\ X_{1} & X_{2} & \dots & X_n \end{bmatrix}
=
\begin{bmatrix} x_{11} & x_{1\tilde{2}} & \dots & x_{1n} \\ x_{21} & x_{22} & \dots & x_{2n} \\ \vdots & \vdots & \ddots & \vdots \\ x_{n1} & x_{n\tilde{3}} & \dots & x_{nn} \end{bmatrix}
\end{verbatim}

.

Тогда матрица Л, определяемая соотношением

\begin{verbatim}
\Lambda = P^{-1}AP,
\end{verbatim}

где \texttt{P\^{}\{-1\}} --- обратная к матрице P, будет диагональной
матрицей вида

\begin{verbatim}
\Lambda =
\begin{bmatrix} \lambda_1 & & & \\ & \lambda_2 & 0 & \\ & & \ddots & \\ & & & \lambda_n \end{bmatrix}
\end{verbatim}

.

Например, у матрицы

\begin{verbatim}
A =
\begin{bmatrix} 0 & 8 \\ 2 & 0 \end{bmatrix}
\end{verbatim}

собственные числа равны \texttt{\textbackslash{}lambda\_\ 1\ =\ 4} и
\texttt{\textbackslash{}lambda\_2\ =\ -4} , а собственные векторы
соответственно

\begin{verbatim}
X_1 =
\begin{bmatrix} 2 \\ 1 \end{bmatrix}
in  X_2 =
\end{verbatim}

\pandocbounded{\includegraphics[keepaspectratio]{images/_page_362_Figure_5.jpeg}}

.

Значит, произведение

\begin{verbatim}
P^{-1}AP =
\begin{bmatrix} 2 & -2 \\ 1 & 1 \end{bmatrix}
^{-1}
\begin{bmatrix} 0 & 8 \\ 2 & 0 \end{bmatrix}
\begin{bmatrix} 2 & -2 \\ 1 & 1 \end{bmatrix}
= \\ =
\begin{bmatrix} 1/4 & 1/2 \\ -1/4 & 1/2 \end{bmatrix}
\begin{bmatrix} 0 & 8 \\ 2 & 0 \end{bmatrix}
\begin{bmatrix} 2 & -2 \\ 1 & 1 \end{bmatrix}
\end{verbatim}

дорівнює діагональній матриці

\begin{verbatim}
\Lambda =
\begin{bmatrix} 4 & 0 \\ 0 & -4 \end{bmatrix}
\end{verbatim}

.

(Читатель может проверить это самостоятельно.)

Тепер можна розв'язати просту систему з двох диференціальних рівнянь у
похідних у похідних (з відповідними початковими умовами).

\section{\texorpdfstring{Лінійне \texttt{u\_t} + Au\_x = 0 системне
рішення.}{Лінійне u\_t + Au\_x = 0 системне рішення.}}\label{ux43bux456ux43dux456ux439ux43dux435-u_t-au_x-0-ux441ux438ux441ux442ux435ux43cux43dux435-ux440ux456ux448ux435ux43dux43dux44f.}

Розглянемо задачу Коші для системи, що містить два рівняння та дві
початкові умови:

(УЧП 1)

\begin{verbatim}
\frac{\partial u_1}{\partial t} + 8 \frac{\partial u_2}{\partial x} = 0
\end{verbatim}

,

\begin{verbatim}
(29.3) (УЧП 2)  \frac{\partial u_2}{\partial t} + 2 \frac{\partial u_1}{\partial x} = 0  ,  -\infty < x < \infty  ,  0 < t < \infty  ,
\end{verbatim}

\begin{verbatim}
(HУ 1)  u_1(x, 0) = f(x)  ,
(HУ 2)  u_2(x, 0) = g(x)  ,  -\infty < x < \infty  .
\end{verbatim}

Ця задача може відповідати визначенню \texttt{u\_1}(x, t) тиску та
густини \texttt{u\_2}(x, t) як функції просторової координати x і часу t
за відомими розподілами цих величин у початковий момент.

Запишемо систему рівнянь у матричній формі

\begin{verbatim}
\begin{bmatrix} \frac{\partial u_1}{\partial t} \\ \frac{\partial u_2}{\partial t} \end{bmatrix}
\end{verbatim}

\begin{itemize}
\tightlist
\item
\end{itemize}

\begin{verbatim}
\begin{bmatrix} 0 & 8 \\ 2 & 0 \end{bmatrix}
\begin{bmatrix} \frac{\partial u_1}{\partial x} \\ \frac{\partial u_2}{\partial x} \end{bmatrix}
=
\begin{bmatrix} 0 \\ 0 \end{bmatrix}
\end{verbatim}

или \texttt{(29\ \textbackslash{}\ 4)}

\begin{verbatim}
u_t + Au_s = 0
\end{verbatim}

где

\begin{verbatim}
A =
\begin{bmatrix} 0 & 8 \\ 2 & 0 \end{bmatrix}
, \quad u_1 =
\begin{bmatrix} \frac{\partial u_1}{\partial t} \\ \frac{\partial u_2}{\partial t} \end{bmatrix}
, \quad u_x =
\begin{bmatrix} \frac{\partial u_1}{\partial x} \\ \frac{\partial u_2}{\partial x} \end{bmatrix}
, \quad 0 =
\begin{bmatrix} 0 \\ 0 \end{bmatrix}
\end{verbatim}

.

Введем новые неизвестные функции

\begin{verbatim}
v =
\begin{bmatrix} v_1 \\ v_2 \end{bmatrix}
\end{verbatim}

с помощью преобразования

\begin{verbatim}
u = Pv
\end{verbatim}

.

де P --- матриця, стовпці якої містять власні вектори матриці A (матриця
P, яку вже знають). Виявляється, що після такого перетворення буде
отримана дуже проста система для визначення v (два рівняння відносно
нового невідомого \texttt{v\_1} і \texttt{v\_2} виявляються
незалежними). Це означає, що \texttt{v\_1} і \texttt{v\_2} легко знайти.
Після \texttt{v\_1} і \texttt{v\_2} знайдено, згідно з формулою

\begin{verbatim}
u = Pv
\end{verbatim}

находятся искомые функции и, и и2.

Спершу ж давайте розглянемо, як виглядає система визначення v.
Диференціюючи дві частини співвідношення u = Pv, отримуємо

(29 5)

\begin{verbatim}
\frac{\partial u}{\partial t} = P \frac{\partial v}{\partial t},
\frac{\partial u}{\partial x} = P \frac{\partial v}{\partial x}.
\end{verbatim}

(Чтобы лучше понять дальнейшее изложение, читателю рекомендуется
записать эти уравнен я в развернутой форме.)

Теперь подставим соотношения (29.5) в систему

\begin{verbatim}
u_t + Au_x = 0
\end{verbatim}

и получим

\begin{verbatim}
Pv_t + APv_x = 0.
\end{verbatim}

Если умножить теперь обе части последнего уравнения на
\texttt{P\^{}\{-1\}} , то получим

\begin{verbatim}
v_t + P^{-1}APv_x = 0,
\end{verbatim}

или

\begin{verbatim}
(29.6) v_t + \Lambda v_x = 0.
\end{verbatim}

Раньше мы уже видели, что для нашей матрицы

\begin{verbatim}
\Lambda =
\begin{bmatrix} 4 & 0 \\ 0 & -4 \end{bmatrix}
\end{verbatim}

.

Запишем (29.6) в развернутой форме

(29.7)

\begin{verbatim}
\frac{\frac{\partial v_1}{\partial t} + 4 \frac{\partial v_1}{\partial x} = 0, \\ \frac{\partial v_2}{\partial t} - 4 \frac{\partial v_2}{\partial x} = 0.
\end{verbatim}

Результатом є система з двох непов'язаних рівнянь, які можна розв'язати
незалежно. Читач має знати, що розв'язки цих рівнянь --- це хвилі, що
рухаються

\begin{verbatim}
v_1(x, t) = \varphi(x-4t),
\end{verbatim}

\texttt{v\_2(x,\ t)\ =\ \textbackslash{}psi(x+4t),} где ф и
ф-произвольные дифференцируемые функции.

Таким образом, мы нашли функции \texttt{v\_1} и \texttt{v\_2} . Для
получения общего решения u нужно вычислить

\begin{verbatim}
u = Pv =
\begin{bmatrix} 2 & -2 \\ 1 & 1 \end{bmatrix}
\begin{bmatrix} v_1 \\ v_2 \end{bmatrix}
=
=
\begin{bmatrix} 2 & -2 \\ 1 & 1 \end{bmatrix}
\begin{bmatrix} \varphi(x - 4t) \\ \psi(x + 4t) \end{bmatrix}
=
=
\begin{bmatrix} 2\varphi(x - 4t) - 2\psi(x + 4t) \\ \varphi(x - 4t) + \psi(x + 4t) \end{bmatrix}
\end{verbatim}

.

Другими словами,

(29.8)

\begin{verbatim}
u_1(x, t) = 2\varphi(x-4t) - 2\psi(x+4t), u_2(x, t) = \varphi(x-4t) + \psi(x+4t).
\end{verbatim}

Следовательно, мы нашли общее решение системы двух уравнений (29.3).

Типичным решением (а их бесконечное число) могут быть, например, функции

\begin{verbatim}
\phi(\xi) = \sin \xi, \ \psi(\xi) = \xi^2
\end{verbatim}

(две произвольные функции)

и, следовательно,

\begin{verbatim}
u_1(x, t) = 2 \sin(x-4t) - 2(x+4t)^2,
\end{verbatim}

\texttt{u\_2(x,\ t)\ =\ \textbackslash{}sin(x-4t)\ +\ (x+4t)\^{}2.}
Подставим теперь общее решение (29.8) в начальные условия

\begin{verbatim}
u_1(x, 0) = f(x),
\end{verbatim}

\texttt{u\_2(x,\ 0)\ =\ g(x),} в результате получаем

\begin{verbatim}
2\varphi(x) - 2\psi(x) = f(x),
\varphi(x) + \psi(x) = g(x),
\end{verbatim}

откуда функци і ф и ф легко находятся:

\begin{verbatim}
\varphi(x) = \frac{1}{4} \left[ f(x) + 2g(x) \right],
\psi(x) = \frac{1}{4} \left[ 2g(x) - f(x) \right].
\end{verbatim}

\textbf{Тепе}рь можно написать решение задачи (29.3) в окончательном
виде

\begin{verbatim}
u_{1}(x, t) = 2\varphi(x-4t) - 2\psi(x+4t) =
= \frac{1}{2} \left[ f(x-4t) + 2g(x-4t) \right] - \frac{1}{2} \left[ 2g(x+4t) - f(x+4t) \right],
\end{verbatim}

(29.9)

\begin{verbatim}
u_{2}(x, t) = \varphi(x-4t) + \psi(x+4t) =
= \frac{1}{4} \left[ f(x-4t) + 2g(x-4t) \right] + \frac{1}{4} \left[ 2g(x+4t) - f(x+4t) \right].
\end{verbatim}

\section{ЗАУВАЖЕННЯ}\label{ux437ux430ux443ux432ux430ux436ux435ux43dux43dux44f-2}

\begin{enumerate}
\def\labelenumi{\arabic{enumi}.}
\item
  Многие численные методы ориентированы на решение систем уравнений и,
  значит, многне программы для ЭВМ написаны так, чтобы решать системы
  уравнений первого порядка. По этой причине, прежде чем воспользоваться
  программои, приходится преобразовать свое уравнение высокого порядка в
  систему уравнений первого порядка.
\item
  Линейную систему
\end{enumerate}

\begin{verbatim}
u_t + A(x, t) u_x + B(x, t) u = 0
\end{verbatim}

можна розв'язати тим самим методом, що й система
\texttt{u\_t\ +\ Au\_x\ =\ 0} . Різниця полягає в тому, що в цьому
випадку матриця P із власних векторів матриці \texttt{A(x,\ t)} буде
функцією від x і t. \# ЗАДАЧІ 1. Запишите систему уравнений (29.2) в
матричной форме

\begin{verbatim}
Au_t + Bu_x + Cu = 0,
\end{verbatim}

где A, B и C-матрицы третьего порядка.

\begin{enumerate}
\def\labelenumi{\arabic{enumi}.}
\setcounter{enumi}{1}
\tightlist
\item
  Найдите собственные векторы и собственные значения матрицы
\end{enumerate}

\begin{verbatim}
A =
\begin{bmatrix} 1 & 1 \\ 4 & 1 \end{bmatrix}
\end{verbatim}

.

\begin{enumerate}
\def\labelenumi{\arabic{enumi}.}
\setcounter{enumi}{2}
\tightlist
\item
  Воспользовавшись результатами задачи 2, решите систему
\end{enumerate}

\begin{verbatim}
\begin{cases} \frac{\partial u_1}{\partial t} + \frac{\partial u_1}{\partial x} + \frac{\partial u_2}{\partial x} = 0, \\ \frac{\partial u_2}{\partial t} + 4 \frac{\partial u_1}{\partial x} + \frac{\partial u_2}{\partial x} = 0. \end{cases}
\end{verbatim}

УКАЗАНИЕ. Запишите сначала систему в матричной форме.

\begin{enumerate}
\def\labelenumi{\arabic{enumi}.}
\setcounter{enumi}{3}
\item
  Проверьте формулы (29.5), записав их в скалярной форме.
\item
  Проверьте, что функции \texttt{u\_1} и \texttt{u\_2} из (29.8)
  удовлетворяют обоим уравнениям (29.3).
\end{enumerate}

\section{КОЛЕБАНИЯ МЕМБРАНЫ (ВОЛНОВОЕ УРАВНЕНИЕ В ПОЛЯРНЫХ
КООРДИНАТАХ)}\label{ux43aux43eux43bux435ux431ux430ux43dux438ux44f-ux43cux435ux43cux431ux440ux430ux43dux44b-ux432ux43eux43bux43dux43eux432ux43eux435-ux443ux440ux430ux432ux43dux435ux43dux438ux435-ux432-ux43fux43eux43bux44fux440ux43dux44bux445-ux43aux43eux43eux440ux434ux438ux43dux430ux442ux430ux445}

МЕТА ЛЕКЦІЇ: Показать, что колебания мембраны можно описать волновым
уравнением

\begin{verbatim}
u_{tt} = c^2 \left( u_{rr} + \frac{1}{r} u_r + \frac{1}{r^2} u_{\theta\theta} \right)
\end{verbatim}

.

Если искать решение этого уравнения в виде стоячих волн:

\begin{verbatim}
u(r, \theta, t) = P(r) \Theta(\theta) T(t),
\end{verbatim}

то оно распадается на три обыкновенных дифференциальных уравнения:

\begin{verbatim}
T'' + \lambda^2 c^2 T = 0
\end{verbatim}

(уравнение гармонических колебаний),
\texttt{r\^{}2\ R\textquotesingle{}\textquotesingle{}\ +\ rR\textquotesingle{}\ +\ (\textbackslash{}lambda\^{}2\ r\^{}2\ -\ n\^{}2)\ R\ =\ 0}
(уравнение Бесселя), (уравнение гармонических колебаний).

Все произведения
\texttt{R(r)\ \textbackslash{}Theta(\textbackslash{}theta)\ T(t)}
решений этих трех обыкновенных дифференциальных уравнений описывают
основные колебания мембраны, а функции
\texttt{R(r)\ \textbackslash{}Theta(\textbackslash{}theta)} описывают
форму мембраны. Чтобы найти колебания мембраны при произвольных
начальных условиях, мы составим такую комбинацию фундаментальных
решений, чтобы начальные условия удовлетворялись.

\texttt{Цель\ этой\ лекции\ ясна—найти\ колебания\ круглой\ мембраны\ с\ заданными\ начальными\ и\ граничными\ условиями.\ Для\ простоты\ будем\ считать,\ что\ радиус\ мембраны\ равен\ единице,\ а\ смещение\ по\ границе\ равно\ нулю.\ Пусть\ \textquotesingle{}u\textbackslash{}left(r,\textbackslash{}theta,t\textbackslash{}right)\textquotesingle{}\ обозначает\ возвышение\ точек\ мембраны\ под\ плоскостью\ равновесия.\ Нам\ необходимо\ решить\ задачу\ (см.\ рис.\ 30.1)}

\begin{verbatim}
\begin{array}{llll} (\text{УЧП}) & u_{tt} = c^2 \left( u_{rr} + \frac{1}{r} u_r + \frac{1}{r^2} u_{\theta\theta} \right), & 0 < r < 1, \ 0 < t < \infty, \\ (\text{ГУ}) & u = 0 & \text{при } r = 1, & 0 < t < \infty, \\ (\text{НУ}) & \begin{cases} u = f(r, \ \theta) \\ u_t = g(r, \ \theta) \end{cases} & \text{при } t = 0. \end{array}
\end{verbatim}

Щоб розв'язати цю задачу, згадаємо проблему коливань скрипкових струн із
лекції 20, яку розв'язали суперпозицією нескінченної кількості простих
коливань. Якщо ви використовуєте-

ся этим же подходом, то решения следует искать в виде

\begin{verbatim}
u(r, \theta, t) = U(r, \theta) T(t)
\end{verbatim}

.

Форма этих колебаний определяется функцией
\texttt{U(r,\ \textbackslash{}theta)} , а характер осцилляций ---
множителем \texttt{T(t)}.

Подставив это представление в волновое уравнение (см. задачу 1), получим
два уравнения

\begin{verbatim}
\Delta U + \lambda^2 U = 0
\end{verbatim}

(уравнение Гельм-гольца),

\begin{verbatim}
T'' + \lambda^2 c^2 T = 0
\end{verbatim}

(уравнение гармонических колебаний),

\pandocbounded{\includegraphics[keepaspectratio]{images/_page_221_Picture_9.jpeg}}

Рис. 30.1. Қолеблющаяся мембрана (гиперболическая смешанная задача).

где

\begin{verbatim}
\Delta U = U_{rr} + \frac{1}{r} U_r + \frac{1}{r^2} U_{\theta\theta}.
\end{verbatim}

Зверніть увагу, що тут ми вимагаємо, щоб константа розділення була
додатною (тому позначимо її \texttt{\textbackslash{}lambda\^{}\{2-1\}}
)). Лише в цьому випадку функції \texttt{T(t)} будуть періодичними.

На следующем шаге мы решим уравнение Гельмгольца, но сначала разберемся
с граничным условием для него. Чтобы найти его, подставим
\texttt{u(r,\ \textbackslash{}theta,\ t)\ =\ U(r,\ \textbackslash{}theta)\ T(t)}
в граничное условие для мембраны

\begin{verbatim}
u(1, \theta, t) = U(1, \theta) T(t) = 0, \quad 0 < t < \infty,
\end{verbatim}

или

\begin{verbatim}
U(1, \theta) = 0.
\end{verbatim}

Следовательно, для того чтобы найти форму фундаментальных колебаний
мембраны, необходимо решить задачу

\begin{verbatim}
\Delta U + \lambda^2 U = 0,
U(1, \theta) = 0.
\begin{cases} \Delta U + \mu u = 0 & \text{в круге,} \\ U = 0 & \text{на границе круга} \end{cases}
\end{verbatim}

при \texttt{\textbackslash{}mu\ \textbackslash{}le\ 0} имеет только
тривиальное решение U=0. --- Прим. ред.

1) Поскольку краевая задача

\texttt{Это\ хорошо\ известная\ эллиптическая\ задача\ на\ собственные\ значения,\ и\ нам\ надо\ найти\ все\ \ \textquotesingle{}\textbackslash{}lambda\textquotesingle{}\ ,\ при\ которых\ она\ имеет\ ненуленые\ решения.\ Решения\ \ \textquotesingle{}U\textbackslash{}left(r,\textbackslash{},\textbackslash{}theta\textbackslash{}right)\textquotesingle{}\ \ описывают\ форму\ фундаментальных\ колебаний\ мембраны,\ а\ собственные\ числа\ \ \textquotesingle{}\textbackslash{}lambda\textquotesingle{}\ \ являются\ нулями\ некоторых\ функций\ Бесселя\ и\ пропорциональны\ частотам\ этих\ колебаний.}

Итак, нам предстоит теперь решить задачу на собственные значения для
уравнения Гельмгольца (эта задача очень важна и сама по себе). Будем
решать ее тем же методом, что и другие линейные однородные уравнения с
нулевыми граничными условиями, а именно методом разделения переменных.

\section{Решение задачи на собственные значения для уравнения
Гельмгольца}\label{ux440ux435ux448ux435ux43dux438ux435-ux437ux430ux434ux430ux447ux438-ux43dux430-ux441ux43eux431ux441ux442ux432ux435ux43dux43dux44bux435-ux437ux43dux430ux447ux435ux43dux438ux44f-ux434ux43bux44f-ux443ux440ux430ux432ux43dux435ux43dux438ux44f-ux433ux435ux43bux44cux43cux433ux43eux43bux44cux446ux430}

Чтобы решить задачу

\begin{verbatim}
\begin{array}{ll} (\mathsf{Y}\mathsf{H}\mathsf{\Pi}) & \Delta U + \lambda^2 U = 0, \\ (\mathsf{\Gamma}\mathsf{Y}) & U(1, \theta) = 0, \end{array}
\end{verbatim}

подставим

\begin{verbatim}
U(r, \theta) = R(r) \Theta(\theta)
\end{verbatim}

в задачу Гельмгольца. Сделав это, получаем

\begin{verbatim}
r^2R'' + rR' + (\lambda^2r^2 - n^2)R = 0
\end{verbatim}

\texttt{(уравнение\ Бесселя),\ \textquotesingle{}R(1)\ =\ 0\textquotesingle{}\ ,\ \textquotesingle{}\textbackslash{}Theta\textquotesingle{}\textquotesingle{}\ +\ n\^{}2\textbackslash{}Theta\ =\ 0\textquotesingle{}\ (Зроби\ все\ сам).\ Зверніть\ увагу,\ що\ ми\ обрали\ константу\ розбиття\ \textquotesingle{}n\^{}2\textquotesingle{}\ ,\ \textquotesingle{}n=0,.\ \textbackslash{}\ 1,\textbackslash{}\ 2,\textbackslash{}\ \textbackslash{}ldots\textquotesingle{}\ тому,\ що\ хочемо,\ щоб\ функції\ \textquotesingle{}\textbackslash{}Theta\textbackslash{}left(\textbackslash{}theta\textbackslash{}right)\textquotesingle{}\ були\ періодичними\ з\ періодом\ \textquotesingle{}2\textbackslash{}pi\textquotesingle{}\ .\ Очевидно,\ що\ форма\ мембрани\ є\ періодичною\ функцією\ \textquotesingle{}\textbackslash{}theta\textquotesingle{}\ .\ Отже,\ щоб\ розв\textquotesingle{}язати\ задачу\ Гельмгольца,\ потрібно\ розв\textquotesingle{}язати\ два\ звичайних\ диференціальні\ рівняння}

\begin{verbatim}
r^2R'' + rR' + (\lambda^2r^2 - n^2) R = 0
\end{verbatim}

\texttt{,\ \textquotesingle{}0\ \textless{}\ r\ \textless{}\ 1\textquotesingle{}\ ,\ \textquotesingle{}R(0)\ \textless{}\ \textbackslash{}infty\textquotesingle{}\ (физическое\ ограничение),\ \textquotesingle{}R(1)\ =\ 0\textquotesingle{}\ ,\ \textquotesingle{}\textbackslash{}Theta\textquotesingle{}\textquotesingle{}\ +\ n\^{}2\textbackslash{}Theta\ =\ 0\textquotesingle{}\ .}

\section{Уравнение
Бесселя}\label{ux443ux440ux430ux432ux43dux435ux43dux438ux435-ux431ux435ux441ux441ux435ux43bux44f}

Уравнение

\begin{verbatim}
r^2R'' + rR' + (\lambda^2r^2 - n^2)R = 0
\end{verbatim}

жорошо известно в теории обыкновенных дифференциальных уравнений. Это
так называемое уравнение Бесселя. Оно имеет два

линейно-независимых решения
\texttt{R\_\{1\}(r)\ =\ A\ J\_\{n\}(\textbackslash{}lambda\ r)} функция
Бесселя \emph{первого} рода \emph{п}-го порядка,
\texttt{R\_\{2\}(r)\ =\ BY\_\{n\}(\textbackslash{}lambda\ r)} функция
Бесселя второго рода n-го порядка,

и, следовательно, общее решение этого уравнения представимо в виде

\begin{verbatim}
R(r) = AJ_n(\lambda r) + BY_n(\lambda r).
\end{verbatim}

Зверніть увагу, що рішення залежать від параметрів n і
\texttt{\textbackslash{}lambda} , включених у рівняння. Деякі графіки
функції Бесселя показані на рис. 30.2. Щоб визначити функції
\texttt{J\_\ n(\textbackslash{}lambda\ r)} і
\texttt{Y\_n(\textbackslash{}lambda\ r)} , можливо

\pandocbounded{\includegraphics[keepaspectratio]{images/_page_223_Figure_8.jpeg}}

\pandocbounded{\includegraphics[keepaspectratio]{images/_page_223_Figure_9.jpeg}}

Рис. 30 2 Функции Бесселя. a --- функции Бесселя первого рода;
\texttt{\textbackslash{}pmb\{\textbackslash{}delta\}} --- функции
Бесселя второго рода.

використовуйте метод Фробеніуса, який дозволяє знаходити \texttt{R(r)} у
вигляді силових рядів. Цей метод містить два лінійно незалежні ряди
ступенів \texttt{J\_n(\textbackslash{}lambda\ r)} та
\texttt{Y\_n(\textbackslash{}lambda\ r)} . Оскільки функція
\texttt{Y\_n(\textbackslash{}lambda\ r)} не обмежена в r=0, ми обираємо
розв'язок у вигляді

\begin{verbatim}
R(r) = AJ_n(\lambda r).
\end{verbatim}

\texttt{Последний\ шаг\ в\ определении\ \textquotesingle{}R\textbackslash{}left(\ r\textbackslash{}right)\textquotesingle{}\ связан\ с\ использованием\ граничного\ условия\ \textquotesingle{}R\textbackslash{}left(\ \textbackslash{}mathbf\{l\}\textbackslash{}right)\ =0\textquotesingle{}\ для\ определения\ всех\ \textquotesingle{}\textbackslash{}lambda\textquotesingle{}\ (в\ данный}

Таблица 30.1

{\def\LTcaptype{none} % do not increment counter
\begin{longtable}[]{@{}lll@{}}
\toprule\noalign{}
Корни & уравнений. & \texttt{J\_n(r)\ =\ 0} \\
\midrule\noalign{}
\endhead
\bottomrule\noalign{}
\endlastfoot
------- & ------------ & -------------- \\
\end{longtable}
}

{\def\LTcaptype{none} % do not increment counter
\begin{longtable}[]{@{}lllllllll@{}}
\toprule\noalign{}
Γ & 1 & n & & & & & & \\
\midrule\noalign{}
\endhead
\bottomrule\noalign{}
\endlastfoot
& & 0 & 1 & 2 & 3 & 4 & & \\
& 1 & 2,40 & 3,83 & 5,13 & 6,38 & 7,59 & & \\
ı & 2 & 5,52 & 7,02 & 8,42 & 9,76 & 11,06 & & \\
ľ & 3 & 8,65 & 10,17 & 11,62 & 13,02 & 14,37 & & \\
& 4 & 11,79 & 13,32 & 14,80 & 16,22 & 17,62 & & \\
& 5 & 14,93 & 16,47 & 17,96 & 19,41 & 20,83 & & \\
& . & • & & • & & • & & \\
& & & • & & • & • & & \\
& • & • & • & • & • & • & & \\
\end{longtable}
}

момент нас не интересует коэффициент A). Подставляя \texttt{R\ (1)} = 0,
в \texttt{R\ (r)} = \texttt{J\_n} ( \texttt{\textbackslash{}lambda\ r}
), получаем \texttt{J\_n} ( \texttt{\textbackslash{}lambda} ) = 0.
Другими словами, для того чтобы функция \texttt{R\ (r)} обращалась в
нуль на границе круга, мы должны выбрать константу разделения так, чтобы
она была корнем уравнения \texttt{J\_n} (r) = 0,
\texttt{\textbackslash{}tau} . e.

\begin{verbatim}
\lambda = k_{nm}
\end{verbatim}

где \texttt{k\_\{nm\}} --- m-й корень уравнения \texttt{J\_n(r)\ =\ 0} .

Таблиці цих коренів добре відомі, але їх можна обчислити на комп'ютері.
Значення перших кількох коренів наведені в Таблиці 1. 30.1. Знаючи ці
корені, ми можемо одразу побудувати розв'язок задачі граничного значення
для рівняння Гельмгольца. Якщо власні значення
\texttt{\textbackslash{}lambda\ =\ k\_\{nm\}} , то власні функції
\texttt{U\_\{\ nm\}(r)} визначаються формулою

\begin{verbatim}
U_{nm}(r, \theta) = J_n(k_{nm}r) [A \sin(n\theta) + B \cos(n\theta)].
\end{verbatim}

На рис. 30.3 изображены некоторые из этих функций для различных значений
n и m. Общая форма функции \texttt{U\_\{nm\}(r,\textbackslash{}theta)}
не зависит от значений констант A и B. Эти константы влияют только на
амплитуду колебаний и расположение линии узлов относительно
\texttt{\textbackslash{}theta=0} .

Каждая функция представляет собой фундаментальное колебание круглой
мембраны с частотой

\begin{verbatim}
f_{nm} = k_{nm}c/2\pi.
\end{verbatim}

Эти частоты мы нашли из решений

\begin{verbatim}
T_{nm}(t) = A \sin(k_{nm}ct) + B \cos(k_{nm}ct)
\end{verbatim}

\section{уравнения}\label{ux443ux440ux430ux432ux43dux435ux43dux438ux44f}

\begin{verbatim}
T'' + k_{nm}^2 c^2 T = 0.
\end{verbatim}

Интересно отметить, что отношение частоты колебания
\texttt{U\_\{nm\}(\textbackslash{}mathbf\{r\},\ \textbackslash{}mathbf\{\textbackslash{}theta\})}
к частоте основного колебания U\_\{nt\}(\ma

\begin{verbatim}
sU(x) = -V \frac{dU}{dx}, \quad 0 < x < \infty,
\begin{cases} u(R_1, \theta) = g_1(\theta), \\ u(R_2, \theta) = g_2(\theta). \end{cases}
U(0) = P/s.
\end{verbatim}

\begin{verbatim}
\begin{cases} 2/3, & n=1, \\0 & \text{для всіх інших } n, \\ d_n=\begin{cases} -2/3, & n=1, \\0 & \text{для всіх інших} n. \end{cases}
\end{verbatim}

\begin{verbatim}
H(\xi) =
\begin{cases} 1, & 0 \le \theta < \pi, \\ 0, & \pi \le \theta < 2\pi. \end{cases}
\begin{cases} 0, & \xi < 0, \\ 1, & \xi \geqslant 0. \end{cases}
\end{verbatim}

\texttt{INL1448X\ (\textbackslash{}Gamma\ Y)\ \ u(1,\ \textquotesingle{}\textbackslash{}varphi\textquotesingle{})\ =\ \textbackslash{}cos(3\textquotesingle{}\textbackslash{}varphi\textquotesingle{})\ \ .}

\texttt{Используя\ тригонометрию,\ попытайтесь\ выразить\ \textquotesingle{}\textbackslash{}cos\{(3\textbackslash{}phi)\}\textquotesingle{}\ через\ \textquotesingle{}\textbackslash{}cos\{\textbackslash{}phi\},\ \textbackslash{},\ \textbackslash{}cos\^{}2\{\textbackslash{}phi\},\ \textbackslash{},\ \textbackslash{}cos\^{}3\{\textbackslash{}phi\},\ \textbackslash{},\ a\textquotesingle{}\ затем\ скомбинируйте\ их\ так,\ чтобы\ по-}

лучилось разложение

\begin{verbatim}
\cos(3\varphi) = a_0 P_0(\cos\varphi) + a_1 P_1(\cos\varphi) + \dots
\end{verbatim}

\begin{enumerate}
\def\labelenumi{\arabic{enumi}.}
\setcounter{enumi}{4}
\tightlist
\item
  Решите задачу Дирихле в шаре
\end{enumerate}

\begin{verbatim}
\Delta u = 0, \quad 0 < r < 1, \quad 0 \le \varphi \le \pi, \quad 0 \le \theta < 2\pi.
u(1, \varphi) =
\begin{cases} 1, & 0 \le \varphi \le \pi/2, \\ -1, & \pi/2 < \varphi \le \pi. \end{cases}
\end{verbatim}

В этой задаче верхняя полусфера нагрета (+1), а нижняя охлаждена (-1).

\begin{enumerate}
\def\labelenumi{\arabic{enumi}.}
\setcounter{enumi}{5}
\tightlist
\item
  Найдите решение задачи Дирихле в шаре
\end{enumerate}

(УЧП)

\begin{verbatim}
\Delta u = 0
\end{verbatim}

\texttt{,\ \textquotesingle{}1\ \textless{}\ r\ \textless{}\ \textbackslash{}infty\textquotesingle{}\ ,\ \textquotesingle{}0\ \textbackslash{}le\ \textbackslash{}varphi\ \textbackslash{}le\ \textbackslash{}pi\textquotesingle{}\ ,\ \textquotesingle{}0\ \textbackslash{}le\ \textbackslash{}theta\ \textless{}\ 2\textbackslash{}pi\textquotesingle{}\ ,\ \textquotesingle{}(\textbackslash{}Gamma\ Y)\textquotesingle{}\ \textquotesingle{}u(1,\ \textbackslash{}varphi)\ =\ 1\ +\ \textbackslash{}cos\ \textbackslash{}varphi\textquotesingle{}\ .}

Проверьте ответ.

Лекция 36

\section{НЕОДНОРОДНАЯ ЗАДАЧА ДИРИХЛЕ (ФУНКЦИЯ
ГРИНА)}\label{ux43dux435ux43eux434ux43dux43eux440ux43eux434ux43dux430ux44f-ux437ux430ux434ux430ux447ux430-ux434ux438ux440ux438ux445ux43bux435-ux444ux443ux43dux43aux446ux438ux44f-ux433ux440ux438ux43dux430}

МЕТА ЛЕКЦІЇ: Показати, як гетерогенну задачу Діріхле можна розв'язати за
допомогою функції Гріна (вихідної функції). У цьому важливому методі
права частина рівняння розглядається як певна вхідна дія і розкладається
на неперервний набір дельтаподібних джерел, розподілених по заданій
площі. Потім визначається відповідь системи на кожне таке джерело
(функція Гріна), і всі відповіді підсумовуються (інтегровані). У
результаті ми отримуємо повне розв'язання проблеми.

Достаточно общей задачей прикладной математики является определение
потенциала в некоторой области пространства, если задано распределение
источников \texttt{f(x,\ y)} енутри этой области. В электростатике
потенциал (в вольтах) в области D обусловлен распределением зарядов с
плотностью \texttt{f(x,\ y)} по этой области. Типичная задача ---
определение потенциала в круге в двумерном случае.

В такой задаче потенциал должен удовлетворять уравнению

Пуассона (например, с нулевыми граничными условиями)

(36.1) (УЧП)

\begin{verbatim}
u_{rr} + \frac{1}{r} u_r + \frac{1}{r^2} u_{\theta\theta} = f(r, \theta), \quad 0 \le r < 1, \quad 0 \le \theta < 2\pi,
\end{verbatim}

\texttt{(ГУ)\ \textquotesingle{}u(1,\ \textbackslash{}theta)\ =\ 0,\ \textbackslash{}quad\ 0\ \textbackslash{}le\ \textbackslash{}theta\ \textbackslash{}le\ 2\textbackslash{}pi.\textquotesingle{}\ Ми\ не\ випадково\ обрали\ нульові\ умови.\ Якщо\ б\ ми\ хотіли\ розв\textquotesingle{}язати\ задачу\ в\ загальному\ випадку,\ з\ ненульовими\ граничними\ умовами\ та\ неоднорідним\ рівнянням,\ то\ внесок\ неоднорідних\ граничних\ умов\ можна\ було\ б\ знайти\ за\ допомогою\ формули\ інтегралу\ Пуассона\ з\ лекції\ 33.}

Чтобы получить хоть какое-то представление о неоднородных
дифференциальных уравнениях, давайте рассмотрим графически решение
следующей задачи для уравнения Пуассона:

(УЧП)

\begin{verbatim}
\Delta u = -q
\end{verbatim}

\texttt{,\ \textquotesingle{}0\ \textless{}\ r\ \textless{}\ 1\textquotesingle{}\ ,\ \textquotesingle{}0\ \textbackslash{}le\ \textbackslash{}theta\ \textbackslash{}le\ 2\textbackslash{}pi\textquotesingle{}\ (\ \textquotesingle{}q\textquotesingle{}\ —положительная\ константа),}

(HY)

\begin{verbatim}
u(1, \theta) = 0, \quad 0 \le \theta \le 2\pi.
\end{verbatim}

\texttt{На\ межі\ потенціал\ (або\ температура,\ якщо\ ви\ бажаєте)\ утримується\ на\ нулі,\ а\ лапласіан\ дорівнює\ q\ у\ всіх\ точках\ області.\ Оскільки\ \textquotesingle{}\textbackslash{}Delta\ u\textbackslash{}left(p\textbackslash{}right)\textquotesingle{}\ слугує\ мірою\ різниці\ між\ функцією\ \textquotesingle{}u\textbackslash{}left(p\textbackslash{}right)\textquotesingle{}\ і\ її\ середнім,\ рівняння\ Пуассона\ стверджує,\ що\ поверхня\ \textquotesingle{}u\textbackslash{}left(r\textbackslash{}right)\textquotesingle{}\ завжди\ опукла\ вниз.\ Інакше\ кажучи,\ вона\ виглядатиме\ як\ тонка\ мембрана\ з\ нерухомим\ краєм,\ яка\ здувається\ зверху\ вниз\ потоком\ повітря.}

Якщо права частина рівняння \texttt{f(x,\ y)} змінюється від точки до
точки, то опуклість також буде змінною функцією точки цієї області.

Перейдем теперь к главной части нашей лекции: определению функции Грина
и решению уравнения (36.1).

Сначала, однако, мы должны ввести понятие потенциала точеных источников
и стоков.

\section{Потенциалы точечных источников и
стоков}\label{ux43fux43eux442ux435ux43dux446ux438ux430ux43bux44b-ux442ux43eux447ux435ux447ux43dux44bux445-ux438ux441ux442ux43eux447ux43dux438ux43aux43eux432-ux438-ux441ux442ux43eux43aux43eux432}

Щоб розв'язати неоднорідне лінійне рівняння, достатньо розв'язати це
рівняння за допомогою точкового джерела, оскільки розв'язок задачі в
загальному випадку можна знайти, підсумувавши внески з точкових джерел.
Наше завдання --- визначити потенціал, який створює точкове джерело (або
поглинач) у певній області. Ці джерела можна інтерпретувати по-різному.
У теорії теплопровідності джерело можна розглядати як точку, де походить
тепло, а дренаж --- як точку, де поглинає тепло. В електростатиці
джерело можна розглядати як окремий позитивний заряд (протон), а дренаж
--- як окремий негативний заряд (електрон). У будь-якому разі, незалежно
від інтерпретації, ми знайдемо двовимірний потенціал і,

создаваемый единичным точечным источником (поиск трехмерного потенциала
мы оставляем читателю в качестве упражнения).

Припустимо, що єдине точкове джерело +q розташоване в початку Clear that
heat (або щось інше)

будет вытекать из области вдоль радиальных линий, и, следовательно, для
везничны полного потока, протекающего через окружность радиуса r можно
записать

Полный поток, вытекающий через окруж-

HOCTD

\begin{verbatim}
= -\int_{0}^{2\pi} u_{r}(r) r d\theta = -2\pi r u_{r}(r).
\end{verbatim}

Но полный поток должен совпадать с количеством тепла, созданного внутри
круга (закон сохранения), т. е.

\begin{verbatim}
-2\pi r u_r(r) = q.
\end{verbatim}

Решая это диф реренциальное уравнение относительно функции
\texttt{u(r)}, получаем Р И С

\pandocbounded{\includegraphics[keepaspectratio]{images/_page_273_Picture_10.jpeg}}

36 1 Радиальный поток тепла от точечного источника; a ---тепло
выделяется в начале координат.

\begin{verbatim}
u(r) = \frac{-q}{2\pi} \ln r = \frac{q}{2\pi} \ln \frac{1}{r}.
\end{verbatim}

(Решение изображено на рис. 362)

В электростатике разность потенциа тов \texttt{u(B)} - \texttt{u(A)}
равна работе, которую необходимо совершить, чтобы переместить единичный
положительный заряд из точки A в точку B (рис. 36.2).

\pandocbounded{\includegraphics[keepaspectratio]{images/_page_273_Figure_15.jpeg}}

Рис. 36 2 Потенциал точечного заряда в двумерном случае.
\texttt{\textbackslash{}alpha} --- работа, необходимая для перемещения
единичного заряда из A в B

Стік можна інтерпретувати як від'ємне джерело, так що стік -q створює
потенціал

\begin{verbatim}
u\left(r\right)=-\frac{q}{2\pi}\ln\frac{1}{r}.
\end{verbatim}

Этим завершается обсуждение потенциалов точечных зарядов. Теперь мы
готовы к решению неоднородного уравнения методом функции Грина.

\section{Уравнение Пуассона в
круге}\label{ux443ux440ux430ux432ux43dux435ux43dux438ux435-ux43fux443ux430ux441ux441ux43eux43dux430-ux432-ux43aux440ux443ux433ux435}

Займемся решением очень важной задачи

(36.2)

\begin{verbatim}
(Y \cup \Pi)
\end{verbatim}

\texttt{\textquotesingle{}u\_\{rr\}\ +\ \textbackslash{}frac\{1\}\{r\}\ u\_r\ +\ \textbackslash{}frac\{1\}\{r\^{}2\}\ u\_\{\textbackslash{}theta\textbackslash{}theta\}\ =\ f(r,\ \textbackslash{}theta),\ \textbackslash{}quad\ 0\ \textless{}\ r\ \textless{}\ 1,\ \textbackslash{}quad\ 0\ \textbackslash{}le\ \textbackslash{}theta\ \textbackslash{}le\ 2\textbackslash{}pi,\textquotesingle{}\ \textquotesingle{}(\textbackslash{}Gamma\ Y)\textquotesingle{}\ \textquotesingle{}u(1,\ \textbackslash{}theta)\ =\ 0,\ \textbackslash{}quad\ 0\ \textbackslash{}le\ \textbackslash{}theta\ \textbackslash{}\ le\ 2\textbackslash{}pi.\textquotesingle{}\ **М**етод\ функции\ Грина\ (метод\ функции\ источника)\ состоит\ из\ двух\ шагов.}

\begin{enumerate}
\def\labelenumi{\arabic{enumi}.}
\item
  Помещаем в точку
  \texttt{(\textbackslash{}rho,\ \textbackslash{}varphi)} точечный
  единичный заряд и находим потенциал
  \texttt{G(r,\ \textbackslash{}theta,\ \textbackslash{}rho,\ \textbackslash{}varphi)}
  , который создается этим зарядом и обращается в нуль на границе.
\item
  Суммируем индивидуальные отклики
  \texttt{G(r,\ \textbackslash{}theta,\ \textbackslash{}rho,\ \textbackslash{}phi)}
  , взвешенные правой частью (плотностью зарядов)
  \texttt{f(r,\ \textbackslash{}theta)} , по всему кругу. В результате
  получаем искомое решение
\end{enumerate}

\begin{verbatim}
u\left(r,\;\theta\right) = \int_{0}^{2\pi} \int_{0}^{1} G\left(r,\;\theta,\;\rho,\;\varphi\right) f\left(\rho,\;\varphi\right) \rho \; d\rho \; d\varphi.
\end{verbatim}

Найдем теперь функцию источника для нашей задачи.

\section{\texorpdfstring{Определение функции источника
\texttt{G(r,\ \textbackslash{}theta,\ \textbackslash{}rho,\ \textbackslash{}phi)}
Замінимо праву частину \texttt{f(r,\textbackslash{}theta)} рівняння
точковим джерелом величини 1, розташованим у деякій довольно обраній
точці \texttt{(\textbackslash{}rho,\ \textbackslash{}varphi)} .
Математично точкове джерело представлене дельта-функцією
\texttt{\textbackslash{}delta(r-\textbackslash{}rho,\ \textbackslash{}theta-\textbackslash{}varphi)\^{}1}
). Ми вважаємо, що дельта-функція дорівнює нулю всюди, крім точки
\texttt{(\textbackslash{}rho,\ \textbackslash{}varphi)} , де
розташований одиничний заряд. У мові сил дельта-функцію можна
інтерпретувати як точкова одиниця сили, прикладена до точки
\texttt{(\textbackslash{}rho,\ \textbackslash{}varphi)} . Наше завдання
--- знайти потенціал точкового заряду, якщо відомо, що він дорівнює нулю
на межі. Функція, яка задовольняє ці умови, називається функцією Гріна
або функцією джерела. Вона дорівнює \emph{відповіді} системи в точці
\texttt{(r,\ \textbackslash{}theta)} на \emph{джерело}, розташоване в
точці \texttt{(\textbackslash{}rho,\ \textbackslash{}varphi)} .
Складність розв'язання цієї проблеми полягає в тому, що бажаний
потенціал має зрівнятися з нулем на кордоні. Якби не ця умова, потенціал
можна було б легко знайти, адже ми вже знаємо,
що}{Определение функции источника G(r, \textbackslash theta, \textbackslash rho, \textbackslash phi) Замінимо праву частину f(r,\textbackslash theta) рівняння точковим джерелом величини 1, розташованим у деякій довольно обраній точці (\textbackslash rho, \textbackslash varphi) . Математично точкове джерело представлене дельта-функцією \textbackslash delta(r-\textbackslash rho, \textbackslash theta-\textbackslash varphi)\^{}1 ). Ми вважаємо, що дельта-функція дорівнює нулю всюди, крім точки (\textbackslash rho, \textbackslash varphi) , де розташований одиничний заряд. У мові сил дельта-функцію можна інтерпретувати як точкова одиниця сили, прикладена до точки (\textbackslash rho, \textbackslash varphi) . Наше завдання --- знайти потенціал точкового заряду, якщо відомо, що він дорівнює нулю на межі. Функція, яка задовольняє ці умови, називається функцією Гріна або функцією джерела. Вона дорівнює відповіді системи в точці (r, \textbackslash theta) на джерело, розташоване в точці (\textbackslash rho, \textbackslash varphi) . Складність розв'язання цієї проблеми полягає в тому, що бажаний потенціал має зрівнятися з нулем на кордоні. Якби не ця умова, потенціал можна було б легко знайти, адже ми вже знаємо, що}}\label{ux43eux43fux440ux435ux434ux435ux43bux435ux43dux438ux435-ux444ux443ux43dux43aux446ux438ux438-ux438ux441ux442ux43eux447ux43dux438ux43aux430-gr-theta-rho-phi-ux437ux430ux43cux456ux43dux438ux43cux43e-ux43fux440ux430ux432ux443-ux447ux430ux441ux442ux438ux43dux443-frtheta-ux440ux456ux432ux43dux44fux43dux43dux44f-ux442ux43eux447ux43aux43eux432ux438ux43c-ux434ux436ux435ux440ux435ux43bux43eux43c-ux432ux435ux43bux438ux447ux438ux43dux438-1-ux440ux43eux437ux442ux430ux448ux43eux432ux430ux43dux438ux43c-ux443-ux434ux435ux44fux43aux456ux439-ux434ux43eux432ux43eux43bux44cux43dux43e-ux43eux431ux440ux430ux43dux456ux439-ux442ux43eux447ux446ux456-rho-varphi-.-ux43cux430ux442ux435ux43cux430ux442ux438ux447ux43dux43e-ux442ux43eux447ux43aux43eux432ux435-ux434ux436ux435ux440ux435ux43bux43e-ux43fux440ux435ux434ux441ux442ux430ux432ux43bux435ux43dux435-ux434ux435ux43bux44cux442ux430-ux444ux443ux43dux43aux446ux456ux454ux44e-deltar-rho-theta-varphi1-.-ux43cux438-ux432ux432ux430ux436ux430ux454ux43cux43e-ux449ux43e-ux434ux435ux43bux44cux442ux430-ux444ux443ux43dux43aux446ux456ux44f-ux434ux43eux440ux456ux432ux43dux44eux454-ux43dux443ux43bux44e-ux432ux441ux44eux434ux438-ux43aux440ux456ux43c-ux442ux43eux447ux43aux438-rho-varphi-ux434ux435-ux440ux43eux437ux442ux430ux448ux43eux432ux430ux43dux438ux439-ux43eux434ux438ux43dux438ux447ux43dux438ux439-ux437ux430ux440ux44fux434.-ux443-ux43cux43eux432ux456-ux441ux438ux43b-ux434ux435ux43bux44cux442ux430-ux444ux443ux43dux43aux446ux456ux44e-ux43cux43eux436ux43dux430-ux456ux43dux442ux435ux440ux43fux440ux435ux442ux443ux432ux430ux442ux438-ux44fux43a-ux442ux43eux447ux43aux43eux432ux430-ux43eux434ux438ux43dux438ux446ux44f-ux441ux438ux43bux438-ux43fux440ux438ux43aux43bux430ux434ux435ux43dux430-ux434ux43e-ux442ux43eux447ux43aux438-rho-varphi-.-ux43dux430ux448ux435-ux437ux430ux432ux434ux430ux43dux43dux44f-ux437ux43dux430ux439ux442ux438-ux43fux43eux442ux435ux43dux446ux456ux430ux43b-ux442ux43eux447ux43aux43eux432ux43eux433ux43e-ux437ux430ux440ux44fux434ux443-ux44fux43aux449ux43e-ux432ux456ux434ux43eux43cux43e-ux449ux43e-ux432ux456ux43d-ux434ux43eux440ux456ux432ux43dux44eux454-ux43dux443ux43bux44e-ux43dux430-ux43cux435ux436ux456.-ux444ux443ux43dux43aux446ux456ux44f-ux44fux43aux430-ux437ux430ux434ux43eux432ux43eux43bux44cux43dux44fux454-ux446ux456-ux443ux43cux43eux432ux438-ux43dux430ux437ux438ux432ux430ux454ux442ux44cux441ux44f-ux444ux443ux43dux43aux446ux456ux454ux44e-ux433ux440ux456ux43dux430-ux430ux431ux43e-ux444ux443ux43dux43aux446ux456ux454ux44e-ux434ux436ux435ux440ux435ux43bux430.-ux432ux43eux43dux430-ux434ux43eux440ux456ux432ux43dux44eux454-ux432ux456ux434ux43fux43eux432ux456ux434ux456-ux441ux438ux441ux442ux435ux43cux438-ux432-ux442ux43eux447ux446ux456-r-theta-ux43dux430-ux434ux436ux435ux440ux435ux43bux43e-ux440ux43eux437ux442ux430ux448ux43eux432ux430ux43dux435-ux432-ux442ux43eux447ux446ux456-rho-varphi-.-ux441ux43aux43bux430ux434ux43dux456ux441ux442ux44c-ux440ux43eux437ux432ux44fux437ux430ux43dux43dux44f-ux446ux456ux454ux457-ux43fux440ux43eux431ux43bux435ux43cux438-ux43fux43eux43bux44fux433ux430ux454-ux432-ux442ux43eux43cux443-ux449ux43e-ux431ux430ux436ux430ux43dux438ux439-ux43fux43eux442ux435ux43dux446ux456ux430ux43b-ux43cux430ux454-ux437ux440ux456ux432ux43dux44fux442ux438ux441ux44f-ux437-ux43dux443ux43bux435ux43c-ux43dux430-ux43aux43eux440ux434ux43eux43dux456.-ux44fux43aux431ux438-ux43dux435-ux446ux44f-ux443ux43cux43eux432ux430-ux43fux43eux442ux435ux43dux446ux456ux430ux43b-ux43cux43eux436ux43dux430-ux431ux443ux43bux43e-ux431-ux43bux435ux433ux43aux43e-ux437ux43dux430ux439ux442ux438-ux430ux434ux436ux435-ux43cux438-ux432ux436ux435-ux437ux43dux430ux454ux43cux43e-ux449ux43e}

См.: Владимиров В. С. Обобщенные функции в математической физике --- М.:
Наука, 1976. --- Прим ред.

потенциал точечного заряда равен

\begin{verbatim}
\frac{1}{2\pi} \ln \left( \frac{1}{r} \right)
\end{verbatim}

,

где r --- расстояние от заряда до точки наблюдения.

Искомый потенциал может иметь следующий физический смысл:

\begin{enumerate}
\def\labelenumi{\arabic{enumi}.}
\tightlist
\item
  В теории теплопроводности --- это равновесная температура внутри
  круга, если тепловой источник размещается в точке (о. ф),
\end{enumerate}

а на границе поддерживается ну-

левая температура.

\begin{itemize}
\tightlist
\item
  \begin{enumerate}
  \def\labelenumi{\arabic{enumi}.}
  \setcounter{enumi}{1}
  \tightlist
  \item
    Він описує відхилення мембрани від положення рівноваги, якщо
    відхилення на межі дорівнюють нулю, а в точці (p.~f) мембрана
    зігнута до дуже великої висоти.
  \end{enumerate}
\item
  \begin{enumerate}
  \def\labelenumi{\arabic{enumi}.}
  \setcounter{enumi}{2}
  \tightlist
  \item
    В електростатиці це розподіл потенціалу всередині кола, якщо
    точковий позитивний заряд розташований у точці
    \texttt{(\textbackslash{}rho,\ \textbackslash{}phi)} , а коло
    заземлене, а його потенціал підтримується на нулі.
  \end{enumerate}
\end{itemize}

\pandocbounded{\includegraphics[keepaspectratio]{images/_page_275_Picture_11.jpeg}}

Единичная окружность РИС. 36 3. Функция Гряна 6 (г. 0. р, ф) точечного
источника, расположенного в точке (р, ф).

Теперь найдем функцию Грина; она будет похожа на функцию, изображенную
на рис. 36.3.

\section{Построение
решения}\label{ux43fux43eux441ux442ux440ux43eux435ux43dux438ux435-ux440ux435ux448ux435ux43dux438ux44f}

ШАГ 1. Поскольку функция

\begin{verbatim}
\frac{1}{2\pi} \ln \frac{1}{R}
\end{verbatim}

является потенциалом поля в точке
\texttt{p\ =\ (r,\ \textbackslash{}varphi)\ \textbackslash{}theta)} ,
которое создано единичным точечным зарядом, помещенным в точку
\texttt{Q\ =\ (\textbackslash{}rho,\ \textbackslash{}phi)} ,
\texttt{(R-paccmonnue\textbackslash{}\ between\textbackslash{}\ points\textbackslash{}\ P\textbackslash{}\ and\textbackslash{}\ Q)}
единственное, что нам осталось сделать---так изменить эту функцию, чтобы
она обратилась в нуль на границе.

ЩАГ 2. Из эксперимента известно, что эквипотенциальные линии точечных
зарядов (положительных или отрицательных) представляют собой окружности
(рис. 36.4).

Основна ідея побудови функції Гріна така: потрібно розмістити інший
(від'ємний) заряд поза колом так, щоб загальний потенціал на колі r = 1
був сталим. Для гему можна відняти цю константу з потенціалу і отримати
нульовий потенціал на межі. Очевидно. Таким чином знайдений потенціал
матиме всі властивості функції
\texttt{G(r,\ \textbackslash{}theta,\ \textbackslash{}rho,\ \textbackslash{}phi)}
. Залишається незрозумілим, у якій точці за межами кола слід розмістити
негативний заряд, щоб потенціал на колі став сталим. Можливо порівняльно
P И С

\pandocbounded{\includegraphics[keepaspectratio]{images/_page_276_Picture_3.jpeg}}

. 36 4. Потенциал поля, образованного двумя противоположно заряженными
частицами; a, b--- линии (окружности) постоянного потенциала.

\pandocbounded{\includegraphics[keepaspectratio]{images/_page_276_Picture_5.jpeg}}

\textbf{РИС.} 36.5. Заряды Q в \texttt{\textbackslash{}overline\{Q\}}
дают постоянный потенциал на окружности r=1.

Легко побачити, що якщо негативний заряд розміщений у точці
\texttt{\textbackslash{}bar\{Q\}\ =\ (\textbackslash{}bar\{\textbackslash{}rho\},\ \textbackslash{}bar\{\textbackslash{}phi\})\ =\ (1/\textbackslash{}rho,\ \textbackslash{}phi)}
, то повний потенціал

\begin{verbatim}
u(r, \theta) = \frac{1}{2\pi} \ln 1/R - \frac{1}{2\pi} \ln 1/\overline{R}
\end{verbatim}

буде сталою на колі r=1 Тут відстані від зарядів до точки (g, c)
позначаються R і \texttt{\textbackslash{}overline\{R\}} (див. рис.
36.5). Можна показати, що величина потенціалу на колі r = 1 визначається
виразом

\begin{verbatim}
-\frac{1}{2\pi}\ln\rho
\end{verbatim}

(положительная константа).

С учетом этого результата строим функцию Грина

(36.3)

\begin{verbatim}
G(r, \theta, \rho, \phi) = \frac{1}{2\pi} \ln 1/R - \frac{1}{2\pi} \ln 1/\overline{R} + \frac{1}{2\pi} \ln \rho,
\end{verbatim}

Потенциал положентельного заряда, расположенного в тенциал границы
постоянного в тенциал границы постоянного в тенциал границы постоянного
в тенциал границы постоянного в тенциал границы постоянного в тенциал
границы постоянного в тенциал границы постоянного в тенциал границы
постоянного в тенциал границы постоянного в тенциал границы постоянного
в тенциал границы постоянного в тенциал границы постоянного в тенциал
границы постоянного в тенциал границы постоянного в тенциал границы
постоянного в тенциал границы постоянного в тенциал границы постоянного
в тенциал границы постоянного в тенциал границы постоянного в тенциал
границы постоянного в тенциал границы постоянного в тенциал границы
постоянного в тенциал границы постоянного в тенциал границы постоянного
в тенциал границы постоянного в тенциал границы постоянного в тенциал
границы постоянного в тенциал границы постоянного в тенциал границы
постоянного в тенциал границы постоянного в тенциал границы постоянного
в тенциал границы постоянного в тенциал границы постоянного в тенциал
границы постоянного в тенциал границы постоянного в тенциал границы
постоянного в тенциал границы постоянного в тенциал границы постоянного
в тенциал границы постоянного в тенциал границы постоянного в тенциал
границы постоянного в тенциал границы постоянного в тенциал границы
постоянного в тенциал границы постоянного в тенциал границы постоянного
в тенциал границы постоянного в тенциал границы постоянного в тенциал
границы постоянного в тенциал границы постоянного в тенциал границы
постоянного в тенциал границы постоянного в тенциал границы постоянного
в тенциал границы постоянного в тенциал границы постоянного в тенциал
границы постоянного в тенциал границы постоянного в тенциал границы
постоянного в тенциал границы постоянного в тенциал границы постоянного
в тенциал границы постоянного в тенциал границы постоянного в тенциал
границы постоян границы постоян границы постоян границы постоян границы
постоян

где

\begin{verbatim}
R = V \frac{r^2 - 2\rho r \cos(\theta - \phi) + \rho^2}{R},
\overline{R} = V \frac{r^2 - 2\frac{r}{\rho}\cos(\theta - \phi) + 1/\rho^2}{r^2 - 2\frac{r}{\rho}\cos(\theta - \phi) + 1/\rho^2}
\end{verbatim}

Формули, що виражають відстань між двома точками в полярних координатах.
Щоб знайти розв'язок початкової проблеми, потрібно накласти всі
імпульсні характеристики системи. Отже, переходимо до останнього кроку.

КРОК 3. Суперпозиція імпульсних характеристик Цей крок досить простий. З
відношення

\begin{verbatim}
u(r, \theta) = \int_{0}^{2\pi} \int_{0}^{1} G(r, \theta, \rho, \varphi) f(\rho, \varphi) \rho d\rho d\varphi
\end{verbatim}

получаем

(36.4)

\begin{verbatim}
u(r, \theta) = \frac{1}{2\pi} \int_{0}^{2\pi} \int_{0}^{1} \ln (\rho \overline{R}/R) f(\rho, \varphi) \rho d\rho d\varphi.
\end{verbatim}

\begin{verbatim}
Формула (36.4) виражає розв'язок задачі Діріхле для рівняння Пуассона всередині кола через функцію Гріна цієї задачі. Якщо щільність заряду  f(r, \theta)  відома, то інтеграл можна обчислити, наприклад, чисельно.
\end{verbatim}

\section{ЗАУВАЖЕННЯ}\label{ux437ux430ux443ux432ux430ux436ux435ux43dux43dux44f-3}

\begin{enumerate}
\def\labelenumi{\arabic{enumi}.}
\tightlist
\item
  Задачу
\end{enumerate}

\begin{verbatim}
\begin{array}{llll} (\text{УЧ\Pi}) & \Delta u = 0, & 0 < r < 1, & 0 \leqslant \theta < 2\pi, \\ (\text{ГУ}) & u(1, \theta) = g(\theta), & 0 \leqslant \theta \leqslant 2\pi, \end{array}
\end{verbatim}

також можна розв'язати за допомогою методу \emph{функції Гріна}. У цьому
випадку рішення фіксується у формі

\begin{verbatim}
u(r, \theta) = \int_{0}^{2\pi} \frac{\partial G}{\partial r} (r, \theta, 1, \varphi) g(\varphi) d\varphi.
\end{verbatim}

\begin{verbatim}
Останнє відношення можна надати більш обчислювально корисну форму, явно виставляючи вираз у  \partial G/\partial r  . Внаслідок цього ми отримуємо інтегральну формулу Пуассона, знайдену раніше в лекції 33.
\end{verbatim}

(36.5)

\begin{verbatim}
u(r, \theta) = \frac{1}{2\pi} \int_{0}^{2\pi} \left[ \frac{1 - r^2}{1 - 2r\cos(\theta - \varphi) + r^2} \right] g(\varphi) d\varphi.
\end{verbatim}

\begin{enumerate}
\def\labelenumi{\arabic{enumi}.}
\setcounter{enumi}{1}
\tightlist
\item
  Решение общей задачи Дирихле
\end{enumerate}

\begin{verbatim}
\begin{array}{llll} (\text{УЧП}) & \Delta u = f(r, \theta), & 0 < r < 1, & 0 \le \theta < 2\pi, \\ (\text{ГУ}) & u(1, \theta) = g(\theta), & 0 < \theta \le 2\pi, \end{array}
\end{verbatim}

находится как сумма решений (36.4) и (36.5).

\begin{enumerate}
\def\labelenumi{\arabic{enumi}.}
\setcounter{enumi}{2}
\item
  Метод функцій Гріна дозволяє отримати розв'язки багатьох задач у
  областях різних форм. Однак для кожної області і (точніше, для кожного
  оператора зліва від граничної умови) для кожного рівняння потрібно
  знайти власну функцію Гріна, і це не завжди легко зробити.
\item
  Щоб дійсно знайти розв'язок за формулою (36.4), у більшості випадків
  необхідно обчислити інтеграли на комп'ютері.
\end{enumerate}

\#TASKS: Знайти потенціал точкового джерела у тривимірному просторі.

\begin{verbatim}
2. Знайти функцію Гріна  G(x, y, \xi, \eta)  для рівняння Лапласа у верхній півплощині y > 0. Іншими словами, знайти потенціал у (x, y) у верхній півплощині, якщо заряд розташований у  (\xi, \eta)  , а потенціал на межі y=0 дорівнює нулю (див. рисунок).
\end{verbatim}

\pandocbounded{\includegraphics[keepaspectratio]{images/_page_278_Figure_11.jpeg}}

\begin{verbatim}
ПРИМІТКА. Якщо від'ємний заряд розмістити в точці  \overline{Q} = (\xi,  -\eta)  , очевидно, що негативний потенціал на прямій y = 0 буде нульовим. Те саме те, що функція Усмішки залежить від потенціалу цих двох зарядів.
\end{verbatim}

\begin{verbatim}
- 3. Використовуючи результати задачі 2, знайдемо розв'язок рівняння Пуассона  \Delta u = -k  у верхній півплощині з умовою нульової межі.
\end{verbatim}

\begin{itemize}
\tightlist
\item
  \begin{enumerate}
  \def\labelenumi{\arabic{enumi}.}
  \setcounter{enumi}{3}
  \tightlist
  \item
    Як би ви побудували функцію Зеленого для першого квадранта
    x\textgreater0,~y\textgreater0 ?
  \end{enumerate}
\item
  \begin{enumerate}
  \def\labelenumi{\arabic{enumi}.}
  \setcounter{enumi}{4}
  \tightlist
  \item
    Інший підхід до розв'язання рівняння Пуассона такий. Припустимо, ви
    хочете розв'язати проблему
  \end{enumerate}
\end{itemize}

(УЧП)

\begin{verbatim}
\Delta u = 1
,  0 < r < 1  ,  0 \le \theta < 2\pi  ,  (\Gamma Y)   u(1, \theta) = \sin \theta  ,  0 \le \theta \le 2\pi  .
\end{verbatim}

Сначала попытайтесь найти частное решение уравнения с помощью
подстановки

\begin{verbatim}
u_p = (r, \theta) = Ar^2
\end{verbatim}

.

\begin{verbatim}
Потім підставляємо цей вираз у рівняння Пуассона і визначаємо константу  A_{\cdot}  числа  u_p  , що знаходиться в частині розв'язку, необхідно пожартувати поза розв'язком з огляду на  u=w+u_p  . Яка межа розуму, яку можна задовольнити функцією w? Знайте  w(r, \theta)  . Як вирішити проблему  u(r, \theta)  ? Зачекай. Retelno proanalizite vidpovid'. Зробіть перетин пеніса шкіри.
\end{verbatim}

\section{Частина 5. ЧИСЕЛЬНІ ТА ПРИБЛИЗНІ
МЕТОДИ}\label{ux447ux430ux441ux442ux438ux43dux430-5.-ux447ux438ux441ux435ux43bux44cux43dux456-ux442ux430-ux43fux440ux438ux431ux43bux438ux437ux43dux456-ux43cux435ux442ux43eux434ux438}

Лекция 37

\section{ЧИСЕЛЬНІ РОЗВ'ЯЗКИ (ЕЛІПТИЧНІ
ЗАДАЧІ)}\label{ux447ux438ux441ux435ux43bux44cux43dux456-ux440ux43eux437ux432ux44fux437ux43aux438-ux435ux43bux456ux43fux442ux438ux447ux43dux456-ux437ux430ux434ux430ux447ux456}

МЕТА ЛЕКЦІЇ: Показати, як диференціальні рівняння в похідних похідних
можна звести до систем алгебраїчних рівнянь, якщо похідні з частковими
похідними замінити на наближення з скінченними різницями. Розв'язання
цієї системи алгебраїчних рівнянь ітеративним методом дозволяє отримати
наближене розв'язок

Надається інформація про пакет ELLPACK, розроблений для розв'язання
еліптичних приладів

задач общего вида на ЭВМ.

На сьогодні ми ознайомилися з кількома методами розв'язання лінійних
диференціальних рівнянь у похідних похідних. Однак більшість рівнянь, з
якими ми працювали, були дуже простими. Форми регіонів, у яких
розв'язувалися проблеми, також були простими. Однак багато завдань не
можна спростити настільки, щоб їх зводити до певного набору шаблонних
завдань. Такі задачі слід розв'язувати приблизно чисельними методами.
Прогрес у сфері високопродуктивних комп'ютерів призвів до створення
нових чисельних методів. Сьогодні такі нелінійні проблеми гідродинаміки,
теорії пружності та теорії потенціалу вже розв'язані, про які ніхто
навіть не думав десять років тому.

Під загальною назвою «чисельні методи» поєднуються кілька різних
підходів до розв'язання задач. Детальний опис усіх цих підходів можна
знайти у книзі {[}4{]} з рекомендованого списку джерел. У цій та
наступних двох лекціях ми покажемо, як розв'язувати еліптичні,
гіперболічні та параболічні рівняння методом скінченної різниці.

Спершу ознайомимося з поняттям скінченних різниць, а потім покажемо, як
їх використовувати для розв'язання задачі Діріхле в квадраті.

\section{Апроксимації скінченних
різниць}\label{ux430ux43fux440ux43eux43aux441ux438ux43cux430ux446ux456ux457-ux441ux43aux456ux43dux447ux435ux43dux43dux438ux445-ux440ux456ux437ux43dux438ux446ux44c}

Вспомним ряд Тейлора для функции \texttt{f(x)}:

\begin{verbatim}
f(x+h) = f(x) + f'(x)h + \frac{f''(x)}{2!}h^2 + \dots
\end{verbatim}

Якщо ми розірвемо цю суперечку на другому терміні, отримаємо

\begin{verbatim}
f(x+h) \cong f(x) + f'(x) h
\end{verbatim}

.

Откуда

\begin{verbatim}
(37.1) f'(x) \cong \frac{f(x+h)-f(x)}{h}.
\end{verbatim}

Вираз праворуч називається похідною правої різниці. Вона апроксимує
першу похідну f'(x) у точці x.

\begin{verbatim}
В разложении Тейлора функции  f(\mathbf{r})  можно заменить h на -h и получить левую разностную производную
\end{verbatim}

(37.2)

\begin{verbatim}
f'(x) \cong \frac{f(x) - f(x - h)}{h}
\end{verbatim}

.

\begin{verbatim}
Buyutan  f(x - h) \cong f(x) - f'(x) h  \nif  f(x + h) \cong f( x) + f'(x) h  ,
\end{verbatim}

получаем центральную разностную производную

(37.3)

\begin{verbatim}
f'(x) \simeq \frac{1}{2h} [f(x+h) - f(x-h)].
\end{verbatim}

Якщо в ряду Тейлора залишиться ще один член, то центральну різницю
похідну для наближення f'\,'(x) можна отримати точно таким самим
способом

(37.4)

\begin{verbatim}
f''(x) \cong \frac{1}{h^2} [f(x+h) - 2f(x) + f(x-h)].
\end{verbatim}

Тепер ми можемо розширити поняття наближення скінченних різниць до
\emph{часткових похідних}. Якщо виходити з розкладу Тейлора функції двох
змінних

\begin{verbatim}
u(x+h, y) = u(x, y) + u_x(x, y)h + u_{xx}(x, y)\frac{h^2}{2!} + \dots,
u(x-h, y) = u(x, y) - u_x(x, y)h + u_{xx}(x, y)\frac{h^2}{2!} - \dots,
\end{verbatim}

можно получить следующие аппроксимации частных производных:

\begin{verbatim}
\begin{split} u_x(x,\ y) & \cong \frac{u(x+h,\ y)-u(x,\ y)}{h}, \\ u_{xx}(x,\ y) & \cong \frac{1}{h^2} \big[ u(x+h,\ y)-2u(x,\ y)+u(x-h,y) \big], \\ u_u(x,\ y) & \cong \frac{u(x,\ y+k)-u(x,\ y)}{k}, \\ u_{uy}(x,\ y) & \cong \frac{1}{h^2} \big[ u(x,\ y+k)-2u(x,\ y)+u(x,\ y-k) \big]. \end{split}
\end{verbatim}

В этих формулах частные производные аппроксимируются правыми,
центральными и левыми разностными производными, но в данной лекции мы
будем пользоваться только центральными аппроксимациями.

\pandocbounded{\includegraphics[keepaspectratio]{images/_page_282_Figure_5.jpeg}}

Рис. 37.1 Сетка для задачи Дирихле внутри квадрата.

Чтобы познакомиться с основными правилами использования
конечно-разностных аппроксимаций, рассмогрим простую задачу Дирихле.

Решение задачи Дирихле методом конечных разностей

Пусть необходимо решить задачу Дирихле

(ЗЧП)

\begin{verbatim}
u_{xx} + u_{yy} = 0
\end{verbatim}

\texttt{,\ \textquotesingle{}0\ \textless{}\ x\ \textless{}\ 1\textquotesingle{}\ ,\ \textquotesingle{}0\ \textless{}\ y\ \textless{}\ 1\textquotesingle{}\ ,\ \textquotesingle{}u\ =\ 0\textquotesingle{}\ ,\ на\ верхней\ и\ боковых\ сторонах\ квадрата,\ \textquotesingle{}u(x,\ 0)\ =\ \textbackslash{}sin(x\textbackslash{}pi)\textquotesingle{}\ ,\ \textquotesingle{}0\ \textbackslash{}le\ \textbackslash{}le\ x\ \textbackslash{}le\ 1\textquotesingle{}\ ,\ \textquotesingle{}y\ =\ 0\textquotesingle{}\ .}

Построим в плоскости (x, y) сетку, как показано на рис. 37.1. Удобно
(особенно, если мы собираемся применять ЭВМ)

непользовать следующие обозначения:

\begin{verbatim}
\begin{split} u\left(x,\,y\right) &= u_{ij}\,,\\ u\left(x,\,y+k\right) &= u_{i+1,j}\,,\\ u\left(x,\,y-k\right) &= u_{i-1,j}\,,\\ u\left(x-h,\,y\right) &= u_{i,j-1}\,,\\ u\left(x+h,\,y\right) &= u_{i,j+1}\,,\\ u_{x}\left(x,\,y\right) &= \frac{1}{2h}\left(u_{i,j+1}-u_{i,j-1}\right)\,,\\ u_{y}\left(x,\,y\right) &= \frac{1}{2k}\left(u_{i+1,j}-u_{i-1,j}\right)\,,\\ u_{xx}\left(x,\,y\right) &= \frac{1}{h^{2}}\left(u_{i,j+1}-2u_{i,j}+u_{i,j-1}\right)\,,\\ u_{yy}\left(x,\,y\right) &= \frac{1}{h^{2}}\left(u_{i+1,j}-2u_{i,j}+u_{i-1,j}\right)\,. \end{split}
\end{verbatim}

Основная идея решения задачи Дирихле основана на замене частных
производных в уравнении Лапласа

\begin{verbatim}
u_{xx} + u_{yy} = 0
\end{verbatim}

их конечно-разностными аппроксимациями. Проделав это и воспользовавшись
компактными обозначениями, получаем следующее разностное уравнение:

\begin{verbatim}
\Delta u = \frac{1}{h^2} (u_{i,j+1} - 2u_{i,j} + u_{i,j-1}) + \frac{1}{k^2} (u_{i+1,j} - 2u_{i,j} + u_{i-1,j}) = 0.
\end{verbatim}

В случае когда h=k, уравнение Лапласа приводится к виду (37.6)
\texttt{(u\_\{i+1,j\}+u\_\{i-1,j\}+u\_\{i,j+1\}+u\_\{i,j+1\}-4u\_\{ij\})=0.}
Разрешив его относительно \texttt{u\_\{II\}} , получим

\begin{verbatim}
u_{i,j} = \frac{1}{4} (u_{i+1,j} + u_{i-1,j} + u_{i,j+1} + u_{i,j-1}).
\end{verbatim}

Зверніть увагу, що в цьому співвідношенні всі значення
\texttt{u\_\{tj\}} беруть для внутрішніх вузлів сітки. Згідно з останнім
співвідношенням, розв'язок \texttt{u\_\{ij\}} апроксимується середнім
значенням розв'язку у чотирьох сусідніх точках. Відповідно, ми можемо
розробити чисельний метод для розв'язання цієї задачі.

Алгоритм численного решения задачи Дирихле (метод Либманна)

ШАГ 1. Присвоим величинам \texttt{u\_\{ij\}} во внутренних узлах
численное вначение, равное среднему всех значений граничных условий.

КРОК 2. Ми перерахуємо значення у всіх внутрішніх точках сітки,
замінивши старе значення на середнє значення чотирьох суміжних точок. Не
дуже важливо, як буде організовано підрахунок процесу, але зазвичай він
здійснюється рядками (або колонками). Після кількох операцій процес
збігається до приблизного розв'язання задачі. Швидкість збіжності цього
алгоритму низька, але її можна збільшити кількома способами.
Рекомендуємо зацікавленим читачам звернутися до книги {[}4{]} зі списку
рекомендованої літератури. \# ЗАУВАЖЕННЯ 1. Если взять сетку с чегырьмя
внутренними точками (m=n=4), то система (37.6) запишется в виде

\begin{verbatim}
-4u_{22} + 0 + \sin(\pi/3) + u_{23} + u_{32} = 0,
-4u_{23} + u_{22} + \sin(2\pi/3) + 0 + u_{33} = 0,
-4u_{32} + 0 + u_{22} + u_{33} + 0 = 0,
-4u_{33} + u_{32} + u_{23} + 0 + 0 = 0.
\end{verbatim}

Це потрібно вирішувати у зв'язку з \texttt{u\_\{22\}} ,
\texttt{u\_\{21\}} , \texttt{u\_\{32\}} та \texttt{u\_\{33\}} .
Розв'язок такої системи можна знайти за допомогою \emph{ітеративного
методу}. Метод Лібмана також є одним із ітеративних методів.

\begin{itemize}
\tightlist
\item
  \begin{enumerate}
  \def\labelenumi{\arabic{enumi}.}
  \setcounter{enumi}{1}
  \tightlist
  \item
    Якщо зменшити кроки сітки на h і k (кількість вузлів сітки
    збільшиться), то отримаємо систему рівнянь типу (37,7), але більших
    розмірів.
  \end{enumerate}
\item
  \begin{enumerate}
  \def\labelenumi{\arabic{enumi}.}
  \setcounter{enumi}{2}
  \tightlist
  \item
    Систему рівнянь (37.7) можна переписати у матричній формі
  \end{enumerate}
\end{itemize}

\begin{verbatim}
\begin{bmatrix} -4 & 1 & 1 \\ 1 & -4 & 0 & 1 \\ 1 & 0 & -4 & 1 \\ 0 & 1 & 1 & -4 \end{bmatrix}
\begin{bmatrix} u_{22} \\ u_{23} \\ u_{32} \\ u_{33} \end{bmatrix}
=
\begin{bmatrix} -\sin{(\pi/3)} \\ -\sin{(2\pi/3)} \\ 0 \\ 0 \end{bmatrix}
=
\begin{bmatrix} -0.86 \\ -0.86 \\ 0 \\ 0 \end{bmatrix}
\end{verbatim}

.

Зі збільшенням кількості рівнянь (наприклад, до 1000) матриця
коефіцієнтів стає розрідженою, тобто містить багато нулів. Системи з
розрідженими матрицями розв'язуються спеціальними методами. Зазвичай
використовуються ітеративні методи, такі як метод Якобі, метод
Гаусса-Зайделя або послідовний метод верхньої релаксації.

\begin{enumerate}
\def\labelenumi{\arabic{enumi}.}
\setcounter{enumi}{3}
\item
  При розв'язанні задачі Неймана ті \emph{похідні}, що входять до
  граничної умови, також слід замінити різницею апроксимаціями.
\item
  Таким чином ви зможете вирішувати проблеми
\end{enumerate}

\begin{enumerate}
\def\labelenumi{\alph{enumi})}
\tightlist
\item
\end{enumerate}

\begin{verbatim}
u_{xx} + u_{yy} = f(x, y)
\end{verbatim}

\begin{enumerate}
\def\labelenumi{\arabic{enumi})}
\setcounter{enumi}{5}
\tightlist
\item
\end{enumerate}

\begin{verbatim}
xu_{xx} + u_{yy} + 2u = \sin(x - y)
\end{verbatim}

\begin{enumerate}
\def\labelenumi{\Alph{enumi})}
\setcounter{enumi}{1}
\tightlist
\item
\end{enumerate}

\begin{verbatim}
\sin xu_{xx} + u_{xy} + 3u = 0
\end{verbatim}

(гетерогенні рівняння з змінними коефіцієнтами), (рівняння зі змінними
коефіцієнтами).

\begin{enumerate}
\def\labelenumi{\arabic{enumi}.}
\setcounter{enumi}{5}
\tightlist
\item
  Якщо область, в якій розв'язується задача, має неправильну форму, її
  можна покрити сіткою і наближити розв'язок у точках, найближчих до
  межі, інтерполюючи граничні умови. Після цього задачу вирішують
  звичайним способом (див. рисунок 37.2).
\end{enumerate}

\pandocbounded{\includegraphics[keepaspectratio]{images/_page_285_Figure_8.jpeg}}

□ новые ГУ находятся интерполяцией ГУ на кривой Рис. 37 2.

\begin{itemize}
\tightlist
\item
  \begin{enumerate}
  \def\labelenumi{\arabic{enumi}.}
  \setcounter{enumi}{6}
  \tightlist
  \item
    Деякі журнали публікують тексти програм для розв'язання
    диференціальних рівнянь у похідних на комп'ютерах. Нижче наведено
    короткий список цих журналів:
  \end{enumerate}
\item
  \begin{enumerate}
  \def\labelenumi{\alph{enumi})}
  \tightlist
  \item
    Транзакції ACM на математичному програмному забезпеченні,
  \end{enumerate}
\item
  \begin{enumerate}
  \def\labelenumi{\arabic{enumi})}
  \setcounter{enumi}{5}
  \tightlist
  \item
    Комп'ютерний журнал,
  \end{enumerate}
\item
  \begin{enumerate}
  \def\labelenumi{\Alph{enumi})}
  \setcounter{enumi}{1}
  \tightlist
  \item
    Numerische Mathematik,
  \end{enumerate}
\item
  \begin{enumerate}
  \def\labelenumi{\alph{enumi})}
  \setcounter{enumi}{17}
  \tightlist
  \item
    БІТ
  \end{enumerate}
\end{itemize}

Крім того, нещодавно було розроблено пакет застосунків ELLPACK для
розв'язання досить широкого класу задач на лінії прикордонних ліній. Цей
пакет дозволить вам розв'язувати різноманітні дво- та тривимірні задачі
у

різні координатні системи для довільних меж і спільних граничних умов.
Користувачу пропонується широкий спектр різних методів розв'язання 1).

\section{\texorpdfstring{ЗАДАЧІ - 1 Отримаємо наближення (39,4) для
другої похідної f'\,`(x) f'\,'(x) = \{\texttt{f(x+h)} 2f(x) +
\texttt{f(x-h)}\}/h\^{}2
.}{ЗАДАЧІ - 1 Отримаємо наближення (39,4) для другої похідної f'\,`(x) f'\,'(x) = \{f(x+h) 2f(x) + f(x-h)\}/h\^{}2 .}}\label{ux437ux430ux434ux430ux447ux456---1-ux43eux442ux440ux438ux43cux430ux454ux43cux43e-ux43dux430ux431ux43bux438ux436ux435ux43dux43dux44f-394-ux434ux43bux44f-ux434ux440ux443ux433ux43eux457-ux43fux43eux445ux456ux434ux43dux43eux457-fx-fx-fxh-2fx-fx-hh2-.}

\begin{itemize}
\tightlist
\item
  \begin{enumerate}
  \def\labelenumi{\arabic{enumi}.}
  \setcounter{enumi}{1}
  \tightlist
  \item
    Виконайте дві ітерації задачі Діріхле (37.5) за ітеративним методом
    Лібмана.
  \end{enumerate}
\item
  \begin{enumerate}
  \def\labelenumi{\arabic{enumi}.}
  \setcounter{enumi}{2}
  \tightlist
  \item
    До якої алгебранічної системи буде зведена задача Діріхле для
    рівнянь Пуассона в квадраті
  \end{enumerate}
\end{itemize}

\texttt{(УЧ\textbackslash{}Pi)}

\begin{verbatim}
u_{xx} + u_{yy} = f(x,y), \quad 0 < x < 1, \quad 0 < y < 1, (\Gamma Y) u(x,y) = g(x,y) на границе, если ее решать ме-
\end{verbatim}

Скінченні відмінності? 4. Розв'язати задачу 3, якщо

(УЧП)

\begin{verbatim}
u_{xx} + u_{yy} + 2u = 0
\end{verbatim}

, 0 \textless{} x \textless{} 1 , 0 \textless{} y \textless{} 1 , (ГУ)
\texttt{u(x,\ y)} = \texttt{g(x,\ y)} на границе.

\begin{enumerate}
\def\labelenumi{\arabic{enumi}.}
\setcounter{enumi}{4}
\tightlist
\item
  Как бы вы решали задачу Неймана внутри квадрата
\end{enumerate}

(УЧП)

\begin{verbatim}
u_{xx} + u_{yy} = 0
\end{verbatim}

, 0 \textless{} x \textless{} 1 , 0 \textless{} y \textless{} 1 , на
верхней, нижней и левой стороне квадрата,

\begin{verbatim}
\begin{cases} u = 0 & \text{квадрата,} \\ \frac{\partial u}{\partial x}(1, y) = 1, & 0 \leqslant y \leqslant 1, \end{cases}
\end{verbatim}

методом конечных разностей?

\begin{enumerate}
\def\labelenumi{\arabic{enumi}.}
\setcounter{enumi}{5}
\tightlist
\item
  Постройте блок-схему алеоритма решения задачи Дирихле в квадрате
\end{enumerate}

\begin{verbatim}
(\text{УЧП})
\end{verbatim}

\begin{verbatim}
u_{xx} + u_{yy} = f(x, y), \quad 0 < x < 1, \quad 0 < y < 1, (\text{ГУ}) u(x, y) = g(x, y) на границе,
\end{verbatim}

\begin{itemize}
\item
  з довільною кількістю вузлів сітки. Якщо ви знаєте мову програмування,
  напишіть програму для виконання цих обчислень.
\item
  \begin{enumerate}
  \def\labelenumi{\arabic{enumi})}
  \tightlist
  \item
    Російською мовою публікуються статті з чисельних методів розв'язання
    диференціальних рівнянь у численних журналах і збірках. Наприклад:
    «Журнал обчислювальної математики та математичної фізики»,
    видавництво «Наука», «Журнал інженерної фізики», видавництво
    Академії наук БРСР, збірки серії «Обчислювальні методи і
    програмування», Видавництво Московського державного університету.
    Багато публікацій, переважно статей, присвячені наборам прикладних
    проблем. Як приклад, можна вказати: Горбунов-Посадов М. М., Карпов
    В. Я, Корягін та ін. Пакет SAFRA. Програмне забезпечення для
    обчислювальних експериментів. У: Пакети прикладних програм.
    Обчислювальний експеримент. Москва, Наука, 1983. Примітка. Переклад.
  \end{enumerate}
\end{itemize}

\section{МЕТАЛЕКЦІЯ ПРО ЯВНІ РІЗНИЦІ СХЕМ: Введіть поняття явної різниці
та покажіть, як її можна використовувати для розв'язання гіперболічних і
параболічних задач. Основна ідея полягає в тому, що після заміни
рівняння
типу}\label{ux43cux435ux442ux430ux43bux435ux43aux446ux456ux44f-ux43fux440ux43e-ux44fux432ux43dux456-ux440ux456ux437ux43dux438ux446ux456-ux441ux445ux435ux43c-ux432ux432ux435ux434ux456ux442ux44c-ux43fux43eux43dux44fux442ux442ux44f-ux44fux432ux43dux43eux457-ux440ux456ux437ux43dux438ux446ux456-ux442ux430-ux43fux43eux43aux430ux436ux456ux442ux44c-ux44fux43a-ux457ux457-ux43cux43eux436ux43dux430-ux432ux438ux43aux43eux440ux438ux441ux442ux43eux432ux443ux432ux430ux442ux438-ux434ux43bux44f-ux440ux43eux437ux432ux44fux437ux430ux43dux43dux44f-ux433ux456ux43fux435ux440ux431ux43eux43bux456ux447ux43dux438ux445-ux456-ux43fux430ux440ux430ux431ux43eux43bux456ux447ux43dux438ux445-ux437ux430ux434ux430ux447.-ux43eux441ux43dux43eux432ux43dux430-ux456ux434ux435ux44f-ux43fux43eux43bux44fux433ux430ux454-ux432-ux442ux43eux43cux443-ux449ux43e-ux43fux456ux441ux43bux44f-ux437ux430ux43cux456ux43dux438-ux440ux456ux432ux43dux44fux43dux43dux44f-ux442ux438ux43fux443}

\begin{verbatim}
u_t = u_{xx}
\end{verbatim}

Її апроксимація скінченної різниці дає формули, які явно виражають
значення розв'язку за один момент часу через значення розв'язку в
попередній момент часу. Таким чином, змішану задачу для параболічного
або гіперболічного рівняння можна розв'язати, послідовно обчислюючи
розв'язки для всіх наступних точок часу.

Явна схема не позбавлена недоліків. Якщо необхідно підвищити точність
наближення похідних, то при зміні сегмента збільшується не лише обсяг
обчислень, а й похибки округлення.

Попередня лекція була присвячена розв'язанню еліптичних задач на межі
лінії (стаціонарних задач). У таких задачах необхідно знайти розв'язок
диференціального рівняння в похідних у заданій області простору, якщо
розв'язок або його похідна задано на межі області визначення Приблизне
розв'язання еліптичних задач зводиться до розв'язку системи алгебраїчних
рівнянь для значень функції у внутрішніх вузлах сітки. Іншими словами,
значення розв'язку у всіх внутрішніх вузлах визначаються одночасно.

\begin{verbatim}
У цій лекції ми розглянемо диференціальні методи розв'язання задач, залежних від часу. Основна ідея полягає в тому, що якщо розв'язок відомий у початковий момент часу, то можливо знайти розв'язок за допомогою схеми підрахунку t=\Delta t,\ 2\Delta t,\ 3\Delta t,\ldots Замінивши часткові похідні часу та просторову змінну на скінченнорізні похідні, можна отримати явні вирази для u_{tj} через значення функції u у попередні моменти часу. Такий процес називається явною схемою підрахунку з тикуванням.
\end{verbatim}

Щоб показати цей метод у дії, розглянемо типову задачу теплопровідності.

\section{Явная схема для уравнения
теплопроводности}\label{ux44fux432ux43dux430ux44f-ux441ux445ux435ux43cux430-ux434ux43bux44f-ux443ux440ux430ux432ux43dux435ux43dux438ux44f-ux442ux435ux43fux43bux43eux43fux440ux43eux432ux43eux434ux43dux43eux441ux442ux438}

Розглянемо проблему теплопровідності у стрижні, початкова температура
якого дорівнює нулю. Нехай температура лівого кінця фіксована, а на
правому кінці відбувається теплообмін із навколишнім середовищем, щоб
тепловий фланець був пропорційний P и с

\pandocbounded{\includegraphics[keepaspectratio]{images/_page_288_Figure_3.jpeg}}

. 38.1. Сетка для уравнения теплопроводности.

\begin{verbatim}
різниця температур між кінцем стрижня і середовищем. Нехай температура середовища задається функцією g\left(t\right) . Іншими словами, ми розв'язуємо проблему
\end{verbatim}

задачу (УЧП)

\begin{verbatim}
u_t = u_{xx}, \quad 0 < x < 1, \quad 0 < t < \infty,
\end{verbatim}

(38.1) (ГУ)

\begin{verbatim}
\begin{cases} u(0, t) = 1, \\ u_x(1, t) = -[u(1, t) - g(t)], \\ u(x, 0) = 0, \quad 0 \le x \le 1. \end{cases}
\end{verbatim}

Щоб розв'язати цю задачу, ми побудуємо прямокутну сітку методом
скінченної різниці, вершини якої визначаються формулами (див. рис. 38.1)

\begin{verbatim}
x_j = jh
,  j = 0, 1, 2, ..., n  ,  y_i = ik  ,  i = 0, 1, 2, ..., m  .
\end{verbatim}

Зверніть увагу, що значення \texttt{u\_\{ij\}} зліва і внизу сітки на
рис. 38.1 відомі з граничних і початкових умов, і наша вадача полягає у
знаходженні решти значень \texttt{u\_\{ij\}} . Щоб розв'язати цю задачу,
ми замінимо часткові похідні рівняння теплопровідності на їхні
наближення з скінченними різницями

\begin{verbatim}
u_{t} = \frac{1}{k} [u(x, t+k) - u(x, t)] = \frac{1}{k} [u_{t+1,j} - u_{t,j}],
u_{xx} = \frac{1}{h^{2}} [u(x+h, t) - 2u(x, t) + u(x-h, t)] = \frac{1}{h^{2}} [u_{t,j+1} - 2_{t,j} + u_{t,j-1}].
\end{verbatim}

Подставим эти выражения в уравнение \texttt{u\_t\ =\ u\_\{xx\}} и
разрешим получившееся уравнение относительно значений функции на верхнем
временном слое. В результате получаем

(38.2)

\begin{verbatim}
u_{t+1,j} = u_{t,j} + \frac{k}{h^2} \left[ u_{t,j+1} - 2u_{t,j} + u_{t,j-1} \right].
\end{verbatim}

Це та формула, яку ми шукаємо, оскільки вона виражає розв'язок у певний
момент часу через розв'язок у попередній момент часу (індекс \emph{i}
означає часову змінну). Фіг. 38.1 підкреслює значення, включені до цієї
формули.

Тепер можна починати розрахунки. Однак спочатку потрібно апроксимувати
похідну в граничній умові праворуч

\begin{verbatim}
a_x(1, t) = -[u(1, t) - g(t)].
\end{verbatim}

В результате аппроксимации получаем

(38.3)

\begin{verbatim}
\frac{1}{h}[u_{t,n} - u_{t,n-1}] = -[u_{t,n} - g_t].
\end{verbatim}

где значения \texttt{g\_i\ =\ g(ik)} известны. Здесь мы заменили
\texttt{u\_x(1,\ t)} левой разностной производной, поскольку правая
разностная производная потребовала бы значений функции за пределами
сетки. Из (38.3) находим

(38.4)

\begin{verbatim}
u_{t,n} = \frac{u_{t,n-1} - hg_t}{1 + h}.
\end{verbatim}

Формулы (38.2) и (38.4) позволяют начать вычисления.

Алгоритм вычислений по явной схеме

\texttt{ШАГ\ 1.\ Находим\ решение\ на\ сеточном\ слое\ \textquotesingle{}t\ =\ \textbackslash{}Delta\ t\textquotesingle{}\ ,\ используя\ явную\ формулу\ (см.\ рис.\ 38.2)}

\begin{verbatim}
u_{2,j} = u_{1,j} + \frac{k}{h^2} [u_{1,j+1} - 2u_{1,j} + u_{1,j-1}], \quad j = 2, 3, \ldots, n-1.
\end{verbatim}

ШАГ 2. Величину \texttt{u\_\{2,n\}} находим по формуле (38.4)

\begin{verbatim}
u_{2,n} = \frac{u_{2,n-1} + hg_2}{1+h}.
\end{verbatim}

\texttt{Після\ виконання\ кроків\ 1\ і\ 2\ ми\ отримуємо\ рішення\ для\ \textquotesingle{}t=\textbackslash{}Delta\ t\textquotesingle{}\ .\ Щоб\ отримати\ розв\textquotesingle{}язок\ у\ \textquotesingle{}t=2\textbackslash{}Delta\ t\textquotesingle{}\ (другий\ рядок\ знизу\ на\ рис.\ 38.2),\ повторіть\ кроки\ 1\ і\ 2,\ рухаючись\ на\ одну\ лінію\ вгору,\ тобто\ збільшуючи\ t\ на\ 1\ і\ використовуючи\ \textquotesingle{}u\_\{t,j\}\textquotesingle{}\ з\ попередньої\ лінії.\ Розв\textquotesingle{}язок\ обчислюється\ таким\ самим\ способом\ у\ наступні\ \textquotesingle{}t=3\textbackslash{}Delta\ t,\textbackslash{}\ 4\textbackslash{}Delta\ t\textbackslash{}dots\textquotesingle{}\ Для\ того\ чтобы\ помочь\ читателю\ в\ проведении\ расчетов\ по\ этому\ методу,\ на\ рис.\ 38.3\ приведена\ достаточно\ подробная\ блоксхема\ алгоритма.\ Те,\ кто\ не\ очень\ хорошо\ знаком\ с\ блок-схемами.}

\pandocbounded{\includegraphics[keepaspectratio]{images/_page_290_Figure_3.jpeg}}

Рис. 38.2. Иллюстрация явной разностной схемы: n--- число узлов сетки
вдоль оси \texttt{x;\textbackslash{}\ m} --- число узлов сетки вдоль оси
\texttt{y;\textbackslash{}\ h=1/(m-1),\textbackslash{}\ k=T/(m-1)} могут
счигать их промежуточным звеном между вычислительными алгоритмами и
подробными программами для ЭВМ. Блок-схемы дают строгое описание
процесса вычислений. \# ЗАУВАЖЕННЯ - Явна схема має серйозний недолік.
Якщо часовий крок достатньо великий порівняно з кроком x, похибки
округлення можуть стати настільки великими, що отриманий розв'язок
втрачає сенс. Відношення кроків у t і x залежить від рівняння та
граничних умов, але в загальному випадку крок у часі має бути значно
меншим за крок у координаті. У книзі {[}4{]} доведено, що для
застосування явної схеми необхідно виконати умову k/h2 ≤ 0,5.
\texttt{-\ 2.\ Існує\ таке\ правило:\ якщо\ зменшити\ кількість\ кроків\ \textquotesingle{}\textbackslash{}Delta\ t\textquotesingle{}\ і\ \textquotesingle{}\textbackslash{}Delta\ x\textquotesingle{}\ ,\ то\ похибка\ у\ наближенні\ часткових\ похідних\ за\ скінченними\ різницями\ також\ зменшиться,\ але\ чим\ менша\ сітка,\ тим\ більше\ обчислень\ потрібно\ виконати,\ і,\ відповідно,\ тим\ більшими\ будуть\ похибки\ округлення.\ Це\ явище\ ілюструється\ на\ рис.\ 38.5.}
- 3. Для гіперболічної задачі

(УЧП)

\begin{verbatim}
u_{tt} = u_{xx}, \quad 0 < x < 1, \quad 0 < t < \infty,
\end{verbatim}

(ГУ)

\begin{verbatim}
\begin{cases} u(0, t) = g_1(t), \\ u(1, t) = g_2(t), \end{cases}
\quad 0 < t < \infty,
\end{verbatim}

(HY)

\begin{verbatim}
\begin{cases} u(x, 0) = \varphi(x), \\ u_t(x, 0) = \psi(x), \end{cases}
\quad 0 \leqslant x \leqslant 1,
\end{verbatim}

Також можливо побудувати явну різницю схеми. Для цього ми апроксимуємо
\texttt{u\_\{tt\}} і \texttt{u\_\{xy\}} центральних різницьких
пропозицій И С

\pandocbounded{\includegraphics[keepaspectratio]{images/_page_291_Figure_4.jpeg}}

. 38.3. Алгоритм решения задачи по явной схеме.

водными

\begin{verbatim}
u_{tt} \simeq \frac{1}{h^2} [u(x, t+k) - 2u(x, t) + u(x, t-k)],
u_{xx} \simeq \frac{1}{h^2} [u(x+h, t) - 2u(x, t) + u(x-h, t)].
\end{verbatim}

а начальное условие--- по схеме

\begin{verbatim}
u_t(x, 0) \approx \frac{1}{k} [u(x, k) - u(x, 0)] = \frac{1}{k} [u(x, k) - \varphi(x)].
\end{verbatim}

В результате для вычисления величины \texttt{u(x,\ t+k)} получаем
следующую явную схему:

(38.5)

\begin{verbatim}
u(x, t+k) = 2u(x, t) - u(x, t-k) + + (k/h)^2 [u(x+h, t) - 2u(x, t) + u(x-h, t)].
\end{verbatim}

З (38.5) видно, що для обчислення розв'язку на заданому часовому шарі
необхідно знати розв'язок на двох попередніх шарах. Отже, щоб почати
рахувати, потрібно використовувати P И С

\pandocbounded{\includegraphics[keepaspectratio]{images/_page_292_Figure_6.jpeg}}

. 38.4. Использование массивов при реализации явной схемы.

начальным условнем на скорость

\begin{verbatim}
\frac{1}{k}\left[u\left(x,\,k\right)-\varphi\left(x\right)\right]=\psi\left(x\right),
\end{verbatim}

\begin{verbatim}
з якої ми виводимо  u(x,  k) = \varphi(x) + k\psi(x)  , тобто цінність розв'язку на  t = \Delta t  . Розв'язок у наступні часи можна знайти за явною формулою (38.5).
\end{verbatim}

\section{ЗАДАЧІ}\label{ux437ux430ux434ux430ux447ux456}

\begin{enumerate}
\def\labelenumi{\arabic{enumi}.}
\tightlist
\item
  Постройте явную конечно-разностную схему для задачи
\end{enumerate}

(УЧП)

\begin{verbatim}
u_t = u_{xx}
,  0 < x < 1  ,  0 < t < \infty  ,
\end{verbatim}

(ГУ)

\begin{verbatim}
\begin{cases} u(0, t) = 0, \\ u(1, t) = 0, \\ u(x, 0) = \sin(\pi x), \quad 0 \le x \le 1. \end{cases}
\end{verbatim}

\begin{verbatim}
(HУ) u(x, 0) = \sin(\pi x), \quad 0 \le x \le 1. Полагая h = \Delta x = 0.1 , найдите решение при t_1 = 0.005 , t_2 = 0.010 , t_3 = 0.015 . Постройте график полученного решения на сетке x = 0; 0.1; 0.2; 0.3; . .; 0.9: 1 при t = 0.015.
\end{verbatim}

\pandocbounded{\includegraphics[keepaspectratio]{images/_page_293_Figure_6.jpeg}}

Фіг. 38.5. Залежність повної похибки від кроку сітки.

\begin{itemize}
\tightlist
\item
  \begin{enumerate}
  \def\labelenumi{\arabic{enumi}.}
  \setcounter{enumi}{1}
  \tightlist
  \item
    Розв'яжіть задачу 1 аналітично (ділячи змінні). Знайдіть значення
    розв'язку у вузлах сітки x=0; 0,1; 0,2; 0,3; \ldots; 0,9; 1 при
    t=0,015. Порівняйте ці результати з результатами чисельного
    розв'язку задачі 1. (Ви можете скористатися калькулятором або
    написати невелику комп'ютерну програму для табуляції аналітичного
    розв'язку.)
  \end{enumerate}
\item
  \begin{enumerate}
  \def\labelenumi{\arabic{enumi}.}
  \setcounter{enumi}{2}
  \tightlist
  \item
    Відповідно до примітки 3, побудуйте блок-схему для розв'язання
    гіперболічної задачі.
  \end{enumerate}
\item
  \begin{enumerate}
  \def\labelenumi{\arabic{enumi}.}
  \setcounter{enumi}{3}
  \tightlist
  \item
    Розв'яжіть задачу 1 з новою граничною умовою при x=1
    \texttt{u\_x}(1, t) = -{[}\texttt{u(1,\ t)} 1{]}. \# СХЕМИ НЕЯВНИХ
    РІЗНИЦЬ (СХЕМА КРАНКА--НІКОЛСОНА)
  \end{enumerate}
\end{itemize}

МЕТАЛЕКЦІЇ: Показати, як задачі, що містять функції, залежні від часу,
можна розв'язати за допомогою неявної різницької схеми. У цій схемі, як
і в попередній схемі, часткові похідні замінюються скінченно-різницею,
але \texttt{u\_\{l+1,\textbackslash{}\ J\}} більше не виражаються явно
через значення на попередніх шарах. Тепер, щоб визначити значення
\texttt{u\_\{l+1,\textbackslash{}\ J\}} , необхідно розв'язати систему
рівнянь. Іншими словами, щоб визначити розв'язок на кожному часовому
nage, необхідно розв'язати систему рівнянь.

Перевага неявних схем над явними полягає в тому, що в неявних схемах
танг сітки можна зробити достатньо великим без страху, що помилки
округлення «знищать» розв'язок.

Для розв'язання параболічної задачі ми використаємо відому схему
\emph{Кранк}-\emph{Ніколсона}.

Як уже згадувалося в попередній лекції, щоб явна схема працювала
надійно, часовий крок сітки має бути невеликим. Зокрема, якщо просте
завдання для рівняння теплопровідності

(39.1)

\begin{verbatim}
\begin{array}{ccccccccccccccccccccccccccccccccccc
\end{verbatim}

\begin{verbatim}
розв'язується за явною схемою, тоді кроки сітки \Delta t і \Delta x повинні задовольняти умову
\end{verbatim}

\begin{verbatim}
\frac{\Delta t}{(\Delta x)^2} \leqslant 0.5.
\end{verbatim}

\begin{verbatim}
Інакше метод буде чисельно нестабільним (а помилки округлення зростуть нескінченно). (Питання чисельної стабільності розглядаються у [53].) Іншими словами, якщо відстань сітки на осі x вибрана \Delta x=0,1 , то часовий крок \Delta t не може бути більшим за 0.5\Delta x^2=0.005 . Це означає, що ви можете пройти від t=0 до t=1, зробивши 200 кроків.
\end{verbatim}

Однак існують такі схеми (неявні схеми), які дозволяють виконувати
обчислення з набагато більшим кроком, тоді як

Правда, обсяг розрахунків, які виконуються на одному кроці, зростає.
Такі схеми можна використовувати для виконання обчислень із досить
великим кроком, але на кожному кроці потрібно розв'язати систему
алгебраїмних рівнянь. Щоб проілюструвати цей метод, розв'яжемо наступну
проблему теплопровідності.

\section{Неявная схема для уравнения
теплопроводности}\label{ux43dux435ux44fux432ux43dux430ux44f-ux441ux445ux435ux43cux430-ux434ux43bux44f-ux443ux440ux430ux432ux43dux435ux43dux438ux44f-ux442ux435ux43fux43bux43eux43fux440ux43eux432ux43eux434ux43dux43eux441ux442ux438}

Рассмотрим задачу

(39 2)

\begin{verbatim}
\begin{array}{ccccccccccccccccccccccccccccccccccc
\end{verbatim}

Використаємо наступні апроксимації скінченних різниць для часткових
похідних \texttt{u\_t} та \texttt{u\_\{xx\}} :

\begin{verbatim}
\begin{split} u_t(x,\ t) &= \frac{1}{k} \big[ u(x,\ t+k) - u(x,\ t) \big], \\ u_{xx}(x,\ t) &= \frac{\lambda}{h^2} \big[ u(x+h,\ t+k) - 2u(x,\ t+k) + u(x-h,\ t+k) \big] + \\ &\quad + \frac{(1-\lambda)}{h^2} \big[ u(x+h,\ t) - 2u(x,\ t) + u(x-h,\ t) \big], \end{split}
\end{verbatim}

\begin{verbatim}
де \lambda взято з сегмента [0,1]. Зверніть увагу, що u_{xx} наближається зваженим середнім центральних різницевих похідних у момент часу t і t+k. \lambda=0.5 отримується звичайне середнє для цих двох центральних похідних, і при \lambda=0.75 одна з різницевих похідних бере вагу 0,75, а інша — 0,25. У \lambda=0 отримується звичайна явна схема, яка була обговорена в попередній лекції.
\end{verbatim}

Після заміни часткових похідних \texttt{u\_t} і \texttt{u\_\{xx\}} у
задачі (39.2) отримуємо різницю (див. рисунок 39.1).

Разностное уравнение

\begin{verbatim}
\frac{1}{k} (u_{i+1, j} - u_{i, j}) = \frac{\lambda}{h^2} (u_{i+1, j+1} - 2u_{i+1, j} + u_{i+i, j-1}) + \frac{(1-\lambda)}{h^2} (u_{i, j+1} - 2u_{i, j} + u_{i, j-1}),
\end{verbatim}

(39.3)

\begin{verbatim}
(\Gamma Y) \begin{cases} u_{i, 1} = 0, & i = 1, 2, ..., m, \\ u_{i, n} = 0, & i = 1, 2, ..., m, \\ u_{i, j} = 1, & j = 2, ..., n-1. \end{cases}
\end{verbatim}

298

Перенесіть усі невідомі значення u з верхнього часового шару (i+1) на
ліву частину рівняння (39.3) і отримайте (39.4)

\begin{verbatim}
-\lambda r u_{i+1, j+1} + (1+2r\lambda) u_{i+1, j} - \lambda r u_{i+1, j-1} = = r(1-\lambda) u_{i, j+1} + [1-2r(1-\lambda)] u_{i, j} + r(1-\lambda) u_{i, j-1},
\end{verbatim}

де вводиться позначення r = k/h\^{}2 . Зверніть увагу, що якщо i
фіксоване, а j змінюється від 2 до n-1, співвідношення (39.4) визначають
систему рівнянь n-2 з n-2 невідомим \texttt{u\_\{t+1,2\}},
\texttt{u\_\{t+1,3\}},

\pandocbounded{\includegraphics[keepaspectratio]{images/_page_296_Figure_5.jpeg}}

Фіг. 39.1. Сітка для схеми неявних різниць.

\pandocbounded{\includegraphics[keepaspectratio]{images/_page_296_Figure_7.jpeg}}

\begin{verbatim}
Фіг. 39.2. Шаблон неявної схеми u_{i+1, 4}, \ldots, u_{i+1, n-1} , які є розв'язком проблеми у внутрішній
\end{verbatim}

\begin{verbatim}
Mesh-вузли на часовому рівні t=(i+1)\Delta t .
\end{verbatim}

Фіг. 39.2 дає візуальне представлення структури кожного рівняння системи
(39.4). Тепер перейдемо до розв'язку системи (39.4).

\section{Алгоритм розв'язання задачі
(39.2)}\label{ux430ux43bux433ux43eux440ux438ux442ux43c-ux440ux43eux437ux432ux44fux437ux430ux43dux43dux44f-ux437ux430ux434ux430ux447ux456-39.2}

\begin{verbatim}
КРОК 1. Виберіть якусь цінність \lambda ( 0 \le \lambda \le 1 ). Якщо \lambda = 0 , то рівняння (39.4) перетворюються на явні формули з лекції 38.
\end{verbatim}

\begin{verbatim}
КРОК 2. Припустимо, наприклад, h=\Delta x=0,2 і k=\Delta t=0,08 (у цьому випадку r=k/h^2=2 ). У цьому випадку сітка містить 6 вузлів уздовж осі (4 внутрішні вузли) (див. рисунок 39.1). Візьмемо параметр ваги \lambda=0,5 (отримана схема називається схемою Кренка-Ніколсона). Відповідно до обчислювальної схеми (її зазвичай називають шаблоном), показаної на рис. На рисунку 39.2, рухаючись зліва направо (j=2,3,4,5) уздовж перших двох шарів (i=1), отримуємо наступні чотири рівняння:
\end{verbatim}

\begin{verbatim}
\begin{array}{l} -u_{21}+3u_{22}-u_{23}=u_{11}-u_{12}+u_{13}=1,\\ -u_{22}+3u_{23}-u_{24}=u_{12}-u_{13}+u_{14}=1,\\ -u_{23}+3u_{24}-u_{25}=u_{13}-u_{14}+u_{15}=1,\\ -u_{24}+3u_{25}-u_{26}=u_{14}-u_{15}+u_{16}=1. \end{array}
\end{verbatim}

Перепишем их в матричной форме

(39.5)

\begin{verbatim}
\begin{bmatrix} 3 & -1 & 0 & 0 \\ -1 & 3 & -1 & 0 \\ 0 & -1 & 3 & -1 \\ 0 & 0 & -1 & 3 \end{bmatrix}
\begin{bmatrix} u_{2\bar{z}} \\ u_{\bar{z}a} \\ u_{24} \\ u_{\bar{z}x} \end{bmatrix}
=
\begin{bmatrix} 1 \\ 1 \\ 1 \\ 1 \end{bmatrix}
\end{verbatim}

.

Матриця цієї системи називається тридіагональною. Щоб розв'язати систему
з трьома діагоналями

\begin{verbatim}
\begin{bmatrix} b_1 & c_1 & 0 & 0 & \dots & & & 0 \\ a_1 & b_2 & c_2 & 0 & \dots & & & 0 \\ 0 & a_{\bar{z}} & b_3 & c_3 & \dots & & & 0 \\ & & & & & & & & \\ \vdots & & & & & & & & \\ 0 & 0 & 0 & \dots & & & & a_{n-1} & b_n \end{bmatrix}
\begin{bmatrix} x_1 \\ x_{\bar{z}} \\ x_{\bar{s}} \\ \vdots \\ x_n \end{bmatrix}
=
\begin{bmatrix} d_t \\ d_{\bar{z}} \\ d_{\bar{s}} \\ \vdots \\ \vdots \\ d_n \end{bmatrix}
\end{verbatim}

,

преобразуем ее в эквивалентную систему вида

\begin{verbatim}
\begin{bmatrix} 1 & c_1^* & 0 & 0 & \dots & 0 \\ 0 & 1 & c_2^* & 0 & \dots & 0 \\ 0 & 0 & 1 & c_2^* & \dots & 0 \\ \vdots & \vdots & & & \vdots & \vdots \\ 0 & 0 & \dots & & 1 \end{bmatrix}
\begin{bmatrix} x_1 \\ x_2 \\ x_3 \\ \vdots \\ x_n \end{bmatrix}
=
\begin{bmatrix} d_1^* \\ d_2^* \\ d_3^* \\ \vdots \\ d_n^* \end{bmatrix}
\end{verbatim}

.

где

\begin{verbatim}
c_1^* = c_1/b_1
,  c_{j+1}^* = \frac{c_{j+1}}{b_{j+1} - a_j c_j^*}  ,  j = 1, 2, \ldots, n-2  , и  d_1^* = d_1/b_1  ,  d_{j+1}^* = \frac{d_{j+1} - a_j d_j^*}{b_{j+1} - a_j c_j^*}  ,  j = 1,  2, \ldots, n-1  .
\end{verbatim}

У цьому перетворенні немає нічого дивного, оскільки друга система
повністю еквівалентна першої. Матриця нової системи організована так, що
цю систему дуже легко розв'язати. Розв'язуючи рівняння послідовно знизу
вгору, отримуємо

\begin{verbatim}
x_n = d_n^*, \quad x_j = d_j^* - c_j^* x_{j+1}, \quad j = n-1, \ n-2, \ \dots, \ 2, \ 1.
\end{verbatim}

\begin{verbatim}
Застосовуючи цей метод до системи чотирьох рівнянь (39,5), отримуємо розв'язок  u_{22}=0,60  ;  u_{23}=0,80  ;  u_{24}=0,80  ;  u_{23}=0,60  . Ми отримали наближені значення розв'язку у внутрішніх точках сітки на  t=\Delta t  . Тепер ми можемо зробити наступний крок у часі, але для цього нам доведеться розв'язати нову систему рівнянь.
\end{verbatim}

У неявній схемі кількість обчислень на кожному кроці більша, ніж на
явному, але хорошу точність можна отримати навіть при значно більшому
кроці.

\section{ЗАВДАННЯ - 1. З диференціального рівняння (39.3) отримуємо
рівняння
(39.4).}\label{ux437ux430ux432ux434ux430ux43dux43dux44f---1.-ux437-ux434ux438ux444ux435ux440ux435ux43dux446ux456ux430ux43bux44cux43dux43eux433ux43e-ux440ux456ux432ux43dux44fux43dux43dux44f-39.3-ux43eux442ux440ux438ux43cux443ux454ux43cux43e-ux440ux456ux432ux43dux44fux43dux43dux44f-39.4.}

\begin{itemize}
\tightlist
\item
  \begin{enumerate}
  \def\labelenumi{\arabic{enumi}.}
  \setcounter{enumi}{1}
  \tightlist
  \item
    Побудуємо неявну схему скінченної дельти для задачі
  \end{enumerate}
\end{itemize}

\begin{verbatim}
\begin{array}{ll} (\mathrm{Y}\mathrm{H}\Pi) & u_t = u_{xx}, \quad 0 < x < 1, \quad 0 < t < \infty, \\ (\Gamma\mathrm{Y}) & \left\{ \begin{array}{ll} u_{\,0}(0, \ t) = 1, \\ u_{\,x}(1, \ t) + u_{\,}(1, \ t) = g_{\,}(t), \\ u_{\,x}(1, \ t) = 0, \quad 0 \leqslant \kappa \leqslant 1. \end{array}
\end{verbatim}

\begin{verbatim}
\right. & 0 < t < \infty, \end{array}
\end{verbatim}

\begin{enumerate}
\def\labelenumi{\arabic{enumi}.}
\setcounter{enumi}{2}
\tightlist
\item
  Как с помощью неявной схемы решить задачу
\end{enumerate}

(УЧП)

\begin{verbatim}
u_t = u_{xx} + u, \quad 0 < x < 1,
\end{verbatim}

(ГУ)

\begin{verbatim}
- 4. Как выглядит шаблон для уравнения (39.4), когда  \lambda = 1  ?
\end{verbatim}

\begin{itemize}
\tightlist
\item
  \begin{enumerate}
  \def\labelenumi{\arabic{enumi}.}
  \setcounter{enumi}{4}
  \tightlist
  \item
    Постройте блок-схему решения уравнения теплопроводности (39.2). Если
    есть возможность, то напишите программу для ЭВМ Было бы очень
    полезно решить эту задачу численно при различных значениях λ с
    начальным условием и (x, 0) = sin (πx). Сравните полученные при
    различных λ численные результаты с аналитическим решением, которое в
    данном случае имеет вид
  \end{enumerate}
\end{itemize}

\pandocbounded{\includegraphics[keepaspectratio]{images/_page_172_Figure_2.jpeg}}

6 Воспользовавшись формулами, приведенными в тексте, решите систему
алгебраических уравнений (39.5).

Лекция 40

\section{СРАВНЕНИЕ АНАЛИТИЧЕСКИХ РЕШЕНИЙ С ЧИСЛЕННЫМИ МЕТА ЛЕКЦІЇ:
Сравнить достоинства и недостатки аналитических и численных методов
решения уравнений с частными производными. Приведена важная
математическая задача об идентификации физических величин
(параметрическая идентификация). В связи с этой задачей рассматривается
один важный пример из
биологии.}\label{ux441ux440ux430ux432ux43dux435ux43dux438ux435-ux430ux43dux430ux43bux438ux442ux438ux447ux435ux441ux43aux438ux445-ux440ux435ux448ux435ux43dux438ux439-ux441-ux447ux438ux441ux43bux435ux43dux43dux44bux43cux438-ux43cux435ux442ux430-ux43bux435ux43aux446ux456ux457-ux441ux440ux430ux432ux43dux438ux442ux44c-ux434ux43eux441ux442ux43eux438ux43dux441ux442ux432ux430-ux438-ux43dux435ux434ux43eux441ux442ux430ux442ux43aux438-ux430ux43dux430ux43bux438ux442ux438ux447ux435ux441ux43aux438ux445-ux438-ux447ux438ux441ux43bux435ux43dux43dux44bux445-ux43cux435ux442ux43eux434ux43eux432-ux440ux435ux448ux435ux43dux438ux44f-ux443ux440ux430ux432ux43dux435ux43dux438ux439-ux441-ux447ux430ux441ux442ux43dux44bux43cux438-ux43fux440ux43eux438ux437ux432ux43eux434ux43dux44bux43cux438.-ux43fux440ux438ux432ux435ux434ux435ux43dux430-ux432ux430ux436ux43dux430ux44f-ux43cux430ux442ux435ux43cux430ux442ux438ux447ux435ux441ux43aux430ux44f-ux437ux430ux434ux430ux447ux430-ux43eux431-ux438ux434ux435ux43dux442ux438ux444ux438ux43aux430ux446ux438ux438-ux444ux438ux437ux438ux447ux435ux441ux43aux438ux445-ux432ux435ux43bux438ux447ux438ux43d-ux43fux430ux440ux430ux43cux435ux442ux440ux438ux447ux435ux441ux43aux430ux44f-ux438ux434ux435ux43dux442ux438ux444ux438ux43aux430ux446ux438ux44f.-ux432-ux441ux432ux44fux437ux438-ux441-ux44dux442ux43eux439-ux437ux430ux434ux430ux447ux435ux439-ux440ux430ux441ux441ux43cux430ux442ux440ux438ux432ux430ux435ux442ux441ux44f-ux43eux434ux438ux43d-ux432ux430ux436ux43dux44bux439-ux43fux440ux438ux43cux435ux440-ux438ux437-ux431ux438ux43eux43bux43eux433ux438ux438.}

Вероятно, пришло время обсудить достоинства и недостатки аналитических и
численных решений уравнений с частными производными. Сначала выясним,
что мы подразумеваем под этими двумя типами решений.

\section{Аналитические
решения}\label{ux430ux43dux430ux43bux438ux442ux438ux447ux435ux441ux43aux438ux435-ux440ux435ux448ux435ux43dux438ux44f}

Под аналитическими решениями мы подразумеваем такие решения, в которых
неизвестная функция и выражена через независимые переменные и параметры
системы в виде формул, бесконечных рядов и интегралов.

\section{Численные
решения}\label{ux447ux438ux441ux43bux435ux43dux43dux44bux435-ux440ux435ux448ux435ux43dux438ux44f}

Під чисельним ми зрозуміємо чисельні розв'язки, отримані чисельно після
приблизної заміни початкового рівняння іншим, простішим рівнянням.
Наприклад, у методі скінченної різниці похідні виражаються через
скінченні різниці, а розв'язок диференціального рівняння в похідних у
похідних апроксимується розв'язком різницького рівняння. Результатом
такої процедури зазвичай є таблиця значень розв'язків, а для деяких
значень пояснювальних змінних.

Таблина 40.1

~~~~~~~~~~~~~~~~~~~~~~~~~~\textbar{} × \textbar{} \textbar{} \textbar{}
\textbar{} \textbar{} \textbar{} \textbar{} \textbar{} \textbar{}
\textbar{} \textbar{}

\textbar---------------------------\textbar-------------\textbar------------------\textbar-----------\textbar----------\textbar----------\textbar-----\textbar-----\textbar-----\textbar-----------\textbar------------------\textbar------------------\textbar{}

~~~~~~~~~~~~~~~~~~~~~~~~~~\textbar{} 0 \textbar{} 0,1 \textbar{} 0,2
\textbar{} 0,3 \textbar{} 0,4 \textbar{} 0,5 \textbar{} 0,6 \textbar{}
0,7 \textbar{} 0,8 \textbar{} 0,9 \textbar{} 1 \textbar{}

00,010,020,03 \textbar{} 100 \textbar{} 10,20,15 \textbar{} 10,34
\textbar{} l: \textbar{} : \textbar{} : \textbar{} ! \textbar{} :
\textbar{} 10,34 \textbar{} 10,20,15 \textbar{} 1000 \textbar{}

; \textbar{} \textbar{} : \textbar{} : \textbar{} : \textbar{} :
\textbar{} : \textbar{} : \textbar{} : \textbar{} : \textbar{} :
\textbar{} : \textbar{}

Теперь, когда мы знаем, что подразумевается под решением каждого из этих
типов, давайте сравним их между собой.

\section{\texorpdfstring{\textbf{Сравнительный анализ численного} и
аналитического
решений}{Сравнительный анализ численного и аналитического решений}}\label{ux441ux440ux430ux432ux43dux438ux442ux435ux43bux44cux43dux44bux439-ux430ux43dux430ux43bux438ux437-ux447ux438ux441ux43bux435ux43dux43dux43eux433ux43e-ux438-ux430ux43dux430ux43bux438ux442ux438ux447ux435ux441ux43aux43eux433ux43e-ux440ux435ux448ux435ux43dux438ux439}

Рассмотрим простейшую смещанную параболическую задачу

(40.1)

\begin{verbatim}
y^2 - x^2 = \text{const},
\end{verbatim}

Розв'язок якого показано на рис. 40.1. Що краще --- мати аналітичне
рішення цієї проблеми?

(40.2)

\begin{verbatim}
y^2 + x^2 = \text{const}.
\end{verbatim}

\begin{verbatim}
или численное, представленное при  \alpha = 1  в табл. 40.1?
\end{verbatim}

Вопрос хороший, а ответ на него зависит от того, что мы хотим делать с
этим решением дальше. Однако каждому из типов решений присущи очевидные
преимущества.

\section{Преимущества аналитического
решения}\label{ux43fux440ux435ux438ux43cux443ux449ux435ux441ux442ux432ux430-ux430ux43dux430ux43bux438ux442ux438ux447ux435ux441ux43aux43eux433ux43e-ux440ux435ux448ux435ux43dux438ux44f}

\begin{enumerate}
\def\labelenumi{\arabic{enumi}.}
\tightlist
\item
  Очевидно, що розв'язок (40.2) є більш інформативним, ніж таблиця
  чисел. Якщо ми хочемо обчислити розв'язок у певній точці (x, t), ми
  можемо зробити це так точно, як хочемо, просто збільшуючи кількість
  членів ряду, які розглядаються. Водночас легко отримати оцінку зверху
  щодо значення допущеної помилки.
\end{enumerate}

\pandocbounded{\includegraphics[keepaspectratio]{images/_page_301_Figure_4.jpeg}}

Рис. 40.1. Решение уравнения теплопроводности (40.1) в различные моменты
времени.

\begin{itemize}
\tightlist
\item
  \begin{enumerate}
  \def\labelenumi{\arabic{enumi}.}
  \setcounter{enumi}{1}
  \tightlist
  \item
    Аналітичне розв'язання завжди дозволяє обчислити значення розв'язку
    в одній точці (x, t) без необхідності обчислювати значення
    розв'язків в інших точках, як це відбувається при розв'язанні задач
    із явними або неявними різницею схем.
  \end{enumerate}
\item
  \begin{enumerate}
  \def\labelenumi{\arabic{enumi}.}
  \setcounter{enumi}{2}
  \tightlist
  \item
    Аналітичне рішення дозволяє визначити розв'язок у будь-якій точці, а
    не лише на вузлах сітки.
  \end{enumerate}
\item
  \begin{enumerate}
  \def\labelenumi{\arabic{enumi}.}
  \setcounter{enumi}{3}
  \tightlist
  \item
    Найважливішою перевагою аналітичного розв'язку для нас є можливість
    простежити вплив фізичних параметрів, початкових і граничних умов на
    природу розв'язку.
  \end{enumerate}
\end{itemize}

Чисельні методи не виявляють цих регулярностей, оскільки дозволяють
знаходити розв'язок лише за заданих параметрів --- початкових і
граничних умовах. Іноді важливо знати взаємозв'язок між параметрами
моделі та розв'язком, особливо коли йдеться про оцінку фізичних
параметрів за типом розв'язку. Наприклад, припустимо, що розв'язок і
визначається експериментально, і ми знаємо аналітичний розв'язок

u = функция от параметров.

Тогда можно поставить вопрос об определении параметров нак функции
входных данных по схеме

Параметры = функция от u = функция исходных данных.

Такая задача называется параметрической идентификацией. Для
параметрической идентификации необходимо уметь решать уравнение с
частными производными. Несколько позже в этой лекции мы приведем пример
параметрической идентификации в биологии, но сначала посмотрим, чем же
хороши численные решения.

\section{Преимущества численных
решений}\label{ux43fux440ux435ux438ux43cux443ux449ux435ux441ux442ux432ux430-ux447ux438ux441ux43bux435ux43dux43dux44bux445-ux440ux435ux448ux435ux43dux438ux439}

Головна перевага чисельних розв'язків полягає в тому, що їх можна
отримати навіть якщо аналітичне рішення неможливо. Майже всі нелінійні
диференціальні рівняння в похідних похідних потрібно розв'язувати
чисельними методами, і більшість реальних моделей фізики, хімії,
біології тощо є нелінійними. Як правило, лінійні моделі наближаються до
нелінійних лише якщо в них відкинути нелінійні члени. Наприклад, ось
кілька нелінійних рівнянь:

\begin{enumerate}
\def\labelenumi{\arabic{enumi})}
\tightlist
\item
  нелинейное волновое уравнение \texttt{u\_\{tt\}\ =\ u\_\{xx\}\ +}f(u)
\end{enumerate}

\begin{verbatim}
;

2) уравнение реакции с нелинейной диффузией  u_t = u_{xx} +f(u)
\end{verbatim}

;

\begin{enumerate}
\def\labelenumi{\arabic{enumi})}
\setcounter{enumi}{2}
\tightlist
\item
  система Ходжкина---Хаксли
\end{enumerate}

\begin{verbatim}
u(x, t) =
\end{verbatim}

Аналитические решения для этих уравнений не известны ни при каких
нелинейных функциях f и \texttt{g\^{}\{(1)\}} . Поэтому общий подход к
решению нелинейных (а в ряде случаев и линейных) задач обычно базируется
на численных решениях.

Тепер давайте розглянемо приклад того, як можна використати аналітичне
рішення для визначення фізичних параметрів. Детальніший опис цих завдань
можна знайти у рекомендованій літературі.

\section{Параметрическая идентификация (в
биологии)}\label{ux43fux430ux440ux430ux43cux435ux442ux440ux438ux447ux435ux441ux43aux430ux44f-ux438ux434ux435ux43dux442ux438ux444ux438ux43aux430ux446ux438ux44f-ux432-ux431ux438ux43eux43bux43eux433ux438ux438}

Припустимо, що біолог намагається визначити швидкість, з якою (K\^{}+)
іони калію дифундують у екзоплазмовому розчині. Якщо ви знаєте цей
коефіцієнт дифузії, можна багато сказати про те, як нервові імпульси
передаються вздовж аксонів. Проблема в тому, що це співвідношення майже
неможливо знайти прямим вимірюванням. Однак можливо знайти математичний

\begin{verbatim}
<sup>1)</sup> Это не так. В качестве примера можно привести уравнение  u_{ii} = u_{xx} - \sin u_x - \Pi pum  , ped.
\end{verbatim}

\begin{verbatim}
связь между концентрацией поташа  u\left(x,\,t\right)  и коэффициентом диффузии  D,\,\,a\,\,  затем по измерениям величины  u\left(x,\,t\right)  определить величину D.
\end{verbatim}

Біологи Ходжкін і Кейз виявили, що після розміщення гігантського аксону
у спеціальному соляному розчині концентрація радіохктивного калію (**2K)
вздовж аксону приблизно описується кривою в початковий момент

\begin{verbatim}
\vec{B} = 2A\xi_x\eta_x + B(\xi_x\eta_y + \xi_y\eta_x) + 2C\xi_y\eta_y = -16x^2y^2
\end{verbatim}

.

Отже, параметри A і A слід обирати так, щоб крива відповідала
експериментальним даним (наприклад, у сенсі методу найменших квадратів).
Також було встановлено, що P И С

\pandocbounded{\includegraphics[keepaspectratio]{images/_page_303_Figure_6.jpeg}}

\begin{verbatim}
. 40.2. Начальная концентрация ионов (  ^{42}K  ): a—теоретическая кривая;  \delta  —экспериментальные точки.
\end{verbatim}

с течением времени ноны поташа под действием диффузии и конвекции
растекаются вдоль аксона. Ходжкин и Кейс предположили, что концентрация
и является решением следующей задачи с начальными данными для уравнения
диффузии:

(40.3)

\begin{verbatim}
\bar{D} = A\xi_{xx} + B\xi_{xy} + C\xi_{yy} + D\xi_x + E\xi_y = -2(x^2 + y^2)
\end{verbatim}

Цю проблему вирішили у лекції 15, переключившись на рухому систему
координат. Якщо читач пам'ятає, загальна схема методу виглядає так.
Спочатку відкидаємо конвективний член і отримуємо розв'язок суто
дифузійної задачі (при V = 0):

\begin{verbatim}
\widetilde{E} = A\eta_{xx} + B\eta_{xy} + C\eta_{yy} + D\eta_x + E\eta_y = 2(y^2 - x^2),
= \frac{A\sqrt{a}}{\sqrt{a+4Dt}} e^{-[x^2/(a+4Dt)]},
\tilde{F} = F = 0
u(x, t) = \frac{A\sqrt{a}}{\sqrt{a+4Dt}}e^{-[(x-Vt)^{2}/(a+4Dt)]}.
\overline{G} = G = 0
SS = \sum_{i=1}^{n} [y_i - (a + bx_i)]^2.
\overline{A}u_{\xi\xi} + \overline{B}u_{\xi\eta} + \overline{C}u_{\eta\eta} + \overline{D}u + \overline{E}u_{\eta} + \overline{F}u = \overline{G}
u_t = \frac{1}{r} \frac{\partial}{\partial r} \left( Dr u_r - s \omega^2 r^2 u \right), \quad 0 < r < 1,
u_{\xi\eta} = \frac{-(x^2 + y^2) u_{\xi} + (y^2 - x^2) u_{\eta}}{8x^2y^3}.
Au_{xx} + Bu_{xy} + Cu_{yy} + Du_x + Eu_y + Fu = G
u_{\xi\eta} = \frac{\eta u_{\xi} - \xi u_{\eta}}{2(\xi^2 - \eta^2)}.
\begin{cases} 0, & t < x/V, \\ P, & t \ge x/V. \end{cases}
\alpha = \alpha (\xi, \eta) = \xi + \eta,
\end{verbatim}

Конечно, все это очень просто! Не сложнее, чем сбрасывать что-нибудь на
ленту транспортера, а затем наблюдать за его движением. Однако все
становится очень интересным, если примесь диффундирует в среде. Чтобы
разобраться в том, что про-исходит с движущейся волной при наличии
диффузии, решим

следующую задачу

(15.2)

З формальною заміною змінних
\texttt{\textbackslash{}xi\ =\ \textbackslash{}xi(x,\ y)} та
\texttt{\textbackslash{}eta\ =\ \textbackslash{}eta(x,\ y)} (41.1)
зводиться до тієї ж форми

(41.2)

где, как обычно, \texttt{H(x)} --- функция Хевисайда Начальное
распределение концентрации изображено на рис 15 4.

\pandocbounded{\includegraphics[keepaspectratio]{images/_page_114_Figure_4.jpeg}}

\begin{verbatim}
Рис. 15.3 Конвективная волна a—передиції фронт примеси,   \delta  —на одну единицу вперед волнового фронта
\end{verbatim}

\begin{verbatim}
Обратим внимание на то, что в задаче (15.2) мы отодвинули границу в —   \infty   (теперь решается задача Коши), так что она не оказывает теперь влияние на оценку решение. Чтобы решить задачу (15.2), можно воспользоваться преобразованием Лапласа Р и с
\end{verbatim}

\pandocbounded{\includegraphics[keepaspectratio]{images/_page_114_Figure_7.jpeg}}

. 15 4 Начальные условия в задаче конвективной диффузии.

\begin{verbatim}
ло переменной t или преобразованием Фурье по переменной x. Однако в этом случае гораздо интереснее поступить совсем по-другому Введем новую систему координат, которая движется вдоль старой со скоростью V. Другими словами, вместо системы координат, привязанной к берегу реки, мы рассмотрим систему координат, которая движется со скоростью фронта примеси (конечно, при налични диффузии фронт будет «смазан»). С точки зрения математики это означает, что мы заменяем пространственную координату x на новую   \xi = x - Vt  .
\end{verbatim}

Теперь ясно, что

\begin{verbatim}
когда   \xi = 0  , мы находимся па фронте распространяющейся примеси.
\end{verbatim}

\begin{verbatim}
когда   \xi = 1  , мы находимся на одну единицу длины впереди фронта.
\end{verbatim}

\begin{verbatim}
когда   \xi = -1  , мы находимся на одну единицу длины позади фронта.
\end{verbatim}

Наша задача --- преобразовать исходную задачу с НУ

\begin{verbatim}
\overline{D} = A\xi_{xx} + B\xi_{xy} + C\xi_{yy} + D\xi_{x} + E\xi_{y},
[\xi_x/\xi_y] = -B/2A
\xi(x, y) = c
\frac{dy}{dx} = -\left[\xi_x/\xi_y\right] = B_t 2A.
u_{\alpha\alpha} - u_{\beta\beta} = \frac{-\beta u_{\alpha} - \alpha u_{\beta}}{2\alpha\beta}.
\overline{B} = 2A\xi_x\eta_x + B(\xi_x\eta_y + \xi_y\eta_x) + 2C\xi_y\eta_y,
\alpha = \xi + \eta = (y^2 - x^2) + (y^2 + x^2) = 2y^3,
\end{verbatim}

1) Завдяки рівності

І з B\^{}2 - 4AC = 0 році , то

\begin{verbatim}
\beta = \xi - \eta = (y^2 - x^2) - (y^2 + x^2) = -2x^2.
\end{verbatim}

\begin{verbatim}
Оскільки  \xi_x/\xi_y = -B \ 2A = -2 \sqrt{A\tilde{C}}/2A = -\sqrt{C/A}  останній вираз для B буде записаний як
\end{verbatim}

\begin{verbatim}
u_{\xi\xi} - u_{\eta\eta} = \Psi(\xi, \, \eta, \, u_{\xi}, \, u, u_{\eta}).
\end{verbatim}

\begin{verbatim}
І оскільки  \overline{A}  дорівнює нулю, то  \overline{B}  , звісно, також дорівнює нулю.
\end{verbatim}

\begin{verbatim}
Оскільки при нашому виборі  \xi  обидва коефіцієнти  \overline{A}  і  \overline{B}  повертаються до нуля, ми можемо взяти будь-яку функцію як  \eta  (якщо вона не пропорційна координаті  \xi  ). У прикладі можна вибрати  \eta = y  .
\end{verbatim}

\begin{verbatim}
Залишається знайти канонічну форму рівняння в нових координатах. Для цього потрібно просто підставити формули  \xi  і  \eta  (41.3) і знайти всі шанси  \overline{A}  ,  \overline{B}  ,  \overline{C}  ,  \overline{D}  ,  \overline{E}  ,  \overline{F}  і  \overline{G}  . На цьому вивчення параболічного випадку завершується, але перш ніж перейти до іншої теми, розглянемо простий приклад.
\end{verbatim}

\section{\texorpdfstring{Перетворення параболічного рівняння
\texttt{u\_\{xx\}} + 2u\_\{xy\} + \texttt{u\_\{yy\}\ =\ 0}у канонічну
форму}{Перетворення параболічного рівняння u\_\{xx\} + 2u\_\{xy\} + u\_\{yy\} = 0у канонічну форму}}\label{ux43fux435ux440ux435ux442ux432ux43eux440ux435ux43dux43dux44f-ux43fux430ux440ux430ux431ux43eux43bux456ux447ux43dux43eux433ux43e-ux440ux456ux432ux43dux44fux43dux43dux44f-u_xx-2u_xy-u_yy-0ux443-ux43aux430ux43dux43eux43dux456ux447ux43dux443-ux444ux43eux440ux43cux443}

Розглянемо рівняння

\begin{verbatim}
3u_{xx} + 7u_{xy} + 2u_{yy} = 0
\end{verbatim}

\begin{verbatim}
Його коефіцієнти мають такі значення: A=1, B=2, C=1, D=E=F=G=0. Отже,  B^2-4AC=0  для всіх значень x і y. Щоб знайти нові координати  \xi  і  \eta  та канонічну форму рівняння, ми рухаємося наступним чином.
\end{verbatim}

КРОК 1. Запишемо \emph{характеристичне рівняння} (тепер лише одне, а не
два)

\begin{verbatim}
u_{\xi\eta} = \Phi(\xi, \eta, u, u_1, u_2).
\end{verbatim}

Розв'язуючи це рівняння відносно y, отримуємо

\begin{verbatim}
u_{\alpha\alpha} - u_{\beta\beta} = \Psi(\alpha, \beta, u, u_{\alpha}, u_{\beta}).
\end{verbatim}

,

\begin{verbatim}
Де — характеристика  \xi = y - x  визначена. Використання цієї характеристики забезпечує досягнення рівності  \overline{A} = 0  . Координату  \eta  можна вибрати довільно (якщо вона не залежить від  \xi  ). Ми оберемо його наступним чином:
\end{verbatim}

\begin{verbatim}
u_{xx} + 4u_{xy} = 0.
\end{verbatim}

.

Нові координати

\begin{verbatim}
u_{tt} = e^{2} (u_{xx} + u_{yy} + u_{zz}),
\end{verbatim}

изображены на рис. 41 1.

\pandocbounded{\includegraphics[keepaspectratio]{images/_page_309_Figure_3.jpeg}}

\begin{verbatim}
Рис. 41.1. Новая система координат   \xi = y - x  ,   \eta = y  .
\end{verbatim}

\begin{verbatim}
ШАГ 2. Осталось найти каноническую форму уравнения. Подставляя   \xi   и   \eta   в формулы для коэффициентов   \overline{A}  ,   \overline{B}  ,   \overline{C}  ,   \overline{D}  ,   \overline{E}  ,   \overline{F}   и   \overline{G}  , получаем
\end{verbatim}

\begin{verbatim}
\overline{A} = 0   (так и должно быть; мы находили   \xi   из условия равенства нулю этого коэффициента),
\end{verbatim}

\begin{verbatim}
\overline{B}=0   (мы уже раньше показали, что этог коэффициент также обращается в нуль),
\end{verbatim}

(41.5)

\begin{verbatim}
\overline{D} = A \xi_{xx} + B \xi_{xy} + C \eta_{yy} + D \xi_x + E \xi_y = 0,
\end{verbatim}

Отже, нове рівняння

\begin{verbatim}
\xi = x - Vt
u(x, t) =
\end{verbatim}

(НУ)

є загальним розв'язком рівняння

Перед нами решение задачи конвективной диффузии (15.2). Его очень легко
интерпретировать, если представить, что мы движемся относительно
графика, изображенного на рис. 15.6.

\pandocbounded{\includegraphics[keepaspectratio]{images/_page_117_Figure_3.jpeg}}

Рис. 15.6. Диффузня из области высокой концентрации в область низкой
концентрации. Чем больше коэффициент диффузии, тем быстрее установится
стационарное значение.

Другими словами, в зависимости от относительной величины D (коэффициента
диффузии) и V (скорости потока) решение движется слева направо со
скоростью V и в то же самое время передний фронт расплывается со
скоростью, определяемой величиной D.

\pandocbounded{\includegraphics[keepaspectratio]{images/_page_117_Figure_6.jpeg}}

Рис. 15.7. Решение задачи конвективнои диффузии. Вещество одновременно
движется и диффундирует.

(На рис. 15.7 показано, как расплывается передний фронт концентрации.)
\# ЗАУВАЖЕННЯ Преобразование координат является важным методом решения
уравнений с частными производными. Выбрав подходящую систему координат,
можно существенно упростить уравнение. \# ЗАДАЧІ 1. Решите задачу Коши:

Але тепер у B\^{}2-4AC\textless0 . Переходячи до нових пояснювальних
змінних, ми хочемо перетворити її у нову форму

\begin{enumerate}
\def\labelenumi{\arabic{enumi}.}
\setcounter{enumi}{1}
\tightlist
\item
  Найдите решение задачи Коши для уравнения конвективной диффузии:
\end{enumerate}

\begin{verbatim}
Це можливо, але лише шляхом введення складних координат. Щоб знайти ці комплексні координати  \xi  і  \eta  , розв'яжемо характеристичні рівняння, так само як у гіперболічному випадку:
\end{verbatim}

Воспользуйтесь преобразованием из лекции 8.

\begin{enumerate}
\def\labelenumi{\arabic{enumi}.}
\setcounter{enumi}{2}
\tightlist
\item
  Найдите решение задачи переноса:
\end{enumerate}

\begin{verbatim}
u_{tt} = c^2 (u_{xx} + u_{yy}), \quad
\frac{dy}{dx} = \sqrt{-4x^2} = 2ix
\\ -\infty < x < \infty, \quad 0 < t < \infty,
u_t = Du_{xx} - Vu_x, \quad -\infty < x < \infty, \quad 0 < t < \infty,
u(x, 0) = e^{-x^2}, \quad -\infty < x < \infty.
y^2 u_{xx} + x^2 u_{yy} = 0,
\frac{dy}{dx} = \frac{B - \sqrt{B^2 - 4AC}}{2A} = -\frac{\sqrt{-4x^2y^2}}{2y^2} = -i\frac{x}{y},
\begin{cases} -\infty < x < \infty, \\ -\infty < y < \infty, \\ -\infty < z < \infty, \end{cases}
u(\xi, \tau) = \frac{1}{2\sqrt{D\pi\tau}} \int_{-\infty}^{+\infty} e^{-\beta^2} e^{-(\xi-\beta)^2/4D\tau} d\beta.
u_{tt} = \alpha^2 u_{xx}
P_{2}
u_{\xi\eta} = \psi(\xi, \eta, u, u_{\xi}, u_{\eta}).
\begin{aligned} u_{tt} &= c^2 \Delta u, \\ (HY) & \begin{cases} u(x, y, z, 0) = 0, \\ u_t(x, y, z, 0) = \psi(x, y, z), \end{cases} \end{aligned}
\xi(x, y) = y^2, \quad \eta(x, y) = x^2.
T \sin \theta_2 - T \sin \theta_1 =
P_{1}
\Delta x \rho u_{tt} = T \left[ u_x \left( x + \Delta x, \ t \right) - u_x \left( x, \ t \right) \right] + \Delta x F \left( x, \ t \right) - \frac{1}{\Delta x \rho u_t} \left( x, \ t \right) - \Delta x \gamma \dot{u} \left( x, \ t \right)
\end{verbatim}

где \texttt{\textbackslash{}overline\{\textbackslash{}psi\}} ---
\emph{среднее} значение начального распределения
\texttt{\textbackslash{}psi} \emph{по сфере} радиуса ct с центром в
точке (x, y, z), т. е.

\begin{verbatim}
u_{tt} = \alpha^{3} u_{xx} - \beta u_{t} - \gamma u + F(x, t).
\end{verbatim}

\pandocbounded{\includegraphics[keepaspectratio]{images/_page_177_Picture_3.jpeg}}

\begin{verbatim}
(16.2) u_{tt} = \alpha^2 u_{xx},
\end{verbatim}

\pandocbounded{\includegraphics[keepaspectratio]{images/_page_177_Picture_7.jpeg}}

\begin{verbatim}
u_{tt} = ku_{xx}
u_{tt} = \frac{\partial}{\partial x} [\alpha^2(x) u_x],
u_h(r, \theta) = \sum_{n=0}^{\infty} r^n [a_n \cos(n\theta) + b_n \sin(n\theta)].
u(x, 0) = f(x)
i_x + Cv_t + Gv = 0, \\ v_x + Li_t + Ri = 0,
u = \frac{\partial}{\partial t} [t\overline{\varphi}].
i_{xx} + Gv_x - CLi_{tt} - CRi_t = 0.
u_p(r, \theta) = Ar^4 + Br^4 \cos(2\theta).
v_{x} = -Li_{t} - Ri,
u(x, t) = \frac{1}{2c} \int_{x-ct}^{x+ct} \varphi(s) ds.
i_{t,s} = CLi_{tt} + (CR + GL)i_t + GRi.
u_p(r, \theta) = -\frac{r^4}{32} - \frac{r^4}{24} \cos(2\theta).
v_{xx} = CLv_{tt} + (CR + GL)v_t + GRv.
u_t(x, t) = \frac{1}{2} [\varphi(x + ct) + \varphi(x - ct)].
v_{tt} = \alpha^2 v_{xx}, \\ i_{tt} = \alpha^2 i_{xx}, \qquad \alpha^2 = 1/CL,
\sum_{n=0}^{\infty} [a_n \cos(n\theta) + b_n \sin(n\theta)] - \frac{1}{32} - \frac{1}{24} \cos(2\theta) = 0.
u_{tt} = u_{xx} - u_t
u_1(r, \theta) = \frac{1}{32} + \frac{1}{24}r^2\cos(2\theta) - \frac{r^4}{24}\cos(2\theta) - \frac{r^4}{32} = \frac{(r^4 - 1)}{32} - \frac{(r^4 - r^2)}{24}\cos(2\theta).
\end{verbatim}

Що ви можете сказати про точність?

\begin{enumerate}
\def\labelenumi{\arabic{enumi}.}
\setcounter{enumi}{3}
\tightlist
\item
  Як отримати послідовність випадкових точок, що лежать всередині
  трикутника T?
\end{enumerate}

\pandocbounded{
\includegraphics[keepaspectratio]{images/_page_319_Figure_13.jpeg}}

\pandocbounded{\includegraphics[keepaspectratio]{images/_page_319_Figure_14.jpeg}}

\begin{enumerate}
\def\labelenumi{\arabic{enumi}.}
\setcounter{enumi}{4}
\tightlist
\item
  Як отримати послідовність випадкових чисел, розподіл яких показаний на
  рисунку?
\end{enumerate}

\pandocbounded{\includegraphics[keepaspectratio]{images/_page_319_Figure_16.jpeg}}

Іншими словами, як отримати послідовність чисел \{0, 1, 2\}, якщо
ймовірність виникнення 0 і 2 дорівнює 0,25, а ймовірність появи 1 ---
0,5?

Лекція 43

\section{РОЗВ'ЯЗАННЯ ДИФЕРЕНЦІАЛЬНИХ РІВНЯНЬ У ПОХІДНИХ ПОХІДНИХ
МОНТЕ-КАРЛО}\label{ux440ux43eux437ux432ux44fux437ux430ux43dux43dux44f-ux434ux438ux444ux435ux440ux435ux43dux446ux456ux430ux43bux44cux43dux438ux445-ux440ux456ux432ux43dux44fux43dux44c-ux443-ux43fux43eux445ux456ux434ux43dux438ux445-ux43fux43eux445ux456ux434ux43dux438ux445-ux43cux43eux43dux442ux435-ux43aux430ux440ux43bux43e}

МЕТА ЛЕКЦІЇ: Показати, як побудувати таку гру, результатом якої є
приблизне розв'язання диференціального рівняння. Одна з цих ігор (модель
випадкової ходьби) веде до скінченної різниці апроксимації задачі
Діріхле в квадраті. Цю модель можна узагальнити для отримання розв'язків
інших задач.

У попередній лекції було сказано, що можливо винайти такі ігри на удачу,
результатом яких є розв'язки (можливо, приблизні) задач для рівнянь з
частковими коефіцієнтами проникнення. У цій лекції ми покажемо, як
побудувати таку гру для наближеного розв'язання задачі Діріхле (PPP)
\texttt{u\_\{xx\}} +
\texttt{u\_\{yy\}\ =\ 0\ ,\ 0\ \textless{}\ x\ \textless{}\ 1,\ 0\ \textless{}\ y\ \textless{}\ 1},

(GI)

\begin{verbatim}
u = u_0 + u_1 = r \cos \theta - \frac{(r^4 - 1)}{32} - \frac{(r^4 - r^2)}{24} \cos (2\theta).
R(A) = g_1 P_A(p_1) + g_2 P_A(p_2) + \dots + g_{12} P_A(p_{12}).
R(A) = 1 \cdot (0.04) + 1 \cdot (0.15) + 1 \cdot (0.03) + 0 \cdot (0.06) + \dot + 0 \cdot (0.04).
R(A) = \frac{1}{4} [R(B) + R(C) + R(D) + R(E)]
u_{i,j} = \frac{1}{4} \left[ u_{i-1,j} + u_{i+1,j} + u_{i,j-1} + u_{i,j+1} \right],
\Delta u = 0, \quad 0 < r < 1 + \frac{1}{4} \sin \theta,
\end{verbatim}

(PPP)

\begin{verbatim}
u \left( 1 + \frac{1}{4} \sin \theta, \ \theta \right) = \cos \theta, \quad 0 \le \theta \le 2\pi.
u_{xx}=u_{l,\ j+1}-2u_{l,\ j}+u_{l,\ j-1},
u\left(1+\frac{1}{4}\sin\theta,\ \theta\right)=\cos\theta
u_{i, j} = \frac{u_{i, j+1} + u_{i, j-1} + \sin x_j (u_{i+1, j} + u_{i-1, j})}{2 (1 + \sin x_j)}.
u(1+\epsilon\sin\theta, \theta) = \cos\theta, \quad 0 \le \epsilon \le 1/4.
\end{verbatim}

(i, j+1)

\begin{verbatim}
u(A) = g_1 P_A(p_1) + g_2 P_A(p_2) + \dots + g_N P_A(p_N),
u_{xx} + u_{yy} = 0
u_{xx} + x^2 u_{yy} = 0
\end{verbatim}

\pandocbounded{\includegraphics[keepaspectratio]{images/_page_350_Figure_2.jpeg}}

\begin{verbatim}
(\text{UBP}) \qquad
u(1+\varepsilon\sin\theta,\ \theta)=u(1,\ \theta)+u_r(1,\ \theta)(\varepsilon\sin\theta)+u_{rr}(1,\ \theta)\frac{(\varepsilon\sin\theta)^2}{21}+\ldots
\end{verbatim}

\begin{enumerate}
\def\labelenumi{\arabic{enumi}.}
\setcounter{enumi}{4}
\tightlist
\item
  Разработайте схему метода Монте-Карло для решения смешанной
  параболической задачи
\end{enumerate}

\begin{verbatim}
\Delta u = 0, \quad 0 < r < 1 + \varepsilon \sin \theta,
\end{verbatim}

\begin{verbatim}
\right. \end{array}
\end{verbatim}

\begin{verbatim}
u(1, \theta) + u_r(1, \theta) (\varepsilon \sin \theta) + u_{rr}(1, \theta) \frac{(\varepsilon \sin \theta)^2}{2!} + \dots = \cos \theta.
J[y] = \int_{b}^{a} F(x, y, y') dx,
(46.8) u = u_0 + \varepsilon u_1 + \varepsilon^2 u_2 + \dots.
J[y] = \int_{0}^{1} [y^{2}(x) + y'^{2}(x)] dx.
P_{0}
\frac{df(x)}{dx} = 0
\begin{cases} \Delta u_{0} = 0, & 0 < r < 1 \text{ (внутри круга),} \\ u_{0}(1, \theta) = \cos \theta, & u_{0}(r, \theta) = r \cos \theta, \end{cases}
T = \int_{0}^{T} dt = \int_{a}^{L} \frac{dt}{ds} ds = \int_{0}^{L} \frac{ds}{v} = \frac{1}{\sqrt{2mg}} \int_{0}^{L} \frac{ds}{\sqrt{y}} = \frac{1}{\sqrt{2mg}} \int_{a}^{b} \sqrt{\frac{1+y'^{2}}{y}} dy
J[y] = \int_{a}^{b} F(x, y, y') dx.
P_{1}
\frac{df(x)}{dx} = 0
\begin{cases} \Delta u_{1} = 0, & 0 < r < 1 \text{ (внутри круга),} \\ u_{1}(1, \theta) = -\sin \theta \frac{\partial u_{0}(1, \theta)}{\partial r} = -\sin \theta \cos \theta \end{cases}
J[y] = \int_{a}^{y} F(x, y, y') dx
J[y] = \int_{a}^{b} F(x, y, y') dx
y(a) = A
J[\overline{y}] \leqslant J[\overline{y} + \varepsilon \eta]
u_u + \frac{1}{4} u_t
\varphi(\varepsilon) = J[\bar{y} + \varepsilon \eta]
u_t = (1+x) u_{xx}, \\ u(x, 0) = q(x), \quad -\infty < x < \infty,
\frac{d\varphi\left(\varepsilon\right)}{d\varepsilon} = \frac{d}{d\varepsilon} J\left[\overline{y} + \varepsilon\eta\right]|_{\varepsilon=0} =
u_1 = (1 + \varepsilon_1) u_{xx}
\end{verbatim}

(Читач має робити це самостійно.) Інтегруючи на частини, ми отримуємо

Оскільки цей інтеграл є нульовим у будь-якій функції
\texttt{\textbackslash{}eta(x)} , що задовольняє граничні умови
\texttt{\textbackslash{}eta(a)\ =\ \textbackslash{}eta(b)\ =\ 0} , ми Р
и с

\pandocbounded{\includegraphics[keepaspectratio]{images/_page_330_Figure_6.jpeg}}

. 44 3. Граф функції
\texttt{J{[}y+{]}\ \textbackslash{}in\ \textbackslash{}eta} {]} в
околиці \texttt{\textbackslash{}varepsilon\ =\ 0} .

Ми приходимо до висновку, що друга частина підоб'єктивного виразу має
дорівнювати нулю, тобто

(44 2)

(Рівняння Ейлера--Лагранжа).

\begin{verbatim}
Рівняння (44 2) називається рівнянням Ейлера-Лагранжа, і хоча загалом воно здається комплексним, коли в нього підставляється конкретна функція F(x, y, y'), воно стає звичайним диференціальним рівнянням другого порядку відносно невідомої функції  \overline{y}(x)  . Отже, щоб визначити мінімізуючу функцію  \overline{y}  , необхідно розв'язати рівняння Ейлера-Лагранжа.
\end{verbatim}

Отже, ми показали, що

\begin{verbatim}
якщо функція y(x) мінімізує функціональний  J[y] = \int_a^b F(x, y, y') dx  (у класі гладких функцій з граничними умовами y(a) = A; y(b) = B), то вона повинна задовольняти рівняння
\end{verbatim}

(з граничною умовою).

(Щоб спростити позначення, ми прибрали лінію над
\texttt{\textbackslash{}bar\{y\}} .) Щоб проілюструвати теорію, наведемо
приклад.

Пошук мінімальної функціональності
\texttt{J{[}y{]}\ =\ \textbackslash{}int\_\{0\}\^{}\{1\}\ {[}y\^{}2\ +\ y\textquotesingle{}\^{}2{]}\ dx}
.

Спробуємо знайти функцію \texttt{y(x)}, яка проходить через точки (0,0)
і (1,1) і мінімізує функціональність

\begin{verbatim}
P_{0}
\end{verbatim}

З контексту зрозуміло, що шукана функція має бути диференційовною,
оскільки інтегративний вираз залежить від y'. Щоб визначити
\texttt{\textbackslash{}bar\{y\}(x)} , запишемо рівняння
Ейлера--Лагранжа (разом із граничною умовою \texttt{y(0)}=0 та
\texttt{y(1)}=1).

\begin{verbatim}
,  (x,  y, z) \in R^3  ,  u = 0   u_t =
\end{verbatim}

\texttt{y(0)\ =\ 0,} \texttt{y(1)\ =\ 1.} Але
\texttt{F(x,\ y,\ y\textquotesingle{})\ =\ y\^{}2\ +\ y\textquotesingle{}\^{}2}
і, отже,

\begin{enumerate}
\def\labelenumi{\arabic{enumi}.}
\setcounter{enumi}{3}
\tightlist
\item
  Решите аналогичную (см. задачу 3) задачу для двумерного уравнения
\end{enumerate}

(УЧП)

(Продифференцировали \texttt{F} по \texttt{y} и
\texttt{y\textquotesingle{}.} )

Уравнение Эйлера --- Лагранжа принимает вид

Зверніть увагу, що в перших двох задачах коефіцієнти є сталими. Читач
повинен розуміти, що параметр b має бути достатньо малим, а необережний
ряд може відрізнятися

\section{ЗАДАЧИ}\label{ux437ux430ux434ux430ux447ux438-1}

\begin{verbatim}
- 1. Подставьте разложение (46.8) в задачу (46.7) и получите последовательность задач P_0,\ P_1,\ P_2,\ \dots . 2. Покажите, что нелинейную задачу
\end{verbatim}

откуда

(НУ)

Решая это простое дифференциальное уравнение с граничными условиями
\texttt{y(0)}=0 и \texttt{y(1)}=1, получаем

\begin{verbatim}
u(r, \theta) = u_0(r, \theta) + \frac{1}{4}u_1(r, \theta)
\end{verbatim}

.

График этой функции изображен на рис. 44.4. Ясно, что любая другая
гладкая функция, удовлетворяющая тем же граничным условиям, даст большее
значение функционала \texttt{J(y)}.

\pandocbounded{\includegraphics[keepaspectratio]{images/_page_332_Figure_5.jpeg}}

Рис. 44 4. Все допустимые гладкие кривые удовлетворяют граничным
условиям \texttt{y(0)} = 0 и \texttt{y(1)} = 1: a---допустимая кривая;
\texttt{\textbackslash{}delta} ---минимизирующая функция
\texttt{y(x)\ =\ 0.42e\^{}x\ -\ 0.42e\^{}\{-x\}} . \# ЗАУВАЖЕННЯ 1.
Уравнение Эйлера---Лагранжа аналогично равенству нулю производной в
дифференциальном исчислении. Если читатель помнит, то на самом деле
равенство нулю производной не является достаточным условием экстремума.
Например, производная функции \texttt{f(x)\ =\ x\^{}3} при x = 0 равна
нулю, но эта точка не является нп точкой минимума, ни точкой максимума.
Так же и с уравнением Эйлера---Лагранжа: оно является только необходимым
условием экстремума, но не достаточным. Значит, экстремум может
достигаться на других функциях. Очень часто, однако, решение уравнения
Эйлера---Лагранжа по существу задачи дает локальный или даже глобальный
минимум.

\begin{enumerate}
\def\labelenumi{\arabic{enumi}.}
\setcounter{enumi}{1}
\tightlist
\item
  Если у (х) минимизирует функционал
\end{enumerate}

\begin{verbatim}
\Delta u = 0,
\end{verbatim}

то площадь поверхности, образованной вращением кривой \texttt{y(x)}
вокруг оси x, будет минимальной (рис. 44.5). Решением уравнения Эйлера
--- Лагранжа в этом случае будет

цепная линия

\begin{verbatim}
q_{xx} + \varphi_{uu} = 0
\end{verbatim}

где константы \texttt{\textbackslash{}alpha} и
\texttt{\textbackslash{}beta} определяются из условий \texttt{y(a)} = A
и \texttt{y(b)} = B (т. е. цепная линия должна проходить через заданные
точки). Решить уравнение Эйлера---Лагранжа в данном случае достаточно
трудно, поэтому мы рекомендуем читателю только проверить, что цепная
линия ему удовлетворяет.

\begin{enumerate}
\def\labelenumi{\arabic{enumi}.}
\setcounter{enumi}{2}
\tightlist
\item
  В этой лекции мы нашли, что функция
\end{enumerate}

\begin{verbatim}
q_{uu} + q_{vv} = 0
\end{verbatim}

доставляет минимум функционалу

\pandocbounded{\includegraphics[keepaspectratio]{images/_page_353_Figure_9.jpeg}}

При этом оказывается, что
\texttt{J{[}\textbackslash{}overline\{y\}{]}\ =\ 0.46} . Если подставить
в функционал любую другую гладкую функцию, график которой проходит через
точки (0, 0) и (1, 1), то значение функционала J{[}y{]} станет больше.

\pandocbounded{\includegraphics[keepaspectratio]{images/_page_333_Picture_8.jpeg}}

Рис. 44 5. Минимальная поверхность вращения.

\begin{enumerate}
\def\labelenumi{\arabic{enumi}.}
\setcounter{enumi}{3}
\tightlist
\item
  Основные законы физики чаще всего формулируются на языке вариационных
  принципов, а не дифференциальных уравнений. В качестве примеров можно
  привести принцип Ферма (свет при распространении из одной точки в
  другую выбирает путь, которому соответствует наименьшее время
  распространения) или принцип Гамильтона (в консервативном поле частица
  движется так, что интеграл действия
\end{enumerate}

\begin{verbatim}
w = f(z)
\end{verbatim}

(кинетическая энергия--- потенциальная энергия) \texttt{dt}

будет минимальным). Таким образом, физические явления развиваются только
так, что эти функционалы принимают минимальные значения.

Основные идеи вариационного исчисления можно распространить на
функционалы, зависящие от функции нескольких независимых переменных,
вида

\pandocbounded{\includegraphics[keepaspectratio]{images/_page_355_Figure_5.jpeg}}

Соответствующее уравнение Эйлера --- Лагранжа для такого функционала
будет уже уравнением с частными производными

\pandocbounded{\includegraphics[keepaspectratio]{images/_page_355_Figure_7.jpeg}}

Такого рода функционалы мы рассмотрим в следующей лекции. Однако общая
философия метода будет иной. Некоторым новым способом мы найдем функцию
\texttt{u(x,\ y)}, которая минимизирует функционал, а поскольку она
является решением уравнения Эйлера --- Лагранжа, то тем самым мы найдем
решение уравнения с частными производными. Другими словами, мы будем
искать решение уравнения с частными производными методом минимизации
функционала. В этой лекции все было наоборот: задачу минимизации
функционала, зависящего от функции одной переменной, мы сводили к
решению уравнения Эйлера --- Лагранжа. Методы решения дифференциальных
уравнений путем минимизации соответствующих функционалов принято
называть прямыми методами вариационного исчисления. Методы минимизации
функционалов путем решения соответствующих уравнений Эйлера --- Лагранжа
называются непрямыми методами вариационного исчисления. Следующая лекция
посвящена хорошо известному прямому методу Ритца. \# ЗАДАЧІ 1. Среди
всех кривых, удовлетворяющих условиям \texttt{y(0)} = 0, и \texttt{y(1)}
= 1, найдите ту, которая минимизирует функционал

\begin{verbatim}
u + iv = (x + iy)^2 = x^2 - y^2 + 2ixy
\end{verbatim}

Как вы интерпретируете получелный результат? Чему равна величина
\texttt{J{[}\textbackslash{}overline\{y\}{]}} ? Каков смысл величины
\texttt{J{[}\textbackslash{}overline\{y\}{]}} ?

\texttt{2.\ Кинетическая\ энергия\ колеблющейся\ материальной\ точки\ определяется\ выражением\ \ \textquotesingle{}KE\ =\ \textbackslash{}frac\{1\}\{2\}\ m\textbackslash{}dot\{y\}\^{}2\textquotesingle{}\ ,\ где\ y\ =\ dy/dt,\ а\ потенциальная\ равна\ \ \textquotesingle{}PE\ =\ \textbackslash{}frac\{1\}\{2\}\ ky\^{}2\textquotesingle{}\ .}

\pandocbounded{\includegraphics[keepaspectratio]{images/_page_335_Picture_2.jpeg}}

Принцип Гамильтона утверждает, что материальная точка будет двигаться
так, что интеграл

\begin{verbatim}
\left\{
\end{verbatim}

будет минимален. Если этот закоп справедлив, то получите
дифференциальное уравнение движения материальной точки.

\begin{enumerate}
\def\labelenumi{\arabic{enumi}.}
\setcounter{enumi}{2}
\tightlist
\item
  Покажите, что в классе гладких функций, удовлетворяющих граничным
  условиям \texttt{y(0)}=0 и \texttt{y(\textbackslash{}pi/2)=1} ,
  минимум функционала
\end{enumerate}

Чи знаєте ви, як виглядає графік цієї функції з сегментом {[}0, 1{]}?
Подумай про це. Зверніть увагу, що синусоїдальне перетворення функції
\texttt{f(x)} визначається лише для додатних цілих чисел n.~Іншими
словами, скінченний синус і косинус перетворюють функції у чисельні
послідовності.

\pandocbounded{\includegraphics[keepaspectratio]{images/_page_185_Figure_3.jpeg}}

РИС. 25.1. Графики функции \texttt{f(x)} = 1 и ее преобразования.

\section{Свойства
преобразований}\label{ux441ux432ux43eux439ux441ux442ux432ux430-ux43fux440ux435ux43eux431ux440ux430ux437ux43eux432ux430ux43dux438ux439-1}

Прежде чем приступить к решению задач, мы должны получить некоторые
свойства этих преобразований.

\begin{verbatim}
Якщо you(x, t) — функція  \partial \theta yx  змінних, і ми виконуємо перетворення змінної x, то
\end{verbatim}

достигается на функции
\texttt{\textbackslash{}overline\{y\}(x)\ =\ \textbackslash{}sin\ x} .
Вычислите \texttt{J(\textbackslash{}sin\ x)} . 4. Исходя из функционала

(Зверніть увагу, що перетворення було виконано на змінній x, і отримана
послідовність залежить лише від часу t.)

А как быть с производными? Ниже приводятся несколько полезных формул

получите уравнение Эйлера --- Лагранжа

\section{ВАРИАЦИОННЫЕ МЕТОДЫ РЕШЕНИЯ УРАВНЕНИЙ С ЧАСТНЫМИ ПРОИЗВОДНЫМИ
МЕТА ЛЕКЦІЇ: Показать, как можно решать дифференциальное уравнение,
рассматривая его как уравнение Эйлера---Лагранжа для некоторого
функционала и находя некоторым новым методом функцию, минимизирующую
этот функционал. Тогда минимизирующая функция будет решением уравнения с
частными производными. Задача, конечно, состоит в том, чтобы найти такой
функционал, для которого исходное уравнение является уравнением
Эйлера---Лагранжа. Хорошо известен результат (теорема о минимуме
энергии), который гласит, что нахождение решения и некоторой
эллиптической краевой
задачи}\label{ux432ux430ux440ux438ux430ux446ux438ux43eux43dux43dux44bux435-ux43cux435ux442ux43eux434ux44b-ux440ux435ux448ux435ux43dux438ux44f-ux443ux440ux430ux432ux43dux435ux43dux438ux439-ux441-ux447ux430ux441ux442ux43dux44bux43cux438-ux43fux440ux43eux438ux437ux432ux43eux434ux43dux44bux43cux438-ux43cux435ux442ux430-ux43bux435ux43aux446ux456ux457-ux43fux43eux43aux430ux437ux430ux442ux44c-ux43aux430ux43a-ux43cux43eux436ux43dux43e-ux440ux435ux448ux430ux442ux44c-ux434ux438ux444ux444ux435ux440ux435ux43dux446ux438ux430ux43bux44cux43dux43eux435-ux443ux440ux430ux432ux43dux435ux43dux438ux435-ux440ux430ux441ux441ux43cux430ux442ux440ux438ux432ux430ux44f-ux435ux433ux43e-ux43aux430ux43a-ux443ux440ux430ux432ux43dux435ux43dux438ux435-ux44dux439ux43bux435ux440ux430ux43bux430ux433ux440ux430ux43dux436ux430-ux434ux43bux44f-ux43dux435ux43aux43eux442ux43eux440ux43eux433ux43e-ux444ux443ux43dux43aux446ux438ux43eux43dux430ux43bux430-ux438-ux43dux430ux445ux43eux434ux44f-ux43dux435ux43aux43eux442ux43eux440ux44bux43c-ux43dux43eux432ux44bux43c-ux43cux435ux442ux43eux434ux43eux43c-ux444ux443ux43dux43aux446ux438ux44e-ux43cux438ux43dux438ux43cux438ux437ux438ux440ux443ux44eux449ux443ux44e-ux44dux442ux43eux442-ux444ux443ux43dux43aux446ux438ux43eux43dux430ux43b.-ux442ux43eux433ux434ux430-ux43cux438ux43dux438ux43cux438ux437ux438ux440ux443ux44eux449ux430ux44f-ux444ux443ux43dux43aux446ux438ux44f-ux431ux443ux434ux435ux442-ux440ux435ux448ux435ux43dux438ux435ux43c-ux443ux440ux430ux432ux43dux435ux43dux438ux44f-ux441-ux447ux430ux441ux442ux43dux44bux43cux438-ux43fux440ux43eux438ux437ux432ux43eux434ux43dux44bux43cux438.-ux437ux430ux434ux430ux447ux430-ux43aux43eux43dux435ux447ux43dux43e-ux441ux43eux441ux442ux43eux438ux442-ux432-ux442ux43eux43c-ux447ux442ux43eux431ux44b-ux43dux430ux439ux442ux438-ux442ux430ux43aux43eux439-ux444ux443ux43dux43aux446ux438ux43eux43dux430ux43b-ux434ux43bux44f-ux43aux43eux442ux43eux440ux43eux433ux43e-ux438ux441ux445ux43eux434ux43dux43eux435-ux443ux440ux430ux432ux43dux435ux43dux438ux435-ux44fux432ux43bux44fux435ux442ux441ux44f-ux443ux440ux430ux432ux43dux435ux43dux438ux435ux43c-ux44dux439ux43bux435ux440ux430ux43bux430ux433ux440ux430ux43dux436ux430.-ux445ux43eux440ux43eux448ux43e-ux438ux437ux432ux435ux441ux442ux435ux43d-ux440ux435ux437ux443ux43bux44cux442ux430ux442-ux442ux435ux43eux440ux435ux43cux430-ux43e-ux43cux438ux43dux438ux43cux443ux43cux435-ux44dux43dux435ux440ux433ux438ux438-ux43aux43eux442ux43eux440ux44bux439-ux433ux43bux430ux441ux438ux442-ux447ux442ux43e-ux43dux430ux445ux43eux436ux434ux435ux43dux438ux435-ux440ux435ux448ux435ux43dux438ux44f-ux438-ux43dux435ux43aux43eux442ux43eux440ux43eux439-ux44dux43bux43bux438ux43fux442ux438ux447ux435ux441ux43aux43eux439-ux43aux440ux430ux435ux432ux43eux439-ux437ux430ux434ux430ux447ux438}

\begin{verbatim}
w = \ln\left\{\frac{z-1}{z+1}\right\}
\end{verbatim}

, в области D , u = 0 , на границе D ,

эквивалентно нахождению такой функции u (тоже обращающейся в нуль на
границе области D), которая минимизирует функционал потенинальной
энергии

\texttt{То\ есть,\ \ \textquotesingle{}\textbackslash{}Delta\ u\ =\ f\textquotesingle{}\ —уравнение\ Эйлера—Лагранжа\ для\ J{[}u{]}.\ Функция,\ приближенно\ минимизирующая\ функционал\ J{[}u{]},\ ищется\ методом\ Ритца;\ тем\ самым\ мы\ получаем\ решение\ (приближениое)\ уравнения\ с\ частными\ производными.\ В\ этой\ лекции\ мы\ познакомимся\ с\ методом\ Ритца\ и\ посмотрим,\ как\ с\ помощью\ этого\ метода\ можно\ и\ и\ и\ и\ и\ и\ и\ уравнения.}

Очень хорошо решать граничные задачи (как, например, задачу о
закрепленной мембране) путем отыскания гладкой поверхности, которая
минимизирует потенциальную энергию мембраны. То есть, если мы
рассматриваем исходное уравнение как уравнение Эйлера --- Лагранжа для
искоторого функционала J{[}u{]}, то решение дифференциального уравнения
можно получить путем минимизации функционала (поскольку минимизирующая
функция функционала является решением соответствующего уравнения

Эйлера --- Лагранжа). В лекции 44 мы рассматривали такие функционалы, у
которых уравнение Эйлера --- Лагранжа было обыкновенным дифференциальным
уравнением. В этой лекции рассматриваются функционалы, для которых
уравнение Эйлера --- Лагранжа относится к классу уравнений с частными
производными. Например для функционала

уравнение Эйлера---Лагранжа (подобно случаю обыкновенного
дифференциального уравнения в последней лекции) имеет вид

и, следовательно, для решения задачи Дирихле в единичном квадрате

(УЧП)

\begin{verbatim}
(\text{УЧП}) \quad u_{tt} = u_{xx} + \sin(\pi x), \quad 0 < x < 1, \quad 0 < t < \infty,
\end{verbatim}

, \texttt{0\ \textless{}\ x\ \textless{}\ 1} ,
\texttt{0\ \textless{}\ y\ \textless{}\ 1} , (ГУ) \texttt{u\ =\ g} , на
границе квадрата,

можно, наоборот, найти такую функцию \texttt{u(x,\ y)}, которая
минимизирует J{[}u{]} и равна g на границе. Вероятно, не стоит
удивляться тому, что функционал

\begin{verbatim}
(\text{ГУ}) \quad
\end{verbatim}

представляет собой потенциальную эпергию мембраны и фактически мы
находим поверхность с минимальной потенциальной энерешей. Вопрос,
конечно, состоит в следующем: если дано дифференциальное уравнение, как
мы находим функционал, выражающий потенциальную энергию решения? Ответ
на этот вопрос дает хорошо известная теорема (теорема о минимуме
энергии), которая утверждает:

Решение и задачи Дирихле

\begin{verbatim}
^{\Lambda}u = f   в области   ^{D}     (ГУ)     ^{U}=0   на границе   ^{D}
\end{verbatim}

является той же функциен, которая минимизирует (среди функций,
удовлетворяющих граничному условию u=0) энергегический функционал

\begin{verbatim}
\begin{cases} u_{tx}, \quad 0 = \begin{cases} 1, \quad -1 < x < 1, \\ 0 \quad \text{в остальных точках,} \end{cases}
\end{verbatim}

\begin{verbatim}
= \ln\left[\frac{z-1}{z+1}\right] + i \arg\left[\frac{z-1}{z+1}\right]
\end{verbatim}

и, следовательно, для решения задачи (45.1) достаточно найти функцию
\texttt{\textbackslash{}overline\{u\}} (среди функций, обращающихся в
нуль на граннце), минимизирующую
\texttt{J\textbackslash{}\{u\textbackslash{}\}} . Этим замечанием
завершается первая часть лекции. В оставшейся части мы покажем, как
находить минимизирующую функцию \texttt{\textbackslash{}overline\{u\}}
для J{[}u{]}, используя метод Ритца.

\section{Метод Ритца минимизации
функционалов}\label{ux43cux435ux442ux43eux434-ux440ux438ux442ux446ux430-ux43cux438ux43dux438ux43cux438ux437ux430ux446ux438ux438-ux444ux443ux43dux43aux446ux438ux43eux43dux430ux43bux43eux432}

Этот метод --- один из многих вариационных методов, описанных в
литературе по вариационному исчислению. Идея, предложенная математиком
В. Ритцем, очень проста. Метод состоит из следующих шагов.

ШАГ 1. Выбираем n и функцию u, минимизирующую функционал

\begin{verbatim}
= \arg\left[\frac{x^2+y^2-1+i2y}{(x+1)^2+y^2}\right] =
\end{verbatim}

ищем в виде

\begin{verbatim}
\begin{cases} u(x, 0) = 1, \\ u_t(x, 0) = 0, \end{cases}
\end{verbatim}

где функции
\texttt{\textbackslash{}phi\_1,\ \textbackslash{}phi\_2,\ \textbackslash{}phi\_3,\ \textbackslash{}ldots,\ \textbackslash{}phi\_n}
принадлежат к классу достаточно хороших функций и все обращаются в нуль
на границе, так что из них можно построить разумное приближение для
решения задачи. Эти функции принято называть пробными функциями.
Типичный набор пробных функций для задачи Дирихле в единичном квадрате
приведен ниже:

\begin{verbatim}
x^{2} + y^{2} = 1,
\end{verbatim}

\texttt{—\ обращаются\ в\ нуль\ на\ границе,\ \ \textquotesingle{}\textbackslash{}varphi\_2(x,\ y)\ =\ x\textbackslash{}varphi\_1(x,\ y),\textquotesingle{}\ \ \ \textquotesingle{}\textbackslash{}varphi\_3(x,\ y)\ =\ y\textbackslash{}varphi\_1(x,\ y),\textquotesingle{}\ \ \ \textquotesingle{}\textbackslash{}varphi\_4(x,\ y)\ =\ x\^{}2\textbackslash{}varphi\_1(x,\ y),\textquotesingle{}\ \ \ \textquotesingle{}\textbackslash{}varphi\_5(x,\ y)\ =\ xy\textbackslash{}varphi\_1(x,\ y),\textquotesingle{}\ \ \ \textquotesingle{}\textbackslash{}varphi\_6(x,\ y)\ =\ y\^{}2\textbackslash{}varphi\_1(x,\ y),\textquotesingle{}\ \ \ \textquotesingle{}\textbackslash{}vdots\textquotesingle{}\ \ \ \textquotesingle{}\textbackslash{}vdots\textquotesingle{}}

Иными словами, первые четыре приближения имеют вид

Решение задачи выполним за три следующих шага:

ШАГ І (Определение подходящего преобразования.)

Поскольку переменная x меняется от 0 до 1, мы воспользуемся конечным
преобразованием, а именно синус-преобразованием; позже станет ясно
почему. Мы могли бы решить эту задачу с помощью преобразования Лапласа
по переменной t (по трудоемкости этот метод не отличается от конечного
синус-преобразования).

ШАГ 2 (Выполнение преобразования).

Для удобства будем использовать обозначение \texttt{S\_n}(t) = S{[}u{]}
. Применим синус-преобразование к исходному уравнению

\begin{verbatim}
\begin{aligned} & u_{ti} = c^{2}u_{xx}, & -\infty < x < +\infty, & 0 < t < \infty, \\ & (\text{НУ}) & \begin{cases} u_{t}(x, 0) = 0, \\ u_{t}(x, 0) = \begin{cases} 1, & -1 < x < 1, \\ 0, & \text{в остальных случаях.} \end{cases} \end{aligned}
\end{verbatim}

и воспользуемся тождеством для синус-преобразования, получим

в координатах (x, t).

\pandocbounded{\includegraphics[keepaspectratio]{images/_page_134_Figure_7.jpeg}}

Рис. 18.4. Решение задачи (18.4) в xt-плоскости.

Задача (18.4) описывает поведение струны, которой сообщается рачальная
единичная скорость на отрезке -1 \textless{} x \textless{} 1. Смещение .
.

описывается формулой Даламбера

\begin{verbatim}
\begin{aligned} & (\text{УЧП}) & \phi_{xx} + \phi_{yy} = 0, & \text{внутри } D, \\ & (\text{ГУ}) & \begin{cases} \varphi(x, y) = 1 & \text{на } x^2 + y^2 = 1, \\ \varphi(x, y) = 2 & \text{на } (x - 1)^2 + y^2 = 9. \end{cases} \end{aligned}
\end{verbatim}

\texttt{является\ теперь\ функцией\ коэффициентов\ \ \textquotesingle{}a\_1a\_2\textbackslash{}dots\ a\_n\textquotesingle{}\ .\ Следовательно,\ для\ того\ чтобы\ найти\ минимум\ функционала\ J,\ приравняем\ нулю\ частные\ производные}

\begin{verbatim}
Aa = b
b_i = -\int_0^1 \int_0^1 f(x, y) \varphi_i(x, y) dx dy,
\end{verbatim}

где

Это решение в различные моменты времени изображено на рис. 18.5.

\pandocbounded{\includegraphics[keepaspectratio]{images/_page_135_Figure_5.jpeg}}

Рис. 18.5. Решение задачи (18.4) в различные моменты времени.

Этим завершается наша интерпретация формулы Даламбера в плоскости
переменных \texttt{\textbackslash{}boldsymbol\{x\}} и t. В оставшейся
части лекции мы займемся решением смешанной задачи для полубесконечной
струны:

(18.6)

\begin{verbatim}
u = \gamma [(x-s)(x-t)-\gamma^2],
\end{verbatim}

Единственное отличие в том, что аппроксимирующие функции
\texttt{\textbackslash{}phi\_1,\ \textbackslash{}phi\_2,\ \textbackslash{}phi\_3,\ \textbackslash{}ldots,\ \textbackslash{}phi\_n}
должны удовлетворять граничным условиям:

\begin{verbatim}
\gamma = 2t/[(x-t)^2 + y^2].
\end{verbatim}

является иравнение Лапласа, но это доказательство проводится аналогично
тому, как это делалось в лекции 44 для функционала, зависящего от
функции одной переменной.

\begin{enumerate}
\def\labelenumi{\arabic{enumi}.}
\setcounter{enumi}{2}
\item
  В книгах по вариационному исчислению показано, как строить
  энергетический функционал не только для уравнения Пуассона с граничным
  условием Дирихле u=0, но и для многих других дифференциальных
  уравнений и многих других типов граничных условий. Следовательно,
  большое число краевых задач можно решить, минимизируя соответствующие
  энергетические функционалы.
\item
  Чем большее значение п мы выбираем, тем меньшее значение
  \texttt{J{[}\textbackslash{}bar\{u\}\_n{]}} получается. Значит, с
  увеличением n растет точность л приближенного решения
  дифференциального уравнения. Один из методов определения и, основан на
  последовательном вычислении функционала \texttt{J{[}u\_n{]}} для
  возрастающих значений n.~Вычисления прекращаются, если с ростом п
  функционал перестает практически уменьшаться.
\item
  Если и велико, то проведение вычислений по методу Ритца требует
  привлечения ЭВМ. На рис. 45.2 (с. 345) приведена блок-
\end{enumerate}

схема решения краевой задачи методом Ритца.

Пользователь должен позаботиться о подпрограмме, которая должна
вычислять значения пробных функций. Такая подпрограмма может иметь вид:

SUBROUTINE BC (X.Y.PHI)

\begin{itemize}
\tightlist
\item
  SUBROUTINE PROVIDED BY THE USER TO C - EVALUATE THE - FUNCTIONS
  PHI(1), PHI(2), \ldots, PHI(N) DIMENSION PHI(20)
\end{itemize}

\begin{verbatim}
PHI(1) = X\*Y\*(1-X)\*(1-Y)
\end{verbatim}

\texttt{PHI(2)\ =\ X\^{}*PHI(1)} (используемые функции)

\begin{verbatim}
PHI(3) = Y\*PHI(1)
\end{verbatim}

PHI(N) = (что получится)RETURN END

\section{\texorpdfstring{\textbf{ЗА}ДАЧИ}{ЗАДАЧИ}}\label{ux437ux430ux434ux430ux447ux438-2}

Найдите энергетический функционал для задачи

(УЧП)

Если теперь преобразовать начальные условия краевой задачи, то получим
начальные условия для обыкновенного дифференциального уравнения

, \texttt{0\ \textless{}\ x\ \textless{}\ 1} ,
\texttt{0\ \textless{}\ y\ \textless{}\ 1} , (ГУ) \texttt{u\ =\ 0} на
границе.

\begin{enumerate}
\def\labelenumi{\arabic{enumi}.}
\setcounter{enumi}{1}
\tightlist
\item
  Как минимизировать функционал
\end{enumerate}

методом Ритца?

УКАЗАНИЕ. Если ввести новую функцию \texttt{z(x)}:

то можно заметить, что она удовлетворяет граничным условиям
\texttt{z(0)} = 0 и \texttt{z(1)} = 0.1

\begin{itemize}
\tightlist
\item
  \begin{enumerate}
  \def\labelenumi{\arabic{enumi}.}
  \setcounter{enumi}{2}
  \tightlist
  \item
    Напишите программу для проведения вычислений по блоксхеме,
    изображенной на рис. 45 2.
  \end{enumerate}
\item
  \begin{enumerate}
  \def\labelenumi{\arabic{enumi}.}
  \setcounter{enumi}{3}
  \tightlist
  \item
    Покажите, что уравнение
  \end{enumerate}
\end{itemize}

\begin{verbatim}
\varphi(r) = a \ln r + b, \quad r = u^2 + v^2.
\end{verbatim}

является уравнением Эйлера --- Лагранжа для функционала

\begin{verbatim}
\varphi(x, y) = 0.57 \ln(u^2 + v^2) + 1,
\end{verbatim}

\begin{enumerate}
\def\labelenumi{\arabic{enumi}.}
\setcounter{enumi}{4}
\tightlist
\item
  Задачу Дирихле
\end{enumerate}

Итак, решим теперь новое семейство задач Кощи:

(ОДУ)

\texttt{где\ A=0.06\ и\ B=0.04.\ Как\ найти\ потенциальную\ энергию\ этого\ решения?\ Мы\ рекомендуем\ читателю\ вспомнить\ синуспреобразование\ и\ самостоятельно\ получить\ решение\ \ \textquotesingle{}u\textbackslash{}left(x,y\textbackslash{}right)\textquotesingle{}\ .}

*) При такой замене возникает особенность в т. x=1. Лучше ввести функцию
\texttt{z(x)} по формуле \texttt{z(x)}=\texttt{y(x)}=x. Прим. ред.

\pandocbounded{\includegraphics[keepaspectratio]{images/_page_343_Figure_2.jpeg}}

Далсе. Приближенное решение
\texttt{u\_n(x,y)\ =\ a\_1\ \textbackslash{}varphi\_1\ +\ a\_2\ \textbackslash{}varphi\_2\ +\ \textbackslash{}dots\ +\ a\_n\ \textbackslash{}varphi\_n}
. В это время можно численно найти эначение функционала
\texttt{J{[}u\_n{]}} . Если необходимо найти значение решения
\texttt{u\_n(x,y)} в некоторой точке, можно написать небольшую программу
эля быполнения отой операции. При необходимости можно повторить
вычисления по этой программе с другими значениями п Рис. 452. Блок-схема
метода Ритца.

\section{РЕШЕНИЕ УРАВНЕНИЙ С ЧАСТНЫМИ ПРОИЗВОДНЫМИ МЕТОДАМИ ТЕОРИИ
ВОЗМУЩЕНИЙ МЕТА ЛЕКЦІЇ: Показать, как различные сложные задачи
(нелинейные уравнения с переменными коэффициентами, области неправильной
формы и т. д.) можно решить методом возмущения более простых задач.
Иными словами, мы хотим показать, как следует изменить решение простой
задачи, чтобы оно давало приближенное решение трудной
задачи.}\label{ux440ux435ux448ux435ux43dux438ux435-ux443ux440ux430ux432ux43dux435ux43dux438ux439-ux441-ux447ux430ux441ux442ux43dux44bux43cux438-ux43fux440ux43eux438ux437ux432ux43eux434ux43dux44bux43cux438-ux43cux435ux442ux43eux434ux430ux43cux438-ux442ux435ux43eux440ux438ux438-ux432ux43eux437ux43cux443ux449ux435ux43dux438ux439-ux43cux435ux442ux430-ux43bux435ux43aux446ux456ux457-ux43fux43eux43aux430ux437ux430ux442ux44c-ux43aux430ux43a-ux440ux430ux437ux43bux438ux447ux43dux44bux435-ux441ux43bux43eux436ux43dux44bux435-ux437ux430ux434ux430ux447ux438-ux43dux435ux43bux438ux43dux435ux439ux43dux44bux435-ux443ux440ux430ux432ux43dux435ux43dux438ux44f-ux441-ux43fux435ux440ux435ux43cux435ux43dux43dux44bux43cux438-ux43aux43eux44dux444ux444ux438ux446ux438ux435ux43dux442ux430ux43cux438-ux43eux431ux43bux430ux441ux442ux438-ux43dux435ux43fux440ux430ux432ux438ux43bux44cux43dux43eux439-ux444ux43eux440ux43cux44b-ux438-ux442.-ux434.-ux43cux43eux436ux43dux43e-ux440ux435ux448ux438ux442ux44c-ux43cux435ux442ux43eux434ux43eux43c-ux432ux43eux437ux43cux443ux449ux435ux43dux438ux44f-ux431ux43eux43bux435ux435-ux43fux440ux43eux441ux442ux44bux445-ux437ux430ux434ux430ux447.-ux438ux43dux44bux43cux438-ux441ux43bux43eux432ux430ux43cux438-ux43cux44b-ux445ux43eux442ux438ux43c-ux43fux43eux43aux430ux437ux430ux442ux44c-ux43aux430ux43a-ux441ux43bux435ux434ux443ux435ux442-ux438ux437ux43cux435ux43dux438ux442ux44c-ux440ux435ux448ux435ux43dux438ux435-ux43fux440ux43eux441ux442ux43eux439-ux437ux430ux434ux430ux447ux438-ux447ux442ux43eux431ux44b-ux43eux43dux43e-ux434ux430ux432ux430ux43bux43e-ux43fux440ux438ux431ux43bux438ux436ux435ux43dux43dux43eux435-ux440ux435ux448ux435ux43dux438ux435-ux442ux440ux443ux434ux43dux43eux439-ux437ux430ux434ux430ux447ux438.}

Например, мы покажем, как нелинейную задачу

(УЧП)

\begin{verbatim}
w = \ln \left[ \frac{z-1}{z+1} \right].
\end{verbatim}

, \texttt{0\ \textless{}\ r\ \textless{}\ 1} ,\\
\texttt{(ГУ)\ \ \textquotesingle{}u(1,\ \textbackslash{}theta)\ =\ \textbackslash{}cos\ \textbackslash{}theta\textquotesingle{}\ ,\ \ \textquotesingle{}0\ \textbackslash{}le\ \textbackslash{}theta\ \textbackslash{}le\ 2\textbackslash{}pi\textquotesingle{}\ .}

Очень часто задачу, решение которой не известно, можно путем
непрерывного изменения каких-то параметров свести к близкой, но легко
решаемой задаче. Тогда и решение исходной задачи будет близко решению
модифицирующей задачи.

Например, уравнение Лапласа

\begin{verbatim}
\begin{cases} 1, & n = 1, \\ 0, & n = 2, 3, \dots, \end{cases}
\end{verbatim}

в помощью семейства уравнений

\begin{verbatim}
,   0 \le \varepsilon \le 1  ,
\right.
,   0 < r < 1  ,   0 \le \theta < 2\pi  ,   u(1, \theta) = g(\theta)  ,   0 \le \theta < 2\pi  ,
\end{verbatim}

\begin{verbatim}
при одних граничных условиях (например,   g(\theta) = \cos \theta  ) имеет два различных вещественных решения, а при других (например,   g(\theta) = A \cos \theta   с A > 20,65) эта задача не имеет ни одного вещественного решения. — Прим. ред.
\end{verbatim}

4) Краевая задача

\begin{verbatim}
можно непрерывно модифицировать в нелинейное уравнение   \Delta u = u^2 = 0   1).
\end{verbatim}

На рис. 46.1 изображена схема решения нелинейной задачи путем сведения
ее к решению уравнения Лапласа2)

\pandocbounded{\includegraphics[keepaspectratio]{images/_page_362_Figure_1.jpeg}}

Рис. 46.1. Схема метода возмущении.

Чтобы найти решение возмущенного уравнения Лапласа

\begin{verbatim}
u(0, t) = f(t),
u_{x}(0, t) = 0.
\end{verbatim}

\pandocbounded{\includegraphics[keepaspectratio]{images/_page_362_Figure_2.jpeg}}

\begin{verbatim}
x + ct = \text{const},
\end{verbatim}

\pandocbounded{\includegraphics[keepaspectratio]{images/_page_362_Figure_5.jpeg}}

\begin{verbatim}
0 < r < 1,     0 \le \theta < 2\pi,     u(1, 0) = 0,     0 \le \theta < 2\pi,
\end{verbatim}

нмеет два решения (одно --- нулсьое, другое --- нетривиальное), в то
время как задача

\pandocbounded{\includegraphics[keepaspectratio]{images/_page_362_Figure_7.jpeg}}

\begin{verbatim}
u(1, \theta) = 0, \quad 0 \le \theta < 2\pi,
\end{verbatim}

\begin{verbatim}
\begin{array}{ccccccccccccccccccccccccccccccccccc
\end{verbatim}

\begin{verbatim}
\begin{array}{ccc} (\text{УЧ}\Pi) & \Delta u = 0, & 0 < r < 1, \\ (\text{ГУ}) & u(1, \theta) = \cos \theta, & 0 \leqslant \theta \leqslant 2\pi, \end{array}
H(x-a) =
u_{tt} = c^2 u_{xx}, \quad 0 < x < L, \quad 0 < t < \infty,
\begin{cases} 0, x < a, \\ 1, x \geqslant a, \end{cases}
u(0, t) = g_1(t),
\end{verbatim}

В результате получаем

\begin{verbatim}
u_x(0, t) = g_1(t),
H(x-a) =
u_x(0, t) - \gamma_1 u(0, t) = g_1(t),
\begin{cases} 0, & x < a, \\ 1, & x \geqslant a, \end{cases}
u(0, t) = g_1(t),
u_x(0, t) = g_1(t),
f(x) = \mathcal{F}^{-1}[F] = \frac{1}{\sqrt{2\pi}} \int_{-\infty}^{+\infty} F(\omega) e^{-i\omega x} d\omega \qquad F(\omega) = \mathcal{F}[f] = \frac{1}{\sqrt{2\pi}} \int_{-\infty}^{+\infty} f(x) e^{-i\omega x} dx
u_x(0, t) - \gamma_1 u(0, t) = g_1(t),
\begin{cases} \Delta u_{1} = -r^{2} \cos^{2} \theta = -\frac{r^{2}}{2} \left[1 + \cos(2\theta)\right] = -\frac{r^{2}}{2} - \frac{r^{2}}{2} \cos(2\theta), \\ u_{1}(1, \theta) = 0. \end{cases}
\begin{cases} 1, |x| < a \\ 0, |x| > a \end{cases}
\Delta u = -\frac{r^2}{2} - \frac{r^3}{2} \cos{(2\theta)}
\end{verbatim}

\section{ЗАМЕЧ
АНИЯ}\label{ux437ux430ux43cux435ux447-ux430ux43dux438ux44f}

\begin{enumerate}
\def\labelenumi{\arabic{enumi}.}
\tightlist
\item
  Для решения задачи методом конечного синус- или косинуспреобразования
  краевые условия при x=0 и x=L должны иметь вид
\end{enumerate}

\begin{verbatim}
u_p(r, \theta) = Ar^4 + Br^4 \cos(2\theta).
\begin{cases} 1, |x| < 1 \\ 0, |x| > 1 \end{cases}
u_p(r, \theta) = -\frac{r^4}{32} - \frac{r^4}{24} \cos(2\theta).
\sum_{n=0}^{\infty} [a_n \cos(n\theta) + b_n \sin(n\theta)] - \frac{1}{32} - \frac{1}{24} \cos(2\theta) = 0.
\end{verbatim}

Отметим, что если величина \texttt{u(0,\ t)} положительна, то величина
\texttt{u\_x}(0, t) также положительна, если \texttt{u(L,\ t)}
положительна, то \texttt{u\_x}(L, t)

\pandocbounded{\includegraphics[keepaspectratio]{images/_page_144_Picture_4.jpeg}}

Рис. 19.6. Возникновение неоднородных граничных условий упругого
закрепления: h---коэффициент упругости пружины.

отрицательна. Однородные граничные условия (19.2) можно записать в виде

(19.3)

\begin{verbatim}
P_{2}
\end{verbatim}

\texttt{f(x)} * \texttt{g(x)}

\begin{verbatim}
\Delta u = 0, \quad 0 < r < 1 + \frac{1}{4} \sin \theta,
(1+x^2)^{-1}
\end{verbatim}

досягається в кінцевій точці сімейства граничних умов виду (див. рисунок
46 3)

\begin{verbatim}
Tu_x(0, t) = -\beta u_t(0, t).
Tu_{x}(0, t) = \varphi[u(0, t)],
\frac{a}{x^2+a^2}
Tu_{x}(0, t) = -hu^{3}(0, t)
\frac{2ax}{(x^2+a^2)^2}
mu_{tt}(L, t) = -ku_x(L, t) + mg.
P_{0}
\end{verbatim}

\pandocbounded{\includegraphics[keepaspectratio]{images/_page_189_Picture_10.jpeg}}

Отже, розв'язуючи кожну з задач Дірікла всередині кола для функцій
\texttt{u\_0} , \texttt{u\_1} , \texttt{u\_2} , отримуємо розв'язок

(46.9)

\begin{enumerate}
\def\labelenumi{\arabic{enumi}.}
\setcounter{enumi}{17}
\tightlist
\item
\end{enumerate}

Задачі Діріхле в деформованій області (зверніть увагу, що тут

\begin{verbatim}
Ми вже взяли  \epsilon = 1/4  :)
\end{verbatim}

Як вправу, читача просять знайти пертурбацію першого порядку
\texttt{u\_1} і перевірити, наскільки добре є приблизне розв'язання

\begin{verbatim}
\begin{cases} 1-|x|, |x| < 1 \\ 0, |x| > 1 \end{cases}
\end{verbatim}

якщо розглянути рівняння з параметром

(46.11)

і побачити його розв'язок у вигляді

Підставляючи цю декомпозицію на (46.11), отримуємо таку послідовність
завдань

\begin{verbatim}
u_{x}(0, t) = \frac{h}{T} [u(0, t) - \theta_{1}(t)]
\frac{2}{\pi} \frac{\sin a\omega}{\omega}
u_{tt} = u_{xx}, \quad 0 < x < 1, \quad 0 < t < \infty,
\frac{2}{\pi} \frac{\sin a\omega}{(\omega^2+a^2)^2}
\Delta u = 0,
q_{uu} + q_{vv} = 0
w = f(z)
\end{verbatim}

\begin{verbatim}
\left\{
\end{verbatim}

\begin{verbatim}
\Rightarrow.
\frac{\pi}{2} e^{-a+\omega}
\end{verbatim}

\begin{verbatim}
\begin{aligned} (Y & \Pi) & \varphi_{xx} + \varphi_{yy} = 0, & -\infty < x < \infty, & 0 < y < \infty, \\ (47.1) & \varphi(x, 0) = \begin{cases} 0, & |x| > 1, \\ 1, & |x| \le 1. \end{cases}
\end{verbatim}

\begin{verbatim}
\frac{\pi}{2} e^{-a+\omega}
w = \ln\left\{\frac{z-1}{z+1}\right\}
\frac{2ax}{(x^2+a^2)^2}
\end{verbatim}

\begin{verbatim}
\begin{cases} (\text{УЧП}) & \varphi_{uu} + \varphi_{vv} = 0, \quad -\infty < u < \infty, \quad 0 < v < \pi, \\ (\text{47.2}) & \begin{cases} \varphi(u, 0) = 0, \\ \varphi(u, \pi) = 1. \end{cases}
\end{verbatim}

\begin{verbatim}
\frac{\cos(ax), |x| > \pi/2a}{0, |x| > \pi/2a}
u(x, t) = \sum_{n=1}^{\infty} c_n X_n(x) T_n(t),
u_{tt} = \alpha^2 u_{xx}, \quad 0 < x < L, \quad 0 < t < \infty,
(\Gamma Y)
\end{verbatim}

Щоб візуально зрозуміти, як працює ця функція, ми радимо читачеві
побудувати її за кількома лініями y=c.

У другому прикладі ми перетворюємо область між двома неконцентричними
колами у кільце.

Задача Діріхле в області між двома некопоцентричними колами

Припустимо, ви хочете знайти потенціал між двома колами

\begin{verbatim}
\cos(ax)
\quad 0 < t < \infty,
\frac{\pi}{2} [\delta(\omega+a) + \delta(\omega-a)]
\end{verbatim}

(HY)

\begin{verbatim}
u = \gamma [(x-s)(x-t)-\gamma^2],
\quad 0 \le x \le L,
u(x, t) = X(x) T(t).
F(\omega) = \frac{?}{\pi} \int_{0}^{\infty} f(x) \sin(\omega x) dx
T'' - \alpha^2 \lambda T = 0,
X'' - \lambda X = 0,
u(x, t) = [C \sin(\beta x) + D \cos(\beta x)][A \sin(\alpha \beta t) + B \cos(\alpha \beta t)]
\beta_n = \frac{n\pi}{L}, \quad n = 0, 1, 2, \ldots
u_n(x, t) = X_n(x) T_n(t) = \sin(n\pi x/L) [a_n \sin(n\pi at/L) + b_n \cos(n\pi at/L)],
\end{verbatim}

f'\,'(0)

\begin{verbatim}
u_n(x, t) = R_n \sin\left(\frac{n\pi x}{L}\right) \cos\left(\frac{n\pi\alpha (t - \delta_n)}{L}\right).
\frac{1}{\alpha} F\left(\frac{\omega}{\alpha}\right)
H(x-a) =
\frac{2\omega}{\pi (\alpha^2 + \omega^2)}
\end{verbatim}

Функція Хевісайду,

\begin{verbatim}
[2/\pi\omega]^{1/2}
\begin{cases} 0, & x < a, \\ 1, & x \geqslant a, \end{cases}
\frac{2}{\pi\omega} [1-\cos(\omega a)]
f(x) = \mathcal{F}^{-1}[F] = \frac{1}{\sqrt{2\pi}} \int_{-\infty}^{+\infty} F(\omega) e^{-i\omega x} d\omega \qquad F(\omega) = \mathcal{F}[f] = \frac{1}{\sqrt{2\pi}} \int_{-\infty}^{+\infty} f(x) e^{-i\omega x} dx
e^{-1}a + \omega
\begin{cases} 1, |x| < a \\ 0, |x| > a \end{cases}
\frac{1}{2}e^{-\alpha}\sin\omega
\delta(x-a)
a_{n} = \frac{2}{n\pi\alpha} \int_{0}^{L} g(x) \sin(n\pi x/L) dx, \\ b_{n} = \frac{2}{L} \int_{0}^{L} f(x) \sin(n\pi x/L) dx.
u(x, t) = \sum_{n=1}^{\infty} \sin(n\pi x/L) \left[ a_n \sin(n\pi \alpha t/L) + b_n \cos(n\pi \alpha t/L) \right],
u(x, t) = \sum_{n=1}^{\infty} b_n \sin(n\pi x/L) \cos(n\pi\alpha t/L)
u(x, \theta) = f(x)
f(x) = \int_{0}^{\infty} F(\omega) \cos(\omega x) d\omega \qquad \qquad F(x) = \frac{2}{\pi} \int_{0}^{\infty} f(x) \cos(\omega x) dx
f(x) = \sum_{n=1}^{\infty} b_n \sin(n\pi x/L).
u_n(x, t) = b_n \sin(n\pi x/L) \cos(n\pi\alpha t/L).
u(x, t) = \sin(\pi x/L)\cos(\pi \alpha t/L) + 0.5\sin(3\pi x/L)\cos(3\pi \alpha t/L) + 0.25\sin(5\pi x/L)\cos(5\pi \alpha t/L).
R_n \sin(n\pi x/L) \cos[n\pi\alpha(t-\delta_n)/L],
\omega_n = \frac{n\pi\alpha}{L} = \frac{n\pi}{L} \sqrt{\frac{T}{\rho}},
S_n = \frac{2}{\pi} \int_0^{\pi} f(x) \sin(nx) dx.
u(x, 0) = \sin(\pi x/L) + 0.5 \sin(3\pi x/L),
y = \frac{\pi}{b - a} (x - a).
u(x, 0) = 0,
f(x) = \sum_{n=1}^{\infty} S_n \sin(nx).
u(x, 0) = \sin(3\pi x/L),
f(x) = \sum_{n=1}^{\infty} S_n \sin(nx)
u(x, 0) =
S_n = \frac{2}{\pi} \int_{0}^{\pi} f(x) \sin(nx) dx
\begin{cases} 2hx, & 0 \le x \le 0.5, \\ 2h(1-x), & 0.5 \le x \le 1. \end{cases}
0 \le x \le \pi
\end{verbatim}

К а к

\pandocbounded{\includegraphics[keepaspectratio]{images/_page_153_Picture_8.jpeg}}

ие движения будет совершать струна, если ее отпустить? 6. Решите задачу
о затухающих колебаниях струны

\begin{verbatim}
n = 1, 2, ...
\frac{2ax}{(x^2+a^2)^2}
\end{verbatim}

f'\,'(x)

Представляется ли вам разумным полученное решение? Удовлетворяет ли оно
уравнению с частными производными, граничным и начальным условиям?

\begin{enumerate}
\def\labelenumi{\arabic{enumi}.}
\setcounter{enumi}{6}
\tightlist
\item
  Как вы будете решать неоднородное уравнение с заданными начальными и
  граничными условиям?
\end{enumerate}

\begin{verbatim}
\frac{\pi}{2} e^{ax}
\cos(ax)
C_n = \frac{2}{\pi} \int_0^{\pi} f(x) \cos(nx) dx.
\frac{\pi}{2} [\delta(\omega+a) + \delta(\omega-a)]
y = \frac{\pi}{b-a}(x-a).
f(x) = \int_{0}^{\infty} F(\omega) \sin(\omega x) d\omega
f(x) = \frac{C_0}{2} + \sum_{n=1}^{\infty} C_n \cos(nx) \qquad C_n = \frac{2}{\pi} \int_{0}^{\pi} f(x) \cos nx \, dx
0 < x < \infty
\end{verbatim}

\begin{enumerate}
\def\labelenumi{\arabic{enumi}.}
\tightlist
\item
\end{enumerate}

\begin{verbatim}
\frac{1}{\alpha} F\left(\frac{\omega}{\alpha}\right)
x^{-1/2}
\end{verbatim}

f'\,'(x)

\begin{enumerate}
\def\labelenumi{\arabic{enumi}.}
\setcounter{enumi}{4}
\tightlist
\item
  \texttt{H(a-x)}
\end{enumerate}

\begin{verbatim}
\begin{array}{ll} u\left(x, \ 0\right) = \sin{(\pi x)}, \\ u_{t}\left(x, \ 0\right) = 0, \end{array}
\end{verbatim}

\begin{verbatim}
6.  x^{-1} 
\end{verbatim}

\begin{enumerate}
\def\labelenumi{\alph{enumi})}
\tightlist
\item
\end{enumerate}

\begin{verbatim}
7.  \frac{x}{x^2+a^2} 
\end{verbatim}

\begin{enumerate}
\def\labelenumi{\arabic{enumi}.}
\setcounter{enumi}{4}
\tightlist
\item
  \texttt{\textbackslash{}cos(mx),\ m=1,\ 2,\ \textbackslash{}dots}
\end{enumerate}

\begin{verbatim}
9.  \arctan \frac{a}{x} 
\end{verbatim}

\texttt{\textbackslash{}mathbf\{\textbackslash{}Pi\}} для яких рівнянь
функція \texttt{u\_1\ +\ u\_2} також буде розв'язком? Які висновки ви
можете зробити зі своїх відповідей?

\begin{enumerate}
\def\labelenumi{\arabic{enumi}.}
\setcounter{enumi}{3}
\tightlist
\item
  Найдите четыре смещанные задачи, сумма решений которых дает решение
  следующей задачи:
\end{enumerate}

Откуда получаем два уравнения

\begin{verbatim}
0 < x < \infty \qquad \qquad 0 < \omega < \infty
\end{verbatim}

j'\,'(x)

\begin{enumerate}
\def\labelenumi{\arabic{enumi}.}
\setcounter{enumi}{4}
\tightlist
\item
  Решите задачу Коши
\end{enumerate}

(OAY)

\begin{verbatim}
S_n = \frac{2}{\pi} \int_0^{\pi} f(x) \sin(nx) dx.
\end{verbatim}

(HY) \texttt{U\_n(0)\ =\ 0.} Можете ли вы проверить полученное решение?

\begin{enumerate}
\def\labelenumi{\arabic{enumi}.}
\setcounter{enumi}{5}
\tightlist
\item
  Предположим, что функции \texttt{u\_1} и \texttt{u\_2} удовлетворяют
  линейным однородным граничным условиям
\end{enumerate}

\begin{verbatim}
y = \frac{\pi}{b - a} (x - a).
\end{verbatim}

\begin{verbatim}
\begin{cases}  u_x(1, t) + h_2 u(1, t) = 0.  Будет ли функция   u_i + u_i   удовлетворять этим условиям? , .
\end{verbatim}

Лекция 27

\section{УРАВНЕНИЯ ПЕРВОГО ПОРЯДКА (МЕТОД
ХАРАКТЕРИСТИК)}\label{ux443ux440ux430ux432ux43dux435ux43dux438ux44f-ux43fux435ux440ux432ux43eux433ux43e-ux43fux43eux440ux44fux434ux43aux430-ux43cux435ux442ux43eux434-ux445ux430ux440ux430ux43aux442ux435ux440ux438ux441ux442ux438ux43a}

ЦЕЛЬ ЛЕКЦИИ: Ввести понятие уравнения с частными производными первого
порядка (до сих пор мы рассматривали уравнения второго порядка) и
познакомиться с важным методом решения задач с начальными
условиями---методом характеристик. Задача, которой мы теперь займемся,
имеет вид

(УЧП)

\begin{verbatim}
f(x) = \sum_{n=1}^{\infty} S_n \sin(nx).
S_n = \frac{2}{\pi} \int_{0}^{\pi} f(x) \sin(nx) dx
\end{verbatim}

\begin{verbatim}
Сразу же отметим, что впервые мы будем решать задачу с переменными коэффициентами. Оказывается, если от координат (x, t) перейти к новым, характеристическим координатам   (s, \tau)  , то наше УЧП превратится в обыкновенное дифференциальное уравнение. Затем можно решить обыкновенное дифференциальное уравнение, т. е. найти   u(s, \tau)  , а на последнем шаге выразить s и t через x и t и получить u(x, t).
\end{verbatim}

Напомним читателю, что, когда мы решали уравнение дифф

Воспользовавшись ортогональностью семейства функции (sin (nnx)) на
отреже {[}0, 1{]}, находим формулы для коэффициентов

(21.6)

\begin{verbatim}
константа   \alpha^2   играла роль коэффициента диффузии, а v—была скоростью среды. Очевидно, что если   \alpha=0   (т. е. диффузии нет), то решение будет смещаться как целое вдоль оси x со скоростью v (у нас осталась только конвекция). Другими словами, если начальное условие   u(x,0)=\varphi(x)  , то соответствующее решение уравнения
\end{verbatim}

\begin{verbatim}
\frac{\pi}{2} e^{ax}
\end{verbatim}

\begin{verbatim}
будет иметь вид   u(x, t) = \varphi(x-vt)  .
\end{verbatim}

Это дает нам основания думать, что решение уравнения первого порядка

\begin{verbatim}
C_n = \frac{2}{\pi} \int_0^{\pi} f(x) \cos(nx) dx.
\end{verbatim}

переносится вдоль потока со скоростью

\begin{verbatim}
y = \frac{\pi}{b-a}(x-a).
\end{verbatim}

Конечно, если a и b---константы, то решение представляет собой бегущую с
постоянной скоростью волну. Если же \texttt{a(x,\ t)} и
\texttt{b(x,\ t)}

\begin{verbatim}
f(x) = \frac{C_0}{2} + \sum_{n=1}^{\infty} C_n \cos(nx).
\end{verbatim}

\begin{verbatim}
РИС. 27.1. Начальное значение в точке x оказывает влияние на решение только вдоль характеристики в плоскости переменных (x,t). Параметр s изменяется вдоль характеристик от нуля до бесконечности, начиная с линии начального условия;   \tau   постоянна вдоль каждой характеристики;   \tau   изменяется вдоль линии начального условия.
\end{verbatim}

зависят от x и t, то скорость потока изменяется как вдоль потока, так и
во времени (читатель сможет увидеть, что начальная кривая сильно
искажается). Здесь так много аналогий с конвекцией!

Вернемся к нашей основной задаче:

(УЧП)

\begin{verbatim}
0 \le x \le \pi \qquad n = 0, 1, 2, \dots
\end{verbatim}

\begin{verbatim}
(НУ)   u(x, 0) = \varphi(x), -\infty < x < \infty.
\end{verbatim}

Решение этого линейного уравнения первого порядка основывается на
следующем физическом факте: начальное условие в некоторой точке x
переносится в tx-плоскости вдоль линии, которая

называется характеристи-кой (см. рис. 27.1)

Этим наше уразнение отличается от других уравнений (таких, как уравнение
теплопроводности \texttt{u\_t\ =\ u\_\{xx\}\ )},для которых начальное
значение в точке x оказывает влияние на решение во всех точках
пространства и во все моменты времени. Если читатель помнит, то
начальное сме- Р И С

f'\,'(x)

\begin{verbatim}
. 27.2. Начальное возмущение   u\left(x,0\right)   в точке порождает две волны. Начальное возмущение струны согласно волновому уравнению распространяется вдоль двух характеристик.
\end{verbatim}

\begin{verbatim}
щение скрипичной струны в точке x оказывает влияние на решение вдоль   \partial syx   линий в tx-плоскости (соответствующих двум бегущим волнам) (см. рис. 27.2).
\end{verbatim}

\begin{verbatim}
Естественно, напрашивается идея ввести две новые координаты s и   \tau   (вместо x и t) так, чтобы
\end{verbatim}

\begin{itemize}
\item
  s изменялась вдоль характеристических кривых,
\item
  т изменялась вдоль начальной кривой (желательно вдоль линии t=0).
\end{itemize}

Рассмотрим сначала новую координату s. Если s выбрать так, что она
удовлетворяет приведенному выше условию, то уравнение

\texttt{2.\ \textquotesingle{}\textbackslash{}frac\{a\_\{0\}\}\{2\}\ +\ \textbackslash{}sum\_\{n=1\}\^{}\{\textbackslash{}infty\}\ a\_\{n\}\ \textbackslash{}cos(nx)\textquotesingle{}\ \textquotesingle{}a\_\{n\}\textquotesingle{}\ 3.\ \textquotesingle{}f(n-x)\textquotesingle{}\ \textquotesingle{}(-1)\^{}\{n\}\ \textbackslash{}frac\{2\}\{\textbackslash{}pi\}\ C\_\{n\}\textquotesingle{}\ 4.\ 1}

превратится в обыкновенное дифференциальное уравнение К о н

\begin{enumerate}
\def\labelenumi{\arabic{enumi}.}
\setcounter{enumi}{4}
\tightlist
\item
  \texttt{\textbackslash{}cos(mx),\ m=1,\ 2,\ \textbackslash{}dots}
\end{enumerate}

\begin{verbatim}
ечно, остается открытым вопрос, как найти эти характеристики. Ответ прост: мы выберем характеристики   \{[x(s), t(s)]: 0 < s < \infty\}   так, чтобы они удовлетворяли системе
\end{verbatim}

\begin{verbatim}
\begin{cases}
\end{verbatim}

\end{document}
